

\chapter{Spatial Reconstruction and Generation}
\label{sec:generation}


\FramedInitial{S}{patial} reconstruction and generation is the productive counterpart to spatial perception and reasoning: rather than analyzing a scene that is given, it synthesizes spatial content, recovering or creating the geometry, appearance, and dynamics of objects, scenes, and worlds from conditions such as text, images, partial observations, or explicit constraints.
Where perception (Chapter~\ref{chapter:3:sec:6}) maps observations into structured spatial representations and reasoning (Chapter~\ref{sec:reasoning}) draws inferences over them, generation closes the loop by mapping intent back into concrete spatial artifacts, a capability central to embodied agents, content creation, simulation, and world modeling.

The chapter is organized along an axis of increasing output structure and dimensionality, from flat images to volumetric objects and scenes, to video, embodied motion, and full four-dimensional space-time content, and finally to world models and to programmatic generation that emits interpretable, executable structure. What holds this range together is a single recurring backbone: an encoder maps conditions into a latent, a generator samples structure, and a decoder or renderer produces the final artifact, specialized to each output domain. Read this way, the subfields are less seven separate problems than one conditional-generation abstraction instantiated against progressively richer codomains.


\section{2D Spatial Generation}
\label{sec:p3_s8_1}

\FramedInitial{T}{he} 2D spatial generation encompasses a broad spectrum of tasks that require synthesizing or editing images under explicit spatial constraints.
These constraints may be \emph{user-specified} (bounding boxes, masks, sketches, keypoints, depth maps) or \emph{derived} from scene understanding (layout graphs, semantic maps, relational structures), and the overarching goal is to produce images that are not only visually realistic but also \emph{faithful to the specified spatial configuration}: objects appearing in the correct locations, at the correct scales and orientations, while remaining globally coherent in lighting, perspective, and style.
2D generation is where the most mature \emph{semantic-to-pixel prior} in contemporary AI is taught to honor \emph{spatial structure}: it is the proving ground on which large pre-trained visual priors are forced to localize objects, follow geometric scaffolds, respect occlusion ordering, and obey region- or instruction-level edits, all without falling out of the implicit grammar of natural images.

The section follows the 2D pipeline in three steps that share one latent-generator framework.
§\ref{subsec:p3_s8_1_1} reviews foundational text-to-image generation, the semantic-to-pixel prior on which everything else here is built, and the three paradigms (adversarial, autoregressive, and diffusion with its flow-matching generalization) that parameterize it.
§\ref{subsec:p3_s8_1_2} adds a structural condition (layouts, masks, edges, depth, pose, scene graphs, region phrases) and asks how it is injected into a frozen or co-trained generator.
§\ref{subsec:p3_s8_1_3} replaces the empty canvas with an existing image and a mask, so that the generator must change what lies inside the mask while preserving what lies outside it.
% A short Summary (§\ref{subsec:p3_s8_1_summary}) compares the three along shared axes.

\subsection{Text-to-image Generation}
\label{subsec:t2i}
\label{subsec:p3_s8_1_1}

Text-to-image (T2I) generation synthesizes an image from a natural-language description.
For the rest of this chapter its importance lies less in the images themselves than in the prior it provides: a high-capacity mapping from compositional language to the visual space of objects, surfaces, layouts, and viewpoints, learned once at web scale and then reused.
The spatial controllers of §\ref{subsec:p3_s8_1_2} and the editors of §\ref{subsec:p3_s8_1_3} impose constraints on this prior rather than replacing it, and the text-to-3D pipelines of §\ref{sec:p3_s8_2} distill it through score distillation or multi-view consistency objectives; whatever relational competence the prior has, or lacks, therefore propagates downstream.
The field has passed through three architectural paradigms, from adversarial training to autoregressive token modeling to iterative denoising, with diffusion and its flow-matching generalization now setting the standard for quality and prompt fidelity~\cite{rombach2022ldm,peebles2023dit}.


\subsubsection{Task Definition}
We formulate text-to-image generation in a paradigm-agnostic form that exposes the structural choices shared by modern systems.
Let $\mathcal{X}$ denote the space of tokenized text prompts and $\mathcal{Y}=[0,1]^{H\times W\times 3}$ the pixel-space image domain.
Almost all modern T2I systems operate not directly on $\mathcal{Y}$ but on a compressed \emph{latent space} $\mathcal{Z}\subset\mathbb{R}^d$ provided by an autoencoder pair
\begin{equation}
E:\mathcal{Y}\rightarrow\mathcal{Z},\qquad D:\mathcal{Z}\rightarrow\mathcal{Y},\qquad D\!\circ\!E\approx \mathrm{id}_{\mathcal{Y}},
\label{eq:t2i_autoencoder}
\end{equation}
where $E$ encodes a training image into a low-dimensional code and $D$ decodes a code back into a pixel image.
A text encoder $\phi:\mathcal{X}\rightarrow\mathcal{E}$ maps the prompt into a conditioning embedding $\phi(\mathbf{x})\in\mathcal{E}$.
The core learning problem is then to train a conditional latent sampler
\begin{equation}
G_\theta:\mathcal{E}\rightarrow\mathcal{Z},\qquad \hat{\mathbf{z}}\sim p_\theta\!\big(\mathbf{z}\,\big|\,\phi(\mathbf{x})\big),
\label{eq:t2i_generator}
\end{equation}
so that, after decoding, the resulting image $\hat{\mathbf{y}}=D(\hat{\mathbf{z}})$ is both photorealistic and semantically aligned with the prompt $\mathbf{x}$.
The three paradigms covered below differ chiefly in how the conditional distribution $p_\theta(\mathbf{z}\mid\phi(\mathbf{x}))$ is parameterized and optimized; the autoencoder pair $(E,D)$ and the text encoder $\phi$ are largely shared design choices across paradigms.
Figure~\ref{fig:p3_s8_1_1_framework} schematizes this paradigm-agnostic pipeline and the substrate that the three paradigms share.

\begin{figure}[!ht]
\centering
\includegraphics[width=\linewidth]{figures/part3/section8/t2i_framework_and_paradigms.pdf}
\caption{The latent-generator view of text-to-image generation. Top: the shared pipeline of Eqs.~\eqref{eq:t2i_autoencoder} and~\eqref{eq:t2i_generator}; the encoder $E$ and the dashed paths are used only in training. Bottom: the sampler $G_\theta$ as (a) a GAN, (b) an autoregressive token model, and (c) diffusion or flow matching. Not every architecture uses a separate autoencoder; the corgi is a schematic illustration.}
\label{fig:p3_s8_1_1_framework}
\end{figure}


\subsubsection{Methodology}

Three generative paradigms, taken in the order in which each dominated T2I, parameterize the sampler $G_\theta$ of Eq.~\eqref{eq:t2i_generator} differently; the bottom row of Figure~\ref{fig:p3_s8_1_1_framework} draws them against the common pipeline.
A second set of advances refines the sampler without redefining it: backbone evolution, sampling acceleration, and the migration of leading systems to proprietary stacks.
Structural conditioning interfaces, the fourth such refinement, are the subject of §\ref{subsec:p3_s8_1_2} and are not repeated here.

\paragraph{Generative Adversarial Networks.}\label{para:p3_s8_1_1_gan}
GAN-based T2I models instantiate $G_\theta$ as a feed-forward generator $G_\theta(\mathbf{n},\phi(\mathbf{x}))\!\to\!\hat{\mathbf{z}}$ driven by an auxiliary noise vector $\mathbf{n}\sim\mathcal{N}(0,I)$; the conditional distribution $p_\theta(\mathbf{z}\mid\phi(\mathbf{x}))$ is defined only implicitly, as the pushforward of $\mathbf{n}$ through $G_\theta$.
Training follows the two-player game of \citet{goodfellow2014gan}, pitting the generator against a discriminator $D_\psi$,
\begin{equation}
\min_\theta \max_\psi\ \mathbb{E}_{\mathbf{y}\sim p_\text{data}}\!\big[\log D_\psi(\mathbf{y},\phi(\mathbf{x}))\big] + \mathbb{E}_{\mathbf{n}}\!\big[\log\!\big(1-D_\psi(D(G_\theta(\mathbf{n},\phi(\mathbf{x}))),\phi(\mathbf{x}))\big)\big],
\label{eq:t2i_gan_objective}
\end{equation}
in which $D_\psi$ scores both photorealism and text-image consistency.
The result is sharp high-frequency detail and single-pass inference, bought at the price of unstable training and mode collapse.
Within this framework, StackGAN~\cite{stackgan2017} generated a low-resolution image first and refined it with a second GAN conditioned on both the text and the coarse image; AttnGAN~\cite{attngan2018} let the generator attend to individual words while rendering different regions, with an auxiliary image-text matching loss (DAMSM) that reduced attribute errors such as wrongly colored objects; XMC-GAN~\cite{xmcgan2021} replaced that loss with cross-modal contrastive learning and improved multi-object prompts.
By 2021 diffusion models had overtaken these systems on COCO, chiefly because they train stably and cover more of the image distribution, and research attention moved on; the GAN emphasis on high-frequency detail survives in hybrid and adversarially fine-tuned samplers.

\paragraph{Autoregressive Models.}\label{para:p3_s8_1_1_ar}
Autoregressive (AR) models discretize the latent space: the autoencoder $E$ is realized as a vector-quantized variational autoencoder (VQ-VAE) whose output is a sequence of discrete tokens $\mathbf{z}=(z_1,\dots,z_N)\in\mathcal{V}^N$ over a learned codebook $\mathcal{V}$, and $G_\theta$ is factorized in causal raster order,
\begin{equation}
p_\theta\!\big(\mathbf{z}\,\big|\,\phi(\mathbf{x})\big) \;=\; \prod_{n=1}^{N} p_\theta\!\big(z_n\,\big|\,z_{<n},\,\phi(\mathbf{x})\big),
\label{eq:t2i_ar_factorisation}
\end{equation}
trained by maximum likelihood on the next-token cross-entropy.
T2I thus becomes sequence-to-sequence translation, and the transformer that performs it inherits the compositional competence of language models: DALL-E~\cite{dalle2021} modeled text and image tokens as a single sequence and showed zero-shot generalization across prompt types.
The cost is decoding latency, one forward pass per token, which grows with resolution.
Parti~\cite{parti2022} scaled the transformer to 20B parameters, generated at $256\times256$, and reached $1024\times1024$ with a convolutional super-resolution stage; Muse~\cite{chang2023muse} abandoned strict sequential order for masked token prediction, filling in many tokens per step over a handful of rounds, and matched diffusion quality on several benchmarks at a fraction of the sampling time.
AR generation remains attractive where images are produced jointly with text (interleaved documents, storyboards) and a single sequence model handles both.

\paragraph{Diffusion Models.}\label{para:p3_s8_1_1_diffusion}
Diffusion models realize $G_\theta$ as iterative denoising rather than a single feed-forward pass.
A fixed \emph{forward} process corrupts the clean latent $\mathbf{z}_0=E(\mathbf{y})$ along a schedule $(\alpha_t,\sigma_t)_{t=0}^{T}$, $\mathbf{z}_t=\alpha_t \mathbf{z}_0+\sigma_t\epsilon$ with $\epsilon\sim\mathcal{N}(0,I)$ and $\mathbf{z}_T\!\approx\!\mathcal{N}(0,I)$; a noise-prediction network $\epsilon_\theta$ learns to invert it under the denoising score-matching objective
\begin{equation}
\min_\theta\ \mathbb{E}_{t,\mathbf{z}_0,\epsilon}\!\left[\left\|\epsilon-\epsilon_\theta\!\big(\mathbf{z}_t,\,t,\,\phi(\mathbf{x})\big)\right\|_2^2\right],
\label{eq:t2i_ldm_objective}
\end{equation}
which is the diffusion-specific realization of the training objective implicit in Eq.~\eqref{eq:t2i_generator}.
At inference the model starts from $\mathbf{z}_T\sim\mathcal{N}(0,I)$, applies $\epsilon_\theta$ for $t=T,\dots,1$, and decodes $\hat{\mathbf{y}}=D(\hat{\mathbf{z}}_0)$; each step corrects a small error, which is why the final samples are both detailed and faithful.
Running this process in the latent space of Eq.~\eqref{eq:t2i_autoencoder} rather than in pixels, as Latent Diffusion Models~\cite{rombach2022ldm} did (\eg, $4\times64\times64$ instead of $3\times512\times512$ for a 512px image), cut memory and compute enough to train on commodity hardware and to release Stable Diffusion openly.
Conditioning the denoiser on embeddings from a pretrained language-image model, as in GLIDE~\cite{glide2022} and DALL$\cdot$E~2~\cite{dalle2_2022}, improved prompt alignment, and classifier-free guidance~\cite{ho2022classifierfree} supplied a single knob that trades diversity against fidelity by extrapolating the conditional prediction away from the unconditional one.
The remaining weakness is sampling cost, dozens of sequential network evaluations per image; higher-order ODE samplers and the distillation methods below address it.
Diffusion or flow-matching cores now underlie most leading systems, with the native multimodal autoregressive stacks discussed under proprietary systems the main exception.

\paragraph{Flow Matching and Rectified Flow.}\label{para:p3_s8_1_1_fm}
Rather than a fourth paradigm, flow matching is diffusion reformulated in continuous time.
It parameterizes $G_\theta$ as a learned \emph{velocity field} $v_\theta(\mathbf{z}_t,t,\phi(\mathbf{x}))$ instead of a noise predictor: training picks the straight interpolation $\mathbf{z}_t=(1-t)\mathbf{z}_0+t\mathbf{z}_1$ between a clean latent $\mathbf{z}_0=E(\mathbf{y})$ and a Gaussian sample $\mathbf{z}_1\!\sim\!\mathcal{N}(0,I)$ and regresses the velocity,
\begin{equation}
\min_\theta\ \mathbb{E}_{t,(\mathbf{z}_0,\mathbf{z}_1)}\!\left[\left\|v_\theta(\mathbf{z}_t,t,\phi(\mathbf{x}))-(\mathbf{z}_1-\mathbf{z}_0)\right\|_2^2\right],
\label{eq:t2i_fm_objective}
\end{equation}
and inference integrates the probability-flow ODE $\mathrm{d}\mathbf{z}_t/\mathrm{d}t=v_\theta(\mathbf{z}_t,t,\phi(\mathbf{x}))$ from $t=1$ back to $t=0$.
The formulation was introduced concurrently by~\citet{lipman2023flow} (Flow Matching) and~\citet{liu2023rectifiedflow} (Rectified Flow); \citet{albergo2023building} gave the unifying stochastic-interpolant view in which diffusion is the special case of a noisier interpolation schedule.
Three consequences matter in practice. Straight-line interpolations yield nearly straight ODE trajectories, so matched quality is reached in noticeably fewer steps (genuine one-step generation still needs a reflow or distillation stage); the objective is a single regression with no noise-schedule design and no per-schedule guidance recalibration; and the loss scales smoothly with model size and resolution.
Stable Diffusion~3~\cite{sd3rft2024} scaled the recipe to T2I on a DiT-style transformer backbone and reported the best GenEval scores among open models together with favorable human preference on Parti-prompts; the Flux family followed the same recipe.
Flow matching composes with every refinement below as readily as a noise-predictor core does and has become the default continuous-time formulation in 2024-onwards systems.

\paragraph{Backbone Evolution: From U-Net to Diffusion Transformers.}
Early diffusion models used the U-Net for its multiscale feature fusion and locality bias.
The Diffusion Transformer (DiT) of~\citet{peebles2023dit} replaced it with a stack of transformer layers over latent patch tokens, giving every denoising step a global receptive field, and showed that FID improves monotonically with depth, width, and token count along predictable scaling laws, reaching FID 2.27 on class-conditional ImageNet at $256^2$.
Attention cost grows quadratically with resolution, so hybrid CNN-transformer designs and efficient attention remain active; the DiT lineage nonetheless underlies SD3 and most 2024-onwards systems, and it points toward a single transformer that generates several modalities, differing only in tokenization and training data.

\paragraph{Sampling Acceleration and Distillation.}
Fifty or more denoising steps rule out interactive use and low-power deployment.
\citet{salimans2022progressive} distilled a $T$-step teacher into a $T/2$-step student and iterated the halving to reach $1$--$4$ steps.
Consistency Models~\cite{song2023consistency} train a network to map a noisy latent at any timestep directly to the clean image, either from scratch with a consistency loss or by distilling a pretrained diffusion model, and Latent Consistency Models~\cite{lcm2023} applied the idea to Stable Diffusion's latent space, matching $50$-step quality in $2$--$4$ steps.
All of these are post-training procedures that leave the architecture untouched, so any strong backbone can be made fast after the fact.

\paragraph{Proprietary and Commercial Systems.}
Leading systems have migrated to proprietary, industrially scaled stacks of three kinds.
Native multimodal models such as OpenAI's GPT~Image~1.5~\cite{openai2025images} and Google's Gemini~3~Pro~Image (Nano~Banana~Pro) run a unified transformer over interleaved text and image tokens and emphasize semantic reasoning and multi-turn editing; latent-diffusion iterations such as Imagen~4 scale the LDM recipe with stronger text encoders and curated data; preference-tuned systems such as Midjourney~v7~\cite{midjourney_v7_alpha} and FLUX.2~[max]~\cite{bfl_flux2_max} fine-tune strong priors on human-preference data (FLUX.2 on the rectified-flow recipe of §\ref{para:p3_s8_1_1_fm}; Midjourney's architecture is undisclosed).
Each reduces to one of the three paradigms plus some of the refinements just listed; because architectural and training detail is withheld, public benchmarking and reproduction still rest on the open SD / DiT stack.

\subsubsection{Benchmarks}
Three independent things are usually conflated under ``benchmark'': the training corpora on which models are fit, the evaluation suites of prompts that probe particular capabilities, and the metrics that score the resulting images.
A model trained on LAION-5B can be evaluated on GenEval prompts and scored by FID, CLIPScore, or an LLM judge, and disagreement between those scores is itself informative.

\paragraph{Training Corpora.}
Corpus scale has driven much of T2I progress.
Early work fit fine-grained sets with ten captions per image, CUB-200-2011 (11,788 bird images)~\cite{wah2011cub} and Oxford-102 Flowers (8,189 images)~\cite{nilsback2008flowers}; MS-COCO Captions~\cite{lin2014microsoft} (123,287 images, five captions each) became the standard general-scene corpus and, for years, the evaluation set as well, a conflation that the suites below were designed to end.
Current models train on web-scale image-text pairs: LAION-5B~\cite{schuhmann2022laion5b} releases 5.85 billion CLIP-filtered pairs, and proprietary corpora of similar scale, less documented, underlie DALL-E~2 and Parti.
Scale brought duplication, demographic bias, unsafe content, and licensing disputes with it, and filtered subsets such as LAION-Aesthetics and curation pipelines are now part of every training recipe.

\paragraph{Evaluation Benchmark Suites.}
An evaluation suite is a curated prompt set plus a verification protocol: prompts are adversarial or capability-stratified, reference images are usually absent, and scoring is automatic (an object detector, a VQA model, or an LLM/MLLM judge) rather than reconstruction-based.
Table~\ref{tab:p3_s8_1_1_eval_benchmarks} lists the main suites since 2022.
The earliest, DrawBench~\cite{saharia2022photorealistic} and PartiPrompts~\cite{parti2022}, are human-pairwise prompt sets released with Imagen and Parti; HRS-Bench~\cite{bakr2023hrsbench} and HEIM~\cite{lee2023heim} are holistic, scoring 13 and 12 axes respectively; TIFA~\cite{hu2023tifa}, GenEval~\cite{ghosh2023geneval}, and T2I-CompBench~\cite{huang2023t2icompbench} introduced automatic verification by VQA, object detection, and BLIP-VQA plus UniDet, and the latter two are now reported alongside FID and CLIPScore in most papers.
For spatial intelligence the salient suite is VISOR~\cite{gokhale2022visor}, which isolates two-object relations (left/right, above/below) over $25{,}280$ prompts and finds that otherwise strong models frequently render the wrong configuration.
DPG-Bench~\cite{hu2024ella} targets dense, long compositional prompts, GenAI-Bench~\cite{li2024genaibench} pairs compositional prompts with VQAScore judging, and WISE~\cite{niu2025wise} asks whether generated images agree with real-world facts (flag colors, landmarks, anatomy), a knowledge axis orthogonal to compositional binding.

\begin{table}[!ht]
\centering
\fontsize{9}{11}\selectfont
\setlength{\tabcolsep}{1.0mm}
\caption{Representative evaluation benchmark suites for text-to-image generation, ordered approximately by year of release. ``Scoring'' indicates the dominant verification mechanism: detector- or VQA-based automatic checking, learned preference models, MLLM/LLM-as-judge rubric scoring, or aggregated human preference. Several suites combine multiple scoring channels.}
\label{tab:p3_s8_1_1_eval_benchmarks}
\begin{tabular}{L{0.17\textwidth} C{0.06\textwidth} C{0.11\textwidth} L{0.33\textwidth} L{0.25\textwidth}}
\toprule
\textbf{Benchmark} & \textbf{Year} & \textbf{\#~Prompts} & \textbf{Evaluation Axis} & \textbf{Scoring} \\
\midrule
DrawBench~\cite{saharia2022photorealistic}       & 2022 & 200   & Broad capability + alignment ($11$ categories) & Human pairwise \\
PartiPrompts~\cite{parti2022}     & 2022 & 1{,}600 & Capability spectrum ($12$ categories) & Human pairwise \\
VISOR~\cite{gokhale2022visor}     & 2022 & 25{,}280 & Two-object spatial relations (left/right, above/below) & Object-detector verification \\
TIFA~\cite{hu2023tifa}            & 2023 & 4{,}081 & Faithfulness (VQA over auto-generated questions) & VQA / LLM-as-judge \\
HRS-Bench~\cite{bakr2023hrsbench} & 2023 & 45{,}000 & Holistic; $13$ skills $\times$ $50$ scenarios & Automatic, mixed \\
HEIM~\cite{lee2023heim}           & 2023 & 62 scenarios & Holistic; $12$ aspects (alignment, quality, knowledge, bias, toxicity, \ldots) & Automatic $+$ human \\
GenEval~\cite{ghosh2023geneval}   & 2023 & 553   & Object-centric: single/two-obj., counting, color, position, attribute & Detector verification \\
T2I-CompBench~\cite{huang2023t2icompbench} & 2023 & 6{,}000 & Compositional: attribute binding, object relationships, complex prompts & BLIP-VQA $+$ UniDet \\
DPG-Bench~\cite{hu2024ella}       & 2024 & 1{,}065 & Dense, long compositional prompts & MLLM-based scoring \\
GenAI-Bench~\cite{li2024genaibench} & 2024 & 1{,}600 & Compositional T2V/T2I prompts & VQAScore (LLM-judge) \\
WISE~\cite{niu2025wise}           & 2025 & 1{,}000 & World-knowledge-informed semantic accuracy & LLM-based grading \\
\bottomrule
\end{tabular}
\end{table}

\paragraph{Evaluation Metrics.}
Metrics fall into four tiers.
Distributional metrics compare the generated and real image populations and say nothing about any individual prompt: the Fr\'echet Inception Distance (FID)~\cite{heusel2017gans} is the Wasserstein-2 distance between Gaussian fits of Inception-V3 features of generated and real images, the Inception Score rewards images that are individually confidently classified yet collectively diverse, and generative precision/recall~\cite{sajjadi2018assessing,kynkaanniemi2019improved} separate quality from coverage through neighborhood tests in feature space.
Alignment metrics score one image against its prompt; CLIPScore~\cite{hessel2021clipscore}, the cosine similarity between the CLIP embeddings of $\hat{\mathbf{y}}$ and $\mathbf{x}$, is cheap and reference-free but inherits CLIP's insensitivity to attribute binding, counting, and spatial relations.
Learned preference models, HPSv2~\cite{wu2023hpsv2}, ImageReward~\cite{xu2023imagereward} (fit to $137$K human comparisons), and PickScore~\cite{kirstain2023pickapic} (fit to the ${\sim}500$K real-user comparisons of Pick-a-Pic), regress a scalar that correlates with average user preference, bundling realism, alignment, and aesthetics; they have become the default automatic ranking while inheriting their raters' biases and degrading out of distribution.

The newest tier uses a vision-language model as judge.
TIFA~\cite{hu2023tifa} has a language model write verification questions from the prompt and a VQA model answer them against the image, the score being the fraction answered correctly; VQAScore~\cite{lin2024evaluating} collapses this to a single yes/no probability under a fine-tuned CLIP-FlanT5 model and reports the best human correlation on compositional suites; VIEScore~\cite{ku2023viescore} and LLMScore~\cite{lu2023llmscore} prompt a multimodal LLM with a structured rubric and return per-axis scores with brief rationales, and direct prompting of frontier VLMs with custom rubrics has become routine for ad-hoc capability probes.
Judge metrics track human preference on compositional and spatial prompts far better than CLIPScore, at the cost of a frontier-model call per image, sensitivity to rubric wording, and drift across judge-model versions that makes numbers from different papers hard to compare.
No single tier suffices, and current papers report a tuple (FID or IS, CLIPScore, one or two preference scores, one judge score) on several suites plus a small human study; cross-metric disagreement is read as a warning rather than averaged away.


\subsubsection{Challenges and Trends}
The weakness that matters most for this handbook is \textbf{spatial-relation faithfulness}, a sub-problem of compositionality.
Models frequently misplace objects relative to one another: a prompt that puts a cup to the left of a laptop is often rendered with the two swapped or stacked, and VISOR~\cite{gokhale2022visor} shows that language-specified relations are frequently ignored outright.
Data-centric remedies help; SPRIGHT~\cite{chatterjee2024spright} re-captions training images with explicit spatial descriptions and measurably improves consistency, but reliable control over where objects appear remains open.
Because layout-conditioned generation (§\ref{subsec:p3_s8_1_2}) and text-to-3D (§\ref{sec:p3_s8_2}) inherit whatever relational competence this prior provides, the limitation is foundational rather than cosmetic.
Fine structure is a second gap: legible text, hands, and faces still fail in otherwise strong models and appear to need architectural or data-specific fixes rather than scale alone.

\textbf{Evaluation reliability} is the third.
Automatic scores correlate imperfectly with human judgment precisely on the compositional, counting, and spatial-relation prompts that matter here; judge-based metrics close part of the gap but add rubric sensitivity and judge-version drift, so headline numbers across papers are not directly comparable.
As GenEval and T2I-CompBench become standard, their prompts risk leaking into web-scraped training corpora, in which case high scores increasingly measure memorization rather than capability; periodic refresh of held-out prompts, blinded judges, and explicit train/test decontamination are the obvious countermeasures and are not yet routine.

\textbf{Societal concerns} follow the data.
Models inherit demographic bias from their corpora and will produce unsafe content on request; prompt classifiers, safety filters, and fine-tuning on curated subsets reduce but do not remove the problem, and photorealistic synthesis raises misinformation and intellectual-property questions that watermarking and data curation address only in part.
The extension of these priors along the time axis, where identity and motion must stay consistent across frames, is the subject of §\ref{sec:p3_s8_3}; the next two subsections stay in the image plane and ask how the prior can be made to obey spatial constraints and local edits.



\subsection{2D Spatial-aware Image Generation}
\label{subsec:p3_s8_1_2}
\emph{2D spatial-aware image generation} synthesizes images that are faithful not only to a text prompt but to explicit spatial constraints in the image plane: where objects appear, how large they are, what shape they take, and how they relate to one another (\emph{left of}, \emph{on}, \emph{behind}).
The constraint may be a set of boxes or a scene layout, an instance or semantic mask, an edge or sketch map, a depth or normal map, 2D keypoints, a scene graph, or a phrase bound to a region, and the pipelines that honor it are almost all diffusion models augmented with a conditioning pathway or a guidance term.
The tension that organizes the subfield is three-way: \emph{adherence} to hard constraints against \emph{perceptual realism} and \emph{diversity}, and both against \emph{compositionality} when many small or occluding objects carry discrete relations.
Depth and normal maps enter here as structural guides for occlusion ordering and surface orientation, not as metric measurements; recovering true distance belongs to reconstruction (§\ref{sec:p3_s8_2}).


\subsubsection{Task Definition}
\label{subsubsec:p3_s8_1_2_task}
We inherit the latent-generator framework of §\ref{subsec:p3_s8_1_1}.
Let $\mathcal{X}$ denote tokenized prompts, $\mathcal{Z}$ the latent space induced by the autoencoder pair $(E,D)$ of Eq.~\eqref{eq:t2i_autoencoder}, and $\phi:\mathcal{X}\!\rightarrow\!\mathcal{E}$ the text encoder.
Spatial-aware generation \emph{augments} this framework with a structural condition $\mathbf{c}\in\mathcal{C}$ drawn from one or more conditioning channels (boxes, masks, edges, depth, pose, scene graphs, region phrases, reference images), encoded by a condition encoder $\kappa:\mathcal{C}\!\rightarrow\!\mathcal{C}'$.
The conditional latent sampler of Eq.~\eqref{eq:t2i_generator} is extended to a joint conditional
\begin{equation}
G_\theta:\mathcal{E}\times\mathcal{C}'\rightarrow\mathcal{Z},\qquad \hat{\mathbf{z}}\sim p_\theta\!\big(\mathbf{z}\,\big|\,\phi(\mathbf{x}),\,\kappa(\mathbf{c})\big),
\label{eq:t2i_spatial_objective}
\end{equation}
and the generated image $\hat{\mathbf{y}}=D(\hat{\mathbf{z}})$ is required to be (i) photorealistic, (ii) semantically aligned with $\mathbf{x}$, and (iii) faithful to the structural condition $\mathbf{c}$ in placement, scale, shape, and pairwise relation.
The unconditional limit $\mathbf{c}=\varnothing$ recovers pure T2I generation as in §\ref{subsec:p3_s8_1_1}; the substance of the present subsection lies in how $\kappa$ is realized for each conditioning channel and how its output is injected into $G_\theta$.
Which channel populates $\mathbf{c}$, and how densely, ranges from pixel-level maps (masks, edges, depth) through coordinate-level scaffolds (boxes, keypoints) to discrete symbols (scene graphs, region phrases); the modality survey below follows that ordering.
Interactive constraints supplied by dragging handles during generation belong to the editing regime of §\ref{subsec:p3_s8_1_3}.


\subsubsection{Methodology}
\label{subsubsec:p3_s8_1_2_methodology}
Two axes organize the methods: the \emph{conditioning modality} (what kind of structural input $\mathbf{c}$ is consumed) and the \emph{injection mechanism} (how $\kappa(\mathbf{c})$ enters the latent sampler $G_\theta$).
They are nearly independent. One mechanism, a ControlNet-style side branch, serves canny edges, pose skeletons, depth, and segmentation alike, and one modality, bounding boxes, is served by side branches, grounded cross-attention, or inference-time attention shaping, so the design space is a matrix of (modality~$\times$~mechanism) cells rather than a list of monolithic methods.
Figure~\ref{fig:p3_s8_1_2_mechanisms} draws the five mechanisms on one backbone silhouette.

\paragraph{Conditioning Modalities.}
\label{para:p3_s8_1_2_modalities}
\textit{Dense pixel-level scaffolds} supply a value at every location, which makes them the easiest channels to follow and the most restrictive on creativity.
Instance and semantic masks drive region-conditioned synthesis, with MultiDiffusion~\cite{multidiffusion} offering a training-free protocol that samples jointly over overlapping mask windows; edges, sketches, and line art, popularized by the canny and soft-edge variants of ControlNet~\cite{zhang2023adding} and the edge head of T2I-Adapter~\cite{t2iadapter2023}, fix only the location of intensity discontinuities and leave texture and appearance to the prompt; depth and normal maps, served by ControlNet's depth variant, impose coarse occlusion ordering and surface orientation without requiring a metric reconstruction.
\textit{Coordinate-level scaffolds} drop to per-region or per-point supervision.
Boxes with labels are the most studied channel, a set of $(\text{box},\text{label})$ tuples fixing location and rough scale, with LayoutDiffusion~\cite{zheng2023layoutdiffusion}, LayoutLLM~\cite{qu2023layoutllm}, and the open-set GLIGEN~\cite{li2023gligen} as representatives; 2D keypoints (typically human pose) are sparser still, a handful of joint coordinates, injected through ControlNet or handled by GLIGEN as a special case of grounding.
\textit{Discrete symbolic structure} carries no pixel- or coordinate-level information at all, and the model must map it onto the image plane itself.
Scene graphs reach back to the sg2im line of~\citet{johnson2018image} and re-emerged in the diffusion era through LLM-emitted layouts, as in LayoutLLM, where the graph is first decoded into a coordinate-level layout; region phrases and referring expressions bind a text span to a box, mask, or free-form polygon and are served by GLIGEN's grounded cross-attention and by phrase-grounding heads attached to SD and SDXL~\cite{sdxl2023}.

\paragraph{Injection Mechanisms.}
\label{para:p3_s8_1_2_mechanisms}
Five mechanism families cover the literature; production pipelines routinely stack a side branch for dense geometric channels, an adapter for style or a reference image, and inference-time guidance to sharpen adherence on adversarial prompts.

\textbf{Side-branch controllers.}
The backbone weights stay frozen and an auxiliary branch mirroring the encoder half of the base U-Net is trained, connected through \emph{zero-convolution} layers whose weights are initialized to zero so that the augmented model reproduces the base model exactly at the start of training and degrades gracefully if the branch is poorly tuned.
ControlNet~\cite{zhang2023adding} is the prototype and the first widely adopted retrofit pathway for non-text structural control: the condition (a canny edge map, an OpenPose skeleton, a depth or normal map, a segmentation mask, or a scribble) is encoded by the side branch in parallel with the noisy latent and injected into the base U-Net through additive residual connections at each scale.
One branch is trained per modality on labeled image--condition pairs, which produced a library of plug-in controllers for SD~1.5 and SDXL~\cite{sdxl2023} that compose with prompt conditioning at inference time without retraining the backbone; most later designs either shrink this recipe or fold it into a unified adapter library.

\textbf{Lightweight adapters.}
Where the side branch spends capacity (half a U-Net per modality), T2I-Adapter~\cite{t2iadapter2023}, developed concurrently with the same goal of plug-in control over a frozen backbone, spends as little as possible: a few residual blocks per modality whose multi-scale features are added directly into the corresponding encoder stages.
Several adapters can be combined at inference, an image conditioned simultaneously on an OpenPose skeleton, a depth map, and a color palette with a scalar weight on each, so the adapter is preferred for multi-channel fusion and fast iteration, and the side branch where maximum single-channel fidelity matters.

\textbf{Grounded cross-attention.}
Side branches and adapters run an external pathway alongside the backbone; the next two families intervene \emph{inside} its attention stack.
GLIGEN~\cite{li2023gligen} inserts a new \emph{gated self-attention} layer between the existing self-attention and cross-attention of each U-Net transformer block; the layer attends jointly over the visual tokens and a set of \emph{grounding tokens} that encode each $(\text{phrase},\text{box})$ pair, a learned scalar gate controls how strongly the grounding signal enters the residual stream, and only the new layer is trained, so the network falls back to ungrounded behavior when no constraint is supplied.
The problem it solves is complementary to the side branch: not how to inject a dense pixel-level scaffold, but how to introduce \emph{open-set} symbolic grounding (any category named in text, localized by a box, keypoints, or a reference image) into a frozen T2I backbone, with several grounding channels interleaved within one mechanism.

\textbf{Layout-aware cross-attention.}
LayoutDiffusion~\cite{zheng2023layoutdiffusion} is the dense-supervised cousin of grounded attention: rather than gating a new layer onto a frozen backbone, it co-trains a layout-aware cross-attention block end to end for layout-to-image (L2I) generation specifically.
A structured layout, a set of $(\text{box},\text{class label})$ tuples, becomes a sequence of region tokens fused with the noisy latent at every denoising step, and an explicit \emph{object--scene} factorization models global scene context separately from per-object features so that object identity, scale, and pairwise relations survive in cluttered scenes.
It reports stronger layout IoU and relation satisfaction than ControlNet or T2I-Adapter on multi-object COCO benchmarks, at the cost of being specialized to box-and-label input rather than the broader modality library.

\textbf{Inference-time guidance.}
The fifth family adds no parameters and modifies only sampling.
BoxDiff~\cite{xie2023boxdiff} is the canonical training-free spatial controller: it shapes the cross-attention maps so that activations of box-localized prompt tokens concentrate inside their boxes while the backbone stays untouched.
MultiDiffusion~\cite{multidiffusion} fuses overlapping mask windows at sampling time, and composing scores from several conditional models or tuning the classifier-free guidance scale per region are related strategies.
Because nothing is retrained, these methods suit rapid prototyping, backbones whose weights are unavailable for fine-tuning, and per-prompt sharpening on top of the trained mechanisms above.

\begin{figure}[!ht]
\centering
\includegraphics[width=\linewidth]{figures/part3/section8/Five_injection_mechanisms_for_spatial-aware_T2I.pdf}
\caption{Five mechanisms for injecting a spatial condition $\kappa(\mathbf{c})$ into a text-to-image denoiser: (a) side-branch control (ControlNet~\cite{zhang2023adding}), (b) lightweight adapters (T2I-Adapter~\cite{t2iadapter2023}), (c) gated grounding (GLIGEN~\cite{li2023gligen}), (d) layout-aware cross-attention (LayoutDiffusion~\cite{zheng2023layoutdiffusion}), and (e) inference-time guidance (BoxDiff~\cite{xie2023boxdiff}). The corgi example is illustrative, not a comparison of outputs.}
\label{fig:p3_s8_1_2_mechanisms}
\end{figure}


\subsubsection{Benchmarks}
\label{subsubsec:p3_s8_1_2_benchmarks}
Two things change relative to §\ref{subsec:p3_s8_1_1}.
The ground truth is a \emph{structural annotation} (a layout, a mask, a keypoint set) rather than a reference image, so reconstruction-style metrics give way to adherence metrics that score how closely the generated image honors $\mathbf{c}$; and every channel admits its own adherence metric, so a meaningful report combines a channel-specific score with the realism and alignment scores inherited from the previous subsection.

\paragraph{Training Corpora.}
Spatial-aware models are typically trained on corpora that pair images with rich structural annotations: COCO Captions~\cite{lin2014microsoft} for layouts and captions; ADE20K~\cite{zhou2017ade20k} ($20$K images, $150$ classes) for semantic and instance masks; Visual Genome~\cite{krishna2017visual} for object boxes, attributes, and pairwise relations supporting scene-graph conditioning; RefCOCO / RefCOCO+ / RefCOCOg~\cite{yu2016modeling} for phrase grounding; and PhraseCut~\cite{wu2020phrasecut} for region-prompt mask supervision.
Additional structure (canny edges, depth, normals, OpenPose skeletons) is typically derived on the fly from these base corpora by running pretrained predictors over each image, producing the paired image--condition data on which ControlNet-style and adapter-style pipelines are trained.

\paragraph{Evaluation Benchmark Suites.}
Most of the suites surveyed in §\ref{subsec:p3_s8_1_1} (Table~\ref{tab:p3_s8_1_1_eval_benchmarks}) include axes that probe spatial-aware generation directly and serve as its primary out-of-the-box benchmarks.
GenEval~\cite{ghosh2023geneval} exposes a position subset that scores whether two objects appear in a specified relation (\emph{left of}, \emph{above}); T2I-CompBench~\cite{huang2023t2icompbench} provides explicit \emph{spatial} and \emph{non-spatial relationship} categories with BLIP-VQA and UniDet scoring; HRS-Bench~\cite{bakr2023hrsbench} includes a \emph{spatial} skill among its thirteen; GenAI-Bench~\cite{li2024genaibench} offers compositional prompts on which spatial relations are a frequent failure of CLIP-based scoring; and WISE~\cite{niu2025wise} measures knowledge-grounded location accuracy.
Beyond these, the L2I literature reports task-specific evaluations on COCO layouts: detection-based AP for box-localized objects alongside FID and CLIPScore.

\paragraph{Constraint-specific Metrics.}
Layout adherence is box IoU between the conditioning layout and the boxes recovered by a frozen detector on $\hat{\mathbf{y}}$, often with per-object center distance; mask adherence uses mask IoU and boundary $F$-score against the conditioning mask; edge and sketch adherence report recall (sometimes precision) of edges re-extracted from $\hat{\mathbf{y}}$; geometry adherence reports depth RMSE and normal angular error where ground-truth geometry exists; pose adherence reports keypoint reprojection error or PCK; and relational adherence reports relation-satisfaction accuracy and per-entity object AP for scene-graph and phrase-grounded settings.
These numbers do not aggregate: averaging layout IoU with mask IoU has no operational meaning, since the two live on different scales and reward different shape regimes, and most papers report a different subset of them, which is the main obstacle to cross-paper comparison in this subfield.


\subsubsection{Challenges and Trends}
\label{subsubsec:p3_s8_1_2_challenges}
Spatial-aware T2I inherits the compositionality and evaluation problems of §\ref{subsec:p3_s8_1_1} and sharpens them.

\textbf{Faithful multi-object composition.}
Crowded scenes with many small or partially occluded objects and hard relations (support, containment, depth ordering) fail in three recognizable ways: \emph{attention dilution}, in which a prompt token's cross-attention spreads across the canvas and no object is well localized; \emph{attribute leakage}, in which color, material, or count attaches to the wrong object (the canonical ``red cube on a blue sphere'' confusion); and \emph{object dropping}, in which small or rare entities in a multi-object layout are simply omitted.
On GenEval the position and counting subsets remain the weakest axes across SD~1.5, SDXL, and SD~3~\cite{sd3rft2024}.
Remedies work at three levels: Attend-and-Excite~\cite{chefer2023attend} reweights cross-attention at inference so that every subject token is strongly activated somewhere in the canvas; StructureDiffusion~\cite{feng2023training} uses prompt syntax to constrain attention structure without training; and LLM-emitted layouts, LayoutGPT~\cite{feng2023layoutgpt} and the multimodal-LLM-driven RPG of~\citet{yang2024mastering}, plan a structured layout from the free-form prompt and hand it to an L2I renderer.
Whether the gap closes through a stronger compositional prior in the backbone or through planning as a separate stage is unresolved, and the two are not exclusive.

\textbf{Constraint satisfaction against creativity and diversity.}
Turning guidance up tightens adherence until, past a model-dependent threshold, edges oversharpen, textures flatten, and variation across seeds collapses; the same curve appears for the ControlNet conditioning weight, GLIGEN's gating strength, and the classifier-free guidance scale, leaving practitioners to tune two or three knobs at once.
Annealing the control weight along the denoising trajectory, per-region or per-token guidance scales, and down-weighting guidance where the base model is already confident are the current responses; how to expose such schedules to end users without overwhelming the interface is open.
Chaining spatial-aware generation with repeated local edits adds a drift problem of its own, treated in §\ref{subsec:p3_s8_1_3}.

\textbf{Geometry-aware realism in 2D.}
Depth and normal conditioning fix coarse occlusion but introduce their own artifacts: \emph{boundary banding} where a sharp depth discontinuity meets a smooth gradient in the generated image, \emph{scale ambiguity} because a monocular depth map carries no absolute scale and two objects at the same depth value can be rendered at incompatible physical sizes, and \emph{lighting incoherence} because the map fixes geometry but not the light field, so shadows disagree with the implied surface orientations.
Better estimators are the first lever, with the diffusion-prior Marigold~\cite{ke2024repurposing} producing sharper and more globally consistent depth than CNN baselines and Depth Anything~\cite{DepthAnything} staying robust on out-of-distribution conditioning images; differentiable geometry proxies that translate depth into shading constraints are the second, explored at the boundary with full 3D synthesis (§\ref{sec:p3_s8_2}) but not yet integrated into general T2I backbones.
Region grounding also changes the safety picture, since a harmful edit can be confined to a footprint that whole-image forensics miss; that concern belongs to editing and is taken up in §\ref{subsec:p3_s8_1_3}.


\subsection{2D Spatial Inpainting \& Editing}
\label{subsec:p3_s8_1_3}

\emph{2D spatial inpainting \& editing} modifies an image inside a specified region while leaving the rest intact: filling holes and occlusions, removing or inserting objects, changing attributes (color, material, texture, text), adjusting layout, or harmonizing pasted content with its new background.
The intent arrives as a mask, a structural hint (edges, sketches), a geometry proxy (depth, normals), or a natural-language instruction, and the tensions are between \textbf{local editability} and \textbf{global consistency} of lighting, perspective, and shadow, between \textbf{constraint adherence} and \textbf{realism}, and between \textbf{semantic plausibility} and \textbf{identity preservation} for person- or product-centric edits.
Editing holds the model to its strongest geometric account: the \emph{unedited} region of $I$ supplies ground-truth pixels that any plausible edit must respect, exposing failure modes that pure synthesis can hide, among them color drift, lighting reflow, identity loss, and attribute leakage across the mask boundary.


\subsubsection{Task Definition}
\label{subsubsec:p3_s8_1_3_task}
We inherit the latent-generator framework of §\ref{subsec:p3_s8_1_1}.
Let $(E,D)$ denote the autoencoder pair of Eq.~\eqref{eq:t2i_autoencoder}, $\phi:\mathcal{X}\!\rightarrow\!\mathcal{E}$ the text encoder, and $G_\theta$ the conditional latent sampler of Eq.~\eqref{eq:t2i_generator}.
Inpainting and editing depart from de-novo generation in one crucial respect: the input is no longer the empty set but an existing image $I\in\mathcal{Y}$, encoded into $\mathbf{z}^{\text{in}}=E(I)$.
A mask $\mathcal{M}\in\{0,1\}^{H\times W}$ designates the editable region, and an editing control set $\mathbf{c}\in\mathcal{C}$ (text prompts, structural cues, exemplar patches, or drag-style handles) populates the conditioning channel, encoded by $\kappa:\mathcal{C}\!\rightarrow\!\mathcal{C}'$.
The conditional latent sampler is extended to
\begin{equation}
G_\theta:\mathcal{Z}\times\{0,1\}^{H\times W}\times\mathcal{E}\times\mathcal{C}'\!\rightarrow\!\mathcal{Z},\qquad
\hat{\mathbf{z}}\sim p_\theta\!\big(\mathbf{z}\,\big|\,\mathbf{z}^{\text{in}},\,\mathcal{M},\,\phi(\mathbf{x}),\,\kappa(\mathbf{c})\big),
\label{eq:t2i_edit_objective}
\end{equation}
and the edited image $\hat{I}=D(\hat{\mathbf{z}})$ is required to (i) realize the requested change \emph{inside} $\mathcal{M}$, (ii) preserve $I$ \emph{outside} $\mathcal{M}$ up to imperceptible deviations, and (iii) remain globally coherent in color, lighting, geometry, and semantics.
Setting $\mathbf{z}^{\text{in}}=\varnothing$ and $\mathcal{M}=\mathbf{1}$ recovers the spatial-aware synthesis of §\ref{subsec:p3_s8_1_2} (Eq.~\eqref{eq:t2i_spatial_objective}); the distinctive substance of editing therefore lies in the interplay between the \emph{preservation} constraint outside $\mathcal{M}$ and the \emph{change} constraint inside it.
The mask may be hard, soft, inferred from a referring expression, or implicit (whole-image instruction editing with $\mathcal{M}=\mathbf{1}$).

\paragraph{Variants by Task Type.}
\label{para:p3_s8_1_3_variants}
The tasks differ in what populates $\mathbf{c}$ and $\mathcal{M}$ and in whether ground-truth pixels exist for evaluation.
Reconstruction-style completion carries no semantic intent: hole filling repairs scratches and occlusions from $(I,\mathcal{M})$ alone, and object removal reconstructs the background behind a specified instance.
Constraint-driven local synthesis adds intent inside $\mathcal{M}$: text-guided local editing follows an instruction scoped to a region (``make the \emph{left cup} red''), structure-guided inpainting takes edges, sketches, or semantic maps as a scaffold for shape and boundary, and geometry-guided inpainting conditions on depth or normals to respect occlusion ordering.
Object-level edits treat the region as an object whose identity or embedding in the scene is the unit of change: insertion with harmonization of color, grain, lighting, shadows, and reflections; attribute and appearance editing that changes material or style while preserving identity (garments, products, logos); and exemplar-guided transfer that drives content or style from a reference patch.
Interactive workflows replace declarative constraints with actions taken at edit time: point- or handle-based reshaping, and multi-round, multi-region sessions in which stability across rounds is itself the requirement.
The metrics appropriate to each are collected in Benchmarks below.


\subsubsection{Methodology}
\label{subsubsec:p3_s8_1_3_methodology}
Methods are organized by \emph{editing mechanism}, that is, by how $G_\theta$ of Eq.~\eqref{eq:t2i_edit_objective} is realized so that change is confined to $\mathcal{M}$ and preservation is enforced outside it.
The mechanism is largely independent of the task type above, since text-guided local editing can be served by masked denoising, by inversion and rewrite, or by an instruction-tuned editor, and one mechanism such as inversion serves both attribute edits and structure-guided ones.
Figure~\ref{fig:p3_s8_1_3_mechanisms} sets the five mechanisms against a single example edit.

\begin{figure}[!ht]
\centering
\includegraphics[width=\linewidth]{figures/part3/section8/Five_editing_mechanisms_for_2D_spatial_inpainting_and_editing.pdf}
\caption{Local 2D editing with context preservation (top) and five editing mechanisms: (a) masked denoising, where $M=1$ marks the editable region; (b) inversion and rewrite; (c) instruction-tuned editing; (d) exemplar / reference-guided editing; and (e) drag-based optimization. The cup images are illustrative, not outputs of the surveyed methods.}
\label{fig:p3_s8_1_3_mechanisms}
\end{figure}

\paragraph{Editing Mechanisms.}
\label{para:p3_s8_1_3_mechanisms}

\textbf{Masked latent denoising.}
The most direct realization of Eq.~\eqref{eq:t2i_edit_objective} runs the diffusion sampler only inside $\mathcal{M}$ while clamping the latent outside $\mathcal{M}$ to a noise-perturbed copy of $\mathbf{z}^{\text{in}}$ at every step.
RePaint~\cite{lugmayr2022repaint} does this with a \emph{pretrained, unconditional} DDPM and no retraining, replacing the known-region pixels with a forward-noised sample of the input at each reverse step and adding a \emph{resampling} schedule that re-injects noise periodically so the in-fill can re-harmonize with its context over several rounds.
Blended Diffusion~\cite{avrahami2022blended} made the same scheme text-driven by interleaving a CLIP-guided gradient step that steers the content inside $\mathcal{M}$ toward a target prompt with the context-clamping step outside it, the first widely cited text-guided local editor.
LaMa~\cite{suvorov2022lama} predates the diffusion line and is instead a feed-forward inpainter whose fast Fourier convolutions give every layer an image-wide receptive field, a Fourier-domain inductive bias that stays robust on the very large holes that defeat purely local CNN backbones.
Region-aware noise schedules and classifier-free guidance are standard add-ons, and ControlNet-inpaint and the inpainting variants of T2I-Adapter (§\ref{para:p3_s8_1_2_mechanisms}) compose with masked denoising by supplying the inside-$\mathcal{M}$ condition.

\textbf{Inversion + rewrite.}
A complementary strategy first recovers the noise trajectory that the model would have used to generate $I$ (\emph{inversion}) and then re-runs the reverse process under an edited prompt while reusing or swapping the attention maps of the original (\emph{rewrite}).
Prompt-to-Prompt~\cite{prompt-to-prompt} reuses or reweights source cross-attention maps during generation under the edited prompt, which preserves spatial structure while changing selected semantic content; for a real image an inversion step first supplies a suitable noise trajectory.
Imagic~\cite{kawar2023imagic} jointly optimizes the text embedding and lightly fine-tunes the model to reconstruct $I$, then interpolates the embedding toward the target prompt, preserving identity-level features (faces, garments) under non-rigid edits.
Null-Text Inversion~\cite{mokady2023nulltext} fixes the practical failure that plain DDIM inversion drifts away from $I$ under classifier-free guidance, by optimizing the \emph{null-text} embeddings used in the inverse pass so that exact reconstruction holds at ordinary guidance scales.
The family needs no full backbone fine-tune for Prompt-to-Prompt or Null-Text and dominates attribute, style, and structure-preserving real-image edits.

\textbf{Instruction-tuned editing.}
A third family avoids inversion entirely by retraining a diffusion editor on paired (before, instruction, after) triplets, mapping $I$ directly to $\hat{I}$ from a natural-language instruction with no explicit mask.
InstructPix2Pix~\cite{brooks2023instructpix2pixlearningfollowimage} built its corpus by pairing GPT-3-generated edit instructions with Prompt-to-Prompt edits of Stable Diffusion samples and fine-tuned Stable Diffusion on the triplets, so that at inference the user supplies only $(I,\text{instruction})$ and the editor localizes the change itself.
MagicBrush~\cite{zhang2024magicbrush} replaced the synthetic data with about 10K manually annotated real-image edits across multi-turn sessions, yielding both a stronger supervised editor and the standard instruction-editing test set.
Emu Edit~\cite{sheynin2024emu} trained a single model on a heterogeneous mixture of \emph{recognition} (segmentation, detection) and \emph{generation} (insertion, removal, style change) tasks with learned task embeddings that route each input, and leads its own benchmark on precision and instruction following.
Instruction editors are the most user-friendly mechanism, no mask and no inversion, at the cost of large supervised retraining.

\textbf{Exemplar / reference-guided editing.}
A visual \emph{exemplar}, an example image or patch encoding the desired content or style, replaces the textual instruction and enters through an auxiliary cross-attention pathway beside the standard text conditioning.
Paint-by-Example~\cite{yang2023paintbyexample} fine-tunes Stable Diffusion to fill $\mathcal{M}$ with the content of an exemplar whose identity is encoded by a CLIP image encoder, with strong exemplar augmentation (cropping, color jitter) at training time to prevent copy-paste behavior and force semantic rather than pixel-level transfer.
IP-Adapter~\cite{ye2023ipadapter} generalizes this into a \emph{plug-in} reference channel: a lightweight decoupled cross-attention head consumes a CLIP image embedding of the exemplar and is inserted into a frozen SD or SDXL backbone, so that any number of reference images can be combined with text prompts and with structural adapters such as ControlNet at inference time.
It is the standard route for subject personalization and style transfer when the desired content is hard to describe in words.

\textbf{Drag-based latent optimization.}
The most interactive family takes user-supplied (source point, target point) pairs and iteratively optimizes the latent so that the features at each source point migrate to the corresponding target while a locality regularizer pins the latent outside the user-marked region.
DragGAN~\cite{pan2023draggan} introduced the recipe on a StyleGAN backbone: a motion-supervision loss tracks the source feature and pulls it toward the target by gradient descent in the $\mathcal{W}^+$ latent space, and a point-tracking step re-localizes the moved feature at each iteration to prevent drift.
DragDiffusion~\cite{shi2023dragdiffusion} ported the recipe to the LDM latent space, running a per-timestep latent-update loop together with an identity-preserving LoRA fine-tune of the U-Net so that both the dragged subject and the unmarked context survive.
Both are inference-time optimizers with latency proportional to the number of drag iterations, and no other family yet offers the same directness of spatial control.

Read across the five families, the mechanisms track the generative substrate: feed-forward inpainters such as LaMa gave way to inversion-based diffusion editors (Prompt-to-Prompt, Imagic), which in turn yielded to supervised instruction-tuned editors (InstructPix2Pix, Emu Edit) as paired data became available at scale.
Orthogonal to that migration, the structural conditioning interfaces of §\ref{subsec:p3_s8_1_2} compose with every mechanism here, so a production pipeline pairs a mechanism with a structural channel rather than choosing one: Blended Diffusion with ControlNet-depth for geometry-aware text-guided edits, or IP-Adapter with DragDiffusion for exemplar-driven drag edits.


\subsubsection{Benchmarks}
\label{subsubsec:p3_s8_1_3_benchmarks}
Editing is evaluated in two regimes: \emph{reference-based}, where a held-out ground-truth region exists (holes synthesized over real images), and \emph{reference-free}, where only constraint satisfaction and human preference are available (instruction editing).
No single suite covers both, so the field keeps a small library of task-specific benchmarks next to the holistic T2I suites of §\ref{subsec:p3_s8_1_1}.

\paragraph{Training Corpora.}
Inpainting draws on Places2~\cite{zhou2017places} and Paris StreetView for scene-level masks, DIV2K for high-resolution patches, and the COCO / ADE20K mask supervision of §\ref{subsec:p3_s8_1_2}~\cite{lin2014microsoft,zhou2017ade20k}.
Constraint-driven local synthesis reuses the structural channels of the previous subsection, RefCOCO / RefCOCO+ / RefCOCOg~\cite{yu2016modeling} and PhraseCut~\cite{wu2020phrasecut} for referring-expression masks, SA-1B masks~\cite{kirillov2023sam}, and COCO edge and depth predictions, with the convention that the mask is supplied or derived from the expression; geometry-guided editing adds depth and normal supervision from NYUv2~\cite{silberman2012nyudv2}, Replica~\cite{straub2019replica}, and synthetic corpora, and insertion with harmonization uses composite datasets such as iHarmony4.
Instruction editing has its own data ecology: the synthetic InstructPix2Pix corpus of Prompt-to-Prompt-generated triplets, the manually annotated MagicBrush triplets~\cite{zhang2024magicbrush}, and the multi-task Emu Edit training mixture~\cite{sheynin2024emu}.

\paragraph{Evaluation Benchmark Suites.}
The holistic suites of §\ref{subsec:p3_s8_1_1} (Table~\ref{tab:p3_s8_1_1_eval_benchmarks}) carry over in part: T2I-CompBench~\cite{huang2023t2icompbench} and GenAI-Bench~\cite{li2024genaibench} score the attribute and compositional fidelity that edited images must also satisfy, and HEIM~\cite{lee2023heim} includes alignment and quality axes that apply to edits.
The dedicated suites are EditBench~\cite{wang2023editbench}, released with Imagen Editor, a curated set of text-guided inpainting prompts over object, attribute, and scene categories that serves as the standard evaluation for text-driven mask inpainting; the MagicBrush test set~\cite{zhang2024magicbrush}, $1{,}053$ manually annotated single- and multi-turn triplets, which is the standard for instruction-tuned editors; and the Emu Edit benchmark~\cite{sheynin2024emu}, which extends instruction editing to per-task subsets (insertion, removal, style change, color change, global change).
Drag-based methods evaluate on small custom test sets curated by~\citet{pan2023draggan} and~\citet{shi2023dragdiffusion}; a community-wide drag-editing benchmark has not yet emerged.

\paragraph{Constraint-specific Metrics.}
Reconstruction quality (PSNR / SSIM / LPIPS against ground truth, with FID / KID for diverse-completion settings) is defined only where held-out content exists.
Region adherence uses mask-restricted CLIPScore or region-level VQAScore to ask whether the edit honors the prompt \emph{inside} $\mathcal{M}$; edit locality measures $\ell_1$ or LPIPS \emph{outside} $\mathcal{M}$ against $I$, lower being better; harmonization is scored by color-statistic divergence, shadow-direction consistency, and the rate at which an off-the-shelf forgery detector flags the composite; identity preservation uses feature-cosine similarity over face or landmark regions for people and CLIP-image similarity for products and characters; and instruction following is paired accuracy, whether $\hat{I}$ honors the instruction better than $I$, judged by human raters or by a VLM over a structured rubric, as is now standard on MagicBrush and Emu Edit.
Change inside $\mathcal{M}$ and preservation outside it are reported together; a method that wins on one and loses on the other has not solved the task.


\subsubsection{Challenges and Trends}
\label{subsubsec:p3_s8_1_3_challenges}
Several problems are specific to editing an existing image.

\textbf{Structure-faithful hallucination in large holes.}
Large irregular holes and rare structures remain the hardest case: inpainters tile a salient local texture across the hole instead of synthesizing appropriate content, straight lines, perspective cues, and silhouettes break at the mask seam, and in-fills that are locally plausible contradict the global scene (a patch of grass inside a room).
Quality degrades sharply with hole size.
LaMa's image-wide Fourier convolutions and RePaint's resampling schedule (§\ref{para:p3_s8_1_3_mechanisms}) are the established responses; retrieval-augmented priors, which search a paired corpus for candidate completions and use them as soft conditioning, are an emerging one at the intersection of editing and retrieval.

\textbf{Edit locality against drift across rounds.}
A single well-calibrated edit does not guarantee a stable session.
Across rounds the unedited background shifts in white balance or saturation (\emph{color drift}), shadows and highlights migrate as if the edit had changed the global illumination (\emph{lighting reflow}), and salient textures from the edited region reappear in nearby unedited regions through attention leakage (\emph{texture replication}); the errors compound even when each round individually meets its locality budget.
Hard attention masking that confines non-edit tokens to $\mathcal{M}$, mask-aware noise schedules that anneal noise inside $\mathcal{M}$ only, and edit-history-aware sampling that conditions on the trajectory of prior edits are the active mitigations.
Inversion-based editors (Null-Text Inversion~\cite{mokady2023nulltext}, Imagic~\cite{kawar2023imagic}) carry a variant of the problem: any residual reconstruction error in the inverted trajectory leaks into the edit and accumulates.

\textbf{Harmonization under complex illumination.}
Inserted content betrays itself when the host scene has structured light: its shadow disagrees with the dominant light direction recoverable from the scene, there is no interreflection between the object and adjacent surfaces, and its albedo is tone-mapped inconsistently with the host, producing the cut-and-paste look that off-the-shelf forgery detectors flag.
Differentiable rendering proxies that turn depth and normals into shading constraints and back-propagate harmonization losses into the editor, illumination maps as an additional conditioning channel, and explicit shadow-synthesis modules are being explored; integration with the 3D-aware methods of §\ref{sec:p3_s8_2} is the natural next step and remains nascent.

\textbf{Language grounding and safety.}
Region-level instructions must be grounded reliably: an edit aimed at one object spills onto neighbors with similar attributes, spatial referring expressions are misparsed (``the left cup'' lands on the right one), and mildly perturbed instructions produce large deviations in the edit trajectory.
Borrowing GLIGEN's gated grounding~\cite{li2023gligen} (§\ref{para:p3_s8_1_2_mechanisms}) to confine edit attention, and using LLM-emitted layouts to harden the grounding before the edit is applied, are the current responses.
The same precision cuts the other way for safety.
A harmful edit (an identity swap inside one bounding box, watermark removal inside one mask) can be confined to a footprint that whole-image forensic detectors miss; large generative models reproduce individual training images under targeted prompting~\cite{somepalli2023diffusion}, which becomes acute when the training set contains copyrighted artwork or identifiable faces and motivates artist-side defenses such as Glaze~\cite{shan2023glaze}, which perturbs source artworks so that downstream fine-tunes mimic a different style; and most pipelines still lack robust per-image watermarking, so a photograph with a small synthesized region cannot be told from a wholly synthetic one.
Watermarking, region-aware safety filters, and licensing frameworks for the underlying corpora are all needed, and none is settled.

\textbf{Evaluation under plausibility ambiguity.}
Ground-truth pixels are often not the only plausible completion: a removed pedestrian leaves a hole that could contain road, curb, or shadow, and any reconstruction metric that compares against one ground truth over-penalizes good alternatives, while PSNR and SSIM rate sharp but contextually wrong in-fills highly.
The field is converging on \emph{paired evaluation}, asking whether $\hat{I}$ honors the instruction better than $I$ under a held-out judge, now standard on MagicBrush~\cite{zhang2024magicbrush} and Emu Edit~\cite{sheynin2024emu}; the judge-prompt sensitivity and benchmark-contamination caveats of §\ref{subsec:p3_s8_1_1} apply here as well.

\textbf{Scalability and interactive latency.}
Inversion-based methods store cross-attention activations at every timestep, multi-step sampling at $4$K and above scales with image area for standard U-Net backbones, and drag editors run an inner optimization loop per operation, so interactive editing lags behind the user's hand by hundreds of milliseconds.
Consistency-style distillation~\cite{song2023consistency,lcm2023} (§\ref{subsec:p3_s8_1_1}) compresses the reverse trajectory to a handful of steps, memory-efficient adaptive patching runs the denoiser only in a window around $\mathcal{M}$, and progressive-refinement schedules return a coarse preview within a perceptual-latency budget and refine in the background.
Closing the gap to the responsiveness of a photo editor is increasingly a systems problem shared with production T2I inference stacks.


% \subsection{Summary}
% \label{subsec:p3_s8_1_summary}

% Table~\ref{tab:p3_s8_1_subfield_comparison} places the three subfields side by side: the conditioning input each consumes, the sampler it instantiates, the mechanisms it surveys, and the benchmarks on which it reports.

% \begin{table}[t]
%     \centering
%     \fontsize{8}{10}\selectfont
%     \setlength{\tabcolsep}{1.0mm}
%     \caption{The three subfields of 2D spatial generation; each sampler is a successive specialization of Eq.~\eqref{eq:t2i_generator}.}
%     \label{tab:p3_s8_1_subfield_comparison}
%     \begin{tabular}{L{0.14\textwidth} L{0.19\textwidth} L{0.19\textwidth} L{0.19\textwidth} L{0.21\textwidth}}
%         \toprule
%         \textbf{Subfield} & \textbf{Input / sampler} & \textbf{Mechanisms} & \textbf{Canonical methods} & \textbf{Benchmarks / metrics} \\
%         \midrule
%         §\ref{subsec:p3_s8_1_1} Text-to-image generation & Text $\mathbf{x}$; $G_\theta:\mathcal{E}\!\to\!\mathcal{Z}$ (Eq.~\eqref{eq:t2i_generator}) & GAN; autoregressive; diffusion / flow matching; backbone, acceleration, and proprietary refinements & StackGAN, AttnGAN; DALL-E, Parti; LDM, Imagen, DiT, SD3 & FID, CLIPScore, HPSv2 / PickScore, VQAScore; GenEval, T2I-CompBench, HRS-Bench, HEIM, WISE \\
%         §\ref{subsec:p3_s8_1_2} Spatial-aware image generation & Text + structure $\mathbf{c}$; $G_\theta:\mathcal{E}\!\times\!\mathcal{C}'\!\to\!\mathcal{Z}$ (Eq.~\eqref{eq:t2i_spatial_objective}) & Side branch; adapter; grounded cross-attention; layout-aware cross-attention; inference-time guidance & ControlNet, T2I-Adapter, GLIGEN, LayoutDiffusion, BoxDiff & Channel adherence (IoU, edge recall, depth RMSE, PCK, relation accuracy); GenEval position; T2I-CompBench spatial \\
%         §\ref{subsec:p3_s8_1_3} Spatial inpainting and editing & Image $I$, mask $\mathcal{M}$, prompt, control; $G_\theta:\mathcal{Z}\!\times\!\{0,1\}^{H\times W}\!\times\!\mathcal{E}\!\times\!\mathcal{C}'\!\to\!\mathcal{Z}$ (Eq.~\eqref{eq:t2i_edit_objective}) & Masked denoising; inversion + rewrite; instruction-tuned; exemplar-guided; drag-based & RePaint, Blended Diffusion; Prompt-to-Prompt, Null-Text; InstructPix2Pix, Emu Edit; IP-Adapter; DragGAN & PSNR / LPIPS, edit locality, harmonization, identity preservation; EditBench, MagicBrush, Emu Edit \\
%         \bottomrule
%     \end{tabular}
% \end{table}

% One move structures the whole section: specialize the conditioning channel of $G_\theta$, from text alone (§\ref{subsec:p3_s8_1_1}) to text plus a structural scaffold (§\ref{subsec:p3_s8_1_2}) to scaffold plus a source image and mask (§\ref{subsec:p3_s8_1_3}).
% Each step is conservative, since the unconditional limit ($\mathbf{c}=\varnothing$ for §\ref{subsec:p3_s8_1_2}; $\mathbf{z}^{\text{in}}=\varnothing,\mathcal{M}=\mathbf{1}$ for §\ref{subsec:p3_s8_1_3}) recovers the preceding subfield, so the three sampler equations form a containment chain rather than a horizontal partition.
% The mechanism axes of the second and third subsections are, moreover, largely orthogonal to the generative core: side branches, adapters, grounded cross-attention, masked denoising, inversion, and drag-based optimization compose with LDM, DiT, and flow-matching backbones interchangeably, and production pipelines routinely stack a spatial-aware injection (Eq.~\eqref{eq:t2i_spatial_objective}) on top of an editing mechanism (Eq.~\eqref{eq:t2i_edit_objective}), so the three subfields are best read as one composable toolbox rather than as three disjoint task families.
% Evaluation has followed the same specialization, from a single FID against COCO for §\ref{subsec:p3_s8_1_1} circa~2020 to a task-specific tuple that adds channel adherence in §\ref{subsec:p3_s8_1_2} and reconstruction quality, edit locality, harmonization, and identity preservation in §\ref{subsec:p3_s8_1_3}.

% The extension of these 2D priors to \emph{volumetric} synthesis (turning a text prompt, an image, or a mask into a textured mesh, a NeRF, or a 3DGS scene) is the subject of §\ref{sec:p3_s8_2}. The conditional latent sampler of Eq.~\eqref{eq:t2i_generator} continues to play a central role there, but now as a \emph{prior to be distilled} via score-distillation sampling, multi-view consistency objectives, or volume-rendering supervision, rather than as the terminal generator.

% \newpage

\section{3D Spatial Reconstruction and Generation}
\label{sec:p3_s8_2}

\FramedInitial{T}{he} 3D spatial reconstruction and generation encompasses the tasks that turn images, video, sensor readings, or high-level prompts into an explicit 3D representation (a textured mesh, a signed-distance or occupancy field, a neural radiance field, or a 3D Gaussian field) of an object, a scene, or a complete navigable environment.
Whereas §\ref{sec:p3_s8_1} taught a 2D semantic-to-pixel before honor spatial structure on the image plane, this section lifts that prior off the plane and into the volumetric substrate on which embodied perception, simulation, and interaction ultimately operate.
This is where appearance and geometry are forced into a single, view-consistent representation: a model must not only produce plausible pixels but place matter correctly in $\mathbb{R}^3$ so that the result can be rendered from arbitrary viewpoints, physically simulated, navigated, and edited. It is the bridge between the flat generative priors of §\ref{sec:p3_s8_1} and the embodied world models, simulators, and digital twins that consume 3D structure downstream.

The five subsections are organized by the direction and source of the map into a shared 3D-representation space $\mathcal{R}$.
§\ref{subsec:p3_s8_2_1} starts with a single asset: 3D object synthesis turns the 2D prior of §\ref{sec:p3_s8_1} into a volumetric generator through 3D-native training, score distillation, multi-view lifting, or feed-forward reconstruction.
§\ref{subsec:p3_s8_2_2} scales that generator to a structured scene of placed objects and a room envelope.
§\ref{subsec:p3_s8_2_3} reverses the direction: reconstruction recovers a representation faithful to sensor observations, and meets synthesis exactly where the observations run out.
§\ref{subsec:p3_s8_2_4} routes through a 360-degree panorama as a structured 2D intermediate that is then lifted to 3D, and §\ref{subsec:p3_s8_2_5} closes the loop by editing an existing representation under user control.
% A Summary (§\ref{subsec:p3_s8_2_summary}) compares the five along shared axes.


\subsection{3D Object Synthesis}
\label{subsec:p3_s8_2_1}


3D object synthesis generates a complete, usable 3D asset (geometry and, optionally, appearance) from \textit{text}, \textit{a single image}, or \textit{a sparse set of views}.
Where reconstruction recovers an object consistent with given observations, synthesis additionally asks for plausible completion of what was never observed, for generative diversity, and for usability in downstream graphics and embodied pipelines (\eg, game engines, simulators, manipulation); the output is a textured mesh, a radiance field, or a set of 3D Gaussians depending on which of fidelity, editability, rendering speed, and toolchain compatibility matters most.
It is also the point at which the \emph{semantic-to-pixel} prior of §\ref{sec:p3_s8_1} (the conditional latent sampler $G_\theta$ of Eq.~\eqref{eq:t2i_generator}) becomes a \emph{semantic-to-3D} prior $G_\theta^{\text{3D}}$: 3D-native generators train on 3D corpora directly, score-distillation pipelines invert the frozen 2D prior under rendering supervision, multi-view pipelines fine-tune it to emit consistent views and lift them, and feed-forward reconstructors and 3D-native latent diffusion exploit the recent object-scale 3D corpora to learn one-shot text- or image-to-3D maps.
Figure~\ref{fig:3D_object_generation} illustrates the text- or single-image-conditioned setting.

\begin{figure}[!ht]
    \centering
    \includegraphics[width=0.89\linewidth]{figures/part3/section8/3D-objects-gen-1.png}
    \caption{Illustration of 3D object generation conditioned on text or a single image (adapted from~\cite{xiang2025trellis}).}
    \label{fig:3D_object_generation}
\end{figure}


\subsubsection{Task Definition}
\label{subsubsec:p3_s8_2_1_task}

We cast 3D object synthesis as a paradigm-agnostic extension of the latent-generator framework of §\ref{subsec:p3_s8_1_1}.
Let $\mathcal{R}$ denote a 3D-representation space taken from $\{\,\textit{mesh},\,\textit{SDF},\,\textit{voxel},\,\textit{triplane},\,\textit{NeRF},\,\text{\textit{3DGS}},\,\textit{structured latent}\,\}$, and let $\mathcal{C}$ denote the conditioning-input space comprising a text prompt $\mathbf{x}\!\in\!\mathcal{X}$ embedded via $\phi$ (as in Eq.~\eqref{eq:t2i_generator}), a single image $I\!\in\!\mathcal{Y}$, or a multi-view set $\{(I_v,\pi_v)\}_{v=1}^V$ pairing images with camera poses $\pi_v$.
A condition encoder $\kappa:\mathcal{C}\!\to\!\mathcal{C}'$ maps each input into a shared conditioning embedding.
The core learning problem is to train a conditional 3D generator
\begin{equation}
G_\theta^{\text{3D}}:\mathcal{C}'\rightarrow \mathcal{R},\qquad \hat{\mathbf{r}}\sim p_\theta^{\text{3D}}\!\big(\mathbf{r}\,\big|\,\kappa(\mathbf{c})\big),
\label{eq:p3_s8_2_1_objective}
\end{equation}
so that the sampled asset $\hat{\mathbf{r}}\in\mathcal{R}$ admits high-fidelity geometry, plausibly completes unobserved regions, and renders consistently across viewpoints.
Two limiting cases locate Eq.~\eqref{eq:p3_s8_2_1_objective} in the literature.
When $\kappa(\mathbf{c})=\varnothing$ the sampler degenerates to \emph{unconditional} 3D generation, the regime of early GAN- and diffusion-based 3D generators trained on category-specific corpora.
At the other extreme, when $\mathbf{c}$ is text-only and $G_\theta^{\text{3D}}$ is realized by distilling the 2D sampler $G_\theta$ of Eq.~\eqref{eq:t2i_generator} (via score-distillation sampling, multi-view consistency objectives, or volume-rendering supervision), one recovers the DreamFusion-style text-to-3D pipeline that has dominated open-vocabulary 3D synthesis since 2022.


\subsubsection{Methodology}
\label{subsubsec:p3_s8_2_1_methodology}

The four paradigms below are four ways of parameterizing and optimizing the sampler $G_\theta^{\text{3D}}$ of Eq.~\eqref{eq:p3_s8_2_1_objective}, taken in the order in which each became prominent.
Figure~\ref{fig:p3_s8_2_1_paradigms} contrasts them, making explicit which components are trained versus frozen and where the 3D representation $\mathcal{R}$ is instantiated.

\begin{figure}[!ht]
    \centering
    \includegraphics[width=0.99\linewidth]{figures/part3/section8/3D_object_synthesis.pdf}
    \caption{Four paradigm families for 3D object synthesis: (P1) 3D-native generation; (P2) distillation from a frozen 2D diffusion prior (SDS); (P3) multi-view generation with a fine-tuned prior, then lifting to 3D; and (P4) learned direct-to-3D routes, by one-pass reconstruction or iterative 3D latent sampling. The families can overlap, notably P1 and the latent route of P4; chair images are illustrative.}
    \label{fig:p3_s8_2_1_paradigms}
\end{figure}

\paragraph{3D-Native Generators.}
\label{para:p3_s8_2_1_native}
The earliest 3D generators instantiate $G_\theta^{\text{3D}}$ as a sampler trained \emph{directly} on a 3D representation, by denoising score matching on voxels, SDFs, or triplanes or, in earlier work, by an adversarial fit to a category-specific 3D corpus $p_{\text{data}}^{\text{3D}}(\mathbf{r})$,
\begin{equation}
\mathcal{L}_{\text{native}}(\theta)
=
\mathbb{E}_{\mathbf{r}_0\sim p_{\text{data}}^{\text{3D}},\,t,\,\epsilon}\Big[\big\|\epsilon - \epsilon_\theta^{\text{3D}}\!\big(\mathbf{r}_t,\,t,\,\kappa(\mathbf{c})\big)\big\|^2\Big],
\quad
\mathbf{r}_t=\alpha_t \mathbf{r}_0 + \sigma_t\,\epsilon,
\label{eq:p3_s8_2_1_native_objective}
\end{equation}
where the noise predictor $\epsilon_\theta^{\text{3D}}$ operates directly on $\mathbf{r}\!\in\!\mathcal{R}$ with no 2D intermediary, so the supervision signal scales with the size of the 3D corpus rather than with 2D image--text pairs.
Because 3D data is heterogeneous, these generators are themselves diverse, having been instantiated on signed distance fields~\cite{gao2022get3d}, voxel grids~\cite{wu2016learning, yang20173d, zhu2018visual}, point clouds~\cite{lin2018learning, huang2020pf}, view sets~\cite{park2017transformation, nguyen2019hologan}, neural radiance fields~\cite{chan2022efficient, chan2021pi, deng2022gram, niemeyer2021giraffe, schwarz2020graf}, triplanes~\cite{chen2023singlestage, shue20233dtriplanetriplane, wang2023rodin}, and 3D Gaussian Splatting~\cite{wewer2024latentsplat, roessle2024l3dg, zou2024triplane}.
The scarcity and monotony of 3D corpora remain the bottleneck: even with sophisticated architectures, generated geometry and texture lack fidelity and stay confined to a handful of canonical categories, which is what pushed the field, around 2022--2023, toward paradigms that use the much larger \emph{2D} corpora as an indirect supervision signal for 3D structure.

\paragraph{2D-Prior Distillation (SDS Family).}
\label{para:p3_s8_2_1_sds}
The second paradigm sidesteps the 3D-data bottleneck by realizing $G_\theta^{\text{3D}}$ as a \emph{per-shape optimization} that distills gradients from a pretrained 2D generator (a CLIP model~\cite{radford2021learning} or a 2D diffusion model~\cite{rombach2022ldm}) trained on web-scale image data~\cite{schuhmann2022laion5b}.
Given the \emph{frozen} 2D noise predictor $\epsilon_\phi$ of Eq.~\eqref{eq:t2i_generator} and a differentiable renderer $\mathrm{Render}(\mathbf{r}_\theta,\pi)$ that images a 3D representation $\mathbf{r}_\theta\!\in\!\mathcal{R}$ under camera pose $\pi$, score distillation updates only the \emph{per-shape} 3D parameters $\theta$ via the gradient
\begin{equation}
\nabla_\theta \mathcal{L}_{\text{SDS}}
=
\mathbb{E}_{t,\,\epsilon,\,\pi}\Big[\,w(t)\,\big(\epsilon_\phi(\mathbf{y}_t^\pi,\,t,\,\phi(\mathbf{x})) - \epsilon\big)\,\tfrac{\partial \mathbf{y}^\pi}{\partial \theta}\,\Big],
\quad
\mathbf{y}^\pi = \mathrm{Render}(\mathbf{r}_\theta,\pi),
\label{eq:p3_s8_2_1_sds_objective}
\end{equation}
where $w(t)$ is a step-dependent weighting and $\mathbf{y}_t^\pi$ is the rendered view perturbed by the diffusion forward process.
The contrast with Eq.~\eqref{eq:p3_s8_2_1_native_objective} is sharp: the 2D prior of §\ref{sec:p3_s8_1} is held \emph{fixed} and no 3D corpus is consumed, while the trainable parameters live in the representation of a \emph{single} asset, hence the per-shape optimization cost paid at inference.
A line of work~\cite{jain2022zerodreamfields, poole2022dreamfusion, lin2023magic3d, wang2023score} applies score-distillation sampling (SDS) or its variants to a NeRF or mesh; DreamFusion~\cite{poole2022dreamfusion} obtains zero-shot text-to-3D by optimizing a NeRF whose multi-view renderings the 2D diffusion prior scores, and GaussianDreamer~\cite{yi2024gaussiandreamer} moves the same recipe onto a 3DGS representation.
These methods need no 3D training data, but because the optimization is rerun for every prompt they are slow, prone to over-saturated color, and subject to the ``multi-face'' (Janus) artifact in which the 2D prior fails to disambiguate the front and back of an object; variational score distillation~\cite{wang2023prolificdreamer} mitigates the over-saturation but not the per-prompt cost.

\paragraph{Multi-View Generation and Lifting.}
\label{para:p3_s8_2_1_mv}
The third paradigm realizes $G_\theta^{\text{3D}}$ as a two-stage pipeline: a fine-tuned 2D diffusion model generates a small set of geometrically consistent views, and a separate reconstructor lifts them to a 3D asset, so that Eq.~\eqref{eq:p3_s8_2_1_objective} factorizes as
\begin{equation}
\hat{\mathbf{r}}
=
\mathcal{F}_\psi\Big(\big\{\hat{\mathbf{y}}_v,\,\pi_v\big\}_{v=1}^V\Big),
\qquad
\hat{\mathbf{y}}_v \sim p_\theta\big(\mathbf{y}_v \,\big|\, \mathbf{y}_{\text{ref}},\,\pi_v,\,\phi(\mathbf{x})\big),
\label{eq:p3_s8_2_1_mv_objective}
\end{equation}
where the view-conditional 2D sampler $p_\theta(\cdot\mid\cdot)$ is a fine-tuned variant of Eq.~\eqref{eq:t2i_generator} that additionally consumes a reference image or text $\mathbf{y}_{\text{ref}}$ and a target camera pose $\pi_v$, and $\mathcal{F}_\psi$ is a (typically feed-forward) multi-view-to-3D reconstructor.
Two differences from Eq.~\eqref{eq:p3_s8_2_1_sds_objective} follow: the 2D prior is \emph{fine-tuned} rather than frozen, and the 3D representation appears only inside the deterministic lift $\mathcal{F}_\psi$, with no per-shape inner loop at inference, which is what lets this paradigm run one to two orders of magnitude faster than SDS at the price of depending on the view-consistency of the synthesized $\hat{\mathbf{y}}_v$.
The first step was to make a pretrained diffusion model pose-aware. Following fine-tuning recipes such as ControlNet~\cite{zhang2023adding} and LoRA~\cite{hu2022lora}, several works~\cite{raj2023dreambooth3d, seo2023let3DFuse} fine-tune for cross-view consistency; Zero-1-to-3~\cite{liu2023zero} opened the open-world single-image setting by training a diffusion model to synthesize a novel view conditioned on an input view and a relative camera pose; MVDream~\cite{shi2023mvdream} and ImageDream~\cite{wang2023imagedream} fix a world-camera coordinate convention from which geometry can be recovered; Magic123~\cite{qian2023magic123} combines the Zero-1-to-3 prior with the 2D prior of Stable Diffusion; One-2-3-45~\cite{liu2023one2345} lifts Zero-1-to-3 views to a mesh; and Consistent123~\cite{lin2024consistent123} and MVD-Fusion~\cite{hu2024mvdfusion} add boundary and depth priors.
Attention then turned to consistency \emph{across} the generated views: SyncDreamer~\cite{liu2023syncdreamer} correlates features across views with a 3D-aware attention, CAT3D and Free3D~\cite{gao2024cat3d, zheng2024free3d} model the multi-view distribution jointly, and Zero123++~\cite{shi2023zero123++} tiles six surrounding views into a single grid image, a trick inherited by One-2-3-45++~\cite{liu2024one2345++}, Instant3D~\cite{li2023instant3d}, InstantMesh~\cite{xu2024instantmesh}, and AR-1-to-3~\cite{zhang2025ar123}.

\paragraph{Feed-Forward Reconstruction and 3D-Native Latent Diffusion.}
\label{para:p3_s8_2_1_lrm}
The fourth group amortizes generation across object-scale 3D corpora rather than optimizing each asset against a frozen 2D prior, in two instantiations.
A \emph{feed-forward reconstructor} trains $G_\theta^{\text{3D}}$ end-to-end to map the conditioning embedding $\kappa(\mathbf{c})$ directly to a 3D representation $\hat{\mathbf{r}}\!\in\!\mathcal{R}$, supervised by differentiable rendering against $V$ ground-truth views $\{\mathbf{y}_v,\pi_v\}$ drawn from the corpus,
\begin{equation}
\mathcal{L}_{\text{LRM}}(\theta)
=
\mathbb{E}_{(\mathbf{c},\,\{\mathbf{y}_v,\pi_v\}_v)\sim p_{\text{data}}^{\text{3D}}}\,
\sum_{v=1}^V \big\|\mathrm{Render}\big(G_\theta^{\text{3D}}(\kappa(\mathbf{c})),\,\pi_v\big) - \mathbf{y}_v\big\|^2.
\label{eq:p3_s8_2_1_lrm_objective}
\end{equation}
A \emph{3D-native latent diffusion} model instead learns a 3D autoencoder pair $(E^{\text{3D}}, D^{\text{3D}})$ with $\mathbf{z}^{\text{3D}}\!=\!E^{\text{3D}}(\mathbf{r})$, trains a denoising model on the compact latent, and decodes via $\hat{\mathbf{r}}=D^{\text{3D}}(\hat{\mathbf{z}}^{\text{3D}})$,
\begin{equation}
\mathcal{L}_{\text{3DL}}(\theta)
=
\mathbb{E}_{\mathbf{z}^{\text{3D}}_0,\,t,\,\epsilon}\Big[\big\|\epsilon - \epsilon_\theta\big(\mathbf{z}^{\text{3D}}_t,\,t,\,\kappa(\mathbf{c})\big)\big\|^2\Big],
\label{eq:p3_s8_2_1_latdiff_objective}
\end{equation}
which is the 2D LDM objective Eq.~\eqref{eq:t2i_ldm_objective} of §\ref{subsec:p3_s8_1_1} with the autoencoder moved from pixels to $\mathcal{R}$.
Both variants avoid the per-shape optimization of Eq.~\eqref{eq:p3_s8_2_1_sds_objective}, but only the reconstructor uses a single prediction pass; latent diffusion still samples iteratively.
The unlocking factor was the arrival of Objaverse and Objaverse-XL~\cite{deitke2023objaverse, deitke2023objaversexl}.
LRM~\cite{hong2023lrm} and its follow-ups~\cite{tang2024lgm, wang2024crm, wang2025moving} show that large transformers trained purely with rendering losses capture object geometry and texture in detail and generalize to unseen images, and a parallel line introduces compact \emph{3D-native latent codes}~\cite{chen20253dtopia, zhang20233dshape2vecset, ye2022structured, yang2025prometheus} whose distribution a latent diffusion model can learn; TRELLIS~\cite{xiang2025trellis}, CLAY~\cite{zhang2024clay}, HiGS-RGBN~\cite{shen2025highGSRGBN}, and Hunyuan3D~\cite{yang2024hunyuan3d} build on this to bring generated assets close to production grade while keeping the open-vocabulary reach that the SDS family pioneered.

Two design choices cut across the four paradigms.
The representation $\mathcal{R}$ (\textit{mesh}, \textit{SDF}, \textit{voxel}, \textit{triplane}, \textit{NeRF}, \textit{3DGS}, or \textit{structured latent}) is largely independent of the paradigm: SDS has been run on NeRFs and on 3DGS, feed-forward reconstructors output triplanes, meshes, or Gaussians, and latent-diffusion stacks decode to almost any explicit representation through a learned VAE.
So is the conditioning modality: the same backbone takes text, a single image, or sparse views by changing only the encoder $\kappa$ and the camera-pose embedding, and production systems chain them (image-to-3D conditioned on a text-to-image preview).
The succession of paradigms itself tracks the accumulation of supervision at each level, from 3D-native corpora to 2D image--text pairs to Zero-1-to-3-style multi-view priors to Objaverse-XL-scale 3D, with the most recent industrial systems sitting in the fourth paradigm while inheriting the multi-view-consistency tricks of the third.


\subsubsection{Benchmarks}
\label{subsubsec:p3_s8_2_1_benchmarks}

Evaluation draws on object-scale training corpora, on dedicated benchmark suites that curate prompts and human ratings, and on metrics that probe alignment, geometry, and appearance separately.

\begin{figure}[!ht]
    \centering
    \includegraphics[width=0.7\linewidth]{figures/part3/section8/objaverse.png}
    \caption{Demonstration of the Objaverse~\cite{deitke2023objaversexl} dataset.}
    \label{fig:objaverse}
\end{figure}

\paragraph{Training Corpora.}
Object-scale 3D corpora drive the feed-forward and latent-diffusion paradigms; Table~\ref{tab:p3_s8_2_1_corpora} lists the ones most commonly used, and Figure~\ref{fig:objaverse} shows representative Objaverse-XL assets.
They span the classical CAD repository ShapeNet~\cite{shapenet2015}, category-bounded real-world scans with accurate geometry and PBR textures (Google Scanned Objects~\cite{downs2022googleGSO}, OmniObject3D~\cite{wu2023omniobject3d} with dense multi-view imagery), the open-vocabulary web-scraped Objaverse and Objaverse-XL~\cite{deitke2023objaverse, deitke2023objaversexl} that unlocked the fourth paradigm, and G-Objaverse~\cite{zuo2024gobjavers}, which adds text-driven synthesized assets between the two.

\begin{table}[t]
    \centering
    \fontsize{8}{10}\selectfont
    \setlength{\tabcolsep}{1.0mm}
    \caption{Representative training corpora for 3D object synthesis. ``Scale'' gives order-of-magnitude object counts as reported in each release.}
    \label{tab:p3_s8_2_1_corpora}
    \begin{tabular}{L{0.22\textwidth} C{0.07\textwidth} C{0.14\textwidth} L{0.20\textwidth} L{0.29\textwidth}}
        \toprule
        \textbf{Dataset} & \textbf{Year} & \textbf{Scale (\#objects)} & \textbf{Acquisition} & \textbf{Primary annotation / modality} \\
        \midrule
        ShapeNet~\cite{shapenet2015} & 2015 & ${\sim}51\mathrm{K}$ & Curated CAD & Semantic category, part labels \\
        Google Scanned Objects~\cite{downs2022googleGSO} & 2022 & ${\sim}1\mathrm{K}$ & Real-world scans & High-fidelity geometry + PBR textures \\
        OmniObject3D~\cite{wu2023omniobject3d} & 2023 & ${\sim}6\mathrm{K}$ & Real-world scans & Dense multi-view + reconstructions \\
        Objaverse~\cite{deitke2023objaverse} & 2023 & ${\sim}800\mathrm{K}$ & Web-scraped & Open-vocabulary text tags \\
        Objaverse-XL~\cite{deitke2023objaversexl} & 2023 & ${\sim}10\mathrm{M}$ & Web-scraped at scale & Open-vocabulary text tags \\
        G-Objaverse~\cite{zuo2024gobjavers} & 2024 & ${\sim}280\mathrm{K}$ & Synthesized + curated & Text-driven assets, rendered views \\
        \bottomrule
    \end{tabular}
\end{table}

\paragraph{Evaluation Benchmark Suites.}
Recent benchmarks have moved to \emph{multi-dimensional} assessment beyond plain text--image alignment.
GT23D-Bench~\cite{cai2024gt23d} curates a high-fidelity dataset and a suite of 3D-aware metrics spanning alignment and visual quality; MATE-3D~\cite{zhang2025mate3d} provides designed prompt categories with large-scale human annotations across several perceptual dimensions and introduces HyperScore as a learned evaluator; T23D-QA~\cite{li2025t23d} adopts a question-answer protocol that scores geometry, texture, and alignment independently; and Gen3DEval~\cite{maiti2025gen3deval} fine-tunes a vision-language model for 3D quality assessment, scoring text fidelity, appearance, and surface quality from rendered surface normals.
Together with the 3D-relevant subsets of the 2D suites collected in §\ref{subsec:p3_s8_1_1} (Table~\ref{tab:p3_s8_1_1_eval_benchmarks}), they form the de facto evaluation protocol for 3D object synthesis.

\paragraph{Constraint-specific Metrics.}
Because an asset is judged from many viewpoints and under different downstream uses, evaluation reports complementary metrics along three axes.
\textbf{Text-3D alignment} renders multi-view images from the asset and measures text--image similarity with CLIP-like encoders, aggregated across views (mean, worst-case, or view-consistency variants); pure CLIP alignment can be misled by texture ``shortcuts'' and does not guarantee geometric correctness, which motivates the 3D-aware alignment scores of GT23D-Bench~\cite{cai2024gt23d} and learned evaluators such as HyperScore~\cite{zhang2025mate3d}.
\textbf{Geometry quality} is measured against ground-truth shapes, where they exist, by Chamfer Distance (CD)~\cite{huang2025spar3d}, F-score at varying thresholds, normal consistency, and volumetric IoU for watertight meshes; for text-to-3D without ground truth it falls back on multi-view depth and normal consistency, silhouette consistency, and learned 3D-quality predictors.
\textbf{Appearance and multi-view consistency} are measured by perceptual distances (LPIPS) between rendered views and references where available, by view-consistency measures (variance across views under fixed lighting and camera trajectories), and by human studies, which recent benchmarks collect as multi-dimensional ratings to calibrate the automatic metrics~\cite{zhang2025mate3d, li2025t23d}.

The paradigm succession of §\ref{subsubsec:p3_s8_2_1_methodology} shows up in two measurable quantities.
Generation time has fallen by about three orders of magnitude: DreamFusion~\cite{poole2022dreamfusion} needs roughly 1.5 hours per asset and Magic3D~\cite{lin2023magic3d} about 40 minutes, whereas LRM~\cite{hong2023lrm} reconstructs a shape in about 5 seconds and InstantMesh~\cite{xu2024instantmesh} completes image-to-3D in around 10 seconds.
Image-to-3D geometry is now routinely scored on Google Scanned Objects~\cite{downs2022googleGSO} with Chamfer Distance (lower is better) and F-score at $\tau{=}0.2$ (higher is better), where InstantMesh~\cite{xu2024instantmesh} reaches CD~$\approx 0.18$ and FS~$\approx 0.88$, ahead of TripoSR~\cite{tochilkin2024triposr} ($0.217$ / $0.843$), CRM~\cite{wang2024crm} ($0.252$ / $0.787$), and LGM~\cite{tang2024lgm} ($0.345$ / $0.671$).
Coarse text-to-3D alignment, by contrast, is saturated: CLIP-style scores no longer separate the strongest systems, so Hunyuan3D~2.0~\cite{hunyuan3d22025tencent} reports ULIP- and Uni3D-based image- and text-to-shape similarity backed by a study in which 50 participants rated 300 unselected results, and learned evaluators (HyperScore on MATE-3D~\cite{zhang2025mate3d}, T23D-QA~\cite{li2025t23d}) score geometry and texture on separate axes.
The reconstruction metrics carry their own caveat: CD and F-score saturate on coarse geometry and can reward visually plausible but physically thin assets, so reported gains are credible mainly where geometric, perceptual, and human ratings agree.
Commercial text- and image-to-3D services now ship these capabilities with PBR export to game engines, alongside research systems such as Meta 3D Gen~\cite{bensadoun2024meta3dgen} and the open-vocabulary 3D-native-latent systems CLAY~\cite{zhang2024clay}, TRELLIS~\cite{xiang2025trellis}, and Hunyuan3D~\cite{yang2024hunyuan3d, hunyuan3d22025tencent}.


\subsubsection{Challenges and Trends}
\label{subsubsec:p3_s8_2_1_challenges}

Four problems remain open, and they depend on one another: the multi-view consistency of the asset, the representation it lands in, whether progress on either can be measured reliably, and the data that ultimately drives all three.

\paragraph{Multi-View and 3D Consistency under Open-Vocabulary Prompts.}
Enforcing consistent geometry and texture across viewpoints while preserving fine detail remains difficult in all four paradigms, and the failures have names.
The \emph{Janus artifact} renders both the front and the back of an object as a frontal face, most acutely in SDS pipelines, where the 2D prior cannot disambiguate camera direction; \emph{view drift} makes adjacent views of a multi-view stack disagree on geometry, a residual concern in Zero-1-to-3-style cascades that SyncDreamer's cross-view attention and Zero123++'s tiled six-view trick only partly remove; and \emph{texture seams} appear where neighboring views' color and lighting cannot be merged into a globally coherent UV map, a persistent weakness of feed-forward reconstructors.
These failures motivate 3D-native backbones with \emph{structured latents}~\cite{xiang2025trellis} that share parameters across views by construction, and benchmarks that score multi-view consistency rather than a single rendered view~\cite{zhang2025mate3d}.

\paragraph{Representation Trade-offs between Fidelity, Editability, and Deployability.}
Even a consistent asset must \emph{land somewhere} a downstream tool can consume, and the choice of representation is not neutral.
Meshes with UV/PBR textures are what game engines and simulators want but lose detail to discretization; radiance fields and Gaussians give high visual fidelity and fast rendering but require conversion, which introduces its own failures: \emph{topology breakage} during marching-cubes-style meshing, \emph{baked illumination} that fuses lighting into albedo when converting from a NeRF or 3DGS, and \emph{texture bleeding} across UV seams.
The response is \emph{unified latent representations} that decode to several formats~\cite{xiang2025trellis}, so that a single generator serves diverse downstream requirements and the choice of $\mathcal{R}$ is deferred to inference.

\paragraph{Scalable, Reliable Evaluation and Human-Aligned Metrics.}
These representation choices complicate evaluation, because the metric panel a paper reports tends to depend on which format its generator targets, and the automatic metrics remain imperfect for geometry plausibility and asset usability.
CLIP-style alignment is fooled by \emph{texture shortcuts} that ignore the underlying shape; Chamfer Distance saturates on coarse geometry while penalizing semantically irrelevant fine detail; and LPIPS on rendered views cannot distinguish a physically plausible asset from a visually plausible cardboard cut-out.
Multi-dimensional evaluators trained on large human-annotation sets (\eg, HyperScore on MATE-3D~\cite{zhang2025mate3d}) and diagnostic benchmarks (\eg, GT23D-Bench~\cite{cai2024gt23d}, T23D-QA~\cite{li2025t23d}) that score alignment, geometry, and texture on independent axes are the current answer, so that progress on one axis is no longer masked by averaging.

\paragraph{Data Scaling and Asset Readiness.}
Underlying consistency, representation, and evaluation alike is the supply of training data and the industrial scale at which it can be exploited.
Large-scale corpora such as Objaverse-XL~\cite{deitke2023objaversexl} and industrial-scale systems such as Hunyuan3D~\cite{hunyuan3d22025tencent, hunyuan3d2025hunyuan3domni} and Meta 3D Gen~\cite{bensadoun2024meta3dgen} point toward \emph{high-resolution, engine-ready assets} with controllable materials, relightable textures, and post-processing for watertightness and topology.
The question has shifted from \emph{whether} a paradigm can be trained at scale, which the fourth paradigm has settled, to \emph{which slice} of the 3D-asset distribution the available corpora still under-represent: articulated objects, semi-transparent or refractive materials, and topologically non-trivial assemblies are visibly thinner in current open corpora than in production pipelines, and closing this gap is the explicit target of recent industrial releases.




\subsection{3D Scene Synthesis}
\label{subsec:p3_s8_2_2}


Scene-level synthesis extends 3D generation from a single asset $\mathbf{r}\!\in\!\mathcal{R}$ to an entire environment: a clutter of objects whose positions, orientations, and pairwise relations together compose a navigable room or open-world scene.
Where §\ref{subsec:p3_s8_2_1} models the conditional distribution of one object given a prompt, scene synthesis must jointly model (i) a variable number of objects, (ii) their global $\mathrm{SE}(3)$ placements, (iii) their pairwise spatial and functional relations, and (iv) the envelope geometry (floor, walls, ceiling) that closes the scene; the output is a structured set rather than a vector in $\mathcal{R}$, and the conditioning broadens from text to layouts, scene graphs, sparse calibrated views, and full camera trajectories.
The result must be navigable, not only renderable, which makes the problem harder than object synthesis along three axes at once: geometric coherence across the clutter, semantic consistency between related objects, and view generalization along an $\mathrm{SE}(3)$ trajectory rather than a fixed look-at orbit.
Four families answer it: \emph{layout- and scene-graph-conditioned} generators that produce the structured set from a symbolic prior; \emph{view-to-3D} synthesizers that regress or sample a continuous radiance or Gaussian field from sparse calibrated views; \emph{view-to-view} synthesizers that skip the explicit 3D representation and learn a conditional map from observed to novel views; and \emph{video-model-based} synthesizers that treat a pretrained video diffusion model as an implicit 3D prior with the camera trajectory as an additional conditioning channel.


\subsubsection{Task Definition}
\label{subsubsec:p3_s8_2_2_task}

We cast 3D scene synthesis as a scene-level extension of the object-synthesis framework of §\ref{subsec:p3_s8_2_1}.
Let $\mathcal{S}$ denote the scene-output space whose elements
$\mathbf{s} = \big(\{\mathbf{r}_i\}_{i=1}^N,\, \{\mathbf{T}_i\}_{i=1}^N,\, \mathbf{G},\, \mathcal{E}_{\text{room}}\big)$
package $N$ object representations $\mathbf{r}_i\!\in\!\mathcal{R}$ (with $\mathcal{R}$ the per-object space of §\ref{subsec:p3_s8_2_1}), their global poses $\mathbf{T}_i\!\in\!\mathrm{SE}(3)$, a relational graph $\mathbf{G}$ over the object set, and an envelope geometry $\mathcal{E}_{\text{room}}$ encoding the room shell (walls, floor, ceiling).
The conditioning space broadens to
$\mathcal{C}^{\text{scene}} \supseteq \{\,\text{text }\mathbf{x},\,\text{layout }\mathbf{L},\,\text{scene graph }\mathbf{G}_{\text{in}},\,\text{sparse views }\{(I_v,\pi_v)\}_v,\,\text{camera trajectory }\{\pi_v\}_v\,\}$,
with a condition encoder $\kappa^{\text{scene}}:\mathcal{C}^{\text{scene}}\!\to\!\mathcal{C}^{\text{scene}\prime}$ mapping each input mode into a shared embedding.
The core learning problem is to train a conditional scene generator
\begin{equation}
G_\theta^{\text{scene}}:\mathcal{C}^{\text{scene}\prime}\rightarrow \mathcal{S},
\qquad
\hat{\mathbf{s}}\sim p_\theta^{\text{scene}}\!\big(\mathbf{s}\,\big|\,\kappa^{\text{scene}}(\mathbf{c})\big),
\label{eq:p3_s8_2_2_objective}
\end{equation}
so that the sampled scene $\hat{\mathbf{s}}\!\in\!\mathcal{S}$ admits coherent inter-object geometry, semantically plausible relations, and view-consistent renderings under arbitrary $\mathrm{SE}(3)$ trajectories.
Two limiting cases fix the scope of Eq.~\eqref{eq:p3_s8_2_2_objective}.
When $N\!=\!1$, $\mathbf{G}\!=\!\varnothing$, and $\mathcal{E}_{\text{room}}\!=\!\varnothing$, the scene degenerates to a single asset and the sampler recovers Eq.~\eqref{eq:p3_s8_2_1_objective}: scene synthesis is a strict generalization of object synthesis.
At the opposite extreme, when $\mathcal{S}$ is instantiated as an ordered sequence of views $\{I_v\}_{v=1}^V$ \emph{without an explicit 3D substrate}, the sampler reduces to a camera-trajectory-conditioned video generator (the regime of §\ref{para:p3_s8_2_2_video} below), so that ``scene synthesis'' and ``consistent video generation'' are two boundary regimes of the same problem.


\subsubsection{Methodology}
\label{subsubsec:p3_s8_2_2_methodology}

Each of the four paradigms is a different factorization of $G_\theta^{\text{scene}}$ in Eq.~\eqref{eq:p3_s8_2_2_objective}, from sparse structural controls (layouts and scene graphs) through continuous 3D representations to video-based synthesis; Figure~\ref{fig:p3_s8_2_2_paradigms} previews them against a common scene.

\begin{figure}[!ht]
    \centering
    \includegraphics[width=\linewidth]{figures/part3/section8/Four_paradigm_families_for_3D_scene_synthesis.pdf}
    \caption{Four paradigm families for 3D scene synthesis (Eq.~\eqref{eq:p3_s8_2_2_objective}): (a) layout- or scene-graph-conditioned placement of objects, (b) view-to-3D reconstruction into a radiance field or 3D Gaussians, (c) view-to-view synthesis with no explicit 3D, and (d) video-model-based synthesis along a camera path. The families are not mutually exclusive, and view consistency does not guarantee a valid 3D asset; room images are illustrative.}
    \label{fig:p3_s8_2_2_paradigms}
\end{figure}

\paragraph{Layout and Scene-graph Conditioned Methods.}
\label{para:p3_s8_2_2_layout}
A major line of research generates 3D environments from high-level structural priors such as layouts $\mathbf{L}$ or scene graphs $\mathbf{G}_{\text{in}}$, instantiating Eq.~\eqref{eq:p3_s8_2_2_objective} by materializing the structured tuple as an explicit object set,
\begin{equation}
\hat{\mathbf{s}} = \big\{(\hat{\mathbf{r}}_i,\,\hat{\mathbf{T}}_i)\big\}_{i=1}^N,
\qquad
\hat{\mathbf{s}}\sim p_\theta\!\Big(\{(\mathbf{r}_i,\mathbf{T}_i)\}_i \,\Big|\, \kappa(\mathbf{L} \text{ or } \mathbf{G}_{\text{in}}),\,\phi(\mathbf{x})\Big),
\label{eq:p3_s8_2_2_layout_objective}
\end{equation}
where the joint over the $N$ slots and their poses is factorized autoregressively (a transformer over object slots) or as a denoising diffusion over the slot tensor.
SceneFormer~\cite{wang2021sceneformer} explored autoregressive modeling of room layouts and furniture placement, and DiffuScene~\cite{tang2024diffuscene} introduced diffusion over object arrangements, improving realism and diversity.
Recent work adds language and visual reasoning for compositional control: LayoutGPT~\cite{feng2023layoutgpt} plans the layout with an LLM before neural scene construction, and LayoutVLM~\cite{sun2025layoutvlm} optimizes layout consistency under multimodal alignment.
Scene-graph approaches (GraphDreamer~\cite{gao2024graphdreamer}, CommonScenes~\cite{zhai2023commonscenes}, EchoScene~\cite{zhai2024echoscene}) use relational reasoning to keep placements plausible, and InstructScene~\cite{lin2024instructscene} adds natural-language conditions for controllable scene creation.
Beyond indoor settings, CityCraft~\cite{deng2024citycraft} scales layout synthesis to the city level, and GALA3D~\cite{zhou2024gala3d} couples layout-guided diffusion with Gaussian-based rendering for photorealistic multi-object environments.

\paragraph{View-to-3D Representation Synthesis.}
\label{para:p3_s8_2_2_view3d}
The second paradigm trades the explicit structured set for a \emph{continuous} 3D representation, recovering the scene as a radiance field or 3D Gaussian collection from sparse calibrated views; the scene tuple collapses into a single field,
\begin{equation}
\hat{\mathcal{R}}^{\text{scene}} = \mathcal{F}_\psi\!\Big(\big\{(I_v,\pi_v)\big\}_{v=1}^V\Big),
\qquad \hat{\mathbf{s}}\equiv \hat{\mathcal{R}}^{\text{scene}},
\label{eq:p3_s8_2_2_view3d_objective}
\end{equation}
where $\mathcal{F}_\psi$ is either a deterministic feed-forward regressor or a generative sampler over the field.
This scales the multi-view-and-lift recipe of §\ref{para:p3_s8_2_1_mv} from one asset to a full scene by replacing the per-object reconstructor with a scene-scale one.
Deterministic regressors such as PixelNeRF~\cite{yu2021pixelnerf}, PixelSplat~\cite{charatan2024pixelsplat}, DepthSplat~\cite{xu2024depthsplat}, and MVSplat~\cite{chen2024mvsplat} / MVSplat360~\cite{chen2024mvsplat360} stay strictly faithful to the input views and are reconstruction methods, treated in §\ref{para:p3_s8_2_3_feedforward}.
The generative variant samples the field instead: DreamScene~\cite{li2024dreamscene} introduces a diffusion prior over Gaussian fields for text-conditioned scene creation, and Generative 3DGS~\cite{Schwarz2025ggs} generates 3D Gaussians for scenes with video priors.

\paragraph{View-to-View Synthesis (Regression and Generative).}
\label{para:p3_s8_2_2_viewview}
The third paradigm skips the explicit 3D representation altogether and formulates novel view synthesis as a direct conditional 2D-to-2D map from observed to target views,
\begin{equation}
\hat{I}_{\pi^*}\sim p_\theta\!\Big(I_{\pi^*}\,\Big|\,\big\{(I_v,\pi_v)\big\}_{v=1}^V,\,\pi^*\Big),
\label{eq:p3_s8_2_2_viewview_objective}
\end{equation}
where $\pi^*$ is the target camera pose and no intermediate 3D representation is materialized.
Regression-based approaches learn a deterministic map: SRT~\cite{sajjadi2022scene} encodes the scene into latent tokens that support interpolation and extrapolation across viewpoints, and LVSM~\cite{jin2025lvsm} predicts novel views with a large feed-forward transformer and minimal 3D inductive bias.
Generative approaches model the conditional distribution over unseen views, which permits stochastic synthesis under ambiguity and large viewpoint changes: Photoconsistent-NVS~\cite{yu2023long} and SEVA~\cite{zhou2025stable} use diffusion to improve diversity and robustness along extended camera trajectories, and CausNVS~\cite{kong2025causnvs} generates views autoregressively and causally, modeling inter-view dependencies explicitly to stay coherent across complex viewpoint transitions.

\paragraph{Video-model Based Scene Synthesis.}
\label{para:p3_s8_2_2_video}
The fourth paradigm adds \emph{temporal} coherence to the viewpoint variety of the view-to-view family by exploiting large video generation models as spatiotemporal priors; the scene tuple is materialized as a frame sequence,
\begin{equation}
\{\hat{I}_v\}_{v=1}^V \sim p_\theta^{\text{video}}\!\Big(\{I_v\}_{v=1}^V \,\Big|\, \{\pi_v\}_{v=1}^V,\,\phi(\mathbf{x})\Big),
\label{eq:p3_s8_2_2_video_objective}
\end{equation}
treating the scene as the implicit 3D structure of a coherent video and the camera trajectory $\{\pi_v\}$ as an explicit conditioning channel.
Where the feed-forward reconstructors of §\ref{para:p3_s8_2_1_lrm} output a static asset in one pass, the video paradigm trades static 3D for temporally coherent multi-frame output that mirrors the navigability requirement of a scene more directly.
CameraCtrl~\cite{he2024cameractrl} and MotionCtrl~\cite{wang2024motionctrl} condition video generation on camera trajectories directly; ViewCrafter~\cite{yu2024viewcrafter} and Gen3C~\cite{ren2025gen3c} add geometric guidance by unprojecting observed views into an intermediate 3D representation and reprojecting it to the target viewpoints, which improves pose fidelity; and CAT4D~\cite{wu2025cat4d} and 4DiM~\cite{watson2024controlling} train jointly on posed static multi-view data and on temporal video data, so that viewpoint variation and temporal dynamics are modeled together for 4D scene synthesis.


\subsubsection{Benchmarks}
\label{subsubsec:p3_s8_2_2_benchmarks}

The four paradigms draw on different training corpora and report on different metric panels, each inherited from a precursor task.

\paragraph{Training Corpora.}
Table~\ref{tab:p3_s8_2_2_corpora} lists the corpora in use.
Layout-conditioned methods rely overwhelmingly on 3D-FRONT~\cite{fu20213dfront}, usually combined with its companion 3D-FUTURE~\cite{fu20213dfuture} furniture library for instance-level geometry.
View-based methods (P2 / P3) draw on real-world scans of residential interiors (ScanNet~\cite{dai2017scannet}, Matterport3D~\cite{chang2017matterport3d}), on photorealistic synthetic corpora with ground-truth depth and normals (Hypersim~\cite{roberts2021hypersim}), and on Replica~\cite{straub2019replica} for compact high-fidelity indoor benchmarks.
Video-model-based methods (P4) inherit the web-scale supervision of their underlying video prior, far larger than any 3D corpus but without explicit 3D supervision, supplemented by camera-pose-annotated subsets distilled from ScanNet, RealEstate10K~\cite{zhou2018stereo}, or hand-labeled trajectories.

\begin{table}[t]
    \centering
    \fontsize{8}{10}\selectfont
    \setlength{\tabcolsep}{1.0mm}
    \caption{Representative training corpora for 3D scene synthesis.}
    \label{tab:p3_s8_2_2_corpora}
    \begin{tabular}{L{0.20\textwidth} C{0.07\textwidth} C{0.14\textwidth} L{0.22\textwidth} L{0.29\textwidth}}
        \toprule
        \textbf{Dataset} & \textbf{Year} & \textbf{Scale} & \textbf{Acquisition} & \textbf{Primary use / paradigm} \\
        \midrule
        ScanNet~\cite{dai2017scannet} & 2017 & ${\sim}1.5\mathrm{K}$ scenes & Real-world RGB-D scans & Sparse-view supervision (P2 / P3) \\
        Matterport3D~\cite{chang2017matterport3d} & 2017 & ${\sim}90$ buildings & Real-world RGB-D scans & Whole-house multi-view (P2 / P3) \\
        Replica~\cite{straub2019replica} & 2019 & ${\sim}18$ scenes & Real-world high-fidelity scans & Photorealistic indoor benchmarks (P2 / P3) \\
        Hypersim~\cite{roberts2021hypersim} & 2021 & ${\sim}460$ scenes & Photorealistic synthetic & Ground-truth depth / normals (P2 / P3) \\
        3D-FRONT~\cite{fu20213dfront} & 2021 & ${\sim}19\mathrm{K}$ rooms & Synthetic indoor layouts & Layout / object-arrangement supervision (P1) \\
        3D-FUTURE~\cite{fu20213dfuture} & 2021 & ${\sim}10\mathrm{K}$ assets & Synthetic furniture library & Companion to 3D-FRONT (P1) \\
        \bottomrule
    \end{tabular}
\end{table}

\paragraph{Evaluation Benchmark Suites.}
Layout-conditioned methods report the metrics introduced with ATISS on 3D-FRONT (object-category frequency, pairwise-relation accuracy, and FID on rendered top-down or perspective views).
View-to-3D and view-to-view methods adopt the NeRF-era panel of PSNR / SSIM / LPIPS on held-out views of ScanNet, RealEstate10K~\cite{zhou2018stereo}, or DL3DV-10K~\cite{ling2024dl3dv}.
Video-model-based methods add FVD and a camera-trajectory-adherence metric, the mean translation and rotation error between requested and recovered camera poses.

\paragraph{Constraint-specific Metrics.}
Four axes mirror the scene tuple of Eq.~\eqref{eq:p3_s8_2_2_objective}: \textbf{layout adherence} (object-count agreement, pairwise-relation accuracy on $\mathbf{G}$, and bounding-box IoU between requested and synthesized placements); \textbf{geometric fidelity} where ground truth exists (Chamfer Distance, F-score, and depth and normal consistency against held-out views); \textbf{photometric quality} (PSNR / SSIM / LPIPS against ground-truth views, FID / KID against held-out scene distributions, and FVD for video outputs); and \textbf{camera-pose adherence}, most relevant for P4 (translation and rotation error of poses recovered by structure-from-motion or a learned pose regressor).
A paper reports the subset of axes its paradigm can be scored on, which is why comparison across paradigms remains difficult.
% (§\ref{subsec:p3_s8_2_summary}).


\subsubsection{Challenges and Trends}
\label{subsubsec:p3_s8_2_2_challenges}

Scene synthesis meets the problems of object synthesis on a structured set: more conditioning variables, more output components, and longer trajectories enter the cost at once.

\paragraph{Cross-Paradigm Inconsistency.}
The four paradigms solve subtly different problems, and each leaves a gap when used alone: a layout-conditioned generator can satisfy the symbolic prior $\mathbf{G}_{\text{in}}$ yet produce a non-traversable static scene, and a view-to-3D generator can produce a navigable Gaussian field that ignores the requested object semantics.
The gaps surface as \emph{layout-adherence drift}, where the generated NeRF or 3DGS no longer satisfies the input layout once the user roams outside the trained view envelope; \emph{object drop}, where small or thin objects from $\mathbf{L}$ are silently omitted during continuous-field synthesis; and \emph{geometry flattening}, where a multi-object composition collapses into a single textured slab when the 2D or video prior dominates.
Hybrids such as GALA3D~\cite{zhou2024gala3d} and ViewCrafter~\cite{yu2024viewcrafter} chain a layout or geometry prior with a Gaussian or video renderer to keep both halves honest, and future systems will more likely compose the four paradigms than choose among them.

\paragraph{Scaling Beyond Indoor Rooms.}
Almost all paradigms remain anchored to furnished indoor rooms because that is where the training corpora live: 3D-FRONT~\cite{fu20213dfront} alone supplies the layout supervision for the dominant P1 methods, and ScanNet, Matterport3D, and Hypersim dominate P2 / P3.
The result is an outdoor and urban scale gap that CityCraft~\cite{deng2024citycraft} is among the few systems to address, a category imbalance with many furniture classes but few terrain, vegetation, or industrial ones, and an open-world domain shift under which methods trained on residential interiors fail on industrial sites, excavations, or invented environments.
Web-scraped open-vocabulary scene corpora and hybrid synthetic-real pipelines that blend procedurally generated environments with captured scans are the active responses, both for breadth and to control the long tail of the scene-class distribution.

\paragraph{Camera-Trajectory Coherence under Long-Horizon Synthesis.}
The video paradigm (§\ref{para:p3_s8_2_2_video}) excels at short camera moves but degrades on long, unbounded trajectories that require persistent world memory.
Long trajectories expose \emph{temporal blur}, where textures shimmer or drift between frames; \emph{scene amnesia}, where objects reappear with different geometry when the camera revisits a location; and \emph{pose-conditioning collapse}, where the model learns to ignore the camera signal in favor of the marginal video distribution.
ViewCrafter~\cite{yu2024viewcrafter}'s unprojection and reprojection scaffold, Gen3C~\cite{ren2025gen3c}'s 3D cache, and CAT4D~\cite{wu2025cat4d}'s joint multi-view and video training all inject an explicit 3D inductive bias back into the video prior, but the trade-off between the video prior's perceptual quality and the geometric prior's rigidity remains an open design tension.

\paragraph{Evaluation.}
A P1 layout-adherence score and a P3 PSNR cannot be ranked against each other, no single dataset evaluates layout faithfulness, geometric fidelity, photometric quality, and pose adherence jointly, and no human-aligned evaluator analogous to HyperScore~\cite{zhang2025mate3d} has been validated on scene outputs.
% ; the cross-subfield picture is drawn in §\ref{subsec:p3_s8_2_summary}.

\paragraph{Trends: Toward 4D and World Models.}
Scene synthesis targets static geometry conditioned on camera poses, but the same formulation is being extended toward 4D synthesis and action-conditioned world modeling: Dreamer~\cite{hafner2019dreamer}, GameNGen~\cite{valevski2024diffusion}, Diamond~\cite{alonso2024diffusionworldmodelingvisual}, RTFM~\cite{worldlabs_rtfm}, and Genie3~\cite{genie3} move from static view or video synthesis toward open-ended environment generation in which spatial structure, temporal evolution, and interaction are modeled together.
From the standpoint of this subsection, scalable, temporally coherent video generation with explicit viewpoint control is the foundation those systems build on; they are treated as world models in §\ref{sec:p3_s8_6}.



\subsection{3D Scene Reconstruction}
\label{subsec:p3_s8_2_3}

3D reconstruction recovers a digital 3D representation of a scene or object from sensor observations such as RGB images, video, depth maps, or point clouds, and has long been a foundational problem in computer vision, graphics, and immersive systems such as augmented reality (AR), virtual reality (VR), and digital twins.
Its difference from the synthesis tasks of §\ref{subsec:p3_s8_2_1} and §\ref{subsec:p3_s8_2_2} is one of \emph{intent}: synthesis is open-ended generation in which plausibility and diversity dominate, whereas reconstruction must explain the input views faithfully before any inference about unobserved regions is attempted.
Reconstruction is where perception becomes geometry: it turns raw sensor streams into an explicit, metric 3D substrate that simulators, AR / VR systems, digital twins, and embodied agents can reason over, manipulate, and render, and it supplies the ground-truth geometry against which the generative tasks of §\ref{sec:p3_s8_2} are ultimately measured.
Three families follow: \emph{iterative reconstruction}, which solves a per-scene optimization against the input views; \emph{feed-forward reconstruction}, which amortizes that cost into a learned map from observations to a 3D representation; and \emph{complete reconstruction}, which generatively fills the occluded regions and so makes the boundary with synthesis explicit.
Figure~\ref{fig:3D_reconstruction} previews representative reconstruction tasks.

\begin{figure}[!ht]
    \centering
    \includegraphics[width=0.99\linewidth]{figures/part3/section8/3D-reconstruction.png}
    \caption{Illustrations of 3D reconstruction, adapted from~\cite{zhang2025advances}.}
    \label{fig:3D_reconstruction}
\end{figure}


\subsubsection{Task Definition}
\label{subsubsec:p3_s8_2_3_task}

We cast 3D scene reconstruction as the observation-conditioned counterpart of the generative frameworks of §\ref{subsec:p3_s8_2_1} and §\ref{subsec:p3_s8_2_2}.
Let $\mathcal{O}$ denote the space of multi-sensor observations.
Each observation $\mathbf{O}\in\mathcal{O}$ is represented as a finite collection of sensor measurements.
Depending on the sensing setup, $\mathbf{O}$ may consist of multi-view RGB images $\{I_v\}_{v=1}^{V}$, posed images $\{(I_v,\pi_v)\}_{v=1}^{V}$, RGB-D or RGB-normal observations $\{(I_v,d_v)\}_{v=1}^{V}$, point clouds, LiDAR scans, or their combinations.
Here, $I_v$ denotes the RGB image from view $v$, $\pi_v$ denotes the corresponding calibrated camera pose, and $d_v$ denotes an associated depth or normal map.
Let $\mathcal{R}^{\text{rec}}$ denote the output 3D-representation space, taken from the same $\mathcal{R}$ as in §\ref{subsec:p3_s8_2_1}, including mesh, SDF, voxel, triplane, NeRF, 3DGS, point map, structured latent.
The core learning problem is to obtain a reconstruction operator
\begin{equation}
G_\theta^{\text{rec}}:\mathcal{O}\rightarrow \mathcal{R}^{\text{rec}},
\qquad
\hat{\mathbf{r}}^{\text{rec}} = G_\theta^{\text{rec}}(\mathbf{O}) \quad \text{(deterministic)}
\quad \text{or} \quad
\hat{\mathbf{r}}^{\text{rec}}\sim p_\theta^{\text{rec}}\!\big(\mathbf{r}^{\text{rec}}\,\big|\,\mathbf{O}\big) \quad \text{(generative)},
\label{eq:p3_s8_2_3_objective}
\end{equation}
that minimizes a reconstruction loss against the input views,
\begin{equation}
\mathcal{L}_{\text{rec}}\big(\hat{\mathbf{r}}^{\text{rec}},\mathbf{O}\big)
\;=\;
\sum_{(I_v,\pi_v)\in\mathbf{O}} \big\|\mathrm{Render}(\hat{\mathbf{r}}^{\text{rec}},\pi_v) - I_v\big\|^2
\;+\; \lambda\,\Omega(\hat{\mathbf{r}}^{\text{rec}}),
\label{eq:p3_s8_2_3_loss}
\end{equation}
where $\Omega$ is an optional geometric or smoothness regularizer and $\mathrm{Render}$ is a differentiable renderer of $\mathcal{R}^{\text{rec}}$.
The structural contrast with §\ref{subsec:p3_s8_2_1} is sharp: there, the conditioning $\kappa(\mathbf{c})$ is a high-level prior (text, single image, multi-view) and the objective is a \emph{generative} likelihood; here, the conditioning $\mathbf{O}$ is sensor data and the objective is an \emph{observation-fidelity} loss.
% Two limiting cases anchor Eq.~\eqref{eq:p3_s8_2_3_objective} to the broader literature.
% First, when $\mathbf{O}$ is a dense calibrated multi-view set, the deterministic form of Eq.~\eqref{eq:p3_s8_2_3_objective} reduces to the classical SfM\,+\,MVS pipeline or to per-scene NeRF / 3DGS optimization: the regime of \S\ref{para:p3_s8_2_3_iterative} below.
% Second, when $\mathbf{O}$ is sparse or even a single image and large portions of the scene must be inferred, the generative form of Eq.~\eqref{eq:p3_s8_2_3_objective} merges into the synthesis regime of Eq.~\eqref{eq:p3_s8_2_1_objective}: this is precisely the regime of \S\ref{para:p3_s8_2_3_complete}, where reconstruction and synthesis become two formulations of the same underlying problem.


\subsubsection{Methodology}
\label{subsubsec:p3_s8_2_3_methodology}

The three paradigms are three parameterizations of the reconstruction operator $G_\theta^{\text{rec}}$ of Eq.~\eqref{eq:p3_s8_2_3_objective}, taken in the order in which each became dominant.
Figure~\ref{fig:p3_s8_2_3_paradigms} contrasts them, making explicit whether the 3D representation is optimized per scene, regressed in a single forward pass, or sampled from a generative posterior over the unobserved regions.

\begin{figure}[!ht]
    \centering
    \includegraphics[width=0.99\linewidth]{figures/part3/section8/3D-scene-reconstruction.pdf}
    \caption{Three paradigm families for 3D scene reconstruction (Eq.~\eqref{eq:p3_s8_2_3_objective}): (P1) iterative per-scene optimization, (P2) feed-forward reconstruction with weights shared across scenes, and (P3) generative completion of unseen (hatched) regions from a learned prior.}
    \label{fig:p3_s8_2_3_paradigms}
\end{figure}

\paragraph{Iterative 3D Reconstruction.}
\label{para:p3_s8_2_3_iterative}
The classical paradigm casts reconstruction as a \emph{per-scene} optimization: there is no learned $\theta$ shared across scenes; the 3D representation $\mathbf{r}$ itself is the optimization variable, and Eq.~\eqref{eq:p3_s8_2_3_loss} is minimized from scratch for each new input,
\begin{equation}
\hat{\mathbf{r}}
\;=\;
\arg\min_{\mathbf{r}\in\mathcal{R}^{\text{rec}}}
\;\sum_{(I_v,\pi_v)\in\mathbf{O}} \big\|\mathrm{Render}(\mathbf{r},\pi_v) - I_v\big\|^2
\;+\; \lambda\,\Omega(\mathbf{r}).
\label{eq:p3_s8_2_3_iterative_objective}
\end{equation}
Classical image-based reconstruction is rooted in Structure-from-Motion (SfM)~\cite{bolles1987epipolar, tomasi1992shape, fitzgibbon1998automatic, schonberger2016structure}, which estimates camera poses from unposed images, followed by Multi-View Stereo (MVS)~\cite{ullman1979interpretation, hartley2003multiple, Agarwal2010BAL, snavely2006photo, furukawa2015multi, schonberger2016pixelwise}, which recovers dense geometry.
Decades of refinement have made SfM\,+\,MVS a mature and widely used pipeline (see~\cite{hartley2003multiple} for a comprehensive treatment), yet it still struggles with textureless surfaces, reflective objects, and large illumination changes.

A more recent line represents the scene as a continuous function parameterized by a neural network: signed distance functions~\cite{park2019deepsdf, chibane2020neural, wang2021neus, yu2022monosdf}, occupancy fields~\cite{mescheder2019occupancy, peng2020convolutional}, and neural radiance fields~\cite{mildenhall2021nerf, barron2021mip, wang2021nerf, barron2022mip}.
The original formulations were expensive to train and render and limited in quality; PlenOctrees~\cite{yu2021plenoctrees}, Plenoxels~\cite{fridovich2022plenoxels}, Instant-NGP~\cite{muller2022instant}, RefNeRF~\cite{verbin2022ref}, TensoRF~\cite{chen2022tensorf}, and ZipNeRF~\cite{barron2023zip} are among the many follow-ups that improved efficiency and fidelity, but all of them still require a large number of input images per scene.
Sparse-view variants add auxiliary information such as regularization~\cite{niemeyer2022regnerf}, depth priors~\cite{roessle2022dense, deng2022depth}, point clouds~\cite{xu2022point}, and surface priors~\cite{oechsle2021unisurf}, or distill a 2D diffusion prior into the per-scene representation, the SDS recipe of §\ref{para:p3_s8_2_1_sds} applied to reconstruction from one or a few images~\cite{melas2023realfusion, metzer2023latent, tang2023make}.
3D Gaussian Splatting~\cite{kerbl20233d} then made per-scene optimization fast, training within minutes and rendering in real time, with subsequent work adding depth regularization, anti-aliasing, mesh extraction, and surface alignment~\cite{chung2024depth, yu2024mip, huang20242d, guedon2024sugar}.
All of these systems, however, work well only over nearby views or a small scene and must be re-optimized from scratch for every new input, which motivates the amortized paradigm.

\paragraph{Feed-forward 3D Reconstruction.}
\label{para:p3_s8_2_3_feedforward}
The feed-forward paradigm amortizes the per-scene optimization of Eq.~\eqref{eq:p3_s8_2_3_iterative_objective} across a learned prior: a network $G_\theta^{\text{rec}}$ is trained once on a large multi-view corpus, then deployed in a \emph{single forward pass} on unseen inputs, minimizing
\begin{equation}
\mathcal{L}_{\text{FF}}(\theta)
\;=\;
\mathbb{E}_{(\mathbf{O},\,\mathbf{r}_{\text{gt}})\sim p_{\text{data}}}\Big[\mathcal{D}\big(G_\theta^{\text{rec}}(\mathbf{O}),\,\mathbf{r}_{\text{gt}}\big)\Big],
\label{eq:p3_s8_2_3_feedforward_objective}
\end{equation}
where $\mathcal{D}$ is a representation-specific distance (Chamfer Distance for point maps, a rendering loss for radiance fields) and $(\mathbf{O},\mathbf{r}_{\text{gt}})$ are observation--asset pairs drawn from a 3D corpus.
This is the reconstruction-side analogue of the view-to-3D paradigm of §\ref{para:p3_s8_2_2_view3d}: sparse calibrated views are mapped feed-forward to a continuous scene representation, but the supervision is a reconstruction loss against ground-truth 3D rather than a rendering loss against held-out views alone.

Early feed-forward work applied 2D networks to individual 3D representations, for depth estimation~\cite{eigen2014depth, liu2015learning}, voxel reconstruction~\cite{choy20163d, wu20153d}, point-cloud reconstruction~\cite{fan2017point, qi2017pointnet}, and mesh deformation~\cite{wang2018pixel2mesh, gkioxari2019mesh}, but remained limited to single objects or small datasets.
Learned components then replaced individual stages of the SfM\,+\,MVS pipeline, keypoint detection and description~\cite{detone2018superpoint, dusmanu2019d2, tyszkiewicz2020disk}, image matching~\cite{sarlin2020superglue, lindenberger2023lightglue}, pose estimation~\cite{kendall2015posenet, yin2018geonet, kendall2017geometric}, and depth prediction~\cite{zhou2017unsupervised, yin2018geonet}, proving more robust and accurate than their hand-designed counterparts and being integrated into the classical pipeline~\cite{vzbontar2016stereo, yao2018mvsnet, yao2019recurrent}.
DUSt3R~\cite{wang2024dust3r} and MASt3R~\cite{mast3r} then reconstructed large scenes from unposed images end to end: similar in spirit to MVS but without explicit keypoint detection, matching, and pose estimation, the model predicts point maps in the reference view directly and recovers camera poses in post-processing, and follow-ups~\cite{tang2025mv, wang2025Cut3R, yang2025fast3r, zhang2025FLARE, wang2025vggt} have improved quality and efficiency at scale.

Beyond geometry alone, a parallel line recovers geometry and appearance together through differentiable rendering.
Scene Representation Networks~\cite{sitzmann2019scene} and Light Field Networks~\cite{sitzmann2021light} introduced ray-based querying for novel-view synthesis and were followed by epipolar attention~\cite{suhail2022light, suhail2022generalizable}, large transformers~\cite{jin2025lvsm, jiang2025rayzer}, and integration with pretrained generative models~\cite{zheng2024free3d, gao2024cat3d}; PixelNeRF~\cite{yu2021pixelnerf} pioneered pixel-aligned features for NeRF-based reconstruction and inspired 3D volume representations~\cite{chen2021mvsnerf, xu2024murf}, feature-matching fusion~\cite{chen2025explicit}, and large transformers~\cite{du2023learning, hong2023lrm}.
3D Gaussian Splatting~\cite{kerbl20233d} has since become the representation of choice for feed-forward reconstruction, in PixelSplat~\cite{charatan2024pixelsplat}, Splatter Image~\cite{szymanowicz2024splatter}, MVSplat~\cite{chen2024mvsplat}, Flash3D~\cite{szymanowicz2024flash3d}, LatentSplat~\cite{wewer2024latentsplat}, Splatt3R~\cite{smart2024splatt3r}, GS-LRM~\cite{zhang2024gs}, NoPoSplat~\cite{ye2024no}, and DepthSplat~\cite{xu2024depthsplat}.
These methods synthesize novel views with high appearance fidelity, but the views remain confined to small camera changes within the observed envelope, which motivates the generative paradigm below.

\paragraph{Complete 3D Reconstruction.}
\label{para:p3_s8_2_3_complete}
A harder regime is to reconstruct a \emph{complete} scene from an observation $\mathbf{O}$ that covers only part of it, which requires inferring both \emph{visible} and \emph{occluded} regions; the operator shifts from a deterministic map to a generative posterior over the unobserved regions,
\begin{equation}
\hat{\mathbf{r}}^{\text{rec}}\sim p_\theta^{\text{rec}}\!\big(\mathbf{r}^{\text{rec}}\,\big|\,\mathbf{O}\big),
\qquad
\mathbf{r}^{\text{rec}}\big|_{\text{visible}}\approx \text{observation-consistent},
\quad
\mathbf{r}^{\text{rec}}\big|_{\text{occluded}}\sim p_\theta(\,\cdot\,\big|\,\mathbf{O}),
\label{eq:p3_s8_2_3_complete_objective}
\end{equation}
which shares the formal structure of the synthesis sampler of Eq.~\eqref{eq:p3_s8_2_1_objective} with the conditioning replaced by a strong observation constraint $\mathbf{O}$.

Complete reconstruction has been studied extensively for single objects~\cite{vahdat2022lion, zhang20233dshape2vecset, zhao2023michelangelo, zhang2024clay, ren2024xcube, xiang2025trellis, tochilkin2024triposr, yang2024hunyuan3d, hunyuan3d22025tencent, li2025triposg} but remains under-explored for large scenes.
Classical scene-scale methods fuse depth maps from multiple views into a volumetric representation~\cite{curless1996volumetric, newcombe2011kinectfusion, whelan2015elasticfusion, dai2017bundlefusion, niessner2013real}, but they are sensitive to depth errors and produce noisy, incomplete reconstructions.
A newer line fine-tunes large pretrained diffusion models~\cite{rombach2022ldm, blattmann2023stable} for the task, conditioning on controlled camera trajectories~\cite{wang2024motionctrl, sargent2024zeronvs, wu2024reconfusion, gao2024cat3d, wallingford2024image, zhou2025stable} or on auxiliary signals~\cite{liu2024reconx, yu2024viewcrafter, chen2024mvsplat360, yu2025wonderworld}; these methods synthesize observation-consistent views but leave the underlying geometry implicit.
A small but growing body of work addresses geometry directly: WVD~\cite{zhang2024world} and Bolt3D~\cite{szymanowicz2025bolt3d} generate a hybrid RGB\,+\,point-map output that recovers geometry and appearance rather than novel views alone, and LaRI~\cite{li2025lari} reconstructs a complete scene from a single image, still limited to small indoor scenes.

Two design choices cut across the three paradigms.
The density of the observation decides which paradigm applies: a dense calibrated multi-view set calls for per-scene optimization, a sparse uncalibrated capture for a feed-forward map, and a single image for a generative posterior, since progressively more of the scene must be inferred rather than observed.
The output representation $\mathcal{R}^{\text{rec}}$, by contrast, is largely independent of the paradigm and driven by the downstream consumer: a point map (DUSt3R), a mesh (classical MVS post-processing), a NeRF (per-scene optimization, PixelNeRF), a 3DGS field (3DGS, PixelSplat), or a hybrid RGB\,+\,point map (Bolt3D, WVD).


\subsubsection{Benchmarks}
\label{subsubsec:p3_s8_2_3_benchmarks}

Datasets span object-centric collections, indoor scenes, and outdoor environments, and partition further by acquisition (synthetic or real-world) and dynamics (static or dynamic); Table~\ref{tab:p3_s8_2_3_corpora} summarizes those most commonly used.

\paragraph{Training Corpora.}
Synthetic corpora such as Objaverse~\cite{deitke2023objaverse} supply large and diverse 3D assets, and MegaSynth~\cite{jiang2025megasynth} and Zeroverse~\cite{xie2024lrmzero} explore procedurally generated data for scalable supervision.
Real-world static corpora include ARKitScenes~\cite{baruch2021arkitscenes} and ScanNet++~\cite{yeshwanth2023scannet++} for indoor capture, dynamic corpora such as KITTI-360~\cite{liao2022kitti360} capture outdoor scenes with temporal variation, and TartanAir~\cite{TartanAir} contains both static and dynamic scenes, which permits evaluation across both settings.

\begin{table}[t]
    \centering
    \fontsize{8}{10}\selectfont
    \setlength{\tabcolsep}{1.0mm}
    \caption{Representative training corpora for 3D scene reconstruction. ``Type'' indicates synthetic / real-world acquisition; ``Dynamics'' indicates whether the corpus contains temporal variation.}
    \label{tab:p3_s8_2_3_corpora}
    \begin{tabular}{L{0.17\textwidth} C{0.06\textwidth} C{0.14\textwidth} L{0.13\textwidth} C{0.09\textwidth} L{0.33\textwidth}}
        \toprule
        \textbf{Dataset} & \textbf{Year} & \textbf{Scale} & \textbf{Type} & \textbf{Dynamics} & \textbf{Primary use} \\
        \midrule
        Objaverse~\cite{deitke2023objaverse} & 2023 & ${\sim}800\mathrm{K}$ objects & Synthetic (CAD) & Static & Object-centric multi-view supervision \\
        Zeroverse~\cite{xie2024lrmzero} & 2024 & Procedural & Synthetic & Static & LRM-zero scalable supervision \\
        MegaSynth~\cite{jiang2025megasynth} & 2025 & Procedural & Synthetic & Static & Scalable multi-view scene supervision \\
        ARKitScenes~\cite{baruch2021arkitscenes} & 2021 & ${\sim}5\mathrm{K}$ scans & Real-world & Static & Indoor RGB-D / mesh supervision \\
        ScanNet++~\cite{yeshwanth2023scannet++} & 2023 & ${\sim}460$ scenes & Real-world & Static & High-fidelity indoor multi-modal scans \\
        KITTI-360~\cite{liao2022kitti360} & 2022 & ${\sim}300\mathrm{K}$ frames & Real-world & Dynamic & Outdoor driving scenes \\
        TartanAir~\cite{TartanAir} & 2020 & ${\sim}1\mathrm{M}$ frames & Synthetic & Mixed & Static + dynamic mixed-domain evaluation \\
        \bottomrule
    \end{tabular}
\end{table}

\paragraph{Evaluation Benchmark Suites.}
Iterative methods (P1) are evaluated on DTU~\cite{Aanaes2016DTU} and Tanks-and-Temples~\cite{knapitsch2017tanks} for geometric fidelity and on the NeRF-era benchmarks LLFF~\cite{mildenhall2019llff}, MipNeRF360~\cite{barron2022mip}, and RealEstate10K~\cite{zhou2018stereo} for novel-view quality.
Feed-forward methods (P2) adopt the same panel together with ScanNet~\cite{dai2017scannet} subsets that supply ground-truth point maps for direct geometry supervision.
Complete-reconstruction methods (P3) have no standardized suite and typically report on the precursor benchmark of their underlying 2D or video prior.

\paragraph{Constraint-specific Metrics.}
Four axes mirror the rendering-versus-geometry decomposition of Eq.~\eqref{eq:p3_s8_2_3_loss}.
\textbf{Photometric quality} is measured by PSNR, SSIM~\cite{wang2004image}, and LPIPS~\cite{zhang2018unreasonable}, covering pixel-level fidelity, structural similarity, and perceptual consistency of novel views.
\textbf{Camera-pose adherence} uses Relative Translation Accuracy (RTA), Relative Rotation Accuracy (RRA), and the Area Under the Curve (AUC) that summarizes pose accuracy across angular thresholds.
\textbf{Geometric fidelity} of point maps and point clouds is measured by Accuracy (precision), Completeness (recall), and Chamfer Distance, where Accuracy reflects geometric precision, Completeness measures surface coverage, and Chamfer Distance aggregates both.
\textbf{Depth quality} for monocular estimation uses the absolute relative error $|y - \hat{y}| / y$ and the inlier ratio $\delta\!<\!T$ (the fraction of predictions within a multiplicative threshold of the ground truth), with median scaling or least-squares fitting applied beforehand because monocular depth is scale-ambiguous; multi-view depth reports Accuracy, Completeness, and Overall, the mean of the first two and equivalent to Chamfer Distance.

The paradigm succession of §\ref{subsubsec:p3_s8_2_3_methodology} is visible in two quantities.
Per-scene optimization has moved from NeRF~\cite{mildenhall2021nerf}, which trains for hours and renders at a few frames per second, to 3D Gaussian Splatting~\cite{kerbl20233d}, which trains a scene in minutes and renders in real time (${\geq}30$~fps at 1080p); the feed-forward paradigm removes per-scene optimization altogether, with DUSt3R / MASt3R~\cite{wang2024dust3r, mast3r} and VGGT~\cite{wang2025vggt} recovering point maps, depth, and camera poses in a single sub-second forward pass and Fast3R~\cite{yang2025fast3r} processing over a thousand images at once.
On pose, learned reconstructors now rival classical pipelines: VGGT~\cite{wang2025vggt} reports a camera-pose AUC@30 of 85.3 on RealEstate10K in pure feed-forward mode and 93.5 with bundle-adjustment refinement.
Two caveats apply. There is no cross-paradigm leaderboard, since P1 is scored on novel-view PSNR / SSIM / LPIPS, P2 adds geometric Accuracy / Completeness / Chamfer Distance and pose AUC, and P3 inherits whichever protocol its 2D or video prior used; and on raw depth-map accuracy feed-forward models still trail tuned MVS pipelines, with quality degrading under free-viewpoint extrapolation (§\ref{subsubsec:p3_s8_2_3_challenges}), so the strongest numbers hold inside the observed view envelope.
Phone-based NeRF and Gaussian-Splatting capture apps and the open-source Nerfstudio stack have brought per-scene reconstruction to commodity devices, while classical photogrammetry suites remain the production baseline for metric accuracy.


\subsubsection{Challenges and Trends}
\label{subsubsec:p3_s8_2_3_challenges}

Four problems span the pipeline end to end: the data the model trains on, the accuracy it can deliver, the quality of the views it renders, and the compute envelope that constrains all three at long input lengths; reconstruction shares the first and third with its synthesis counterparts but sharpens the other two because it commits to recovering ground truth from raw sensor data.

\paragraph{Limited Modality in Datasets.}
Most reconstruction and view-synthesis datasets cover a narrow band of modalities.
Widely used benchmarks such as RealEstate10K~\cite{zhou2018stereo} and MipNeRF360~\cite{barron2022mip} are RGB-only, without depth, LiDAR, or semantic annotation, and even Objaverse-XL~\cite{deitke2023objaversexl}, with 10.2M objects, is primarily synthetic mesh data without real-world sensor modalities.
Merging datasets of different modalities is the usual workaround, at the cost of domain shift and annotation inconsistency, and the gap is widest for dynamic scenes: datasets with synchronized RGB, depth, optical flow, and 3D motion remain far fewer than static ones, and most dynamic datasets record either camera motion or object motion rather than both with a full sensor suite.
The open problem is a scalable, modality-rich corpus that combines the diversity of Objaverse~\cite{deitke2023objaverse} with the multi-sensor completeness of ScanNet++~\cite{yeshwanth2023scannet++}, collected under one protocol so that aligned RGB, depth, and semantic signals supervise the same model.

\paragraph{Reconstruction Accuracy.}
Feed-forward models still trail traditional MVS pipelines~\cite{zhang2023geomvsnet, zhang2025advances} on raw depth-map accuracy, because MVS exploits known camera parameters and a hypothesized depth set to build cost volumes from which depth is predicted.
Whether feed-forward reconstructors~\cite{yang2025fast3r, wang2025vggt} learn an approximation of such cost volumes implicitly in their attention layers is an open question; injecting explicit camera parameters or other priors into the backbone, for instance through Diffusion Transformers~\cite{peebles2023dit}, is one route to closing the gap.

\paragraph{Free-viewpoint Rendering.}
The accuracy gap is most visible at \emph{free-viewpoint} rendering, where the model must produce novel views far from the input camera distribution.
Disocclusions, geometric uncertainty, and limited generalization conspire: regions never observed come out blurred, ghosted, or geometrically wrong, because current methods rely on local consistency and cannot infer plausible content for occluded areas, and view-dependent effects and complex light transport demand reasoning beyond interpolation.
Hallucinating missing detail while keeping novel views consistent is exactly the generative regime of §\ref{para:p3_s8_2_3_complete}.

\paragraph{Long Context Input.}
Reconstruction from many views meets a hard computational ceiling: methods for geometry reasoning and novel-view synthesis rely on full attention, whose compute grows quadratically with the number of tokens; exact-attention kernels such as FlashAttention-2~\cite{dao2023flashattention2} remove the quadratic memory of the attention matrix but not the quadratic compute.
VGGT~\cite{wang2025vggt}, for example, needs about 21\,GB of GPU memory for 100 input frames and about 41\,GB for 200, so inference over long sequences and training on many views quickly exceed a single GPU.
Recurrent alternatives such as CUT3R~\cite{wang2025Cut3R} integrate new views incrementally while maintaining a memory state, which keeps inference memory near constant, but they forget earlier views and degrade as the number of inputs grows.
Reasoning efficiently over hundreds or thousands of views while keeping memory and compute manageable remains open, with recurrent and sparse-attention architectures the main candidates.



\subsection{3D Panorama Synthesis}
\label{subsec:p3_s8_2_4}

3D panorama synthesis occupies a distinctive position in spatial generation: rather than producing a 3D representation directly, it generates a complete 360-degree omnidirectional image as a structured \emph{intermediate} from which a full 3D environment can then be lifted.
The panorama gives embodied agents, AR / VR pipelines, and world models a holistic spatial context with minimal occlusion anchored at one observation point, and it bridges the mature 2D generative prior of §\ref{sec:p3_s8_1} with the 3D substrate on which the rest of §\ref{sec:p3_s8_2} operates.
Four method families trace the path from a single omnidirectional snapshot to an explorable world: panoramic image synthesis, panoramic video generation, panoramic 3D reconstruction, and panorama-based 3D scene generation.


\subsubsection{Task Definition}
\label{subsubsec:p3_s8_2_4_task}

Panorama synthesis generates a complete 360-degree environment conditioned on text, a reference panorama or NFoV image, a spatial layout, or a specified viewpoint; unlike object-centric or narrow-field-of-view (NFoV) generation, it covers the whole environment around the observer at once rather than a fragment of it.
Formally, we cast it as the omnidirectional-intermediate counterpart of the generative frameworks of §\ref{subsec:p3_s8_2_1} and §\ref{subsec:p3_s8_2_2}.
Let $\mathcal{P}$ denote the space of panoramic images (functions defined on the sphere $\mathbb{S}^2$ and encoded through a projection operator $\Pi\in\{\text{ERP},\text{cubemap},\text{tangent},\text{fisheye}\}$ onto a processable 2D array) and let $\mathbf{c}\in\mathcal{C}\supseteq\{\text{text},\text{reference panorama / NFoV image},\text{layout},\text{viewpoint / camera trajectory}\}$ be the conditioning.
Panorama synthesis first learns a panoramic generator and then optionally lifts its output into 3D,
\begin{equation}
G_\theta^{\text{pano}}:\mathcal{C}'\rightarrow\mathcal{P},
\qquad
\hat{P}\sim p_\theta^{\text{pano}}\!\big(P\,\big|\,\kappa(\mathbf{c})\big),
\qquad
\hat{\mathbf{r}}^{\text{rec}} = \Lambda\big(\hat{P}\big),
\label{eq:p3_s8_2_4_objective}
\end{equation}
where $\kappa$ is the condition encoder of §\ref{subsec:p3_s8_2_1}, $\hat{P}\in\mathcal{P}$ is the synthesized omnidirectional image, and $\Lambda:\mathcal{P}\rightarrow\mathcal{R}^{\text{rec}}$ is an optional \emph{lifting operator} mapping the panorama into the 3D-representation space $\mathcal{R}^{\text{rec}}$ of §\ref{subsec:p3_s8_2_3} (drawn from $\mathcal{R}$ of §\ref{subsec:p3_s8_2_1}, \eg NeRF, 3DGS, mesh, point map).
Equation~\eqref{eq:p3_s8_2_4_objective} makes the role of the panorama as a structured \emph{intermediate} explicit: the hard generative work is performed on the sphere by $G_\theta^{\text{pano}}$ (inheriting the 2D generator $G_\theta$ of Eq.~\eqref{eq:t2i_generator} with its outputs defined on $\mathbb{S}^2$ rather than the plane), while $\Lambda$ reuses the reconstruction machinery of §\ref{subsec:p3_s8_2_3}.
This is also what separates panorama synthesis from the scene synthesis of §\ref{subsec:p3_s8_2_2}, which produces 3D representations directly from arbitrary viewpoints distributed over $\mathrm{SE}(3)$: here every direction around one observation point is represented coherently first, and 3D structure is recovered from that complete view.
Two limiting cases locate Eq.~\eqref{eq:p3_s8_2_4_objective} within §\ref{sec:p3_s8_2}.
When $\Lambda=\mathrm{id}$ (no lifting), the task collapses to omnidirectional 2D image generation, a spatially aware variant of the text-to-image sampler of §\ref{sec:p3_s8_1} defined on the sphere rather than the plane.
When $\mathbf{c}$ is a set of sparse posed panoramas and $\Lambda$ is a per-scene optimization, Eq.~\eqref{eq:p3_s8_2_4_objective} reduces to panoramic reconstruction, recovering the observation-conditioned operator of Eq.~\eqref{eq:p3_s8_2_3_objective}.

\paragraph{Panoramic Image Formats.}
The projection $\Pi$ fixes how the sphere is laid out for a 2D network and which distortions the generator must absorb.
Equirectangular projection (ERP) maps the sphere onto a rectangle in latitude--longitude coordinates; it stretches content near the poles but is continuous and directly compatible with standard 2D pipelines and convolutional networks, which has made it the dominant format for learning-based synthesis.
Cubemap projection represents the sphere as six perspective faces (front, back, left, right, top, bottom), removing polar distortion at the cost of discontinuities along face boundaries, and is the native format of real-time rendering and VR / AR systems, where each face is processed as an ordinary perspective image.
Tangent images and perspective patches decompose the panorama into overlapping tangent-plane projections so that models trained on ordinary perspective images can be reused.
Fisheye and dual-fisheye projections capture 180--220 degrees per lens with radial distortion; most consumer 360-degree cameras capture two hemispheres this way and stitch them, so dual-fisheye is the native acquisition format even when methods operate on ERP.


\subsubsection{Methodology}
\label{subsubsec:p3_s8_2_4_methodology}

The four families are four instantiations of the panoramic generator $G_\theta^{\text{pano}}$ and lifting operator $\Lambda$ of Eq.~\eqref{eq:p3_s8_2_4_objective}: image synthesis stops at a static panorama, video generation extends it over time, panoramic reconstruction applies $\Lambda$ to observed panoramas, and panorama-based scene generation composes the two.
Figure~\ref{fig:p3_s8_2_4_paradigms} schematizes them and the panorama-as-intermediate pipeline that connects them.

\begin{figure}[!ht]
    \centering
    \includegraphics[width=0.99\linewidth]{figures/part3/section8/Four_method_families_for_3D_panorama_synthesis.pdf}
    \caption{The four method families of 3D panorama synthesis (Eq.~\eqref{eq:p3_s8_2_4_objective}), linked by a panoramic intermediate (top): (a) panoramic image synthesis, (b) panoramic video generation, (c) panoramic 3D reconstruction, and (d) panorama-based 3D scene generation. Room imagery is illustrative.}
    \label{fig:p3_s8_2_4_paradigms}
\end{figure}

\paragraph{Panoramic Image Synthesis.}
Two task formulations differ in their conditioning: \emph{text-to-panorama generation}, which leans on the prior of a pretrained generator and is judged on semantic alignment and diversity, and \emph{image-conditioned panorama generation}, which expands an NFoV image or an incomplete panorama by inpainting or outpainting and must preserve the content, style, and panoramic distortion of its input.
Formally, this family is the base case of Eq.~\eqref{eq:p3_s8_2_4_objective} with the lifting operator disabled: it samples a single static omnidirectional image and stops,
\begin{equation}
\hat{P} \sim p_\theta^{\text{pano}}\!\big(P \mid \kappa(\mathbf{c})\big),
\qquad
\Lambda = \mathrm{id},
\label{eq:p3_s8_2_4_image}
\end{equation}
where the generative work is carried out under a chosen projection $\Pi\in\{\text{ERP},\text{cubemap},\text{tangent}\}$ that dictates how spherical distortion is handled; Figure~\ref{fig:p3_s8_2_4_paradigms}(a) depicts this static-panorama pathway.

Early approaches~\cite{chen2022text2light, may2023cubegan} relied on GANs~\cite{goodfellow2014gan} and were constrained by limited training data and mode collapse; diffusion models~\cite{ho2020denoising, rombach2022ldm} brought diversity and fidelity, and current methods divide by projection format and architecture.
The three strategies below are three realizations of the sampler of Eq.~\eqref{eq:p3_s8_2_4_image}, differing in how the spherical output is factorized before being mapped onto the sphere,
\begin{equation}
\hat{P} =
\begin{cases}
\Pi_{\text{ERP}}^{-1}(\hat{Z}),\ \ \hat{Z}\sim p_\theta\big(Z\mid\kappa(\mathbf{c})\big) & \text{(direct ERP)},\\[3pt]
\Pi_{\text{cube}}^{-1}\big(\{\hat{F}_i\}_{i=1}^{6}\big),\ \ \{\hat{F}_i\}\sim p_\theta\big(\{F_i\}\mid\kappa(\mathbf{c})\big) & \text{(cubemap)},\\[3pt]
\mathcal{S}\big(\{\hat{I}_k\}_{k=1}^{K}\big),\ \ \hat{I}_k\sim p_\theta^{\text{2D}}\big(I\mid\kappa(\mathbf{c})\big) & \text{(perspective + stitch)},
\end{cases}
\label{eq:p3_s8_2_4_image_proj}
\end{equation}
where $\Pi^{-1}$ lifts the generated representation onto the sphere and $\mathcal{S}$ is a spherical stitch-and-blend operator. The first two cases synthesize the panorama \emph{directly} (and so require panoramic training data), whereas the third reuses a frozen perspective generator $p_\theta^{\text{2D}}$ and pushes the cross-view-consistency burden into $\mathcal{S}$.

$\blacktriangleright$ \textit{Equirectangular projection (ERP) methods.} The direct-ERP case generates on the latitude--longitude grid and must handle polar distortion and the wrap-around seam architecturally.
For text-to-panorama generation, DiffPano~\cite{ye2024diffpano} fine-tunes a diffusion model with LoRA on ERP, SMGD~\cite{sun2025spherical} guides the diffusion process with spherical manifold convolutions that respect geometry across the sphere, and PAR~\cite{wang2025conditional} replaces diffusion with autoregressive modeling to avoid the mismatch between ERP and the noise assumptions of standard diffusion.
The seam is treated at the denoising or decoding stage: PanoDiff~\cite{panodiff} uses geometric augmentation when completing a panorama from unregistered NFoV images, PanoDiffusion~\cite{wupanodiffusion} generates depth jointly and rotates the latent circularly to enforce left--right continuity, StitchDiffusion~\cite{stitchdiffusion} re-denoises the wrap-around region at every step together with global cropping, and PanFusion~\cite{zhang2024taming} pairs an ERP branch with a multi-view perspective branch to draw on both representations.
For image-conditioned generation, OmniDreamer~\cite{omnidreamer2022} introduced mask-based outpainting from NFoV images with GANs, AOG~\cite{aog} and ODIS~\cite{odis} explore alternative conditioning for inpainting and outpainting, and PanoDecouple~\cite{zheng2025panorama} separates distortion guidance from content completion to counter the ``visual cheating'' in which a model buys image quality at the cost of geometric accuracy.

$\blacktriangleright$ \textit{Cubemap-based methods.} The second case of Eq.~\eqref{eq:p3_s8_2_4_image_proj} generates six perspective faces, which removes polar distortion but moves the problem to the face boundaries.
CubeDiff~\cite{kalischek2025cubediff} introduced the cubemap representation for text-to-panorama generation, and DreamCube~\cite{huang2025dreamcube} synchronizes the six planes to suppress discontinuous seams and inconsistent color tone; overlapping the fields of view of adjacent faces reduces seams further but duplicates generation in the overlap and lowers the effective resolution.
JoPano~\cite{feng2025jopano} models the faces jointly with a DiT and a Joint-Face Adapter, applies Poisson blending along the boundaries, and introduces the seam metrics (Seam-SSIM, Seam-Sobel) used in §\ref{subsubsec:p3_s8_2_4_benchmarks}.

$\blacktriangleright$ \textit{Perspective image-based methods.} The third case leaves a perspective text-to-image model untouched and pushes the consistency burden into the stitch operator $\mathcal{S}$, with or without panorama-specific training.
\emph{Training-free methods} apply a pretrained perspective model directly: MultiDiffusion~\cite{multidiffusion} denoises overlapping patches independently and blends them, PanoGen~\cite{li2023panogen} outpaints recursively from an initial view until the sphere is covered, DiffCollage~\cite{zhang2023diffcollage} casts the panorama as a factor graph whose factor nodes are portions and whose variable nodes are their overlaps so that portions are generated in parallel rather than autoregressively, PanoFree~\cite{panofree} iterates warping and inpainting over multi-view images, and SphereDiff~\cite{park2025spherediff} moves MultiDiffusion onto a spherical latent with dynamic spherical sampling and distortion-aware weighted averaging so that quality holds up to the poles.
These methods ride on advances in perspective generation but struggle with global consistency across the full sphere.
\emph{Training-based methods} fine-tune on panoramic data: MVDiffusion~\cite{tang2023MVDiffusion} generates eight perspective views with consistent geometry and reprojects them into a closed-loop panorama, and TanDiT~\cite{tandit2025} decomposes the panorama into a structured grid of tangent-plane images that a Diffusion Transformer generates jointly in one denoising pass under spatial attention.


\paragraph{Panoramic Video Generation.}
Extending synthesis over time adds temporal coherence to omnidirectional consistency.
Formally, this family extends the static generator of Eq.~\eqref{eq:p3_s8_2_4_image} over time, sampling a sequence under a joint spatial-and-temporal coherence constraint,
\begin{equation}
\{\hat{P}_t\}_{t=1}^{T} \sim p_\theta^{\text{pano}}\!\big(\{P_t\}_{t=1}^{T} \mid \kappa(\mathbf{c})\big),
\qquad
\text{s.t.}\;\; \hat{P}_t \approx \mathcal{W}^{\text{t}}(\hat{P}_{t-1}),
\label{eq:p3_s8_2_4_video}
\end{equation}
where $\mathcal{W}^{\text{t}}$ enforces frame-to-frame motion consistency on top of the per-frame 360-degree wrap-around continuity; Figure~\ref{fig:p3_s8_2_4_paradigms}(b) depicts this temporally-extended pathway.
The methods share this generator but differ in the conditioning argument $\mathbf{c}$ supplied to Eq.~\eqref{eq:p3_s8_2_4_video}, which separates them into three input modalities,
\begin{equation}
\mathbf{c}=
\begin{cases}
\mathbf{x} & \text{(text-to-video)},\\
P_0 & \text{(image-to-video)},\\
\{(I_t,\mathcal{M}_t)\}_{t=1}^{T} & \text{(perspective-to-video)},
\end{cases}
\label{eq:p3_s8_2_4_video_cond}
\end{equation}
where $\mathbf{x}$ is a text prompt, $P_0$ a static panorama anchor reused as the first frame, and $(I_t,\mathcal{M}_t)$ an NFoV frame paired with its visibility mask. The difficulty grows down the list: as the conditioning becomes spatially sparser, the model must hallucinate progressively more of the sphere while preserving the temporal coherence of Eq.~\eqref{eq:p3_s8_2_4_video}.

$\blacktriangleright$ \textit{Text-to-panoramic video generation.} 360DVD~\cite{wang2024360dvd} adapts a pretrained text-to-video model to the ERP domain through adapter modules, keeping wrap-around consistency across frames but struggling to balance text alignment against spherical geometry and to vary style under several control conditions.
TiP4GEN~\cite{xing2025tip4gen} adopts a dual-branch panoramic-and-perspective architecture in the spirit of PanFusion and lifts the generated video into a geometry-aligned 4D Gaussian scene.
PanoWan~\cite{xia2025panowan} and PanoDiT~\cite{zhang2025panodit} fine-tune large video diffusion transformers on panoramic video data and gain motion quality from scale, and PanoLora~\cite{dong2025panolora} treats the task as a projection adaptation from perspective views that LoRA can learn with little data.
Training-free alternatives carry the static-panorama machinery into time: SphereDiff~\cite{park2025spherediff} extends its spherical latent to video, and DynamicScaler~\cite{liu2025dynamicscaler} denoises through a seamless rotating window with an offset-shifting scheme and Global Motion Guidance, which yields arbitrary output resolution at constant memory.
Text alone, however, offers little control over layout and motion, which motivates the richer conditioning below.

$\blacktriangleright$ \textit{Image-to-panoramic video generation.} Animating a static panorama anchors the video spatially and asks for plausible omnidirectional motion that respects the input geometry.
4K4DGen~\cite{li20244k4dgen} animates user-selected regions with an image-to-video model under overlapping windows, at the cost of fixed-window artifacts and limited motion.
ViewPoint~\cite{fang2025panoramic} introduces a panorama representation that keeps global spatial continuity while Pano-Perspective attention transfers a perspective video prior, and HoloTime~\cite{zhou2025holotime} reconstructs the result as 4D Gaussian Splatting so that the animated scene can be explored under viewpoint changes.
GenEx~\cite{lu2024genex} generates a 3D-consistent imaginative environment from a single image as a panoramic video stream with loop consistency over long trajectories, which makes it usable for embodied exploration.
For explicit control, PanFlow~\cite{panflow} conditions on synthetic optical flow as a scene-agnostic motion signal and warps noise spherically so that camera and object motion can be controlled separately while the seam stays consistent; OmniDrag~\cite{li2026omnidrag} adds drag-style interaction at scene and object level through a spherical motion estimator; and PanoWorld-X~\cite{yin2025panoworld} reprojects ERP features onto the sphere inside a Sphere-Aware Diffusion Transformer to follow diverse camera trajectories.

$\blacktriangleright$ \textit{Perspective-to-panoramic video generation.} Expanding an NFoV video to a full 360-degree video requires coherent hallucination in space and time at once, under visibility masks that change with camera motion.
VidPanos~\cite{ma2024vidpanos} frames the task as space-time outpainting and uses a generative video model to fill the regions a panning camera never observed; Imagine360~\cite{tan2024imagine360} combines a global panoramic branch with local perspective priors and an elevation-aware design for the moving masks; and Beyond the Frame~\cite{beyondtheframe} adds geometry-aware modules and motion-aware conditioning to improve realism and motion consistency in the outpainted regions.


\paragraph{Panoramic 3D Reconstruction and Generation.}
Reconstruction bypasses the generator of Eq.~\eqref{eq:p3_s8_2_4_objective}: rather than generating $\hat{P}$, it takes \emph{observed} posed panoramas and applies the lifting operator directly,
\begin{equation}
\hat{\mathbf{r}}^{\text{rec}} = \Lambda\big(\{P_v, \pi_v\}_{v=1}^{V}\big),
\label{eq:p3_s8_2_4_recon}
\end{equation}
where $\Lambda$ is realized either as a per-scene optimization (iterative) or as an amortized network (feed-forward), mirroring the split of §\ref{subsec:p3_s8_2_3}; Figure~\ref{fig:p3_s8_2_4_paradigms}(c) depicts this reconstruction pathway.
Generation then composes the two directions, feeding a synthesized panorama into the same $\Lambda$.

$\blacktriangleright$ \textit{Iterative optimization-based panoramic reconstruction.} Realizing $\Lambda$ of Eq.~\eqref{eq:p3_s8_2_4_recon} as a \emph{per-scene} optimization,
\begin{equation}
\hat{\mathbf{r}}^{\text{rec}} = \arg\min_{\mathbf{r}} \sum_{v} \big\| \mathrm{Render}(\mathbf{r},\pi_v) - P_v \big\|^2 + \lambda\,\Omega(\mathbf{r}),
\label{eq:p3_s8_2_4_recon_iter}
\end{equation}
re-solved from scratch for each scene, these methods fit a radiance field or an explicit representation to multi-view panoramas, the panoramic analogue of §\ref{para:p3_s8_2_3_iterative}.
NeRF-based work includes 360Roam~\cite{huang2022360roam} for geometry-aware indoor reconstruction, 360FusionNeRF~\cite{kulkarni2023fusion} for novel views from a single panorama under geometric and semantic supervision, EgoNeRF~\cite{choi2023egonerf} with spherical grids for egocentric capture, and OmniNeRF~\cite{gu2022omninerf}, which processes raw dual-fisheye input and its lens distortion.
Gaussian-splatting alternatives are faster: 360GS~\cite{bai2024360gs} projects Gaussians onto spherical tangent planes, ODGS~\cite{lee2024odgs} rasterizes on tangent planes with elevation-aware densification, OP43DGS~\cite{wang2024op43dgs} corrects the projection bias of the local affine approximation across camera models, OmniGS~\cite{li2024omnigs} exploits omnidirectional optical properties, and Seam360GS~\cite{seam360gs} calibrates dual-fisheye input by jointly optimizing translation vectors and learnable angular-distortion embeddings.
Scene4U~\cite{scene4u} reconstructs from a single panorama in depth layers: an open-vocabulary segmenter and an LLM decompose it semantically, diffusion inpainting restores occluded regions, and 3D Gaussians are fitted per layer.

$\blacktriangleright$ \textit{Feed-forward generalization-based panoramic reconstruction.} Amortizing $\Lambda$ of Eq.~\eqref{eq:p3_s8_2_4_recon} into a trained network applied in a single forward pass,
\begin{equation}
\hat{\mathbf{r}}^{\text{rec}} = \Lambda_\psi\big(\{P_v, \pi_v\}_{v=1}^{V}\big),
\label{eq:p3_s8_2_4_recon_ff}
\end{equation}
where $\Lambda_\psi$ is trained once over a panoramic corpus (the analogue of §\ref{para:p3_s8_2_3_feedforward}), trades the per-scene cost of Eq.~\eqref{eq:p3_s8_2_4_recon_iter} for a learned prior that reconstructs from sparse panoramas without scene-specific training.
PanoGRF~\cite{chen2023panogrf} pioneered feed-forward reconstruction of wide-baseline panoramas with depth guidance; PanoSplatt3R~\cite{panosplatt3r} extends DUSt3R~\cite{wang2024dust3r} to unposed wide-baseline panoramas through RoPE rolling, a rotary positional embedding rolled across attention heads so that ERP periodicity and correspondences across the seam are modeled; and PanSplat~\cite{pansplat} builds spherical 3D Gaussian pyramids on a Fibonacci lattice to reach 4K novel views with feed-forward inference and cubemap rendering for memory efficiency.

$\blacktriangleright$ \textit{Panorama-based 3D generation and world modeling.} This family realizes the full composition of Eq.~\eqref{eq:p3_s8_2_4_objective}: it first \emph{generates} the panorama and then lifts it,
\begin{equation}
\hat{\mathbf{r}}^{\text{rec}} = \Lambda\big(\hat{P}\big),
\qquad
\hat{P} \sim p_\theta^{\text{pano}}\!\big(P \mid \kappa(\mathbf{c})\big),
\label{eq:p3_s8_2_4_gen}
\end{equation}
so that the synthesized omnidirectional image of Eq.~\eqref{eq:p3_s8_2_4_image} becomes the structured intermediate fed to the reconstruction operator of Eq.~\eqref{eq:p3_s8_2_4_recon}; Figure~\ref{fig:p3_s8_2_4_paradigms}(d) depicts this generate-then-lift pathway.
Text $\rightarrow$ panorama $\rightarrow$ 3D pipelines generate the omnidirectional image with a diffusion model~\cite{ye2024diffpano,kalischek2025cubediff} before reconstructing a NeRF, Gaussian, or mesh scene.
LayerPano3D~\cite{yang2024layerpano3d} parses the panorama into hierarchical depth layers and reconstructs each with 3D Gaussian Splatting, DreamScene360~\cite{zhou2025dreamscene360} fits panoramic Gaussians for free-roaming text-conditioned scenes, Matrix-3D~\cite{yang2025matrix} couples trajectory-guided panoramic video generation with a dual-pathway reconstruction that conditions on scene-mesh renders, and HunyuanWorld~\cite{team2025hunyuanworld} treats the panorama as a 360-degree world proxy for text- and image-conditioned explorable scenes with mesh export and disentangled object representations.
The panoramic video generators GenEx~\cite{lu2024genex} and PanoWorld-X~\cite{yin2025panoworld} (above) turn the same intermediate into an action- or trajectory-conditioned observation space, which is where panorama synthesis meets the world models of §\ref{sec:p3_s8_6}.


\subsubsection{Benchmarks}
\label{subsubsec:p3_s8_2_4_benchmarks}

Panoramic datasets must supply geometric consistency over the full sphere, a declared projection format, and, for reconstruction, camera poses and depth; Table~\ref{tab:p3_s8_2_4_corpora} lists the corpora in use, grouped by modality.

\paragraph{Training Corpora.}
For static synthesis, Structured3D~\cite{zheng2020structured3d} (synthetic indoor scenes with controlled geometry) and SUN360~\cite{xiao2012recognizing} (real-world indoor and outdoor panoramas) are the standard benchmarks for distortion handling and seam artifacts, in the splits adopted by JoPano, PanoDecouple, and DreamCube~\cite{feng2025jopano,zheng2025panorama,huang2025dreamcube}; Matterport3D~\cite{chang2017matterport3d}, with RGB-D panoramas, camera poses, and mesh geometry across 90 buildings, is the most widely used corpus for both generation and reconstruction~\cite{pansplat,yang2024layerpano3d,team2025hunyuanworld,tang2023MVDiffusion,zhang2024taming}, and Upright360~\cite{yang2024layerpano3d} filters Matterport3D, web, and Blockade Labs panoramas with GeoCalib~\cite{veicht2024geocalib} to near-zero pitch and roll for realistic indoor synthesis.
Panoramic video with geometric annotation has been the bottleneck.
360+x~\cite{chen2023360plusx} and WEB360~\cite{wang2024360dvd} are small real-world collections (the latter captioned for text conditioning), 360-1M~\cite{wallingford2024image} scales to over a million videos for world-model pretraining, PanFlow~\cite{panflow} adds optical flow and camera poses on motion-rich sequences, PANOVID~\cite{xia2025panowan} aggregates these with Imagine360~\cite{tan2024imagine360} and public VR videos but lacks poses and depth, and Matrix-Pano~\cite{yang2025matrix} fills that gap with Unreal Engine 5 renders carrying depth maps and camera trajectories; embodied work renders action-conditioned observations in Unreal Engine directly~\cite{lu2024genex,yin2025panoworld}.
For reconstruction, Replica~\cite{straub2019replica} and Residential~\cite{habtegebrial2022somsi} serve as synthetic indoor benchmarks at baselines of about 1.0\,m and 0.3\,m respectively~\cite{pansplat}, 360Loc~\cite{huang2024360loc} supplies posed panorama sequences used for 4K fine-tuning~\cite{pansplat}, and the scarcity of posed real panoramic video led PanSplat~\cite{pansplat} to capture its own dual-fisheye footage with poses from visual SLAM.

\begin{table}[t]
    \centering
    \fontsize{8}{10}\selectfont
    \setlength{\tabcolsep}{1.0mm}
    \caption{Representative training corpora for 3D panorama synthesis, grouped by primary modality. ``Modality'' indicates the data type and synthetic / real acquisition. Structured3D and SUN360 scales are those of the splits used by recent panorama generators (SUN360 in full has $67{,}583$ panoramas).}
    \label{tab:p3_s8_2_4_corpora}
    \begin{tabular}{L{0.17\textwidth} C{0.06\textwidth} C{0.15\textwidth} L{0.22\textwidth} L{0.32\textwidth}}
        \toprule
        \textbf{Dataset} & \textbf{Year} & \textbf{Scale} & \textbf{Modality} & \textbf{Primary use} \\
        \midrule
        Structured3D~\cite{zheng2020structured3d} & 2020 & ${\sim}21\mathrm{K}$ panoramas & Synthetic image & Indoor synthesis; distortion / seam benchmark \\
        SUN360~\cite{xiao2012recognizing} & 2012 & ${\sim}29\mathrm{K}$ panoramas & Real image & Real-scene generation benchmark \\
        Matterport3D~\cite{chang2017matterport3d} & 2017 & $10{,}800$ RGB-D panoramas & Real image + depth & Generation and reconstruction supervision \\
        Upright360~\cite{yang2024layerpano3d} & 2024 & ${\sim}9.4\mathrm{K}$ panoramas & Mixed image & Geometrically upright indoor synthesis \\
        \midrule
        360+x~\cite{chen2023360plusx} & 2024 & $2{,}152$ videos & Real multimodal video & Panoptic scene understanding \\
        WEB360~\cite{wang2024360dvd} & 2024 & ${\sim}2\mathrm{K}$ videos & Real captioned video & Text-to-panoramic video \\
        360-1M~\cite{wallingford2024image} & 2024 & $>1\mathrm{M}$ videos & Real video & Large-scale world-model pretraining \\
        PANOVID~\cite{xia2025panowan} & 2025 & Aggregated & Real video & Text-to-panoramic video training \\
        Matrix-Pano~\cite{yang2025matrix} & 2025 & $116{,}759$ sequences & Synthetic video + depth + trajectory & Trajectory-conditioned generation \\
        \midrule
        Replica~\cite{straub2019replica} & 2019 & $18$ scenes & Synthetic (rendered scans) & Indoor reconstruction benchmark \\
        Residential~\cite{habtegebrial2022somsi} & 2022 & Indoor scenes & Synthetic image & Indoor reconstruction benchmark \\
        360Loc~\cite{huang2024360loc} & 2024 & $18$ posed sequences & Real reconstruction & 4K reconstruction fine-tuning \\
        \bottomrule
    \end{tabular}
\end{table}

\paragraph{Evaluation Metrics.}
Panorama evaluation combines general image and video quality metrics with panorama-specific measures of distortion, seam consistency, and spherical topology.

$\blacktriangleright$ \textit{Image fidelity and semantic alignment.} FID measures distributional similarity between generated and real panoramas~\cite{feng2025jopano,yang2024layerpano3d,zheng2025panorama,tandit2025}; PanFusion~\cite{zhang2024taming} shows that FAED (Fréchet Auto-Encoder Distance), computed in the latent space of an auto-encoder trained on the target data, tracks panorama quality more closely than Inception features do.
Text-to-panorama generation reports CLIP-Score between prompt and output~\cite{team2025hunyuanworld,feng2025jopano,yang2024layerpano3d,tandit2025}, image-conditioned tasks report CLIP-I between input and output~\cite{team2025hunyuanworld}, and no-reference quality is scored with BRISQUE and NIQE~\cite{team2025hunyuanworld,yang2024layerpano3d,yang2025matrix}, complemented by aesthetic scorers and human preference studies, for instance the Average User Ranking protocol of LayerPano3D~\cite{yang2024layerpano3d}.
Layout-conditioned synthesis measures adherence by 3D and 2D IoU against layouts recovered with HorizonNet~\cite{sun2019horizonnet}, following PanFusion~\cite{zhang2024taming}.

$\blacktriangleright$ \textit{Panorama-specific consistency.} Projection-induced distortion and boundary discontinuities call for metrics of their own.
JoPano~\cite{feng2025jopano} scores cubemap seams with Seam-SSIM (SSIM between adjacent edge pixels of neighboring faces) and Seam-Sobel (gradient discontinuity along the seam); PanoDecouple~\cite{zheng2025panorama} introduces Distort-CLIP, a CLIP variant trained to detect panoramic distortion, so that the ``visual cheating'' of models that improve image quality while degrading geometry becomes measurable; and loop consistency across the 360-degree wrap-around is assessed by PSNR in overlapping regions~\cite{tang2023MVDiffusion} or across the left--right boundary~\cite{panflow,xia2025panowan}.

$\blacktriangleright$ \textit{Reconstruction and novel-view synthesis.} Because pixel density on the sphere is non-uniform, WS-PSNR (Weighted Spherical PSNR) weights pixels by latitude and replaces plain PSNR for equirectangular outputs~\cite{pansplat,chen2023panogrf}, complemented by SSIM and LPIPS~\cite{yang2024layerpano3d,pansplat,yang2025matrix}.

$\blacktriangleright$ \textit{Video quality and spatiotemporal dynamics.} Panoramic videos are often projected to several perspective views (\eg, four views at azimuths $\{0\degree, 90\degree, 180\degree, 270\degree\}$ with a $90\degree$ field of view) so that human raters and perspective-domain tools can be used~\cite{team2025hunyuanworld,yang2024layerpano3d,tan2024imagine360}.
Q-Align~\cite{wu2023qalign}, a large multimodal model trained to score visual quality on discrete text-defined levels, is increasingly used for images and videos alike~\cite{team2025hunyuanworld,tan2024imagine360}; FVD~\cite{unterthiner2018fvd}, computed in the feature space of an inflated 3D ConvNet, is the standard distributional metric for video~\cite{panflow,xia2025panowan}; VBench~\cite{huang2023vbench} adds imaging quality, aesthetic quality, motion smoothness, and subject consistency over projected views; and trajectory-conditioned generation is scored by frame-to-frame consistency and by the range of camera motion it sustains~\cite{yin2025panoworld}.
Two limits remain. Ground truth is scarce for dynamic panoramic video with depth and camera annotation, so evaluation falls back on projected perspective views or user studies; and world-level properties such as loop consistency over long trajectories and action-conditioned spatial coherence have no dedicated benchmark yet.


\subsubsection{Challenges and Trends}
\label{subsubsec:p3_s8_2_4_challenges}

\paragraph{Unique Challenges in Panoramic Synthesis.}
Modeling a complete omnidirectional environment imposes constraints that NFoV synthesis never meets.
The first is \textbf{distortion-aware learning}: ERP stretches content near the poles, so visual information is non-uniformly distributed over the image and convolutional kernels designed for translation invariance on the plane struggle with the varying distortion, which motivates spherical manifold convolutions and latitude-aware architectures~\cite{sun2025spherical,zheng2025panorama}.
The second is \textbf{seamless boundary consistency}: wrap-around coherence across the longitudinal boundary must be preserved to avoid visible seams, typically through circular padding or latent rotation during denoising~\cite{panflow,xia2025panowan}.
The third is \textbf{scaling structured generation} to 4K and beyond while keeping global geometric and semantic consistency over the whole sphere, which pushes toward memory-efficient attention and hierarchical spherical representations~\cite{pansplat,liu2025dynamicscaler}.

\paragraph{From Panorama to World Models.}
The clearest trend is the move from static spherical snapshots toward \textbf{interactive 3D world modeling}, in which the panorama is the interface through which an embodied agent reasons about a complete environment rather than a fragment of one.
The systems that lift a generated panorama into an explorable world (HunyuanWorld~\cite{team2025hunyuanworld}, Matrix-3D~\cite{yang2025matrix}) and World Labs' Marble~\cite{Feifei2025Marble} show what panoramic priors make possible, and \textbf{temporal panoramic modeling} supplies the 4D half: generating temporally coherent panoramic video requires motion dynamics that respect spherical topology and occlusion boundaries~\cite{panflow,yin2025panoworld}.
Integrated with interactive systems, such panoramic world simulations become a basis for embodied interaction, long-range spatial reasoning, and environment simulation in robotics and immersive media.

\paragraph{Representation Innovations.}
Efficiency and fidelity both hinge on the omnidirectional representation.
\textbf{Efficient omnidirectional representations} such as the Fibonacci-lattice Gaussian pyramids of PanSplat~\cite{pansplat} render at high fidelity without the redundancy of uniform grids, and \textbf{multi-modal panoramic fusion}, which binds 360-degree RGB to depth, semantic labels, and geometric proxies such as surface normals~\cite{yang2024layerpano3d,scene4u}, makes the synthesized environment both more realistic and more editable, closing part of the gap between 2D generative priors and full 3D scene intelligence.




\subsection{3D Scene Editing}
\label{subsec:p3_s8_2_5}

The preceding subsections build 3D content; 3D scene editing takes an existing asset and modifies it under user control while preserving everything the edit should leave untouched.
Such incremental control is exactly what spatial intelligence asks of a 3D world model: an agent that must correct a mistake, personalize an environment, or keep a scene consistent as its task unfolds cannot afford to regenerate the world from scratch every time it changes.
Technically, 3D editing lifts the 2D inpainting-and-editing machinery of §\ref{subsec:p3_s8_1_3} from the image plane to a multi-view-consistent 3D representation, adding a cross-view coherence requirement that has no analogue in the 2D case; the five method families below are distinguished by how the edit signal is injected into that representation.


\begin{figure}[!ht]
    \centering
    \includegraphics[width=0.99\linewidth]{figures/part3/section8/3d_editing_tasks.png}
    \caption{Demonstrations of 3D editing tasks. (a) appearance editing, (b) object transformation, (c) shape deformation, (d) scene inpainting, and (e) creative editing. Figure adapted from~\cite{zhu2026survey3D}.}
    \label{fig:3D_editing_tasks}
\end{figure}

\subsubsection{Task Definition}
\label{subsubsec:p3_s8_2_5_task}

3D scene (and object) editing modifies an existing 3D asset or reconstructed scene according to a user intent (typically language or reference images, as in Figure~\ref{fig:3D_editing_tasks}) while preserving 3D consistency for novel-view rendering.
Unlike text-to-3D generation, editing operates on a given representation (\eg, NeRF, 3D Gaussian Splatting, or meshes) and requires (i) locality and identity preservation for unedited regions, (ii) multi-view (and potentially temporal) coherence, and (iii) controllability over the edited attributes, whether geometry, appearance, material, or semantics.

Formally, 3D editing is the 3D-lifted counterpart of the 2D inpainting-and-editing operator of §\ref{subsec:p3_s8_1_3} (Eq.~\eqref{eq:t2i_edit_objective}).
Given an input 3D representation $\mathbf{r}^{\text{in}}\in\mathcal{R}$ drawn from the same representation space $\mathcal{R}$ as §\ref{subsec:p3_s8_2_1}, an edit intent $\mathbf{c}\in\mathcal{C}\supseteq\{\text{text instruction},\text{reference image},\text{mask / region},\text{scribble}\}$, and an optional 3D region of interest $\mathcal{M}\subseteq\mathbb{R}^3$, the goal is an editing operator
\begin{equation}
E_\theta:\mathcal{R}\times\mathcal{C}'\rightarrow\mathcal{R},
\qquad
\hat{\mathbf{r}} = E_\theta\big(\mathbf{r}^{\text{in}},\,\kappa(\mathbf{c})\big),
\label{eq:p3_s8_2_5_objective}
\end{equation}
trained or optimized to minimize an objective that couples three competing terms,
\begin{equation}
\mathcal{L}_{\text{edit}}
=
\underbrace{\mathcal{L}_{\text{intent}}\big(\hat{\mathbf{r}},\mathbf{c};\mathcal{M}\big)}_{\text{edit-region adherence}}
\;+\;
\lambda_{\text{p}}\underbrace{\mathcal{L}_{\text{preserve}}\big(\hat{\mathbf{r}},\mathbf{r}^{\text{in}};\mathbb{R}^3\!\setminus\!\mathcal{M}\big)}_{\text{identity preservation}}
\;+\;
\lambda_{\text{c}}\underbrace{\mathcal{L}_{\text{consist}}\big(\{\mathrm{Render}(\hat{\mathbf{r}},\pi_v)\}_v\big)}_{\text{multi-view coherence}},
\label{eq:p3_s8_2_5_loss}
\end{equation}
where $\kappa$ is the condition encoder of §\ref{subsec:p3_s8_2_1}, $\mathrm{Render}$ is the differentiable renderer of §\ref{subsec:p3_s8_2_3}, and $\lambda_{\text{p}},\lambda_{\text{c}}$ trade off preservation and consistency against edit strength.
The limiting case that anchors Eq.~\eqref{eq:p3_s8_2_5_objective} to its 2D origin is instructive: when $\mathcal{R}$ collapses to a single image ($V=1$, no rendering), $E_\theta$ reduces exactly to the 2D editing operator of Eq.~\eqref{eq:t2i_edit_objective} and the multi-view-consistency term vanishes, so 3D editing is 2D editing plus a cross-view coherence constraint.


\subsubsection{Methodology}
\label{subsubsec:p3_s8_2_5_methodology}

Recent 3D object and scene editing methods can be categorized by \emph{how} the edit signal of Eq.~\eqref{eq:p3_s8_2_5_objective} is injected into the 3D representation, from lifting 2D edits view by view, through enforcing cross-view consistency, to operating directly on 3D primitives, reasoning with language models, and amortizing the edit into a single feed-forward pass.
Figure~\ref{fig:p3_s8_2_5_mechanisms} schematizes the five families and how each realizes the operator $E_\theta$.

\begin{figure}[!ht]
    \centering
    \includegraphics[width=0.99\linewidth]{figures/part3/section8/Five_method_families_for_3D_scene_editing.pdf}
    \caption{Five method families for 3D scene editing: (a) 2D-lifted iterative editing, (b) multi-view-consistent guidance, (c) 3D-native primitive editing, (d) LLM-integrated editing, and (e) efficient and training-free editing. The chair edit (top) is illustrative, not a measured result.}
    \label{fig:p3_s8_2_5_mechanisms}
\end{figure}

\paragraph{2D-lifted Iterative Editing.}
A dominant paradigm inherits from instruction-based NeRF editing~\cite{bao2023sine,khalid2024latenteditor,mikaeili2023sked,zhou2023repaintnerf}: the method iterates between rendering views from the current scene, applying a 2D diffusion-based editor to them, and updating the 3D representation to match the edited guidance under geometric regularization.
Formally, this family instantiates $E_\theta$ of Eq.~\eqref{eq:p3_s8_2_5_objective} as a \emph{per-scene} optimization that drives the rendered views toward the output of a \emph{frozen} 2D editor $\mathcal{E}_{\text{2D}}$ (the InstructPix2Pix-style operator of §\ref{subsec:p3_s8_1_3}),
\begin{equation}
\hat{\mathbf{r}} = \arg\min_{\mathbf{r}\in\mathcal{R}} \sum_{v} \big\| \mathrm{Render}(\mathbf{r},\pi_v) - \mathcal{E}_{\text{2D}}\!\big(\mathrm{Render}(\mathbf{r}^{\text{in}},\pi_v),\,\mathbf{c}\big) \big\|^2 + \lambda\,\Omega(\mathbf{r}),
\label{eq:p3_s8_2_5_lift}
\end{equation}
so that the 2D edits are \emph{lifted} into the 3D representation and the $\arg\min$ is re-solved from scratch for every scene; Figure~\ref{fig:p3_s8_2_5_mechanisms}(a) depicts this render-edit-update loop.
Instruct-NeRF2NeRF~\cite{haque2023instructnerf2nerf}, EditNeRF~\cite{liu2021editingnerf}, Free-Editor~\cite{karim2024freeeditor}, and CTRL-D~\cite{he2025ctrld} lift instruction-driven 2D edits into the radiance field this way, and with the rise of 3D Gaussian Splatting (3DGS)~\cite{kerbl20233d} the same principle moved to Gaussian primitives: Instruct-GS2GS~\cite{vachha2024instructgs2gs} iteratively edits the training images with InstructPix2Pix~\cite{brooks2023instructpix2pixlearningfollowimage} and re-optimizes the 3DGS to reflect the instruction while preserving 3D structure~\cite{wang2024viewedit,wynn2025morpheus}.
The recipe is simple and borrows the full strength of 2D editors, but it is slow, prone to view-dependent artifacts, and only as good as the stability of the 2D editor and the coverage of the rendered guidance views.

\paragraph{Multi-view-consistent Guidance.}
Because independently edited views can disagree (inconsistent textures, hallucinated geometry), recent work~\cite{chen2024dge} enforces multi-view consistency at the \emph{guidance} level rather than relying on the 3D update to reconcile conflicting edits.
Formally, rather than editing each view independently as in Eq.~\eqref{eq:p3_s8_2_5_lift}, this family makes the guidance images mutually consistent before the lift,
\begin{equation}
\{\hat{I}_v\} = \arg\min_{\{I_v\}} \sum_{v} \mathcal{L}_{\text{intent}}(I_v,\mathbf{c}) \;+\; \beta\!\!\sum_{(u,v)} \big\| I_u - \mathcal{W}_{u\to v}(I_v) \big\|^2,
\label{eq:p3_s8_2_5_mvc}
\end{equation}
where $\mathcal{W}_{u\to v}$ warps view $v$ into view $u$ using geometry or depth derived from the current 3D estimate, penalizing cross-view disagreement; the consolidated $\{\hat{I}_v\}$ then drive the 3D update of Eq.~\eqref{eq:p3_s8_2_5_lift}, as in Figure~\ref{fig:p3_s8_2_5_mechanisms}(b).
GaussCtrl~\cite{wu2024gaussctrl} edits several rendered views in a coupled fashion rather than sequentially, improving multi-view consistency for text-driven 3DGS editing~\cite{wang2024viewedit,xia2025towardsConsistent}, and VcEdit~\cite{wang2024viewedit} integrates the 3DGS into the image-editing process itself with consistency modules that counter cross-view disagreement and mode collapse.
A complementary line transfers a \emph{single-view} edit to other views via geometry: \citet{gomel2024diffusionBasedAttention} extracts diffusion attention features from a reference edit, warps them across views using depth from the Gaussian splats, and injects the warped features to propagate a coherent edit.
These methods preserve 3D coherence more robustly but still struggle with large non-rigid shape changes or severe occlusion, where reliable correspondence and geometry are hard to come by.

\paragraph{3D-native Primitive Editing.}
To reduce reliance on unstable 2D guidance and enable interactive use, 3D-native approaches operate directly on explicit scene primitives.
Formally, this family bypasses the rendering loop of Eq.~\eqref{eq:p3_s8_2_5_lift} entirely: a grounding or segmentation operator $\mathcal{G}$ selects the target subset of primitives and a direct geometric transform $\mathcal{T}$ applies the edit,
\begin{equation}
\hat{\mathbf{r}} = \mathcal{T}\!\big(\mathbf{r}^{\text{in}};\,\mathcal{G}(\mathbf{r}^{\text{in}},\mathbf{c})\big), \qquad \mathcal{T}\in\{\text{add},\,\text{delete},\,\text{move},\,\text{recolor},\,\text{replace}\},
\label{eq:p3_s8_2_5_native}
\end{equation}
acting on explicit Gaussians or mesh facets so that the unedited primitives are preserved \emph{by construction}: the identity-preservation term of Eq.~\eqref{eq:p3_s8_2_5_loss} is satisfied structurally rather than by a penalty; Figure~\ref{fig:p3_s8_2_5_mechanisms}(c) shows this pathway.
3DSceneEditor~\cite{yan20243dsceneeditor} runs a fully 3D pipeline, combining instance segmentation and zero-shot grounding to localize the target and then performing Gaussian-level addition, deletion, relocation, recoloring, and replacement with strong controllability in real time, and InterGSEdit~\cite{wen2025intergsedit} adds user-in-the-loop key-view selection and a geometry-consistent attention prior for iterative, preference-guided 3DGS editing with control over the edit degree.
More broadly, GaussianEditor and 3DSceneEditor~\cite{chen2024gaussianeditor,yan20243dsceneeditor} impose higher-level structure on the Gaussians, grouping them by object identity or semantic coherence (typically through zero-shot grounding or segmentation~\cite{hu2024semanticanythinggaussians,cen2025segmentany3dgaussians} followed by anchoring) so that an edit can be formulated at the level of an instance rather than of individual primitives, which localizes the target more precisely and protects the background.

\paragraph{LLM-integrated Editing.}
Large language models have been integrated into 3D editing pipelines not merely as text encoders but as high-level reasoning and control modules.
Formally, the LLM decomposes the instruction into an ordered program of geometry-aware operators applied to the asset,
\begin{equation}
\hat{\mathbf{r}} = (o_K \circ \cdots \circ o_1)(\mathbf{r}^{\text{in}}), \qquad (o_1,\dots,o_K) = \mathrm{LLM}\big(\mathbf{c},\,\mathbf{r}^{\text{in}}\big),
\label{eq:p3_s8_2_5_llm}
\end{equation}
so that high-level reasoning (attribute inference, edit planning, dependency resolution) is separated from the low-level operators $o_k$, often the primitive transforms of Eq.~\eqref{eq:p3_s8_2_5_native}, that it dispatches; Figure~\ref{fig:p3_s8_2_5_mechanisms}(d) depicts this plan-then-execute pathway.
ShapeLLM-Omni~\cite{ye2025shapellmOmni} casts 3D editing as language-grounded reasoning in which the LLM performs shape understanding, attribute inference, and edit planning before invoking geometry-aware operators, and LL3M~\cite{lu2025ll3m} extends this to multimodal 3D manipulation requiring compositional understanding of geometry, appearance, and spatial relations.
A complementary line~\cite{ma2025feedforward3DEditing} uses the LLM to guide or parameterize image-to-3D or text-to-3D backbones in a feed-forward setting rather than a per-scene optimization, and LLMs also serve as coordination agents for global consistency~\cite{huang2024blenderalchemy,zhang2025scenelanguage}: \citet{xia2025towardsConsistent} shows that reasoning over object identities, relations, and edit dependencies reduces the conflicts that arise in purely diffusion- or optimization-based pipelines when several objects and attributes must be edited at once.

\paragraph{Efficient and Training-free Editing.}
A final trend reduces test-time cost through feed-forward reconstruction or reuse of pretrained models~\cite{erkocc2025preditor3d,li2025voxhammer}. These strategies are distinct: a training-free editing pipeline can still involve iterative diffusion inference. The amortized feed-forward case evaluates the editing operator of Eq.~\eqref{eq:p3_s8_2_5_objective} directly,
\begin{equation}
\hat{\mathbf{r}} = E_\theta\big(\mathbf{r}^{\text{in}},\,\kappa(\mathbf{c})\big) \quad \text{(amortized, no per-scene optimization)},
\label{eq:p3_s8_2_5_ff}
\end{equation}
eliminating the per-scene $\arg\min$ of Eq.~\eqref{eq:p3_s8_2_5_lift} by amortizing inference into a trained network. Training-free alternatives instead reuse pretrained components, potentially through single-view guidance, mesh proxies, or iterative native-3D latent editing; Figure~\ref{fig:p3_s8_2_5_mechanisms}(e) distinguishes these routes from the per-scene optimization loop in panel~(a).
Free-Editor~\cite{karim2024freeeditor} edits a NeRF zero-shot from text without scene-specific retraining, countering multi-view style inconsistency with a ``single-view editing'' scheme and an Edit Transformer; on the scene-operation side (insert, replace, delete), FreeEdit~\cite{madhavaram2025towardsTrainingFree} uses a mesh proxy to perform training-free edits guided by foundation models at near real-time rates in complex indoor scenes.
These methods improve usability but typically trade away edit fidelity and coverage, especially for fine-grained geometry, and depend strongly on robust grounding and localization in 3D.


\subsubsection{Benchmarks}
\label{subsubsec:p3_s8_2_5_benchmarks}

Several benchmarks have recently been introduced for systematic evaluation.
3D-Alpaca-Editing~\cite{ye2025shapellmOmni} focuses on instruction-following 3D editing, assessing whether a model interprets natural-language commands correctly while preserving unedited content.
CMD~\cite{li2025cmd} emphasizes compositional and multi-step instructions, where reasoning and planning matter.
Edit3D-Bench~\cite{li2025voxhammer} is a general-purpose suite covering geometry, appearance, and object-level operations, with 3D edit regions annotated on 100 source models (300 edits), which allows fair comparison across editing paradigms.
3DEditVerse~\cite{xia2025towardsConsistent} is a large-scale paired benchmark built for robustness and generalization across varied objects, scenes, and intents.
Together they standardize evaluation in terms of instruction compliance, edit locality, and scalability; Table~\ref{tab:p3_s8_2_5_datasets} summarizes their scale and the evaluation axes each one annotates.

\begin{table}[!ht]
\centering
\fontsize{8}{10}\selectfont
\setlength{\tabcolsep}{1.0mm}
\caption{Representative 3D editing benchmarks. ``Edit region'' indicates whether a localized edit mask / region is annotated; ``Consistency'' and ``Harmony'' indicate whether multi-view consistency and edit harmonization are evaluated. The criteria follow the four requirements set out with 3DEditVerse~\cite{xia2025towardsConsistent}.}
\label{tab:p3_s8_2_5_datasets}
\begin{tabular}{L{0.24\textwidth} C{0.11\textwidth} C{0.10\textwidth} C{0.11\textwidth} C{0.12\textwidth} C{0.12\textwidth} C{0.10\textwidth}}
\toprule
\textbf{Dataset} & \textbf{Training size} & \textbf{Testing size} & \textbf{Edit region} & \textbf{Train / test split} & \textbf{Consistency} & \textbf{Harmony} \\
\midrule
3D-Alpaca-Editing~\cite{ye2025shapellmOmni} & ${\sim}70{,}000$ & -- & \cmark & \crossmark & \crossmark & \cmark \\
CMD~\cite{li2025cmd} & ${\sim}40{,}000$ & 50 & \crossmark & \cmark & \cmark & \crossmark \\
Edit3D-Bench~\cite{li2025voxhammer} & -- & 300 & \cmark & \crossmark & \cmark & \cmark \\
3DEditVerse~\cite{xia2025towardsConsistent} & 116,309 & 1,500 & \cmark & \cmark & \cmark & \cmark \\
\bottomrule
\end{tabular}
\end{table}

\paragraph{Evaluation Metrics.}
Evaluation considers geometric fidelity, appearance quality, and semantic alignment with the instruction.
Geometry is measured by Chamfer Distance (CD)~\cite{ma2025feedforward3DEditing}, lower being better, and by F-score~\cite{xia2025towardsConsistent}, which captures surface overlap at a distance threshold.
Texture and appearance are assessed by comparing renderings of the edited representation against target or inpainted images with SSIM~\cite{wang2004image}, LPIPS~\cite{zhang2018unreasonable}, and DreamSim~\cite{fu2023dreamsim}~\cite{gao20253dmeshedit}.
Visual--text alignment is scored with CLIP- or DINO-based embedding similarity~\cite{shu2025gaussedit}, and GPTEval3D~\cite{wu2024gpt4vt23dgen} uses GPT-4V to compare multi-view renderings of edited shapes on text alignment, 3D plausibility, and texture detail.
To measure the edit relative to the input rather than in isolation, directional CLIP (CLIP$_{\text{dir}}$)~\cite{wynn2025morpheus} checks that the change in image embedding, computed over multi-view renderings of input and output, follows the change in text embedding.

Edit cost is the clearest axis separating the five families.
Instruct-NeRF2NeRF~\cite{haque2023instructnerf2nerf} takes on the order of 30--90 minutes per edit; GaussianEditor~\cite{chen2024gaussianeditor} brings this under 20 minutes by operating on explicit Gaussians; PrEditor3D~\cite{erkocc2025preditor3d} edits a shape in a few minutes, and VoxHammer~\cite{li2025voxhammer} performs precise local edits training-free in a native 3D latent space.
On quality the field reports DINO-I and CLIP-T on Edit3D-Bench~\cite{li2025voxhammer}, directional CLIP~\cite{wynn2025morpheus}, and GPT-4V judgment~\cite{wu2024gpt4vt23dgen}; the 3D-native and feed-forward families lead on locality and speed, while 2D-lifted methods keep an edge on open-vocabulary appearance edits where a strong 2D editor does the work.
The speed of feed-forward editing comes at a cost in fine-grained geometry control and a dependence on robust 3D grounding, taken up below.


\subsubsection{Challenges and Trends}
\label{subsubsec:p3_s8_2_5_challenges}

Several problems remain genuinely open and define the current frontier of 3D editing.

\paragraph{Geometry-level, not merely appearance, editing.}
Most current methods edit color and texture well but treat geometry as fixed: precise, instruction-driven changes to shape, including thin, transparent, or high-frequency structures, remain unreliable, and edits that should alter the underlying surface often only repaint it. Closing this gap requires editors that operate on geometry and appearance jointly rather than deferring geometry to a frozen reconstruction.

\paragraph{Large non-rigid and structural edits.}
Adding, removing, or substantially deforming structure while preserving the identity of the rest of the asset is still difficult: drag-driven, as-rigid-as-possible Gaussian editors~\cite{han2025arapgs} handle moderate deformation, but large non-rigid motion and topology changes frequently lose original detail. Robust structural and compositional editing, including relation-aware multi-object edits, remains an active frontier.

\paragraph{Material--lighting disentanglement and relightable editing.}
Because most representations bake illumination into appearance, an edit that changes an object's material or color is not guaranteed to stay consistent when the scene is relit. Relightable representations that decompose geometry, BRDF, and incident light~\cite{gao2024relightable} are a prerequisite for edits that survive relighting and for inserting edited content into new lighting, yet stable material--light decomposition under editing is not solved.

\paragraph{The consistency bottleneck has moved to 3D grounding.}
Now that native-3D editing handles appearance coherence, the limiting factor is localization: methods such as VoxHammer~\cite{li2025voxhammer} are highly sensitive to 3D mask accuracy and degrade rapidly when the target region is imprecisely grounded. Robust open-vocabulary 3D grounding and segmentation~\cite{hu2024semanticanythinggaussians,cen2025segmentany3dgaussians}, ideally without user-provided masks~\cite{xia2025towardsConsistent}, is therefore the new gate on edit quality, and it is most acute in dynamic scenes, where the edited region must be tracked with temporal identity preserved across frames~\cite{he2025ctrld}.

Two pressures cut across all four. Evaluation remains immature: the benchmarks of §\ref{subsubsec:p3_s8_2_5_benchmarks} are small (Edit3D-Bench annotates 100 source models and 300 edits) and no metric panel is yet standard, so progress on the harder edits above is difficult to measure reliably. And the strongest forward trend is unification, bringing generation, segmentation, grounding, and editing into a single structured pipeline~\cite{park2026generative} with interfaces that go beyond text to reference images~\cite{shu2025gaussedit}, masks~\cite{barda2025instant3dit}, sketches~\cite{mikaeili2023sked}, and object-level operations, so that an agent can specify and verify a geometry-, material-, and relighting-consistent edit rather than only a repaint.



% \subsection{Summary}
% \label{subsec:p3_s8_2_summary}

% Table~\ref{tab:p3_s8_2_subfield_comparison} places the five subfields side by side: the input each consumes, the operator it instantiates over the shared 3D-representation space $\mathcal{R}$, the paradigms it surveys, the methods anchoring each, and the benchmarks on which it reports.

% \begin{table}[t]
%     \centering
%     \fontsize{8}{10}\selectfont
%     \setlength{\tabcolsep}{1.0mm}
%     \caption{The five subfields of 3D spatial reconstruction and generation; each operator acts on the shared 3D-representation space $\mathcal{R}$.}
%     \label{tab:p3_s8_2_subfield_comparison}
%     \begin{tabular}{L{0.13\textwidth} L{0.20\textwidth} L{0.18\textwidth} L{0.19\textwidth} L{0.22\textwidth}}
%         \toprule
%         \textbf{Subfield} & \textbf{Input / operator} & \textbf{Dominant paradigms} & \textbf{Canonical methods} & \textbf{Benchmarks / metrics} \\
%         \midrule
%         §\ref{subsec:p3_s8_2_1} 3D object synthesis & Text / image / multi-view $\to$ asset; $G_\theta^{\text{3D}}:\mathcal{C}'\!\to\!\mathcal{R}$ (Eq.~\eqref{eq:p3_s8_2_1_objective}) & 3D-native generation; SDS distillation; multi-view generation and lifting; feed-forward reconstruction and 3D latent diffusion & GET3D; DreamFusion; Zero-1-to-3, MVDream; LRM, TRELLIS, Hunyuan3D & GT23D-Bench, MATE-3D, Gen3DEval; CLIP alignment, Chamfer / F-score, LPIPS \\
%         §\ref{subsec:p3_s8_2_2} 3D scene synthesis & Text / layout / graph / sparse views / trajectory $\to$ scene tuple; $G_\theta^{\text{scene}}$ (Eq.~\eqref{eq:p3_s8_2_2_objective}) & Layout / scene-graph conditioning; view-to-3D; view-to-view; video-model-based & DiffuScene, GALA3D; DreamScene; LVSM, SEVA; CameraCtrl, ViewCrafter & 3D-FRONT, ScanNet, RealEstate10K; layout adherence, PSNR / LPIPS, FVD, pose error \\
%         §\ref{subsec:p3_s8_2_3} 3D scene reconstruction & Posed or unposed images, depth, point clouds $\to$ 3D; $G_\theta^{\text{rec}}:\mathcal{O}\!\to\!\mathcal{R}^{\text{rec}}$ (Eq.~\eqref{eq:p3_s8_2_3_objective}) & Iterative (SfM--MVS, NeRF, 3DGS); feed-forward; generative complete reconstruction & NeRF, 3DGS; DUSt3R, MASt3R, VGGT; ReconFusion, Bolt3D & DTU, Tanks-and-Temples, Mip-NeRF360; PSNR / SSIM / LPIPS, Chamfer, pose AUC \\
%         §\ref{subsec:p3_s8_2_4} 3D panorama synthesis & Text / NFoV image / layout / viewpoint $\to$ panorama $\to$ 3D; $G_\theta^{\text{pano}}$ $+$ lift $\Lambda$ (Eq.~\eqref{eq:p3_s8_2_4_objective}) & Panoramic image synthesis; panoramic video; panoramic reconstruction; panorama-based scene generation & PanFusion, CubeDiff; 360DVD, Imagine360; PanSplat; LayerPano3D, HunyuanWorld & Structured3D, SUN360, Matterport3D; FID / FAED, WS-PSNR, seam and loop consistency \\
%         §\ref{subsec:p3_s8_2_5} 3D scene editing & Existing asset $+$ edit intent $\to$ edited asset; $E_\theta:\mathcal{R}\!\times\!\mathcal{C}'\!\to\!\mathcal{R}$ (Eq.~\eqref{eq:p3_s8_2_5_objective}) & 2D-lifted iterative; multi-view-consistent guidance; 3D-native primitive; LLM-integrated; feed-forward / training-free & Instruct-NeRF2NeRF; GaussCtrl, VcEdit; 3DSceneEditor; ShapeLLM-Omni; VoxHammer & Edit3D-Bench, 3DEditVerse; Chamfer / F-score, CLIP$_{\text{dir}}$, GPTEval3D \\
%         \bottomrule
%     \end{tabular}
% \end{table}

% The subfields partition not by output domain, since all five produce an element of the same $\mathcal{R}$, but by the direction and source of the map: synthesis (§\ref{subsec:p3_s8_2_1}, §\ref{subsec:p3_s8_2_2}) maps high-level priors into $\mathcal{R}$, reconstruction (§\ref{subsec:p3_s8_2_3}) maps sensor observations into it, panorama synthesis (§\ref{subsec:p3_s8_2_4}) inserts a 2D intermediate before the map, and editing (§\ref{subsec:p3_s8_2_5}) maps $\mathcal{R}$ back into itself.
% Synthesis and reconstruction are not disjoint. As observations grow sparse, the reconstruction operator of Eq.~\eqref{eq:p3_s8_2_3_objective} must hallucinate the unobserved regions and collapses into the generative posterior of Eq.~\eqref{eq:p3_s8_2_1_objective}, so the complete-reconstruction regime of §\ref{subsec:p3_s8_2_3} and the generative regime of §\ref{subsec:p3_s8_2_1} approach one problem from opposite sides, and the formal distinction between the two operators dissolves into a single estimator that consumes whatever observations are available and falls back on a learned prior wherever they are not.
% Across all five subfields the 2D prior of §\ref{sec:p3_s8_1} is reused as a prior to be distilled rather than as a terminal generator (score distillation, multi-view and panorama diffusion, frozen 2D editors lifted into 3D), and all inherit the same failure mode, the 2D$\leftrightarrow$3D consistency gap (Janus artifacts, view drift, cubemap seams, edit flicker) whose suppression motivates the multi-view-consistency machinery surveyed throughout the section.
% Evaluation, finally, has not followed the convergence seen in §\ref{sec:p3_s8_1}: no cross-subfield suite exists, each subfield reports a paradigm-specific tuple over geometry, appearance, and constraint adherence, and a layout-adherence score from §\ref{subsec:p3_s8_2_2} cannot be ranked against a novel-view PSNR from §\ref{subsec:p3_s8_2_3}; a human-aligned protocol that spans the subfields is the clearest open need.

% The temporal extension of these volumetric priors (generating spatially \emph{and} temporally coherent 3D content, where a camera or the scene itself moves through time) is the subject of video-based spatial generation in §\ref{sec:p3_s8_3}.

% \newpage
\section{Video Spatial Generation}
\label{sec:p3_s8_3}

\FramedInitial{V}{ideo} spatial generation is the task of synthesizing a temporally coherent image sequence (a video clip) from a high-level prompt together with optional spatial, temporal, and other conditioning signals.
It has advanced substantially in recent years, evolving from early generative adversarial network (GAN)-based approaches~\cite{goodfellow2014gan} through autoregressive token models to the diffusion- and diffusion-transformer-based systems that now dominate the field~\cite{polyak2024movie,hong2022cogvideo,wan2025wan,sora2024}.
This rapid progress has stimulated a wide range of downstream tasks, extending beyond basic text-to-video generation to encompass fine-grained control mechanisms (spatial, \eg, image, pose, depth, landmark, sketch, and bounding box; temporal, \eg, optical flow, trajectory, camera, and motion; and other conditional controls, \eg, style, audio, bird's-eye view) as well as sophisticated video editing capabilities.

Within the arc of §\ref{sec:generation}, video spatial generation is the \emph{temporal} extension of the priors built in the preceding two sections.
Whereas §\ref{sec:p3_s8_1} learned a 2D semantic-to-pixel prior on a single image plane and §\ref{sec:p3_s8_2} lifted that prior into a static volumetric representation, this section adds the time axis: the model must place matter and appearance correctly \emph{and} make them evolve plausibly as the camera or the scene itself moves through time.
Appearance, geometry, and dynamics must now be \emph{jointly} consistent: a generated clip is judged not only frame by frame but on object permanence, motion plausibility, and multi-view coherence under camera motion, which makes video the natural substrate for the dynamic world models and simulators that consume spatial structure downstream.

Formally, we cast every system in this section as a conditional spatio-temporal sampler that extends the latent-generator framework of §\ref{subsec:p3_s8_1_1}.
Let $\mathcal{V}=[0,1]^{T\times H\times W\times 3}$ be the space of $T$-frame clips and $\mathcal{Z}^{\text{vid}}$ a compressed spatio-temporal latent space supplied by a video autoencoder that extends the image autoencoder of Eq.~\eqref{eq:t2i_autoencoder} along the time axis.
Writing $\phi$ for the text encoder of Eq.~\eqref{eq:t2i_generator} and $\kappa:\mathcal{C}\!\to\!\mathcal{C}'$ for a generic encoder of an auxiliary control signal $\mathbf{c}\in\mathcal{C}$, the core learning problem is to train
\begin{equation}
G_\theta^{\text{vid}}:\mathcal{E}\times\mathcal{C}'\rightarrow\mathcal{Z}^{\text{vid}},\qquad
\hat{\mathbf{z}}^{1:T}\sim p_\theta\!\big(\mathbf{z}^{1:T}\,\big|\,\phi(\mathbf{x}),\,\kappa(\mathbf{c})\big),
\label{eq:p3_s8_3_objective}
\end{equation}
so that the decoded clip $\hat{\mathbf{v}}=D(\hat{\mathbf{z}}^{1:T})$ is photorealistic, temporally coherent, and faithful to both the prompt $\mathbf{x}$ and the control $\mathbf{c}$.
Two limiting cases anchor the framework: with no auxiliary control ($\kappa=\varnothing$) Eq.~\eqref{eq:p3_s8_3_objective} reduces to text-to-video generation, and with a single frame ($T=1$) it recovers the image generator of Eq.~\eqref{eq:t2i_generator} exactly, so video spatial generation is precisely the image prior of §\ref{sec:p3_s8_1} unrolled in time and equipped with a richer conditioning interface.

The subsections are organized by the axis of control each introduces into Eq.~\eqref{eq:p3_s8_3_objective}.
Text-to-video generation (§\ref{subsec:p3_s8_3_1}) is the base case without auxiliary control; spatial control (§\ref{subsec:p3_s8_3_2}) adds a per-frame scaffold such as a pose or depth sequence; temporal control (§\ref{subsec:p3_s8_3_3}) adds a motion or camera signal that governs dynamics rather than appearance; §\ref{subsec:p3_s8_3_4} covers the remaining modalities and the frameworks that fuse several at once; and video editing (§\ref{subsec:p3_s8_3_5}) modifies an existing clip, lifting the 2D editing of §\ref{subsec:p3_s8_1_3} into time.
§\ref{subsec:p3_s8_3_6} compares the five and discusses the challenges they share.



\subsection{Text-to-Video Generation}
\label{subsec:p3_s8_3_1}

Text-to-video (T2V) generation~\cite{sora2024,polyak2024movie,hong2022cogvideo,wan2025wan} synthesizes photorealistic, temporally coherent video from a natural-language description.
It extends text-to-image (T2I) synthesis along time, so the model must produce plausible motion and state transitions as well as high-fidelity frames, and it does so with no explicit spatial scaffold: object permanence, motion plausibility, and view coherence under camera motion cannot be supplied from outside but must \emph{emerge} from the generator itself.
That is why the unconditional T2V model is the substrate on which the explicit control axes of §\ref{subsec:p3_s8_3_2}--§\ref{subsec:p3_s8_3_5} are built.
Its architecture has followed the image case, from generative adversarial networks (GANs) through autoregressive models to the diffusion models that dominate today, with progress driven as much by data as by architecture~\cite{chen2024panda70mcaptioning70mvideos,NEURIPS2024_57f6683e}.

\subsubsection{Task Definition}
\label{subsubsec:p3_s8_3_1_task}
Text-to-video generation is the base case of Eq.~\eqref{eq:p3_s8_3_objective} with no auxiliary control ($\kappa=\varnothing$): the generator samples a latent from the prompt embedding alone,
\begin{equation}
\hat{\mathbf{z}}^{1:T}\sim p_\theta\!\big(\mathbf{z}^{1:T}\,\big|\,\phi(\mathbf{x})\big),
\label{eq:p3_s8_3_1_objective}
\end{equation}
which is exactly the latent sampler $G_\theta$ of Eq.~\eqref{eq:t2i_generator} unrolled over $T$ frames in the spatio-temporal latent space $\mathcal{Z}^{\text{vid}}$.

\subsubsection{Methodology}
\label{subsubsec:p3_s8_3_1_methodology}
The three paradigms factorize the clip distribution $p_\theta(\mathbf{z}^{1:T}\mid\phi(\mathbf{x}))$ adversarially, causally over discrete tokens, or by iterative spatio-temporal denoising, and each changes its image counterpart of §\ref{subsec:p3_s8_1_1} in one structural respect; Figure~\ref{fig:p3_s8_3_1_paradigms} contrasts how each handles the temporal dimension.

\begin{figure}[!ht]
\centering
\includegraphics[width=\linewidth]{figures/part3/section8/text-to-video_generation.pdf}
\caption{Three mechanisms for text-to-video generation: (a) a GAN with spatial and temporal critics, a representative rather than universal design; (b) autoregressive prediction of discrete space--time tokens, with pale cells as generated context and the blue cell as the next token; and (c) diffusion with factorized spatial / temporal attention, where $s$ is the sampling step and $t$ is video time. The ball sequence is illustrative.}
\label{fig:p3_s8_3_1_paradigms}
\end{figure}

\noindent \textbf{Generative Adversarial Networks (GANs)}. The first T2V systems extended the image GANs of §\ref{para:p3_s8_1_1_gan} to clips. The video-specific change is in the discriminator: rather than scoring a single image, it must judge both per-frame realism \emph{and} the plausibility of motion across the sequence, which is naturally realized as a two-critic objective
\begin{equation}
\min_\theta\ \max_{\psi_s,\psi_t}\ \mathbb{E}_{\mathbf{v}\sim p_\text{data}}\!\big[\log D_{\psi_s}(\mathbf{v})+\log D_{\psi_t}(\mathbf{v})\big] + \mathbb{E}_{\mathbf{n}}\!\big[\log\!\big(1-D_{\psi_s}(\hat{\mathbf{v}})\big)+\log\!\big(1-D_{\psi_t}(\hat{\mathbf{v}})\big)\big],
\label{eq:p3_s8_3_1_gan}
\end{equation}
where $\hat{\mathbf{v}}=D\big(G_\theta(\mathbf{n},\phi(\mathbf{x}))\big)$ is the decoded clip, the spatial critic $D_{\psi_s}$ scores individual frames, and the temporal critic $D_{\psi_t}$ scores motion coherence over the clip, the term absent from the single-discriminator image objective of Eq.~\eqref{eq:t2i_gan_objective}. TGAN~\cite{saito2017temporal} and TGAN-v2~\cite{saito2020train} used a temporal generator to produce latent motion vectors from which an image generator rendered frames; MoCoGAN~\cite{tulyakov2018mocogan} and MoCoGAN-HD~\cite{tian2021good} disentangled content from motion, generating the motion code recurrently; DVD-GAN~\cite{clark2019adversarial} instantiated Eq.~\eqref{eq:p3_s8_3_1_gan} directly with separate spatial and temporal discriminators; and DIGAN~\cite{yu2022generating} used implicit neural representations for longer sequences. Even with architectures such as StyleGAN-V~\cite{skorokhodov2022stylegan}, GANs suffered from flicker, training instability, and mode collapse, which limited the diversity of their motion and pushed the field toward other paradigms.


\noindent \textbf{Autoregressive Models}. Autoregressive (AR) models recast T2V as the sequential prediction of §\ref{para:p3_s8_1_1_ar}, tokenizing a clip with a 3D Vector Quantized Variational Autoencoder (3D~VQ-VAE) and predicting each token from its predecessors with a large transformer. The video-specific change is the \emph{order} of the factorization: where the image model of Eq.~\eqref{eq:t2i_ar_factorisation} runs over the spatial tokens of a single frame, the video model factorizes over space \emph{and} time,
\begin{equation}
p_\theta\!\big(\mathbf{z}^{1:T}\,\big|\,\phi(\mathbf{x})\big)\;=\;\prod_{t=1}^{T}\prod_{n=1}^{N} p_\theta\!\big(z_{t,n}\,\big|\,\mathbf{z}^{<t},\,z_{t,<n},\,\phi(\mathbf{x})\big),
\label{eq:p3_s8_3_1_ar}
\end{equation}
so that each frame is conditioned on all previously generated frames $\mathbf{z}^{<t}$. This inter-frame causal dependency is the source of both the paradigm's strength (global temporal coherence) and its characteristic failure, the accumulation of per-step error over long roll-outs. VideoGPT~\cite{yan2021videogpt} demonstrated the approach for unconditional generation, N\"{U}WA~\cite{wu2022nuwa} and CogVideo~\cite{hong2022cogvideo} adapted it to text-to-video synthesis, and TATS~\cite{ge2022long} addressed quality degradation in long sequences. The paradigm has since resurged in large multimodal backbones: Emu3~\cite{wang2024emu3} generates video inside a language-model backbone by treating frames as discrete visual tokens, and CausVid~\cite{yin2025slow} and MAGI-1~\cite{ai2025magi1autoregressivevideogeneration} scale to long-form generation with autoregressive diffusion transformers. Sequential decoding nevertheless remains expensive and prone to the error accumulation that Eq.~\eqref{eq:p3_s8_3_1_ar} makes explicit.


\noindent \textbf{Diffusion Models}. Diffusion models, which now dominate T2V, lift the latent-denoising objective of §\ref{para:p3_s8_1_1_diffusion} from a single image to the clip latent $\mathbf{z}^{1:T}$. The forward and reverse processes are those of Eq.~\eqref{eq:t2i_ldm_objective}; the video-specific change is that the denoiser $\epsilon_\theta$ operates on a spatio-temporal tensor and is trained to predict the noise added to the whole clip at sampling step $s$,
\begin{equation}
\min_\theta\ \mathbb{E}_{s,\,\mathbf{z}_0^{1:T},\,\epsilon}\!\left[\big\|\epsilon-\epsilon_\theta\!\big(\mathbf{z}_s^{1:T},\,s,\,\phi(\mathbf{x})\big)\big\|_2^2\right],
\label{eq:p3_s8_3_1_diff}
\end{equation}
where $\epsilon_\theta$ models dependencies across both space and time, and $s$ is distinct from the frame index $t$. One widely used design \emph{factorizes} attention into spatial attention within each frame and temporal attention across frames, as illustrated in Figure~\ref{fig:p3_s8_3_1_paradigms}(c); other architectures use joint spacetime attention.
VDM~\cite{ho2022video} applied the idea in pixel space with a 3D U-Net; the breakthrough was moving diffusion into the compressed latent space of a video-specific autoencoder, as in VideoLDM~\cite{blattmann2023align}, LVDM~\cite{he2022latent}, and the open-source Stable Video Diffusion (SVD)~\cite{blattmann2023stable}, which cut training and inference cost enough to make high-resolution generation practical.

A second recipe adapts pretrained T2I models. Imagen Video~\cite{ho2022imagen} and Make-A-Video~\cite{singer2022make} freeze a T2I model and insert trainable temporal modules (\eg, temporal attention layers or 1D convolutions) between its spatial layers, which are then trained on video to learn motion while the image prior is inherited intact. VideoCrafter~\cite{chen2023videocrafter1,chen2024videocrafter2} and LaVie~\cite{wang2025lavie} built high-quality generators on this recipe, Tune-A-Video~\cite{wu2023tune} adapted it to one-shot tuning on a single clip, and AnimateDiff~\cite{guo2023animatediff} turned the temporal modules into plug-and-play motion modules that animate personalized T2I models without further fine-tuning.

Diffusion Transformers (DiTs) have since replaced the U-Net backbone, operating on ``spacetime patches'' (unified tokens representing blocks of video) and scaling well. Sora~\cite{sora2024} demonstrated long-duration (up to one minute), high-resolution clips and presented the approach as a step toward world simulation, and a wave of DiT models followed, including Open-Sora~\cite{zheng2024open}, LATTE~\cite{ma2024latte}, CogVideoX~\cite{yang2024cogvideox}, StepVideo~\cite{ma2025stepvideot2vtechnicalreportpractice}, Movie Gen~\cite{polyak2024movie}, Wan~\cite{wan2025wan}, and Veo~\cite{DeepMindVeo}.
Large open DiT foundation models have substantially narrowed the gap to proprietary systems. HunyuanVideo~\cite{kong2024hunyuanvideo}, at $13$B parameters the largest open T2V model at release, couples a ``dual-stream-to-single-stream'' full-attention DiT (text and video tokens processed in separate streams before fusion) with a multimodal-LLM text encoder and a 3D causal VAE, and reports human-evaluated quality competitive with leading closed systems. Mochi~1~\cite{genmo2024mochi} introduces an asymmetric DiT (AsymmDiT) that allocates most capacity to the visual stream, and LTX-Video~\cite{hacohen2024ltxvideo} uses a high-compression video VAE to run latent video diffusion in near real time. NVIDIA's Cosmos~\cite{nvidia2025cosmos} reframes the video generator as a world-foundation model for ``Physical AI'', trained so that embodied and robotics systems can use it as a learned simulator, which ties video generation to the world models of §\ref{sec:p3_s8_6}.
Across these shifts the T2V trajectory has trailed the T2I trajectory of §\ref{subsec:p3_s8_1_1} by roughly a year, each image paradigm being inflated along the time axis once it matured.

\subsubsection{Benchmarks}
\label{subsubsec:p3_s8_3_1_benchmarks}
Video adds a requirement at every level of evaluation: corpora must pair clips with captions at scale, suites must stratify by temporal phenomena (motion, object permanence, dynamic degree), and metrics must measure cross-frame coherence as well as per-frame quality.

\paragraph{Training Corpora.}
Caption--clip scale has been the principal driver of T2V quality, and high-quality paired video--text data is far scarcer than image--text data. Early work repurposed action-recognition corpora such as UCF-101~\cite{soomro2012ucf101} and Kinetics~\cite{carreira2017quo}, which supply short, class-labeled clips but no captions. WebVid-10M~\cite{bain2021frozen}, the first web-scale captioned corpus, underpinned most 2022--2023 systems despite noisy alt-text captions and watermarks; later corpora raised caption density and clip length, with Panda-70M~\cite{chen2024panda70mcaptioning70mvideos} captioned automatically by cross-modality teachers, MiraData~\cite{NEURIPS2024_57f6683e} targeting long clips with dense structured captions, and Koala-36M~\cite{wang2025koala36m} explicitly optimizing the consistency between caption and content, the property that most directly limits a generator's controllability. Table~\ref{tab:p3_s8_3_1_corpora} gives their scale.

\begin{table}[!ht]
\centering
\caption{Representative training corpora for text-to-video generation. Action-recognition sets (top) supply class labels rather than captions; web-scale captioned corpora (bottom) supply the clip--text pairs on which modern generators are fit. MiraData scale is that of its largest release.}
\label{tab:p3_s8_3_1_corpora}
\fontsize{8}{10}\selectfont
\setlength{\tabcolsep}{1.0mm}
\begin{tabular}{L{0.22\textwidth} C{0.08\textwidth} L{0.24\textwidth} L{0.38\textwidth}}
\toprule
\textbf{Corpus} & \textbf{Year} & \textbf{Scale} & \textbf{Annotation / Notes} \\
\midrule
UCF-101~\cite{soomro2012ucf101} & 2012 & 13{,}320 clips & 101 action classes; labels, no captions \\
Kinetics~\cite{carreira2017quo} & 2017 & $\sim$650K clips & Action classes; labels, no captions \\
WebVid-10M~\cite{bain2021frozen} & 2021 & 10M clip--text pairs & Web alt-text captions; noisy, watermarked \\
Panda-70M~\cite{chen2024panda70mcaptioning70mvideos} & 2024 & 70M clip--text pairs & Auto-captioned by cross-modality teachers \\
MiraData~\cite{NEURIPS2024_57f6683e} & 2024 & 330K clips & Long-duration; dense structured captions \\
Koala-36M~\cite{wang2025koala36m} & 2025 & 36M clip--text pairs & Fine-grained structured captions; condition--content consistency \\
\bottomrule
\end{tabular}
\end{table}

\paragraph{Evaluation Benchmark Suites.}
Classical T2V evaluation reported zero-shot distributional scores on the held-out splits of UCF-101 and MSR-VTT~\cite{xu2016msrvtt}, which probe frame quality and coarse text relevance but say nothing about temporal or compositional faithfulness. Recent suites are purpose-built: VBench~\cite{huang2023vbench} decomposes quality into fine-grained dimensions (subject and background consistency, motion smoothness, dynamic degree, temporal flicker, per-frame aesthetic and imaging quality), each with an automatic protocol validated against human preference; EvalCrafter~\cite{liu2024evalcrafter} aggregates objective metrics over prompts derived from real-world usage; and T2V-CompBench~\cite{sun2024t2vcompbench} targets compositional faithfulness over time (attribute binding, spatial relations, motion binding, interactions), the temporal analogue of the T2I-CompBench probe of §\ref{subsec:p3_s8_1_1}. Table~\ref{tab:p3_s8_3_1_eval} lists them.

\begin{table}[!ht]
\centering
\caption{Representative evaluation resources for text-to-video generation. Distributional resources (top) reuse captioned-clip splits for zero-shot FVD/IS/CLIPSIM; multi-dimensional suites (bottom) score temporal coherence, dynamics, and compositional faithfulness directly.}
\label{tab:p3_s8_3_1_eval}
\fontsize{8}{10}\selectfont
\setlength{\tabcolsep}{1.0mm}
\begin{tabular}{L{0.22\textwidth} C{0.10\textwidth} L{0.36\textwidth} L{0.24\textwidth}}
\toprule
\textbf{Suite} & \textbf{Year} & \textbf{Evaluation Axis} & \textbf{Scoring} \\
\midrule
UCF-101 / MSR-VTT~\cite{soomro2012ucf101,xu2016msrvtt} & 2012--16 & Frame quality $+$ coarse text relevance (MSR-VTT: 10K captioned clips) & Zero-shot FVD / IS / CLIPSIM \\
VBench~\cite{huang2023vbench} & 2024 & 16 dimensions (consistency, motion, dynamics, flicker, quality) & Per-dimension automatic, human-validated \\
EvalCrafter~\cite{liu2024evalcrafter} & 2024 & 17 objective metrics over 700 structured prompts & Aggregated automatic, fitted to human opinion \\
T2V-CompBench~\cite{sun2024t2vcompbench} & 2024 & Compositional: attribute/motion binding, relations, interactions & MLLM- / detector-based \\
\bottomrule
\end{tabular}
\end{table}

\paragraph{Evaluation Metrics.}
The metric that defines video evaluation is the Fr\'echet Video Distance (FVD)~\cite{unterthiner2018fvd}, the spatio-temporal analogue of FID: it compares distributions of features extracted from generated and real clips by a pretrained \emph{video} backbone (I3D), and is therefore sensitive to motion realism in a way that per-frame FID~\cite{heusel2017gans} and Inception Score are not. Text--clip alignment is measured by CLIPSIM, the per-frame CLIPScore~\cite{hessel2021clipscore} averaged over the clip, which captures topical relevance but is blind to whether the \emph{described} motion actually occurs. Because none of these scalars captures coherence or dynamics, papers report several VBench dimensions alongside them, which exposes a trade-off a single number hides: \emph{dynamic degree} (how much genuine motion a clip contains) against \emph{temporal consistency} (how stable it stays), since a near-static clip trivially maximizes consistency while failing the prompt. Headline claims are corroborated with human preference.

\subsubsection{Challenges and Trends}
\label{subsubsec:p3_s8_3_1_challenges}
Three problems belong to the unconditional base case. \textbf{Long-horizon coherence from text alone} is the first: with no scaffold, object permanence and motion plausibility must emerge implicitly, and all three paradigms degrade as clips lengthen (GANs flicker, autoregressive roll-outs accumulate the error of Eq.~\eqref{eq:p3_s8_3_1_ar}, diffusion models lose consistency beyond their trained window), so coherence beyond a few seconds remains the binding constraint. \textbf{Motion realism without a control signal} is the second: prompts under-specify dynamics, so the generator must sample from a highly multi-modal motion posterior, and the easy optimum is near-static or physically implausible motion that still scores well on per-frame metrics. The \textbf{video--text data bottleneck} is the third: paired data is scarce relative to images, web captions are short and weakly aligned, and this caps both controllability and the reliability of the automatic captioning used to scale corpora, which is what the dense-caption efforts (Panda-70M, MiraData, Koala-36M) address. The explicit control interfaces of the following subsections relieve the first two by moving structure and dynamics from the prompt into the input.


\subsection{Spatial Controlled Video Generation}
\label{subsec:p3_s8_3_2}

While textual descriptions provide high-level semantic guidance, they often lack the granularity required for precise control over object placement, scene geometry, or character appearance. Spatial-controlled video generation addresses this limitation by augmenting the generation process with explicit spatial conditioning signals, affording users more direct and deterministic influence over the synthesized visual output.
The geometry, articulation, and layout that the text-to-video base case of §\ref{subsec:p3_s8_3_1} had to \emph{hallucinate} from the prompt alone are here \emph{externalized} into a per-frame scaffold: the spatial structure becomes an input the model must honor rather than an emergent property it must invent, which is precisely what makes the output controllable and verifiable.

\subsubsection{Task Definition}
\label{subsubsec:p3_s8_3_2_task}
Spatial control instantiates Eq.~\eqref{eq:p3_s8_3_objective} with a control encoder $\kappa_{\text{s}}$ acting on a spatial condition $\mathbf{c}_{\text{s}}^{1:T}$ (a static or sparse reference image, a pose sequence, a depth-map sequence, facial landmarks, a sketch sequence, or a bounding-box track),
\begin{equation}
\hat{\mathbf{z}}^{1:T}\sim p_\theta\!\big(\mathbf{z}^{1:T}\,\big|\,\phi(\mathbf{x}),\,\kappa_{\text{s}}(\mathbf{c}_{\text{s}}^{1:T})\big),
\label{eq:p3_s8_3_2_objective}
\end{equation}
so that the scaffold pins down the per-frame spatial layout while the prompt $\mathbf{x}$ and the temporal modules supply appearance and motion. The conditioning channels surveyed below differ only in what $\mathbf{c}_{\text{s}}$ encodes (appearance for an image, articulation for pose or landmarks, 3D structure for depth, free-form layout for a sketch, or object location for a bounding box), and recover the text-to-video base case of Eq.~\eqref{eq:p3_s8_3_1_objective} exactly when $\mathbf{c}_{\text{s}}=\varnothing$.

\subsubsection{Methodology}
\label{subsubsec:p3_s8_3_2_methodology}
The six channels carry different spatial information, while the injection mechanism depends on the model. A common design is a ControlNet-style adapter~\cite{zhang2023adding}: a trainable control branch processes the encoded condition and injects a residual into intermediate features of a video backbone, schematically at layer $\ell$
\begin{equation}
h_\ell' = h_\ell + \Delta h_\ell,
\qquad
\Delta h_\ell = \mathcal{A}_{\psi,\ell}\!\big(\kappa_{\text{s}}(\mathbf{c}_{\text{s}}^{1:T})\big),
\label{eq:p3_s8_3_2_inject}
\end{equation}
where $h_\ell$ denotes backbone features and $\mathcal{A}_{\psi,\ell}$ a trainable control branch (dependence on the noisy latent, diffusion step, and other conditioning is omitted, and the backbone may be frozen while the adapter trains); latent concatenation and attention-based conditioning are the other common implementations.
Figure~\ref{fig:p3_s8_3_2_mechanisms} shows the six channels alongside this adapter pathway: a reference image acts as a static appearance anchor, the other five are usually frame-indexed sequences, and compatible channels can be combined.

\begin{figure}[!ht]
\centering
\includegraphics[width=\linewidth]{figures/part3/section8/spatial_control_channels_for_video_generation.pdf}
\caption{Six spatial-control channels for video generation, (a)--(f), which can be selected or combined, and an example injection route (right) in which a trainable adapter adds residual features to a frozen video backbone (Eq.~\eqref{eq:p3_s8_3_2_inject}). Frames and control maps are illustrative.}
\label{fig:p3_s8_3_2_mechanisms}
\end{figure}

\noindent\textbf{Image Control.} A reference image (Figure~\ref{fig:p3_s8_3_2_mechanisms}a) specifies the style, content, or initial state of a video, which is the basis of animating a still and of keeping a subject's appearance consistent. I2VGen-XL~\cite{zhang2023i2vgen}, Stable Video Diffusion (SVD)~\cite{blattmann2023stable}, VideoCrafter1~\cite{chen2023videocrafter1}, and Lumiere~\cite{bar2024lumiere} use the reference image as an initial condition, often by concatenating its latent with the initial noise, and ConsistI2V~\cite{ren2024consisti2v} refines this for consistency. The model then generates a plausible motion sequence that evolves from the image, typically under a text prompt, as in DynamiCrafter~\cite{xing2024dynamicrafter}.

\noindent\textbf{Pose Control.} Pose-guided generation (Figure~\ref{fig:p3_s8_3_2_mechanisms}b) animates articulated subjects, predominantly humans, along a predefined sequence of skeletal poses, the video counterpart of pose-guided image generation~\cite{ju2023humansd,zhang2023adding,Ju_2023_CVPR} and the core of avatar animation and character generation. The pose keypoints (\eg, from MMPose~\cite{mmpose2020}) are encoded and injected at each denoising step, frequently through a ControlNet-style adapter~\cite{zhang2023adding}. MagicAnimate~\cite{xu2024magicanimate}, Animate Anyone~\cite{hu2024animate}, Follow-Your-Pose~\cite{ma2024followyourpose}, and DreamPose~\cite{karras2023dreampose} pair a reference image for appearance with a pose sequence for motion, DisPose~\cite{li2024dispose} disentangles the pose guidance, and Champ~\cite{zhu2024champ} and MagicPose~\cite{chang2023magicpose} target high-fidelity motion, so that the output keeps the subject's identity while following the target poses.

\noindent\textbf{Depth Control.} A sequence of depth maps (Figure~\ref{fig:p3_s8_3_2_mechanisms}c), estimated from an existing video or rendered from a 3D model, specifies dense 3D structure and so enforces consistent object layout and realistic occlusion. ControlVideo~\cite{zhao2025controlvideo}, Make-Your-Video~\cite{xing2024make}, Control-A-Video~\cite{chen2023control}, and ControlNeXt~\cite{peng2024controlnext} adapt the ControlNet framework to video, processing each frame's depth map with a trainable adapter that guides the denoiser; DreamDance~\cite{pang2024dreamdance} and SparseCtrl~\cite{guo2024sparsectrl} also use depth for dynamic scene generation.

\noindent\textbf{Landmark Control.} Facial landmarks (Figure~\ref{fig:p3_s8_3_2_mechanisms}d) give fine-grained control over expression and are the standard interface for talking-head animation. EchoMimic~\cite{chen2024echomimic}, HunyuanPortrait~\cite{xu2025hunyuanportrait}, and LivePortrait~\cite{guo2024liveportrait} animate a static portrait from landmarks, often themselves driven by audio, so that lip movements, blinks, and subtle expressions stay synchronized; TASTE-Rob~\cite{zhao2025taste}, Takin-ADA~\cite{lin2024takin}, and Follow-Your-Emoji~\cite{ma2024follow} extend landmarks to more expressive control.

\noindent\textbf{Sketch Control.} Sketches (Figure~\ref{fig:p3_s8_3_2_mechanisms}e) let users define compositional layout and motion with simple line drawings, an intuitive interface for animation. ToonCrafter~\cite{xing2024tooncrafter}, VidSketch~\cite{jiang2025vidsketch}, AniDoc~\cite{meng2025anidoc}, and SketchBetween~\cite{loftsdottir2022sketchbetween} take a sequence of sketches to guide interpolation or generation, encoding the line art and injecting it as a condition; LayerAnimate~\cite{yang2025layeranimate} animates hand-drawn characters and scenes layer by layer, and MagicStick~\cite{ma2025magicstick} uses sketch-like structural cues for editing.

\noindent\textbf{Bounding Box (BBox) Control.} Bounding-box tracks (Figure~\ref{fig:p3_s8_3_2_mechanisms}f) specify where objects are and how they move over time, and dominate autonomous-driving simulation: MagicDrive~\cite{gao2023magicdrive}, DriveDreamer~\cite{wang2024drivedreamer}, DriveScape~\cite{wu2024drivescape}, and DreamForge~\cite{mei2024dreamforge} use 3D boxes to control the position and motion of vehicles and pedestrians. The box information enters either through cross-attention, as in DiVE~\cite{jiang2024dive}, or as a spatial feature map, and Boximator~\cite{wang2024boximator} and Ctrl-V~\cite{luo2025ctrl} generate and control object motion from box sequences.

Nearly every structural channel (pose, depth, sketch, box) is paired with an appearance reference image, since structure alone cannot fix identity or texture, and the channels cluster by application: image and landmark control in portrait and talking-head synthesis, pose and sketch control in character animation, depth and box control in geometry-faithful and driving-simulation settings.

\subsubsection{Benchmarks}
\label{subsubsec:p3_s8_3_2_benchmarks}
Unlike the text-to-video base case of §\ref{subsec:p3_s8_3_1}, spatial control has no single canonical corpus or suite: each channel is trained and evaluated on data specialized to its application. Models are typically built on the general web-scale corpora of §\ref{subsec:p3_s8_3_1} (WebVid-10M~\cite{bain2021frozen} and successors) for appearance priors, then specialized on task data: human-animation channels on dance- and fashion-video collections paired with extracted pose keypoints (\eg, via OpenPose~\cite{cao2017realtime} or MMPose~\cite{mmpose2020}), talking-head landmark control on portrait-video sets, and bounding-box driving simulation on annotated driving datasets such as nuScenes~\cite{caesar2020nuscenes}, whose 3D box and trajectory labels supply the control signal directly.

\paragraph{Control-fidelity Metrics.}
Spatial control is judged on two axes simultaneously: the generic video-quality metrics of §\ref{subsec:p3_s8_3_1} (FVD~\cite{unterthiner2018fvd}, CLIPSIM, and VBench~\cite{huang2023vbench} dimensions) and channel-specific \emph{adherence} metrics that measure how faithfully the output respects $\mathbf{c}_{\text{s}}$. The latter are necessarily per-channel: pose and landmark control report keypoint accuracy (re-detected-keypoint distance to the target skeleton or face mesh); depth control reports depth consistency between the re-estimated and the conditioning maps; image control reports identity and appearance preservation (face-/subject-similarity to the reference); and bounding-box control reports detection mAP and trajectory IoU of the generated objects against the conditioning boxes. Neither axis suffices alone, so methods report a quality-plus-adherence tuple: a model can maximize adherence by ignoring the prompt (rigid, texture-collapsed output) or maximize quality by ignoring the control. The same holds for every controlled axis that follows.

\subsubsection{Challenges and Trends}
\label{subsubsec:p3_s8_3_2_challenges}
Single-channel spatial control has three problems of its own; those that arise once several conditions are fused are discussed in §\ref{subsec:p3_s8_3_6}. Because the structural scaffold and the appearance reference pass through the same denoiser, strong control leaks structure into texture: pose and depth conditioning induce identity drift or bake the scaffold's silhouette into the surface, and disentangling what moves from what it looks like (as DisPose~\cite{li2024dispose} attempts for pose) remains open. Dense scaffolds (image, depth) over-constrain and suppress plausible motion, whereas sparse ones (boxes, coarse sketches) leave appearance and fine geometry ambiguous, so no channel is uniformly preferable and the right granularity depends on the task. And adapters trained on, say, frontal human dance transfer poorly to unusual viewpoints or non-human subjects, which is the data cost of per-channel specialization and one motivation for the multi-condition frameworks of §\ref{subsec:p3_s8_3_4}.


\subsection{Temporal Controlled Video Generation}
\label{subsec:p3_s8_3_3}

Temporal control regulates how motion and timing evolve across frames, the lever behind cinematic camera work, narrative pacing, and believable physical interaction. It is the complement of the spatial scaffold of §\ref{subsec:p3_s8_3_2}: where spatial control fixes \emph{where things are} in each frame, temporal control fixes \emph{how they move between} frames, externalizing the dynamics (object motion, camera path, timing) that the text-to-video base case of §\ref{subsec:p3_s8_3_1} otherwise had to sample blindly from a highly multi-modal motion posterior.

\subsubsection{Task Definition}
\label{subsubsec:p3_s8_3_3_task}
Temporal control instantiates Eq.~\eqref{eq:p3_s8_3_objective} with an encoder $\kappa_{\text{t}}$ of a motion or timing signal $\mathbf{m}^{1:T}$ (an optical-flow field, an object trajectory, a camera-pose path, or a learned motion embedding),
\begin{equation}
\hat{\mathbf{z}}^{1:T}\sim p_\theta\!\big(\mathbf{z}^{1:T}\,\big|\,\phi(\mathbf{x}),\,\kappa_{\text{t}}(\mathbf{m}^{1:T})\big).
\label{eq:p3_s8_3_3_objective}
\end{equation}
The contrast with spatial control (Eq.~\eqref{eq:p3_s8_3_2_objective}) is that $\mathbf{m}^{1:T}$ constrains the \emph{inter-frame dynamics} (how content moves between frames) rather than the per-frame appearance, so $\kappa_{\text{t}}$ typically steers the temporal-attention or latent-warping pathway of the denoiser. The channels below differ in how the motion is specified: densely (optical flow), as object paths (trajectory), as virtual-camera parameters (camera), or as a transferable behavioral style (motion).

\subsubsection{Methodology}
\label{subsubsec:p3_s8_3_3_methodology}
The four channels specify inter-frame dynamics through different signals, and no single injection architecture serves them all. One implementation uses a trainable temporal-control branch $\mathcal{T}_\psi$ to modulate features of a pretrained video backbone,
\begin{equation}
h_\ell' = h_\ell + \Delta h_\ell,
\qquad
\Delta h_\ell=\mathcal{T}_{\psi,\ell}\!\big(\kappa_{\text{t}}(\mathbf{m}^{1:T})\big),
\label{eq:p3_s8_3_3_inject}
\end{equation}
with dependence on noisy latents, diffusion step, and text omitted.
More broadly, two routes recur: explicit transformations that use the motion to warp noise, latents, or attention features, and encoded conditioning that steers generation through adapters, attention, or learned motion features. They are not tied to particular channels (flow and trajectories can drive learned conditioning, and camera control can use geometry-aware warping), and Figure~\ref{fig:p3_s8_3_3_mechanisms} keeps the four input channels separate from the two routes.

\begin{figure}[!ht]
\centering
\includegraphics[width=\linewidth]{figures/part3/section8/temporal_control_channels_for_video_generation.pdf}
\caption{Four temporal-control channels for video generation, (a)--(d), and two alternative ways of applying them (right): warping noise or latents, and encoding the signal as conditioning features. Frames are illustrative.}
\label{fig:p3_s8_3_3_mechanisms}
\end{figure}

\noindent\textbf{Flow Control.} The most explicit motion signal (Figure~\ref{fig:p3_s8_3_3_mechanisms}a) is an optical-flow field, which records the per-pixel displacement between consecutive frames and so specifies dynamics at the finest granularity; conditioning on a sequence of such fields yields specific, coherent motion at the cost of authoring or estimating a dense signal. Motion-I2V~\cite{shi2024motion} integrates a motion-field predictor into its diffusion framework so that generation follows the supplied flow, I2VControl~\cite{feng2024i2vcontrol} disentangles flow for unified control, and MOFA-Video~\cite{niu2024mofa} animates a still image from user-painted motion fields. Go-with-the-Flow~\cite{burgert2025gowiththeflow} dispenses with the adapter altogether, warping the diffusion \emph{noise prior} along the target flow so that the motion constraint lives in the sampling process rather than in an added branch.

\noindent\textbf{Trajectory Control.} Sparsifying the dense field into a few object paths (Figure~\ref{fig:p3_s8_3_3_mechanisms}b) trades pixel-level precision for an interface a user can actually author: a sequence of 2D or 3D coordinates, or a path drawn interactively, directs character movement or object kinematics, entering the denoiser either by modifying the latents or by steering cross-attention maps. MotionCanvas~\cite{xing2025motioncanvas}, MagicMotion~\cite{li2025magicmotion}, and Tora~\cite{zhang2025tora} consume a predefined curve, FreeTraj~\cite{qiu2024freetraj} relaxes this to free-form trajectories, Peekaboo~\cite{jain2024peekaboo} also schedules when objects appear and disappear along the path, and DragNUWA~\cite{yin2023dragnuwa} and DragAnything~\cite{wu2024draganything} let the user drag regions through the frame interactively. MotionCtrl~\cite{wang2024motionctrl} handles object trajectory and camera motion in one controller, treating an object path and a camera path as two facets of the same inter-frame displacement.

\noindent\textbf{Camera Control.} Moving that displacement from objects to the observer (Figure~\ref{fig:p3_s8_3_3_mechanisms}c), camera-guided generation manipulates the virtual camera's extrinsics (translation, rotation) and intrinsics (\eg, focal length) to produce pans, tilts, zooms, and dolly moves. CameraCtrl~\cite{he2024cameractrl} and CameraCtrl~II~\cite{he2025cameractrl} encode the camera trajectory as a conditioning feature; Direct-a-Video~\cite{yang2024direct}, ViewCrafter~\cite{yu2024viewcrafter}, and Cavia~\cite{xu2024cavia} couple it with novel-view synthesis; and 3DTrajMaster~\cite{fu20243dtrajmaster}, MotionFlow~\cite{meral2024motionflow}, and CineMaster~\cite{wang2025cinemaster} control camera and objects jointly in 3D. Two directions now lead. ReCamMaster~\cite{bai2025recammaster} and TrajectoryCrafter~\cite{yu2025trajectorycrafter} re-render an \emph{existing} monocular clip under a new camera path rather than generating from scratch, and Diffusion-as-Shader~\cite{gu2025diffusionshader3dawarevideo} and GS-DiT~\cite{bian2025gsditadvancingvideogeneration} route the camera signal through a 3D proxy (a rendered shader pass or a pseudo-4D Gaussian field), which sharpens control by making it geometry-aware.

\noindent\textbf{Motion Control.} At the abstract end (Figure~\ref{fig:p3_s8_3_3_mechanisms}d), the signal is no longer a geometric displacement but a \emph{behavioral concept} (walking, dancing, jumping) transferred from a reference video as a pose sequence or a learned motion embedding. Unlike the pose control of §\ref{subsec:p3_s8_3_2}, it captures the dynamic style and cadence of a behavior rather than a frame-exact skeleton, so the central technical problem is disentanglement: extracting the motion while discarding the reference's appearance. MotionDirector~\cite{zhao2024motiondirector}, MotionClone~\cite{ling2024motionclone}, and VMC~\cite{jeong2024vmc} learn this motion--appearance separation; Motion Prompting~\cite{geng2025motion}, MotionAgent~\cite{liao2025motionagent}, and MotionMatcher~\cite{wu2025motionmatcher} refine its granularity through trajectory- or feature-level motion descriptors; and DreamMotion~\cite{jeong2024dreammotion} and Customize-A-Video~\cite{ren2024customize} apply a reference clip's dynamics to a new subject or scene fixed by the text prompt.

The channels trade precision against ease of specification: dense flow gives pixel-accurate control but is laborious to author, a motion embedding is effortless to supply but steers the result only loosely, and the most usable systems therefore pair an abstract signal (a drawn path, a reference clip) with a learned predictor that densifies it into the flow- or latent-level displacement the backbone actually consumes.

\subsubsection{Benchmarks}
\label{subsubsec:p3_s8_3_3_benchmarks}
As with spatial control (§\ref{subsec:p3_s8_3_2}), temporal control has no single canonical suite: each channel is evaluated against data specialized to its signal.
Models build on the general corpora of §\ref{subsec:p3_s8_3_1} for appearance and motion priors, then evaluate on task-specific resources: camera control on the posed-video collection RealEstate10K~\cite{zhou2018stereo}, whose per-frame camera poses supply both the conditioning signal and the evaluation ground truth, increasingly supplemented by the larger and more scene-diverse DL3DV-10K~\cite{ling2024dl3dv} to test generalization beyond real-estate walkthroughs; flow control against the MPI-Sintel optical-flow benchmark~\cite{sintel} or flow re-estimated from generated clips with RAFT~\cite{teed2020raft}; and motion- and trajectory-control quality through the motion-smoothness and dynamic-degree dimensions of VBench~\cite{huang2023vbench}.

\paragraph{Control-fidelity Metrics.}
Alongside the generic video-quality metrics of §\ref{subsec:p3_s8_3_1}, each channel measures adherence by re-extracting the control signal from the generated clip and comparing it with the conditioning input; the estimator and the error differ per channel.

\noindent\textit{Flow adherence.} The optical flow of the generated clip is re-estimated with RAFT~\cite{teed2020raft} and compared to the conditioning flow by the \emph{end-point error} (EPE), the per-pixel $L_2$ displacement difference averaged over pixels and frames. A complementary \emph{flow-warping error} warps each generated frame onto the next with the \emph{target} flow and reports the photometric residual, which tests whether the rendered motion realizes the commanded displacement rather than merely correlating with it.

\noindent\textit{Trajectory adherence.} The controlled object is re-localized in the output with a dense point tracker such as CoTracker~\cite{karaev2024cotracker} or a detector, and its tracked center is compared to the target path by the mean $L_2$ distance, normalized by frame size and averaged over time and instances; this is the object motion-control error reported by Tora~\cite{zhang2025tora} and DragNUWA~\cite{yin2023dragnuwa}, and it penalizes both spatial deviation and mistimed arrival.

\noindent\textit{Camera adherence.} The camera trajectory of the generated clip is recovered by structure-from-motion (COLMAP~\cite{schonberger2016structure}) and aligned to the target up to a similarity transform, since pose is recoverable only up to scale. Following CameraCtrl~\cite{he2024cameractrl}, a \emph{rotation error} (the geodesic angle between estimated and target rotations, summed over frames) and a \emph{translation error} (the $L_2$ distance between scale-normalized camera centers) are reported; failure to reconstruct a stable trajectory at all is itself diagnostic of an incoherent output.

\noindent\textit{Motion-transfer fidelity.} With no ground-truth video, fidelity is measured from two sides: a \emph{motion-similarity} term compares the dynamics of output and reference (the correlation of tracked point trajectories or temporal feature statistics), and an \emph{appearance-divergence} term (CLIP similarity to the new prompt together with low feature similarity to the reference's appearance) confirms that the subject was re-synthesized rather than copied.

As for spatial control, these adherence errors are reported together with generic quality, since a near-static clip trivially minimizes EPE and pose error while ignoring the prompt.

\subsubsection{Challenges and Trends}
\label{subsubsec:p3_s8_3_3_challenges}
Three problems are specific to temporal control; fusing it with other conditions is discussed in §\ref{subsec:p3_s8_3_6}. Because the motion signal passes through the same backbone that renders appearance, a transferred motion often drags the source clip's identity or texture with it, the leakage that MotionDirector~\cite{zhao2024motiondirector} and VMC~\cite{jeong2024vmc} try to disentangle and that remains unsolved for large appearance gaps between reference and target. Dense flow over-constrains and suppresses generative freedom, whereas a single trajectory or a short reference clip under-specifies the global dynamics, which is why methods increasingly interpose a learned predictor that densifies an abstract signal into a flow- or latent-level field. And a commanded camera path or object trajectory can demand views or dynamics inconsistent with the scene geometry, producing parallax errors and object distortion; 3D-aware camera control (ViewCrafter~\cite{yu2024viewcrafter}, CineMaster~\cite{wang2025cinemaster}) routes the control through an explicit geometric proxy for exactly this reason, which links temporal control back to the volumetric representations of §\ref{sec:p3_s8_2}.





\subsection{Other Controlled Video Generation}
\label{subsec:p3_s8_3_4}

Beyond the spatial scaffolds of §\ref{subsec:p3_s8_3_2} and the motion signals of §\ref{subsec:p3_s8_3_3}, a video can be steered by further, application-specific conditions: an artistic style reference, a handful of draggable points, a bird's-eye-view road layout, a target text string to render, or a synchronizing audio track.
Each externalizes an attribute that text specifies only loosely, extending controllable generation to creative, interactive, and domain-specific settings; together they raise the question the single-axis interfaces leave open, how to honor \emph{several} conditions at once, which the universal frameworks at the end of this subsection try to answer.

\subsubsection{Task Definition}
\label{subsubsec:p3_s8_3_4_task}
Taken one at a time, each ``other'' modality instantiates the generic conditional sampler of Eq.~\eqref{eq:p3_s8_3_objective} with an encoder $\kappa_{\text{o}}$ of a heterogeneous control signal $\mathbf{c}$ (a style reference, a set of draggable points, a BEV layout, a target text string, or an audio track). The distinctive task of this subsection, however, is to honor \emph{several} such conditions simultaneously: universal and multi-task frameworks fuse $K$ encoded signals into a single sequence,
\begin{equation}
\hat{\mathbf{z}}^{1:T}\sim p_\theta\!\big(\mathbf{z}^{1:T}\,\big|\,\phi(\mathbf{x}),\,\kappa_1(\mathbf{c}_1),\dots,\kappa_K(\mathbf{c}_K)\big),
\label{eq:p3_s8_3_4_objective}
\end{equation}
so that the $K$ encoded conditions are processed jointly rather than one axis at a time, with $K=1$ recovering the single-modality case and $\kappa_k=\varnothing$ for all $k$ recovering the text-to-video base case of §\ref{subsec:p3_s8_3_1}.

\subsubsection{Methodology}
\label{subsubsec:p3_s8_3_4_methodology}
The methods fall into single-modality channels (style, points, BEV, text rendering, audio) and multi-task frameworks that combine compatible conditions. Writing $u_k=\kappa_k(\mathbf c_k)$, the common interface is
\begin{equation}
\hat{\mathbf z}^{1:T}=G_\theta\!\big(\mathbf n,\phi(\mathbf x);\mathcal F_\psi(u_1,\dots,u_K)\big),
\label{eq:p3_s8_3_4_inject}
\end{equation}
where $\mathcal F_\psi$ is a model-specific conditioning interface that may act at several layers rather than produce one fixed-size embedding. Three implementations are representative and can be combined: summed adapter residuals $\sum_k\Delta h_k$, concatenation of the encoded channels with the model input or features, and unified attention over condition tokens. Figure~\ref{fig:p3_s8_3_4_mechanisms} shows the five input modalities and these three fusion designs; $K=1$ denotes one selected modality and $K>1$ a supported combination.

\begin{figure}[!ht]
\centering
\includegraphics[width=\linewidth]{figures/part3/section8/other_control_modalities_for_video_generation.pdf}
\caption{Further control modalities for video generation, (a)--(e), each with its own encoder $E_k$ implementing $\kappa_k$, and (f) alternative ways of fusing $K\geq 1$ encoded conditions. The watercolor driving sequence is illustrative.}
\label{fig:p3_s8_3_4_mechanisms}
\end{figure}

\noindent\textbf{Style Control.} Style-guided generation transfers the artistic style of a reference image or video to a target video. Per-frame style transfer flickers; diffusion-based approaches such as Rerender A Video~\cite{yang2023rerender}, StyleCrafter~\cite{liu2023stylecrafter}, FRESCO~\cite{yang2024fresco}, and StyleMaster~\cite{ye2025stylemaster} instead integrate the style into generation, often by guiding it with the CLIP embedding of the style reference so that the style holds across frames, and UniVST~\cite{song2024univst} generalizes this into a universal framework. FrameGuidance~\cite{jang2025frameguidancetrainingfreeguidance} and OmniVDiff~\cite{xi2025omnivdiffomnicontrollablevideo} add training-free frame-level guidance and omni-controllable diffusion, respectively.

\noindent\textbf{Point Control.} Point control fixes \emph{motion} through a sparse, interactive interface: the user drags a few points on an object from an initial frame to a new position in a later frame, and the model infers a plausible motion for the whole object. Drag-A-Video~\cite{teng2023drag} introduced the approach by optimizing latent codes to satisfy the point constraints, Track4Gen~\cite{jeong2025track4genteachingvideodiffusion} focuses on tracking the user-specified points robustly, and the 3D-aware camera controllers of §\ref{subsec:p3_s8_3_3}~\cite{gu2025diffusionshader3dawarevideo,bian2025gsditadvancingvideogeneration} also accept sparse point input for non-rigid manipulation.

\noindent\textbf{Bird's-Eye View (BEV) Control.} BEV-guided generation serves autonomous-driving simulation: a sequence of top-down semantic maps dictates the road layout, the positions of other vehicles, and the movements of pedestrians. DriveDreamer-2~\cite{zhao2024drivedreamer2llmenhancedworldmodels}, MagicDrive-V2~\cite{gao2025magicdrivev2highresolutionlongvideo}, DiVE~\cite{jiang2024dive}, DreamForge~\cite{mei2024dreamforge}, and Panacea~\cite{wen2023panaceapanoramiccontrollablevideo} take these maps as input and generate first-person driving video consistent with the specified layout and dynamics.

\noindent\textbf{Text Rendering.} Many generators cannot render coherent, legible text inside a scene, so a dedicated channel treats the target string itself as the control signal. Text-Animator~\cite{liu2024textanimatorcontrollablevisualtext} adds text-embedding injection and refinement modules that improve the structure and temporal stability of rendered text, and Wan~\cite{wan2025wan} reports being the first video foundation model able to render legible Chinese and English text.

\noindent\textbf{Audio Control.} Audio-controlled generation synthesizes video synchronized with an input audio signal, in two sub-domains.
\emph{Voice-guided generation} produces talking heads: EMO~\cite{tian2024emoemoteportraitalive}, VASA-1~\cite{xu2024vasa1lifelikeaudiodriventalking}, DreamTalk~\cite{ma2023dreamtalk}, and AniPortrait~\cite{wei2024aniportrait} take a reference image and a speech clip, extract audio features with a pretrained encoder, and drive facial motion with accurate lip-sync and expression; Say Anything~\cite{ma2025sayanythingaudiodrivenlipsynchronization}, MuseTalk~\cite{zhang2025musetalkrealtimehighfidelityvideo}, and CAFE-TALK~\cite{chen2025cafetalkgenerating3dtalking} focus on high-fidelity lip synchronization, and V-Express~\cite{wang2024vexpressconditionaldropoutprogressive} and MegActor~\cite{yang2024megactorharnesspowerraw} balance audio against other control signals.
\emph{Sound-guided generation} follows more general audio such as music or ambient sound, using its rhythm, pitch, and emotional tone to drive visual elements, motion, and scene transitions: TA2V~\cite{10420468}, MotionCraft~\cite{aira2024motioncraftphysicsbasedzeroshotvideo}, AV-Link~\cite{hajiali2025avlinktemporallyaligneddiffusionfeatures}, and ASVA~\cite{linz2024asva} generate music visualizations or scenes that follow the energy and structure of a track, and MagicInfinite~\cite{yi2025magicinfinitegeneratinginfinitetalking} and Sound2Sight~\cite{cherian2020sound2sightgeneratingvisualdynamics} explore more open-ended audio-driven generation.

\noindent\textbf{Universal and Multi-Task Control.} The single-modality channels each add one encoder; universal frameworks realize the fusion of Eq.~\eqref{eq:p3_s8_3_4_objective} in a single model, along two lines. Autoregressive models built on multimodal large language models generate video from discrete visual tokens, as in VideoPoet~\cite{kondratyuk2024videopoetlargelanguagemodel} and Emu3~\cite{wang2024emu3}. Diffusion models are being adapted to many conditions at once: VideoComposer~\cite{wang2023videocomposercompositionalvideosynthesis} integrates depth, pose, motion, and other conditions by channel concatenation, OmniHuman~\cite{lin2025omnihuman1rethinkingscalinguponestage} generates human video under multiple controls by mixing cross-attention and concatenation, and FullDiT~\cite{ju2025fullditmultitaskvideogenerative} processes all conditions with unified attention, the third realization of $\mathcal{F}_\psi$.

\subsubsection{Benchmarks}
\label{subsubsec:p3_s8_3_4_benchmarks}
Being heterogeneous, the ``other'' modalities share no single suite: each is evaluated against data and metrics specialized to its signal, with the general corpora of §\ref{subsec:p3_s8_3_1} supplying the appearance and motion priors. BEV control is evaluated on driving collections such as nuScenes~\cite{caesar2020nuscenes}, whose paired BEV layouts and front-camera footage furnish both the conditioning signal and the ground truth. Voice-guided talking-head generation is evaluated on portrait-video collections (HDTF~\cite{zhang2021hdtf} and VFHQ~\cite{xie2022vfhq} are standard) and sound-guided generation on audio--video corpora with aligned music or ambient tracks. Style and point control, having no canonical benchmark, are assessed on the general corpora plus qualitative and user studies.

\paragraph{Control-fidelity Metrics.}
Beyond the generic video-quality metrics of §\ref{subsec:p3_s8_3_1} (FVD~\cite{unterthiner2018fvd}, CLIPSIM, and VBench~\cite{huang2023vbench} dimensions), each modality adds an \emph{adherence} metric tailored to its signal, reported together with generic quality as in §\ref{subsec:p3_s8_3_2}.
\emph{Style adherence} combines a CLIP style-similarity to the reference with a temporal-consistency term that penalizes flicker.
\emph{Point adherence} re-localizes the dragged keypoints in the output and reports their $L_2$ deviation from the target path, as in the trajectory metric of §\ref{subsec:p3_s8_3_3}.
\emph{BEV adherence} renders or detects objects in the generated clip and reports detection-mAP and layout-IoU against the input map.
\emph{Text adherence} runs OCR on the rendered frames and reports recognition accuracy and legibility.
\emph{Audio adherence} measures lip-sync quality with SyncNet-style confidence and landmark error for voice, and audio--visual synchrony for sound.

\subsubsection{Challenges and Trends}
\label{subsubsec:p3_s8_3_4_challenges}
The central difficulty here, fusing conditions of incommensurate dimensionality and strength without one washing out or dominating the others, is shared with the spatial and temporal axes and discussed in §\ref{subsec:p3_s8_3_6}. Two problems are specific to these modalities. Text rendering must keep glyphs legible and stable at once, and per-frame legibility and temporal coherence pull against each other as the camera or the text moves. BEV control covers a narrow domain: models trained on a handful of driving datasets inherit their sensor rigs, cities, and traffic statistics, and generalize poorly beyond them.


\subsection{Video Editing}
\label{subsec:p3_s8_3_5}

Video editing takes an \emph{existing} clip and modifies it under user instruction, rather than generating one from scratch (§\ref{subsec:p3_s8_3_1}) or steering that generation with controls (§\ref{subsec:p3_s8_3_2}--§\ref{subsec:p3_s8_3_4}). It is the temporal lift of the image editing of §\ref{subsec:p3_s8_1_3}, and its defining demand is selective preservation: the operator must change only the targeted content while holding the source's geometry, layout, and motion fixed across every frame, so the central technical problem is to balance editability against spatio-temporal consistency without artifacts or unintended drift.

\subsubsection{Task Definition}
\label{subsubsec:p3_s8_3_5_task}
Given a source clip encoded as $\mathbf{z}^{\text{in},1:T}=E(\mathbf{v}^{\text{in}})$ and an edit instruction $\mathbf{c}$, an editing operator
\begin{equation}
\hat{\mathbf{z}}^{1:T}=E_\theta^{\text{vid}}\!\big(\mathbf{z}^{\text{in},1:T},\,\phi(\mathbf{x}),\,\kappa(\mathbf{c})\big)
\label{eq:p3_s8_3_5_objective}
\end{equation}
returns an edited clip that applies the requested change while preserving unedited content and the source motion, recovering the image editing operator of Eq.~\eqref{eq:t2i_edit_objective} exactly at $T=1$. Eq.~\eqref{eq:p3_s8_3_5_objective} encodes two competing objectives: \emph{editability} (the output must realize the instruction $\mathbf{c}$) and \emph{preservation} (the unedited regions and the source dynamics must survive untouched).

\subsubsection{Methodology}
\label{subsubsec:p3_s8_3_5_methodology}
The five paradigms are overlapping choices of training strategy and preservation mechanism: training-free editing keeps the weights fixed, one-shot tuning updates selected parameters on the source clip, trained instruction editors learn across editing data, attention-based editing reuses or modifies source attention, and structural guidance conditions the edit on source-derived depth, edges, or poses. Many of them share a diffusion recipe that inverts the source clip and re-denoises it under the instruction; writing $\mathrm{Inv}$ for DDIM inversion of the pretrained or adapted backbone $\epsilon_\theta^{0}$ of §\ref{subsec:p3_s8_3_1},
\begin{equation}
\mathbf{z}_S^{1:T}=\mathrm{Inv}\big(\mathbf{z}^{\text{in},1:T};\,\phi(\mathbf{x}^{\text{src}})\big),\qquad
\hat{\mathbf{z}}^{1:T}=\mathrm{Denoise}\big(\mathbf{z}_S^{1:T};\,\phi(\mathbf{x}),\,\kappa(\mathbf{c})\big),
\label{eq:p3_s8_3_5_invert}
\end{equation}
where $S$ denotes the terminal diffusion step, distinct from the clip length $T$. Inversion supplies a source-correlated initialization; how much of the source survives then depends on the sampling, feature reuse, and guidance applied during re-denoising. Figure~\ref{fig:p3_s8_3_5_mechanisms} schematizes the five paradigms, with panels (a)--(e) keyed to the paragraphs below.

\begin{figure}[!th]
\centering
\includegraphics[width=\linewidth]{figures/part3/section8/five_video_editing_paradigms.pdf}
\caption{Five overlapping video-editing paradigms, (a)--(e), applied to one task (top): restyling a driving clip while preserving its layout and motion. The dashed link in (c) marks a paired training example; scene sequences are illustrative.}
\label{fig:p3_s8_3_5_mechanisms}
\end{figure}

\noindent\textbf{Training-Free (Zero-Shot) Methods.} The lightest realization of Eq.~\eqref{eq:p3_s8_3_5_invert} (Figure~\ref{fig:p3_s8_3_5_mechanisms}a) edits a video without any fine-tuning, analogous to zero-shot image editing~\cite{ju2023directinversionboostingdiffusionbased}: DDIM inversion encodes the source into latent codes and noise maps, and the denoising process is rerun on them under the edited prompt. Text2Video-Zero~\cite{Khachatryan_2023_ICCV} and Pix2Video~\cite{ceylan2023pix2videovideoeditingusing} established the category, and vid2vid-zero~\cite{wang2024zeroshotvideoeditingusing}, I2VEdit~\cite{ouyang2024i2veditfirstframeguidedvideoediting}, RAVE~\cite{kara2024rave}, and FastVideoEdit~\cite{zhang2024fastvideoeditleveragingconsistencymodels} add latent-code warping and feature injection to improve temporal consistency without fine-tuning.

\noindent\textbf{One-Shot Tuning Methods.} Trading a little training for sharper fidelity (Figure~\ref{fig:p3_s8_3_5_mechanisms}b), Tune-A-Video~\cite{wu2023tune} fine-tunes a limited subset of a pretrained text-to-image diffusion model's parameters on the frames of a single source video, so that the model learns its content and motion and can regenerate it under a new prompt. Dreamix~\cite{molad2023dreamixvideodiffusionmodels} builds on one-shot fine-tuning to improve both editability and fidelity to the original motion.

\noindent\textbf{Training-Based Methods.} The limited control and quality of inference-time editing have motivated editors that learn a reusable operator from editing data (Figure~\ref{fig:p3_s8_3_5_mechanisms}c), particularly for instruction-based editing. InsV2V~\cite{cheng2023consistentvideotovideotransferusing} extends the InstructPix2Pix~\cite{brooks2023instructpix2pixlearningfollowimage} framework directly to video; GenProp~\cite{Liu_2025_CVPR} propagates an edit from a user-modified first frame to the rest of the sequence; and VideoPainter~\cite{videopainter} uses a video inpainting model as an intermediary. UNIC~\cite{ye2025unicunifiedincontextvideo} handles many editing tasks in one model by concatenating the conditions into a sequence with task-aware positional embeddings, and EditVerse~\cite{ju2025editverseunifyingimagevideo} unifies video generation and editing by representing text, image, and video as one token sequence processed with self-attention, which permits in-context learning and flexible inputs.

\noindent\textbf{Attention-Based Editing.} To protect preservation within the inference-time recipe (Figure~\ref{fig:p3_s8_3_5_mechanisms}d), many techniques manipulate the cross-attention maps directly: following Prompt-to-Prompt in images, they blend the attention maps of the original and edited prompts so that the inverted source structure of Eq.~\eqref{eq:p3_s8_3_5_invert} survives the re-denoising. FateZero~\cite{qi2023fatezerofusingattentionszeroshot} and Video-P2P~\cite{liu2023videop2p} were early examples, and Edit-A-Video~\cite{shin2023editavideosinglevideoediting}, TokenFlow~\cite{tokenflow2023}, UniEdit~\cite{bai2024uniedit}, and AnyV2V~\cite{ku2024anyv2vtuningfreeframeworkvideotovideo} improve temporal consistency at the token or feature level, enabling localized edits such as changing an object's material or color while the scene structure is kept.

\noindent\textbf{Editing with Structural Guidance.} At the explicitly guided end (Figure~\ref{fig:p3_s8_3_5_mechanisms}e), editing is conditioned on spatial signals extracted from the source video, the conditioning machinery of §\ref{subsec:p3_s8_3_2} repurposed for preservation rather than generation. Rerender A Video~\cite{yang2023rerender} feeds the source's depth or edge maps to a ControlNet-like module during editing so that geometry and layout stay intact, and MagicEdit~\cite{liew2023magicedit} accepts a variety of spatial inputs (\eg, pose, sketches, segmentation maps) for high structural fidelity, which suits style transfer and object replacement, where preserving structure is what prevents distortion and keeps motion coherent.

The paradigms trade cost against control: training-free editing avoids per-clip weight updates but may still need iterative inference, whereas a trained instruction editor is expensive to build but generalizes across edits.

\subsubsection{Benchmarks}
\label{subsubsec:p3_s8_3_5_benchmarks}
Video editing has no single canonical suite, and evaluation has moved from shared source-video collections toward purpose-built, fine-grained editing benchmarks. Early evaluations report on DAVIS~\cite{pont20172017}, the de facto source set for object-centric edits, and on the LOVEU text-guided video-editing (TGVE) benchmark~\cite{wu2023loveutgve}, which supplies prompt--clip pairs with human ratings; VPBench, released with VideoPainter~\cite{videopainter}, adds masked clips with dense captions for inpainting-style edits. The most recent suites target instruction-guided editing and pair richer data with model-based scoring. FiVE~\cite{li2025five} contributes $74$ real clips (drawn from DAVIS) and $26$ generated clips across six fine-grained editing types, with $420$ object-level prompt pairs and masks, evaluated under a $15$-metric protocol spanning background preservation, text--video similarity, temporal consistency, video quality, and runtime. IVEBench~\cite{chen2026ivebench} scales this to $600$ source videos spanning seven semantic dimensions and clip lengths from $32$ to $1{,}024$ frames, across eight task categories ($35$ subcategories) whose prompts are generated by LLMs and refined by experts, and adopts a three-dimensional protocol (video quality, instruction compliance, and video fidelity) that fuses traditional metrics with multimodal-LLM-based assessment.

\paragraph{Editing-fidelity Metrics.}
Editing is judged on three axes that mirror the editability--preservation tension of Eq.~\eqref{eq:p3_s8_3_5_objective}. The classic metrics are, for \emph{edit adherence}, the CLIP similarity between each edited frame and the target prompt; for \emph{temporal consistency}, frame-to-frame CLIP similarity together with a warping error that flow-aligns adjacent frames and reports the photometric residual; and for \emph{structure and background preservation}, PSNR / SSIM / LPIPS on the unedited regions against the source, atop the generic FVD~\cite{unterthiner2018fvd} quality score. These scalars correlate only loosely with human judgment on fine-grained edits (an object recolored but mislocalized can still score well), so the latest benchmarks add vision-language judges: FiVE-Acc~\cite{li2025five} uses a VLM to test whether the intended fine-grained edit succeeded, and IVEBench~\cite{chen2026ivebench} scores instruction compliance and fidelity with multimodal LLMs alongside the traditional metrics. Results are reported as an edit-versus-preservation pair, since ignoring the instruction trivially maximizes preservation.

\subsubsection{Challenges and Trends}
\label{subsubsec:p3_s8_3_5_challenges}
Three problems are specific to editing. Pushing the edit harder erodes the unedited regions and the source motion, while protecting them too strongly suppresses the requested change; this editability--preservation trade-off of Eq.~\eqref{eq:p3_s8_3_5_objective} is what every paradigm above negotiates differently. The invert-and-re-denoise recipe of Eq.~\eqref{eq:p3_s8_3_5_invert} is imperfect, and its reconstruction error compounds over long clips, so edits that are clean over a few seconds degrade over minute-scale video. And training-based editors follow only the edit types seen in their training pairs, so a flexible instruction editor that handles arbitrary edits remains open; the unified generate-and-edit models UNIC~\cite{ye2025unicunifiedincontextvideo} and EditVerse~\cite{ju2025editverseunifyingimagevideo}, which treat the source clip as one stream among many in shared attention, are the current attempt.







\subsection{Summary and Critical Challenges and Future Directions}
\label{subsec:p3_s8_3_6}
\label{subsec:p3_s8_3_summary}

Table~\ref{tab:p3_s8_3_subfield_comparison} places the five control axes side by side: the control signal each introduces, the operator it instantiates over the shared sampler of Eq.~\eqref{eq:p3_s8_3_objective}, the paradigms it surveys, the methods anchoring each, and the benchmarks on which it reports.

\begin{table}[t]
    \centering
    \fontsize{8}{10}\selectfont
    \setlength{\tabcolsep}{1.0mm}
    \caption{The five control axes of video spatial generation; each operator specializes the spatio-temporal sampler of Eq.~\eqref{eq:p3_s8_3_objective}.}
    \label{tab:p3_s8_3_subfield_comparison}
    \begin{tabular}{L{0.12\textwidth} L{0.22\textwidth} L{0.18\textwidth} L{0.19\textwidth} L{0.21\textwidth}}
        \toprule
        \textbf{Control axis} & \textbf{Control signal / operator} & \textbf{Dominant paradigms} & \textbf{Canonical methods} & \textbf{Benchmarks / metrics} \\
        \midrule
        §\ref{subsec:p3_s8_3_1} Text-to-video & Text only, $\kappa=\varnothing$; $G_\theta^{\text{vid}}:\mathcal{E}\!\to\!\mathcal{Z}^{\text{vid}}$ (Eq.~\eqref{eq:p3_s8_3_1_objective}) & GAN; autoregressive token; diffusion / diffusion transformer & CogVideo, Sora, Wan, HunyuanVideo, Cosmos & UCF-101, MSR-VTT, VBench, EvalCrafter; FVD, CLIPSIM \\
        §\ref{subsec:p3_s8_3_2} Spatial control & Per-frame scaffold $\kappa_{\text{s}}(\mathbf{c}_{\text{s}}^{1:T})$ (Eq.~\eqref{eq:p3_s8_3_2_objective}); residual injection (Eq.~\eqref{eq:p3_s8_3_2_inject}) & Image, pose, depth, landmark, sketch, bounding box & Animate Anyone, MagicAnimate, ControlVideo, EchoMimic, MagicDrive & nuScenes, pose-annotated video; keypoint / depth / identity / detection adherence \\
        §\ref{subsec:p3_s8_3_3} Temporal control & Motion signal $\kappa_{\text{t}}(\mathbf{m}^{1:T})$ (Eq.~\eqref{eq:p3_s8_3_3_objective}); temporal-branch injection (Eq.~\eqref{eq:p3_s8_3_3_inject}) & Flow, trajectory, camera, motion; latent warping or encoded conditioning & Motion-I2V, Tora, DragNUWA, CameraCtrl, MotionDirector & RealEstate10K, DL3DV-10K, MPI-Sintel; EPE, path / pose error, transfer fidelity \\
        §\ref{subsec:p3_s8_3_4} Other control & Heterogeneous $\kappa_{\text{o}}$ fused by $\mathcal{F}_\psi$ (Eqs.~\eqref{eq:p3_s8_3_4_objective}, \eqref{eq:p3_s8_3_4_inject}) & Style, point, BEV, text, audio; summed adapters, concatenation, unified attention & StyleCrafter, DriveDreamer-2, EMO, VideoComposer, FullDiT & nuScenes, HDTF, VFHQ; style / point / BEV / OCR / lip-sync adherence \\
        §\ref{subsec:p3_s8_3_5} Video editing & Existing clip $+$ edit intent $\to$ edited clip; $E_\theta^{\text{vid}}$ (Eq.~\eqref{eq:p3_s8_3_5_objective}) & Training-free, one-shot tuning, trained instruction editors, attention reuse, structural guidance & Tune-A-Video, TokenFlow, InsV2V, UNIC, EditVerse & DAVIS, LOVEU-TGVE, FiVE, IVEBench; edit-vs-preservation tuple, VLM / MLLM judges \\
        \bottomrule
    \end{tabular}
\end{table}

The axes partition not by output domain, since all five decode a clip from the same $\mathcal{Z}^{\text{vid}}$, but by the control fed to one shared sampler: text-to-video is the $\kappa=\varnothing$ limit, spatial, temporal, and other control progressively enrich $\kappa$, and editing alone reverses the direction of the map, taking an existing clip in place of pure noise.
The injection mechanism converges across axes: the spatial residual branch (Eq.~\eqref{eq:p3_s8_3_2_inject}), the temporal branch (Eq.~\eqref{eq:p3_s8_3_3_inject}), and the multi-condition interface (Eq.~\eqref{eq:p3_s8_3_4_inject}) are one idea, a trainable control pathway added to a pretrained backbone, generalized from one condition to many, and the most recent systems (VideoComposer, FullDiT, UNIC, EditVerse) go further by folding every axis, editing included, into a single token sequence processed by shared attention.
Evaluation has not converged: no suite spans the five axes, each reports generic quality together with a channel-specific adherence score, and the newest editing benchmarks (FiVE, IVEBench) add vision-language and multimodal-LLM judges to improve human alignment.

Video generation inherits every difficulty of the 2D and 3D priors of §\ref{sec:p3_s8_1} and §\ref{sec:p3_s8_2} and adds the time axis, which turns several already-hard problems into open ones; the four below cut across all five axes.

\noindent\textbf{Temporal coherence and long-horizon drift.} The defining difficulty of video over image generation is keeping content stable across frames. The signature artifact is \emph{flicker}, the high-frequency appearance instability between adjacent frames that plagues per-frame and early GAN-based methods (§\ref{subsec:p3_s8_3_1}); over longer spans it gives way to \emph{identity drift}, the gradual mutation of a subject's appearance, color, or fine texture across a clip; and in autoregressive roll-outs to \emph{error accumulation}, the compounding of small per-step mistakes that caps usable clip length. Latent video autoencoders with temporal compression~\cite{blattmann2023stable,blattmann2023align} and dedicated spatio-temporal attention suppress flicker, while autoregressive--diffusion hybrids such as CausVid~\cite{yin2025slow} and MAGI-1~\cite{ai2025magi1autoregressivevideogeneration} combine sequential consistency with diffusion fidelity to push toward minute-scale generation; coherence over arbitrarily long horizons remains unsolved and is the binding constraint on the world-model applications that motivate the section.

\noindent\textbf{Controllability versus fidelity under multi-condition fusion.} As §\ref{subsec:p3_s8_3_2}--§\ref{subsec:p3_s8_3_4} stack heterogeneous conditions (spatial scaffolds, camera paths, audio, style) onto Eq.~\eqref{eq:p3_s8_3_objective}, the conditions begin to conflict. Over-strong conditioning yields rigidity and texture collapse; weak conditioning ignores the control; and naive fusion produces \emph{control-signal dilution} (summed adapters wash one another out), \emph{condition leakage} across frames, and \emph{collapse to the dominant condition} when one modality (\eg, audio in talking-head generation) overwhelms the rest. Compositional and universal frameworks (VideoComposer~\cite{wang2023videocomposercompositionalvideosynthesis}, OmniHuman~\cite{lin2025omnihuman1rethinkingscalinguponestage}, and the unified-attention design of FullDiT~\cite{ju2025fullditmultitaskvideogenerative}) replace per-axis injection with joint token fusion and task-aware positional embeddings, but a principled account of how to weight and arbitrate competing conditions is still missing.

\noindent\textbf{Physical and 3D plausibility.} Even a temporally coherent, well-controlled clip may depict an impossible world, which is the bridge back to §\ref{sec:p3_s8_2}. Generated motion frequently violates object permanence (objects teleport, morph, or vanish), contact and rigid-body dynamics (rigid objects warp non-rigidly; feet slide; occlusions resolve implausibly), and \emph{multi-view consistency} under camera motion, where the same scene is rendered inconsistently from successive viewpoints. Systems presented as steps toward ``world simulators'', Sora~\cite{sora2024} among them, still exhibit these failures on closer inspection. Camera-control methods~\cite{he2024cameractrl,yu2024viewcrafter} and explicitly 3D-aware video diffusion (Diffusion-as-Shader~\cite{gu2025diffusionshader3dawarevideo} and GS-DiT~\cite{bian2025gsditadvancingvideogeneration}) narrow the gap by conditioning generation on geometry, which suggests that the more promising route to physical plausibility runs through the volumetric representations of §\ref{sec:p3_s8_2} rather than through scale alone.

\noindent\textbf{Compute, latency, and evaluation.} Underlying all three is cost: a spatio-temporal latent is an order of magnitude larger than an image latent, so long, high-resolution clips strain memory and make iterative sampling slow, which in turn limits the context length available for coherence and the conditioning budget available for control. Consistency distillation and few-step samplers, together with sparse and causal attention, are the active responses. Evaluation compounds the problem: distributional scores such as FVD and per-frame FID are blind to temporal and control adherence, and automatic metrics correlate weakly with human judgment, so reported progress is hard to compare across methods; benchmarks and human-aligned protocols that jointly score fidelity, temporal coherence, and per-axis control adherence are a central near-term need.

Coherence and physical plausibility both demand longer temporal context, which inflates compute, and every added condition enlarges the evaluation surface, so progress on any one of the four is bounded by the others. Together they mark the last step of the progression traced across §\ref{sec:generation}, from the static 2D prior of §\ref{sec:p3_s8_1} through the static volumetric prior of §\ref{sec:p3_s8_2} to dynamic, navigable models that consume spatial structure over time.

% \newpage


\section{Motion Generation}
\label{sec:p3_s8_4}

\FramedInitial{M}{otion} generation is the temporal, embodied output channel of spatial generation: where §\ref{sec:p3_s8_1}–§\ref{sec:p3_s8_3} synthesize static images, volumetric assets, and controllable video, motion generation produces the \emph{articulated dynamics} that an avatar, character, or robot can actually enact.
A generated motion is a time series of body poses that must be believable in its dynamics, physically plausible (balance, ground contact, non-penetration), and faithful to its driving signal, whether a text description, audio, a scene, or a key frame.
The conditional-sampler template that turns a prompt into pixels (Eq.~\eqref{eq:t2i_generator}) is here instantiated over a sequence of poses in a structured kinematic space, so geometric priors meet physical constraints: the output must respect the biomechanics of an articulated body as well as the semantics of the condition.

The subsections are indexed by the moving agent and then by one cross-cutting operation: human motion (§\ref{subsec:p3_s8_4_1}), animal motion with its cross-species morphology (§\ref{subsec:p3_s8_4_2}), vehicle motion planning, where the body is rigid and the difficulty moves into the multi-agent traffic scene (§\ref{subsec:p3_s8_4_3}), and motion editing (§\ref{subsec:p3_s8_4_4}).
% , with a comparison in §\ref{subsec:p3_s8_4_summary}.




\subsection{Human Motion Generation}
\label{subsec:p3_s8_4_1}
Human motion generation produces natural, varied sequences of human body poses from intentions, environmental stimuli, or explicit conditions, and supports animation, virtual characters in film and games, AR/VR, and human--robot interaction.

\subsubsection{Task Definition}
\label{subsubsec:p3_s8_4_1_task}
Let a motion be a temporal sequence of body poses $\mathbf{m}^{1:T}=(\mathbf{m}^1,\dots,\mathbf{m}^T)$, where each frame $\mathbf{m}^t\in\mathcal{P}$ is a pose in a chosen representation space $\mathcal{P}$ and $\mathcal{M}=\mathcal{P}^T$ is the resulting motion space. Writing $\phi$ for the text encoder of Eq.~\eqref{eq:t2i_generator} and $\kappa:\mathcal{C}\!\to\!\mathcal{C}'$ for an encoder of a driving signal $\mathbf{c}\in\mathcal{C}$ (an action label, a text description, a music or speech track, or a scene), human motion generation learns a conditional sampler
\begin{equation}
\label{eq:p3_s8_4_1_objective}
G_\theta^{\text{mot}}:\mathcal{C}'\!\to\!\mathcal{M},\qquad
\hat{\mathbf{m}}^{1:T}\sim p_\theta\!\left(\mathbf{m}^{1:T}\mid\kappa(\mathbf{c})\right).
\end{equation}
This is the conditional-sampler template of Eq.~\eqref{eq:t2i_generator} instantiated over a temporal pose sequence rather than an image latent: with no driving signal ($\kappa=\varnothing$) it reduces to unconditional motion synthesis, and the choice of $\mathbf{c}$ indexes the five conditioning settings surveyed in §\ref{subsubsec:p3_s8_4_1_method}.

Figure~\ref{fig:p3_s8_4_1_repr} illustrates the pose representations that instantiate $\mathcal{P}$. Broadly, representations fall into keypoint-based and rotation-based categories, and the two are convertible: joint rotations map to spatial keypoints via forward kinematics (FK), and keypoints recover joint rotations via inverse kinematics (IK).

\begin{figure}
    \centering
    \includegraphics[width=\linewidth]{figures/part3/section8/human_motion_repr.png}
    \includegraphics[width=\linewidth]{figures/part3/section8/human_motion_repr_smpl.png}
    \caption{Human motion representations. Top: keypoint-based, as (a) joint positions, (b) joints grouped by kinematic chain, and (c) dense surface markers; and rotation-based, as (d) per-joint rotations in $\mathrm{SO}(3)$ on a kinematic tree. Bottom: the SMPL body model, from (a) the template mesh and skinning weights through (b) shape and (c) pose blend shapes to (d) the skinned, posed mesh.}
    \label{fig:p3_s8_4_1_repr}
\end{figure}

\paragraph{Keypoint-based Representation.} The body is modeled as a collection of keypoints, each a distinct anatomical landmark (joints or salient surface locations) given by its 2D or 3D coordinates in a chosen frame (pixel space or world space), and motion is encoded as a temporal series of such configurations.
Keypoint representations are directly available from many motion-capture systems and are highly interpretable, but many downstream uses (character animation, robotic control) require converting keypoints to joint rotations through the IK problem.
Some works instead augment or replace sparse skeletal keypoints with dense body markers distributed across the body surface~\cite{WuWZZHYT22, TaheriCBT22}, which carry richer information about body shape, limb twisting, and fine geometric detail.

\paragraph{Rotation-based Representation.} A pose is represented by joint angles, the rotations of body segments relative to their parent in a hierarchical kinematic tree. Most studies model these in three dimensions, treating each non-root joint as an element of $\mathrm{SO}(3)$, parameterized as Euler angles, axis--angle vectors, or unit quaternions.

\paragraph{Skinned Multi-Person Linear Model.} Beyond pure joint rotation, statistical mesh models jointly capture pose and body shape, including non-rigid deformations associated with articulation.
The Skinned Multi-Person Linear (SMPL) model~\cite{loper2015smpl} defines a template mesh with fixed topology (6{,}890 vertices, 23 joint rotations plus a root) and uses linear blend skinning augmented by learned blend shapes: identity (shape) blend shapes for inter-subject variation and pose blend shapes for rotation-dependent deformations such as muscle bulging and twisting.



\subsubsection{Methodology}
\label{subsubsec:p3_s8_4_1_method}
Methods for human motion generation are organized along two axes. The first is the \emph{learning paradigm}: regression-based techniques learn a deterministic mapping from conditions to a single target motion, whereas generative approaches model the full conditional distribution of motion using GANs, VAEs, normalizing flows, or diffusion models (the generative paradigms themselves are detailed in §\ref{subsec:p3_s8_1_1}).
Formally, a regression model realizes Eq.~\eqref{eq:p3_s8_4_1_objective} as a point estimate,
\begin{equation}
\label{eq:p3_s8_4_1_regression}
\hat{\mathbf{m}}^{1:T}=f_\theta\!\left(\kappa(\mathbf{c})\right),\qquad
\min_\theta\ \big\|\hat{\mathbf{m}}^{1:T}-\mathbf{m}^{1:T}\big\|,
\end{equation}
which is simple and stable but collapses the inherently one-to-many condition--motion relationship onto an averaged trajectory. Generative models instead sample $\hat{\mathbf{m}}^{1:T}\!\sim\!p_\theta(\mathbf{m}^{1:T}\mid\kappa(\mathbf{c}))$; the now-dominant diffusion instantiation trains a denoiser $G_\theta$ to recover the clean motion at every noise level $n$,
\begin{equation}
\label{eq:p3_s8_4_1_diffusion}
\mathcal{L}=\mathbb{E}_{\mathbf{m}^{1:T},\,n}\,\big\|\,\mathbf{m}^{1:T}-G_\theta(\mathbf{m}^{1:T}_n,\,n,\,\kappa(\mathbf{c}))\,\big\|^2,
\end{equation}
which preserves diversity and supports flexible conditioning, at higher sampling cost.
The second axis is the \emph{conditioning modality} $\mathbf{c}$, which indexes the five settings below; each setting instantiates the same Eq.~\eqref{eq:p3_s8_4_1_objective} and has migrated, broadly, from regression and adversarial models toward diffusion. Figure~\ref{fig:p3_s8_4_1_methods} previews this two-axis organization.



\begin{figure}
    \centering
    \includegraphics[width=\linewidth]{figures/part3/section8/two_axis_human_motion_generation.pdf}
    \caption{Two-axis organization of human motion generation: the conditioning modality (left), encoded by $\kappa$, and the learning paradigm (right), either a regressor that returns one estimate (Eq.~\eqref{eq:p3_s8_4_1_regression}) or a generative model that returns alternative samples, shown with diffusion (Eq.~\eqref{eq:p3_s8_4_1_diffusion}). Branches are alternatives, not consecutive stages.}
    \label{fig:p3_s8_4_1_methods}
\end{figure}

\paragraph{Action-to-Motion.} Given a discrete action label (\eg, ``walk'', ``throw''), these methods generate a corresponding sequence. Early work used GANs with graph convolution and self-attention (SA-GAN~\cite{YuZLYC20}, Kinetic-GAN~\cite{DegardinNLBY022}) for realism and diversity, while VAE/CVAE frameworks (Action2Motion~\cite{GuoZWZSDG020}, ACTOR~\cite{PetrovichBV21}) produced probabilistic, controllable motions. More recent models adopt contrastive latent learning (ODMO~\cite{LuZLR22}), discrete autoregressive sequences (PoseGPT~\cite{LucasBWR22}), and diffusion (MDM~\cite{TevetRGSCB23}, the latent-diffusion hybrid MLD~\cite{ChenJLHFCY23}). As most methods excel only at single actions, MultiAct~\cite{LeeML23} recurrently composes multiple action motions into long sequences via a canonicalization scheme.

\paragraph{Text-to-Motion.} Conditioning on free-form language exploits richer semantic cues than fixed labels allow. Early methods used GANs (Text2Action~\cite{AhnHCYO18}) or joint text--motion embeddings (JL2P~\cite{AhujaM19}, two-stream encoder--decoders~\cite{GhoshCOTS21}); later ones added variational modeling and attention (TEMOS~\cite{PetrovichBV22}, TEACH~\cite{AthanasiouPBV22}).
For zero-shot generalization, MotionCLIP~\cite{TevetGHBC22} and AvatarCLIP~\cite{HongZPCYL22} align the motion latent space with CLIP.
The current frontier combines VQ-VAE tokenization~\cite{GuoZWC22}, autoregressive transformers (T2M-GPT~\cite{zhang2023generating}), and diffusion (FLAME~\cite{KimKC23}), with multimodal conditioning extensions such as MoFusion~\cite{DabralMGT23} fusing text and audio. More recent masked-generative models (MoMask~\cite{GuoMJW024}) push quality and inference speed further by predicting quantized motion tokens with a bidirectional, iteratively-refined masking schedule rather than left-to-right.

\paragraph{Music-to-Dance.} The goal is coherent, expressive dance aligned to a music track. Early approaches framed this as supervised sequence-to-sequence translation via LSTM autoencoders~\cite{TangJM18}, transformers (AI Choreographer~\cite{LiYRK21}), and multi-person frameworks (GroupDancer~\cite{Wang0WXCMCW22}, GDancer~\cite{LePDTT023}).
Because music--dance is inherently one-to-many, generative models model the diverse distribution: MNET~\cite{0001OKT022} with style conditioning, Transflower~\cite{PerezHBHOA21}, Bailando~\cite{SiyaoYGLW0L022} (VQ-VAE + GPT), and EDGE~\cite{TsengCL23} (diffusion). Graph-based methods (ChoreoMaster~\cite{ChenTLZGZ021}, PC-Dance~\cite{GaoPZSZ22}, ChoreoGraph~\cite{AuCJG22}) use motion graphs and shared music--motion embeddings for rhythmic coherence, while curriculum learning, dynamics-prior denoising, and hierarchical pose/motif/choreography modeling target long-term, genre-consistent full-song dances~\cite{HuangHWSZJ21, SunWHLJH22, AristidouYACSC23}.

\paragraph{Speech-to-Gesture.} Co-speech gesture synthesis generates upper-body or full gesture sequences from speech (audio, sometimes text).
Early adversarial methods modeled person-specific gesture patterns~\cite{GinosarBKCOM19}; subsequent work disentangled shared versus motion-specific variation~\cite{KucherenkoHHKK19} and decomposed body, hands, and face~\cite{HabibieXMLSPET21}.
Recent methods employ diffusion~\cite{ZhuLLQLY23, AlexandersonNBH23, AoZL23}, VQ-VAE quantization~\cite{YiLLCWBTB23}, and style or speaker embeddings~\cite{AhujaLNM20} to handle diversity, identity, rhythm, and expression more robustly.

\paragraph{Scene-to-Motion.} Here the condition is a 3D scene, and most methods adopt a multi-stage pipeline that predicts a goal or interaction anchor, plans a trajectory, then generates motion along it. Early works predict a goal from 2D/RGB scene inputs and synthesize motion through trajectory-guided VAEs~\cite{CaoGMCVM20, WangYDL21}, sometimes with scene-aware discriminators. SAMP~\cite{HassanCVSYZB21} models 3D paths and human--object interactions; later methods refine interaction with per-joint contact maps~\cite{D00010HS22} or path-search ($A^\ast$) for trajectory diversity~\cite{0003RLYLD22}, and GOAL~\cite{TaheriCBT22} and SAGA~\cite{WuWZZHYT22} target grasping via two-step goal-to-motion VAEs.
Some assume given goals and stitch sub-goal motions hierarchically for long-term synthesis~\cite{WangXXL021}, while refinement-based frameworks such as CIRCLE~\cite{AraujoLVAWGCL23} initialize by interpolation and refine with scene-geometry-aware modules.


\subsubsection{Benchmarks}
\label{subsubsec:p3_s8_4_1_benchmarks}
Human motion datasets group naturally by conditioning modality; Table~\ref{tab:p3_s8_4_1_datasets} gives their representation, scale, and condition.
Text- and action-conditioned corpora dominate, led by the language-annotated KIT Motion-Language~\cite{PlappertMA16} and HumanML3D~\cite{GuoZZ0JL022}, BABEL~\cite{PunnakkalCAQB21}'s sequence- and frame-level labels over the AMASS~\cite{mahmood2019amass} aggregate, and the action taxonomies of UESTC~\cite{JiXYSSZ18}, NTU-RGB+D~120~\cite{LiuSPWDK20}, and HumanAct12~\cite{GuoZWZSDG020}.
Audio-conditioned resources split into music-to-dance sets (AIST++~\cite{LiYRK21}, AIOZ-GDANCE~\cite{LePDTT023}, and several genre-specific collections) and co-speech gesture sets (Trinity~\cite{FerstlM18}, BEAT~\cite{LiuZIPLZBZ22}, PATS~\cite{AhujaLNM20}), many built by fitting parametric body models to in-the-wild video.
Scene-conditioned datasets (PROX~\cite{HassanCTB19}, GRAB~\cite{TaheriGBT20}, HUMANISE~\cite{WangCLZLH22}, CIRCLE~\cite{AraujoLVAWGCL23}) capture human--scene and human--object interaction with synchronized geometry, from Vicon and IMU rigs to synthetic engines such as GTA-IM~\cite{CaoGMCVM20}. Across all three groups, the cost of high-quality capture forces a trade-off between scale and annotation fidelity.

\begin{table}[htbp]
\centering
\fontsize{8}{10}\selectfont
\setlength{\tabcolsep}{1.2mm}
\caption{Benchmarks for human motion generation, grouped by condition (none, text / action, music, speech, scene/object).}
\label{tab:p3_s8_4_1_datasets}
\begin{tabular}{L{0.18\textwidth} C{0.115\textwidth} C{0.07\textwidth} C{0.08\textwidth} C{0.06\textwidth} C{0.06\textwidth} C{0.13\textwidth} L{0.17\textwidth}}
\toprule
\textbf{Dataset} & \textbf{Representation} & \textbf{Subjects} & \textbf{Sequences} & \textbf{Frames} & \textbf{Length} & \textbf{Condition} & \textbf{Remarks} \\
\midrule
Human3.6M \cite{IonescuPOS14} & Keypoint & 11 & - & 3.6M & 5h & - & 15 actions \\
AMASS \cite{mahmood2019amass} & Rotation & 344 & 11,265 & - & 40h & - & Unifies 15 marker-based MoCap datasets \\
HuMMan \cite{CaiRZLYWFGYPHZL22} & Rotation & 1000 & 400K & 60M & - & - & 500 actions \\
KIT Motion Language \cite{PlappertMA16} & Keypoint & 111 & 3,911 & - & 10.3h & Text & 6.3K text descriptions \\
UESTC \cite{JiXYSSZ18} & Keypoint & 118 & 25.6K & - & 83h & Action & 40 action classes \\
NTU-RGB+D 120 \cite{LiuSPWDK20} & Keypoint & 106 & 114.5K & - & 74h & Action & 120 action classes \\
HumanAct12 \cite{GuoZWZSDG020} & Keypoint & 12 & 1,191 & 90K & 6h & Action & 12 action classes \\
BABEL \cite{PunnakkalCAQB21} & Rotation & 344 & - & - & 43.3h & Text & 260 action classes \\
HumanML3D \cite{GuoZZ0JL022} & Keypoint + Rotation & 344 & 14.6K & - & 28.5h & Text & 44.9K text descriptions \\
DanceWithMelody \cite{TangJM18} & Keypoint & - & 61 & 907K & 1.6h & Music & 4 genres \\
Lee \cite{LeeY0WLYK19} & Keypoint & - & 361K & - & 71h & Music & 3 genres \\
Huang \cite{HuangHWSZJ21} & Keypoint & - & 790 & - & 12h & Music & 3 genres\\
AIST++ \cite{LiYRK21} & Rotation & 30 & 1,408 & 10.1M & 5.2h & Music & 10 genres \\
PMSD \cite{PerezHBHOA21} & Keypoint & 8 & - & - & 3.8h & Music & 4 genres \\
ShaderMotion \cite{PerezHBHOA21} & Keypoint & 8 & - & - & 10.2h & Music & 2 genres \\
Chen \cite{ChenTLZGZ021} & Rotation & - & - & 160K & 9.9h & Music & 9 genres \\
PhantomDance \cite{LiZZS22} & Rotation & - & 260 & 795K & 3.7h & Music & 13 genres \\
MMD-ARC \cite{GaoPZSZ22} & Rotation & - & 213 & - & 11.3h & Music & - \\
MDC \cite{Wang0WXCMCW22} & Rotation & - & 798 & - & 3.5h & Music & 2 genres \\
Aristidou \cite{AristidouYACSC23} & Rotation & 32 & - & - & 2.4h & Music & 3 genres \\
AIOZ-GDANCE \cite{LePDTT023} & Rotation & >4000 & - & - & 16.7h & Music & 7 styles, 16 genres \\
Trinity \cite{FerstlM18} & Keypoint & 1 & 23 & - & 4.1h & Speech & Casual talks\\
TED-Gesture \cite{YoonKJLKL19} & Keypoint & - & 1,295 & - & 52.7h & Text & TED talks \\
Speech2Gesture \cite{GinosarBKCOM19} & Keypoint & 10 & - & - & 144h & Speech & TV shows, lectures \\
TED-Gesture++ \cite{YoonCLJLKL20} & Keypoint & - & 1,766 & - & 97.0h & Speech, Text & Extension of \cite{YoonKJLKL19} \\
PATS \cite{AhujaLNM20} & Keypoint & 25 & - & - & 251h & Speech, Text & Extension of \cite{GinosarBKCOM19} \\
Speech2Gesture-3D \cite{HabibieXMLSPET21} & Keypoint & 6 & - & - & 33h & Speech & Videos from \cite{GinosarBKCOM19}\\
BEAT \cite{LiuZIPLZBZ22} & Rotation & 30 & 2,580 & 30M & 76h & Speech, Text, Emotion & 8 emotions, 4 languages \\
Chinese Gesture \cite{AoGLCL22} & Rotation & 5 & - & - & 4h & Speech, Text & Chinese \\
ZEGGS \cite{GhorbaniFHTC23} & Rotation & 1 & 67 & - & 2.3h & Speech, Style & 19 styles \\
SHOW \cite{YiLLCWBTB23} & Rotation & - & - & - & 27h & Speech & Videos from \cite{GinosarBKCOM19} \\
WBHM \cite{ManderyTDVA15} & Rotation & 43 & 3,704 & 691K & 7.68h & Object & 41 objects \\
PiGraph \cite{Hanrahan16a} & Keypoint & 5 & 63 & 0.1M & 2h & Scene, Object & 30 scenes, 19 objects \\
PROX \cite{HassanCTB19} & Rotation & 20 & 60 & 0.1M & 1h & Scene, Object & 12 indoor scenes\\
i3DB \cite{MonszpartGCYM19} & Keypoint & 1 & - & - & - & Scene, Object & 15 scenes\\
GTA-IM \cite{CaoGMCVM20} & Keypoint & 50 & 119 & 1M & - & Scene & Synthetic, 10 indoor scenes \\
GRAB \cite{TaheriGBT20} & Rotation & 10 & 1,334 & 1.6M & - & Object & 51 objects \\
HPS \cite{GuzovMSP21} & Rotation & 7 & - & 300K & - & Scene & 8 large scenes\\
SAMP \cite{HassanCVSYZB21} & Rotation & 1 & - & 185K & 0.83h & Scene, Object & 7 objects\\
COUCH \cite{ZhangBSGP22} & Rotation & 6 & >500 & - & 3h & Scene, Chairs & 3 chairs\\
HUMANISE \cite{WangCLZLH22} & Rotation & - & 19.6K & 1.2M & - & Scene, Object, Text & 643 scenes \\
CIRCLE \cite{AraujoLVAWGCL23} & Rotation & 5 & >7K & 4.3M & 10h & Scene & 9 scenes \\
\bottomrule
\end{tabular}
\end{table}

\paragraph{Evaluation Metrics.} Evaluating synthesized motion is hard because the condition--motion map is one-to-many, judgments are partly subjective, and quality depends on high-level conditional cues, so the literature reports an ensemble of metrics along three axes.
\textbf{Fidelity} captures naturalness, smoothness, and plausibility: distance-based measures such as NPSS~\cite{GopalakrishnanM19} and NDMS~\cite{TankeZG21}, keypoint accuracy (PCK~\cite{YangR13}), motion-space statistics (AVE, Hellinger distance over speed histograms~\cite{yang2023qpgesture}, Fréchet Distance over pose rotations~\cite{HassanCVSYZB21, 0003RLYLD22}, PSKL~\cite{RuizGM19}), feature-space distributional distances led by FID (with IS~\cite{AhujaLNM20}, MMD~\cite{YuZLYC20}, and related scores), and physical-plausibility proxies for foot sliding and foot--ground contact (FCS~\cite{TsengCL23}, non-collision ratios~\cite{WuWZZHYT22}, contact accuracy~\cite{ZhangT22}), the last still lacking a standardized definition.
\textbf{Diversity} guards against the long-sequence ``freezing'' failure via intra-motion variation and freezing rate~\cite{SunWHLJH22}, inter-motion manifold coverage, and multimodality metrics (n-APD, ADE/FDE and their multimodal variants~\cite{YuanK20}) computed by sampling many motions per condition.
\textbf{Condition consistency} measures alignment to the driving signal: recognition accuracy and R-Precision for text/action, feature-level text--motion distance (Multimodal Distance~\cite{GuoZZ0JL022}, mCLIP and MID~\cite{KimKC23}), audio--motion beat metrics (beat coverage, hit rate~\cite{LeeY0WLYK19}, beat alignment~\cite{LiYRK21}, SRGR~\cite{LiuZIPLZBZ22}, Semantic Score~\cite{AoZL23}), and, for scene conditioning, non-collision and human--scene contact scores~\cite{WuWZZHYT22, ZhangBSGP22}. No single metric is decisive, so methods are compared on a task-appropriate composition of the three.


\subsubsection{Challenges and Trends}
\label{subsubsec:p3_s8_4_1_challenges}
\textbf{Data scarcity} is the most basic constraint: high-quality motion capture is expensive and slow, so dataset scale trades against annotation fidelity, and heterogeneous pose representations and conditioning signals fragment the data that exist; fusing heterogeneous sources through weak supervision and multimodal foundation models is the main response.
\textbf{Semantic alignment} is the second. Motion encodes culture, communication, and context, not only mechanics, and a model that fits joint trajectories while missing their relation to language, music, or environment produces motion that is kinematically clean but perceptually empty; semantic priors from language and foundation models help, but how to measure the result is unresolved.
That is the \textbf{evaluation} problem: every objective metric above has blind spots, none replaces a user study, and the field still needs metrics that are both interpretable and better correlated with human perception.
\textbf{Controllability} matters for practical use as much as in image generation, and joint masking, attribute prompts, and style conditioning are moving toward editing tools with textual, spatial, and temporal handles (§\ref{subsec:p3_s8_4_4}).
\textbf{Interactivity} remains underexplored: most methods animate one person in a static scene, and socially coupled groups (conversation, group dance) and motion grounded in dynamic, actionable scenes need both new corpora and new models.





\subsection{Animal Motion Generation}
\label{subsec:p3_s8_4_2}
Most motion-generation work targets the human body and ignores the diversity of animal behavior, although synthesizing animal motion matters for wildlife conservation, ecology, and biomechanics.
Relative to §\ref{subsec:p3_s8_4_1}, animal motion is harder on two fronts that the human setting never confronts: species vary widely in skeletal topology and morphology, so a single fixed-skeleton representation no longer suffices, and the same textual description (``walk stealthily'') maps to radically different gaits across species.
These two pressures, \ie, data scarcity and morphological heterogeneity, shape every method below, which is why animal motion generation borrows the architectures, datasets, and training recipes of human motion generation while adding species-conditioning machinery on top.



\subsubsection{Task Definition}
\label{subsubsec:p3_s8_4_2_task}
Animal motion generation extends the human sampler of Eq.~\eqref{eq:p3_s8_4_1_objective} with an explicit species (morphology) descriptor. Let $s\in\mathcal{S}$ index a species and let $\mathcal{P}_s$ be its species-dependent pose space (its skeleton topology and joint count depend on $s$), so that $\mathcal{M}_s=\mathcal{P}_s^{T}$ is the corresponding motion space. With $\kappa$ encoding a driving signal $\mathbf{c}\in\mathcal{C}$ (an action class, a text description, an anchor pose, or past motion), the task is to learn
\begin{equation}
\label{eq:p3_s8_4_2_objective}
G_\theta^{\text{amot}}:\mathcal{C}'\!\times\!\mathcal{S}\!\to\!\mathcal{M}_s,\qquad
\hat{\mathbf{m}}^{1:T}\sim p_\theta\!\left(\mathbf{m}^{1:T}\mid\kappa(\mathbf{c}),\,s\right).
\end{equation}
Restricting $\mathcal{S}$ to the single human morphology recovers Eq.~\eqref{eq:p3_s8_4_1_objective} exactly; the open difficulty is that $s$ ranges over a heterogeneous, long-tailed set of skeletons for which paired text--motion data are scarce, so the species condition must do double duty as both a generation control and a regularizer that shares statistical strength across morphologies.


\subsubsection{Methodology}
\label{subsubsec:p3_s8_4_2_method}
Almost all animal motion generation is text-driven, and the methods share a common response to the data-scarcity and morphology problems of Eq.~\eqref{eq:p3_s8_4_2_objective}: transfer representations learned from abundant human motion, then specialize them to a species through \emph{feature modulation}. Concretely, given a backbone feature $\mathbf{h}$ at some layer, species-aware modulation applies a FiLM-style affine transform conditioned on $s$,
\begin{equation}
\label{eq:p3_s8_4_2_modulation}
\tilde{\mathbf{h}}=\gamma(s)\odot\mathbf{h}+\beta(s),
\end{equation}
where $\gamma(s),\beta(s)$ are predicted from the species descriptor, allowing one shared network to adapt its internal representation to diverse anatomical forms. The methods differ in how they obtain the shared representation and how strictly they enforce cross-species plausibility, as Figure~\ref{fig:p3_s8_4_2_methods} summarizes.

\textbf{OmniMotionGPT}~\cite{YangZSWXHBQ024} tackles limited animal data by jointly training motion autoencoders on human and animal corpora and aligning their latent encodings to text via CLIP embeddings; it relies most directly on transfer, and at inference it still depends on human-motion priors, which prevents truly text-only animal generation.
\textbf{Motion Avatar}~\cite{0006W0CZH0F0024} uses SinMDM~\cite{RaabLTABC24} for data augmentation in its ZooGen pipeline (the route by which the Zoo-300K dataset is built) and extends the masked-motion modeling of MoMask~\cite{GuoMJW024} to animals, adding an LLM planner that orchestrates joint motion and 3D-avatar mesh generation; it makes little use of species-specific structure, which is what Eq.~\eqref{eq:p3_s8_4_2_modulation} addresses.
\textbf{AniMo}~\cite{WangRZW25} makes that remedy central, injecting species-aware modulation into a joint-aware spatiotemporal tokenizer and learning, through masked modeling, a mapping from text plus species attributes to motion tokens; its AniMo4D dataset (78{,}149 sequences over 114 species) yields markedly better morphological and behavioral fidelity than transfer-only baselines.
\textbf{X-MoGen}~\cite{abs-2508-05162} pushes toward a single model spanning humans and animals: it embeds motion into a shared space with graph-based VAEs and latent autoencoders, generates via masked modeling, and adds a morphological-consistency module that enforces skeletal plausibility across species, supported by the UniMo4D corpus (${\sim}115$ species, ${\sim}119$K sequences). OmniMotionGPT and Motion Avatar still lean on human priors, whereas AniMo and X-MoGen promote species morphology to a first-class, explicitly regularized conditioning variable, so the four together trace a progression from transfer toward native cross-species modeling.

\begin{figure}[!ht]
    \centering
    \includegraphics[width=\linewidth]{figures/part3/section8/cross_species_animal_motion_generation.pdf}
    \caption{Abstract cross-species motion generation: human and animal corpora support a shared representation, species-dependent modulation (Eq.~\eqref{eq:p3_s8_4_2_modulation}) adapts features, and an action-conditioned sampler produces motion for the selected skeleton. Dog, bird, and horse sequences illustrate alternative target morphologies, not simultaneous outputs. The dashed constraint denotes optional morphology regularization; implementations vary across methods.}
    \label{fig:p3_s8_4_2_methods}
\end{figure}


\subsubsection{Benchmarks}
\label{subsubsec:p3_s8_4_2_benchmarks}
Animal motion resources fall into two groups (Table~\ref{tab:p3_s8_4_2_datasets}).
\textbf{Capture and pose-estimation corpora} supply ground-truth articulated motion but no language: RGBD-Dog~\cite{Kearney0PKC20} and 3DDogs-Wild~\cite{abs-2406-14412} pair RGB-D or RGB video with marker-based mocap for canine pose, and DeformingThings4D~\cite{LiTTZN21} provides synthetic 4D sequences with dense mesh-deformation annotations across humanoid and animal categories.
\textbf{Text-to-motion corpora} drive recent progress: AnimalML3D~\cite{YangZSWXHBQ024} is the first text--animal-motion dataset, and the field then scales sharply through Zoo-300K~\cite{0006W0CZH0F0024} (built by the ZooGen pipeline with SinMDM and LLM annotators), AniMo4D~\cite{WangRZW25} (normalized to a 30-joint cross-species skeleton), and UniMo4D~\cite{abs-2508-05162}, which merges AniMo4D with HumanML3D and KIT-ML. Even the largest text-conditioned sets remain smaller and more morphologically fragmented than their human counterparts of §\ref{subsubsec:p3_s8_4_1_benchmarks}, which is what forces the transfer-and-modulate methodology above.



\begin{table}[!ht]
\centering
\fontsize{8}{10}\selectfont
\setlength{\tabcolsep}{1.0mm}
\caption{Benchmarks for animal motion generation. ``Subjects / species'' counts individual animals for the canine capture sets and species or categories elsewhere; ``Texts / frames'' gives text descriptions for text-conditioned corpora and frames for capture corpora.}
\label{tab:p3_s8_4_2_datasets}
\begin{tabular}{L{0.21\textwidth} L{0.13\textwidth} L{0.13\textwidth} C{0.15\textwidth} C{0.14\textwidth} C{0.14\textwidth}}
\toprule
\textbf{Dataset} & \textbf{Type} & \textbf{Modality} & \textbf{Subjects / species} & \textbf{Sequences} & \textbf{Texts / frames} \\
\midrule
RGBD-Dog \cite{Kearney0PKC20} & Canine & RGB-D, mocap & 5 dogs & 5 motion types & 8,346 frames \\
3DDogs-Wild \cite{abs-2406-14412} & Canine & RGB, mocap & 37 dogs & 143 recordings & - \\
DeformingThings4D \cite{LiTTZN21} & Animal / humanoid & 4D mesh & 31 categories & 1,972 & 122,365 frames \\
AnimalML3D \cite{YangZSWXHBQ024} & Animal / humanoid & 4D & 36 identities & 1,240 & - \\
Zoo-300K \cite{0006W0CZH0F0024} & Animal & 3D mesh & 65 categories & ${\sim}$300K text--motion pairs & - \\
AniMo4D \cite{WangRZW25} & Animal & 4D & 114 species & 78,149 & 185,435 texts \\
UniMo4D \cite{abs-2508-05162} & Animal + human & 4D & ${\sim}$115 species & ${\sim}$119K & - \\
\bottomrule
\end{tabular}
\end{table}

\paragraph{Evaluation Metrics.} Animal motion evaluation reuses the three-axis human protocol of §\ref{subsubsec:p3_s8_4_1_benchmarks}, adapted for species-specific biomechanics.
\emph{Fidelity} is measured by FID against real movement distributions and by MPJPE against ground-truth poses, augmented with physical-plausibility checks for gait-cycle regularity, ground contact, and the absence of foot sliding or limb penetration.
\emph{Condition alignment} carries over R-Precision and Multimodal Distance to score text-/action-to-motion correspondence, and \emph{diversity / multimodality} metrics quantify behavioral variety within and across species. The animal-specific addition is perceptual and biomechanical judgment of naturalness, balance, and anatomical consistency, which matters more for quadrupeds and multi-limbed morphologies than for the single human skeleton.



\subsubsection{Challenges and Trends}
\label{subsubsec:p3_s8_4_2_challenges}
\textbf{Physical realism} is the most basic gap: most methods do not simulate ground-contact forces, quadruped balance, or the fluid dynamics of swimming and flight, so artifacts the human setting has largely learned to suppress reappear, and differently for every body plan.
\textbf{Behavioral diversity} is harder than for humans because one instruction such as ``walk stealthily'' realizes into different gaits, postures, and speeds across species, so the conditional distribution to be captured is far wider while the data to capture it are scarcer.
\textbf{Unseen species and radical morphologies} remain out of reach: species-aware modulation (Eq.~\eqref{eq:p3_s8_4_2_modulation}) and morphological-consistency constraints generalize across the skeletons seen in training, but insects and winged or aquatic creatures violate the shared-skeleton, fixed-topology assumptions on which current methods rest.
AniMo and X-MoGen already move toward treating morphology as a variable rather than a template, and topology-flexible representations are the direction this points.




\subsection{Vehicle Motion Planning}
\label{subsec:p3_s8_4_3}
Where §\ref{subsec:p3_s8_4_1} and §\ref{subsec:p3_s8_4_2} generate the motion of an articulated body, vehicle motion planning generates the trajectory of a (near-)rigid agent moving through a traffic scene.
The body is far simpler, since a vehicle has essentially no internal articulation, but the difficulty moves outward into the environment: a planned trajectory must obey vehicle dynamics, stay on the drivable surface dictated by an HD map, and remain socially compatible with the future motion of many surrounding agents whose intentions are unobserved.
This is the setting where motion generation becomes inseparable from \emph{multi-agent} spatial reasoning and physical feasibility: trajectory forecasting (predicting how others will move) and planning (deciding how the ego vehicle should move) are two readings of the same generative object.

\subsubsection{Task Definition}
\label{subsubsec:p3_s8_4_3_task}
Let a trajectory be a sequence of agent states $\boldsymbol{\tau}^{1:H}=(\mathbf{s}^1,\dots,\mathbf{s}^H)$ over a horizon $H$, where each state $\mathbf{s}^t$ collects position, heading, and velocity.
The driving signal is now a structured scene context $\mathbf{c}=(\mathbf{s}^{-P:0},\,\mathcal{G},\,\mathcal{A})$ comprising the agent's own past states, an HD-map / lane graph $\mathcal{G}$, and the observed histories of neighboring agents $\mathcal{A}$.
With $\kappa$ encoding this context, vehicle motion planning learns
\begin{equation}
\label{eq:p3_s8_4_3_objective}
G_\theta^{\text{veh}}:\mathcal{C}'\!\to\!\mathcal{T},\qquad
\hat{\boldsymbol{\tau}}^{1:H}\sim p_\theta\!\left(\boldsymbol{\tau}^{1:H}\mid\kappa(\mathbf{c})\right)
\quad\text{s.t.}\quad
\mathbf{s}^{t+1}=f_{\text{dyn}}(\mathbf{s}^{t},\mathbf{u}^{t}),
\end{equation}
the conditional-sampler template of Eq.~\eqref{eq:p3_s8_4_1_objective} with the pose sequence replaced by a state trajectory and an added kinematic-feasibility constraint $f_{\text{dyn}}$ (\eg, a bicycle model) that distinguishes this setting from articulated-body generation.
Dropping $\mathcal{A}$ recovers single-agent forecasting; supplying a goal or reward in place of $\mathbf{c}$ turns prediction into planning.


\subsubsection{Methodology}
\label{subsubsec:p3_s8_4_3_method}
Prediction methods vary along two axes: how the scene context is \emph{encoded} and how the multimodal future is \emph{decoded}.
Early encoders rasterize the scene into a bird's-eye-view image, but the field has converged on \emph{vectorized} encoders that represent lanes and agent tracks as polylines: VectorNet~\cite{gao2020vectornet} and LaneGCN~\cite{liang2020lanegcn} build graph representations of the HD map, and transformer encoders such as Wayformer~\cite{nayakanti2023wayformer} and the Motion Transformer (MTR)~\cite{shi2022mtr} attend jointly over map and agents; graph-attention predictors~\cite{ZhangFWH22} model spatial and temporal interaction in the same spirit.
On the decoding side, the defining property is multimodality, since a single context admits many plausible futures, so models predict a set of $K$ modes with probabilities,
\begin{equation}
\label{eq:p3_s8_4_3_multimodal}
p_\theta(\boldsymbol{\tau}\mid\kappa(\mathbf{c}))=\sum_{k=1}^{K}\pi_k\,\mathcal{N}\!\left(\boldsymbol{\tau};\,\mu_k(\kappa(\mathbf{c})),\,\Sigma_k\right),
\end{equation}
rather than the single point estimate of Eq.~\eqref{eq:p3_s8_4_1_regression}, which here would average distinct maneuvers (turning \vs going straight) into an infeasible mean.
The earliest learned predictors modeled social interaction directly in agent-coordinate space, Social-LSTM~\cite{alahi2016social} through pooled recurrent states and Social-GAN~\cite{gupta2018social} through adversarial sampling, and Trajectron++~\cite{salzmann2020trajectron} added dynamic feasibility to a graph-structured latent-variable model.
Goal-based decoders make the modes explicit: DenseTNT~\cite{gu2021densetnt} samples dense candidate endpoints on the map and scores trajectories conditioned on them, exposing the maneuver structure that Eq.~\eqref{eq:p3_s8_4_3_multimodal} captures only implicitly.
Generative decoders replace the fixed mixture with a learned sampler: MotionDiffuser~\cite{jiang2023motiondiffuser} denoises whole multi-agent trajectory sets jointly.

Planning turns the same machinery on the ego vehicle. ChauffeurNet~\cite{bansal2019chauffeurnet} learns the ego policy by imitating expert driving and augments the demonstrations with synthesized perturbations and worst cases, so that the policy stays robust when it drifts off the data manifold.
Figure~\ref{fig:p3_s8_4_3_methods} situates these families along the encode/decode axes.

\begin{figure}[!ht]
    \centering
    \includegraphics[width=\linewidth]{figures/part3/section8/vehicle_motion_planning.pdf}
    \caption{Vehicle motion planning as conditional trajectory generation: a vectorized encoder represents map polylines and agent interactions, and a multimodal decoder generates candidate futures (Eq.~\eqref{eq:p3_s8_4_3_multimodal}). Gray intensity indicates mode weight; blue marks a plan selected using a goal or cost and feasibility constraints, not necessarily the highest-probability mode. Road geometry and trajectories are illustrative.}
    \label{fig:p3_s8_4_3_methods}
\end{figure}

\subsubsection{Benchmarks}
\label{subsubsec:p3_s8_4_3_benchmarks}
Progress is measured on large real-world driving logs that pair agent tracks with HD maps: nuScenes~\cite{caesar2020nuscenes}, Argoverse and Argoverse~2~\cite{Argoverse-ChangLSSBHW0LRH19, wilson2023argoverse}, and the Waymo Open Motion Dataset~\cite{ettinger2021womd}, the last two offering the scale and interaction density that current transformer models require.
Evaluation centers on the multimodal-aware displacement metrics that match Eq.~\eqref{eq:p3_s8_4_3_multimodal}: minimum average and final displacement error over the $K$ predicted modes (minADE$_K$, minFDE$_K$), miss rate (the fraction of scenes where no mode lands within a tolerance of the ground truth), and the soft mean-average-precision of the Waymo protocol, which jointly scores accuracy and the calibration of mode probabilities $\pi_k$.
Off-road and collision rates further test the feasibility constraint $f_{\text{dyn}}$ that purely displacement-based scores ignore.

\subsubsection{Challenges and Trends}
\label{subsubsec:p3_s8_4_3_challenges}
A planner must assign trustworthy probabilities to the futures it enumerates, not merely list them: an over-confident missed mode of Eq.~\eqref{eq:p3_s8_4_3_multimodal} is a safety failure rather than a metric error, so \textbf{calibrated multimodality} is the first open problem.
\textbf{Feasibility and safety} are the second: displacement metrics reward trajectories a vehicle cannot physically execute, so feasibility has to be enforced inside the model ($f_{\text{dyn}}$) rather than checked afterward, and the long tail of near-collision scenarios is exactly where data are scarcest.
The third is \textbf{closed-loop distribution shift}: models trained to imitate logged, open-loop futures drift once their own predictions feed back into the scene, the difficulty ChauffeurNet's worst-case synthesis was built to expose.
Diffusion- and transformer-based scene-level samplers that forecast all agents jointly and run closed-loop are the current direction, and they connect directly to the controllable-video and world-model lines of §\ref{sec:p3_s8_3} and §\ref{sec:p3_s8_6}, in which a driving world model both predicts and is conditioned on the ego plan.







\subsection{Motion Editing}
\label{subsec:p3_s8_4_4}
The subsections above \emph{synthesize} motion from scratch; motion editing instead modifies an \emph{existing} motion so that it satisfies a new constraint while leaving the rest intact. It is the temporal, articulated analogue of the image editing of §\ref{subsec:p3_s8_1_3} and the video editing of §\ref{subsec:p3_s8_3_5}, and its demand is again selective preservation: change a targeted aspect (a trajectory, a joint, a style, a semantic action) while holding the body's identity, timing, and untouched portions fixed, which makes editing the controllability layer that closes the generate-and-refine loop for motion.



\subsubsection{Task Definition}
\label{subsubsec:p3_s8_4_4_task}
Given a source motion $\mathbf{m}^{\text{in}}\!\in\!\mathcal{M}$, an edit region $\mathbf{M}\!\in\!\{0,1\}^{T\times J}$ over time $T$ and joints $J$, and an edit control $\mathbf{c}$ (a text instruction, target keyframes, a desired root trajectory, or a target joint), motion editing learns an operator
\begin{equation}
\label{eq:p3_s8_4_4_objective}
E_\theta:\mathcal{M}\!\times\!\{0,1\}^{T\times J}\!\times\!\mathcal{C}'\!\to\!\mathcal{M},\qquad
\hat{\mathbf{m}}=E_\theta\!\left(\mathbf{m}^{\text{in}},\mathbf{M},\kappa(\mathbf{c})\right),
\end{equation}
which lifts the image-editing operator of Eq.~\eqref{eq:t2i_edit_objective} into the motion domain. As in image and video editing, Eq.~\eqref{eq:p3_s8_4_4_objective} encodes two competing objectives, \emph{editability} (the output must realize the control $\mathbf{c}$ inside $\mathbf{M}$) and \emph{preservation} (it must match $\mathbf{m}^{\text{in}}$ outside $\mathbf{M}$ and remain a plausible, smooth motion), and methods differ chiefly in how they trade these off.



\subsubsection{Methodology}
\label{subsubsec:p3_s8_4_4_method}
Because a pretrained motion diffusion model already samples plausible motion, most editors realize Eq.~\eqref{eq:p3_s8_4_4_objective} without retraining, by \emph{conditioning the sampler} of Eq.~\eqref{eq:p3_s8_4_1_diffusion} rather than learning a dedicated $E_\theta$. The simplest mechanism is \textbf{masked infilling}: at every denoising step the observed entries are overwritten with the (noised) source while only the masked entries are generated,
\begin{equation}
\label{eq:p3_s8_4_4_inpaint}
\hat{\mathbf{m}}=(1-\mathbf{M})\odot\mathbf{m}^{\text{in}}+\mathbf{M}\odot G_\theta\!\left(\mathbf{m}_n,\,n,\,\kappa(\mathbf{c})\right),
\end{equation}
which directly enforces preservation by construction. This is the basis of motion in-betweening, where $1-\mathbf{M}$ marks given keyframes: robust in-betweening~\cite{harvey2020robust} predates the diffusion era, and CondMDI~\cite{cohan2024condmdi} casts flexible keyframe-conditioned in-betweening as masked diffusion, while PriorMDM~\cite{shafir2024priormdm} reuses a frozen motion-diffusion prior for editing and composition tasks such as joining and in-filling. A second family handles \emph{spatial} controls that masking alone cannot express (a target root path, or a specific joint at specific times) through classifier-style guidance: GMD~\cite{karunratanakul2023gmd} and OmniControl~\cite{xie2024omnicontrol} steer the denoiser toward arbitrary spatial constraints at any joint and time. A third, optimization-based family treats the initial diffusion noise itself as the free variable, solving $\min_{\mathbf{z}}\,\mathcal{L}_{\text{edit}}(\mathcal{D}(\mathbf{z}))$ so that the decoded motion satisfies an arbitrary differentiable criterion: DNO~\cite{karunratanakul2024dno} shows this single mechanism serves editing, in-betweening, and refinement as one universal motion prior. Finally, \emph{semantic} editors take a natural-language instruction in place of a geometric mask: MotionFix~\cite{athanasiou2024motionfix} learns text-driven edits from paired source--target motions, and LLM-driven approaches~\cite{goel2024iterative} decompose an instruction into a program of primitive edits applied iteratively. Figure~\ref{fig:p3_s8_4_4_methods} groups these four mechanisms by how they obtain $E_\theta$, paralleling the editing taxonomies of §\ref{subsec:p3_s8_1_3} and §\ref{subsec:p3_s8_3_5}.

\begin{figure}[!ht]
    \centering
    \includegraphics[width=\linewidth]{figures/part3/section8/four_motion_editing_mechanisms.pdf}
    \caption{Four motion-editing mechanisms: (a) reinsert known states while infilling missing frames; (b) guide denoising with root-path or joint targets; (c) optimize initial noise through a fixed decoder using an edit loss; and (d) edit source motion from text, using a learned editor or iterative primitive edits. Rendered people illustrate edits rather than measured model outputs.}
    \label{fig:p3_s8_4_4_methods}
\end{figure}

\subsubsection{Benchmarks}
\label{subsubsec:p3_s8_4_4_benchmarks}
Motion editing inherits the corpora of §\ref{subsubsec:p3_s8_4_1_benchmarks} (HumanML3D~\cite{GuoZZ0JL022} and BABEL~\cite{PunnakkalCAQB21} over AMASS), used in editing protocols where keyframes, joints, or trajectories are held out as the constraint, together with the dedicated MotionFix~\cite{athanasiou2024motionfix} benchmark of paired source--edit--instruction triples that supplies the ground truth most editing tasks lack.
Evaluation must report the editability--preservation pair jointly: editability through control adherence (instruction-conditioned R-Precision, target-joint or trajectory error) and overall realism (FID, foot-skating), and preservation through the deviation from the source outside the edit region.
As in video editing, single scalars correlate weakly with perceived edit quality, so reporting a tuple, and increasingly a learned or VLM-style judge of instruction compliance, is becoming standard.

\subsubsection{Challenges and Trends}
\label{subsubsec:p3_s8_4_4_challenges}
The central tension is the \textbf{editability--preservation trade-off} of Eq.~\eqref{eq:p3_s8_4_4_objective}: pushing harder on a control propagates unwanted change through the kinematic chain (editing one arm disturbs balance and foot contact), so a good editor must respect the body's global coordination, not just the masked entries.
This is compounded by \textbf{sparse-control ambiguity}: a single keyframe or a root path under-determines the full-body motion, and the model must fill the gap plausibly, which the guidance and noise-optimization families address with differing robustness.
Underlying both is an \textbf{evaluation} gap sharper than in synthesis: because the ``correct'' edit is rarely unique, ground-truth-based metrics under-reward valid alternatives, motivating the paired-triplet and judge-based protocols above.
The trend mirrors the rest of §\ref{sec:p3_s8_4} and §\ref{sec:p3_s8_3}: from per-mechanism, training-free editors toward unified models that generate, control, and edit motion within a single conditional sampler, increasingly driven by natural-language instructions rather than geometric masks.








% \subsection{Summary}
% \label{subsec:p3_s8_4_summary}
% Table~\ref{tab:p3_s8_4_subfield_comparison} compares the four subfields as specializations of the human-motion sampler of Eq.~\eqref{eq:p3_s8_4_1_objective}: by the moving object and its condition, the sampler or operator each instantiates, the dominant method families, and the benchmarks on which they report.

% \begin{table}[!ht]
% \centering
% \fontsize{8}{10}\selectfont
% \setlength{\tabcolsep}{1.0mm}
% \caption{The four subfields of motion generation; each sampler or operator specializes Eq.~\eqref{eq:p3_s8_4_1_objective}.}
% \label{tab:p3_s8_4_subfield_comparison}
% \begin{tabular}{L{0.13\textwidth} L{0.19\textwidth} L{0.16\textwidth} L{0.22\textwidth} L{0.22\textwidth}}
% \toprule
% \textbf{Subfield} & \textbf{Moving object \& condition} & \textbf{Sampler / operator} & \textbf{Dominant method families} & \textbf{Benchmarks / metrics} \\
% \midrule
% §\ref{subsec:p3_s8_4_1} Human motion & Articulated body; action / text / music / speech / scene & $G_\theta^{\text{mot}}$ (Eq.~\eqref{eq:p3_s8_4_1_objective}) & Regression or generative (GAN / VAE / diffusion); per-modality conditioning & HumanML3D, KIT-ML, AIST++, BEAT; FID, R-Precision, diversity, beat / contact \\
% §\ref{subsec:p3_s8_4_2} Animal motion & Multi-species body; text $+$ species $s$ & $G_\theta^{\text{amot}}$ (Eq.~\eqref{eq:p3_s8_4_2_objective}) & Human-prior transfer $+$ species-aware modulation; morphological consistency & AnimalML3D, Zoo-300K, AniMo4D; FID, MPJPE, biomechanical plausibility \\
% §\ref{subsec:p3_s8_4_3} Vehicle planning & Rigid agent state; HD map $+$ agent histories & $G_\theta^{\text{veh}}$, $K$-mode (Eqs.~\eqref{eq:p3_s8_4_3_objective},~\eqref{eq:p3_s8_4_3_multimodal}) & Vectorized encoders; goal-based and generative decoders; imitation planners & nuScenes, Argoverse 1/2, WOMD; minADE$_K$, minFDE$_K$, miss rate, soft-mAP \\
% §\ref{subsec:p3_s8_4_4} Motion editing & Existing motion $+$ edit region / control & $E_\theta$ (Eq.~\eqref{eq:p3_s8_4_4_objective}) & Masked infilling, spatial guidance, noise optimization, instruction editing & HumanML3D / BABEL protocols, MotionFix; editability $+$ preservation tuple \\
% \bottomrule
% \end{tabular}
% \end{table}

% The four subfields partition by \emph{what is held fixed}: human, animal, and vehicle generation fix nothing but the condition and sample a fresh trajectory, whereas editing fixes the source motion and changes only a region, standing to the other three exactly as image editing (§\ref{subsec:p3_s8_1_3}) stands to image generation. Underneath that partition the dominant failure mode is shared and physical: whether it surfaces as foot-skating and penetration in human and animal motion, kinematic infeasibility in vehicle trajectories, or balance disruption when one limb is edited, every subfield is limited by how well the generator respects the biomechanics and dynamics of the moving body, which is why physics- and contact-aware modeling recurs in all four Challenges discussions. Evaluation stays fragmented, each subfield reporting its own tuple (fidelity, diversity, condition consistency; calibrated multimodality; editability and preservation), and all are migrating toward learned or perceptual judges as scalar metrics prove weakly correlated with human assessment. Motion generation is thus the point in §\ref{sec:generation} where generative spatial modeling must answer to physical law as well as semantics, which leads directly to the dynamic 4D generation of §\ref{sec:p3_s8_5}.


% \newpage
\section{4D Spatial Generation}
\label{sec:p3_s8_5}




\FramedInitial{G}{enerative} modeling has moved outward from static images~\cite{rombach2022ldm,saharia2022photorealistic} to video~\cite{chen2024videocrafter2,chen2023videocrafter1,ho2022video,harvey2022flexible}, multi-view imagery~\cite{shi2023mvdream,liu2023zero,shi2023zero123++,zhao2024sg,zhao2024chase}, and static 3D assets~\cite{wang2023score,raj2023dreambooth3d}, yet each stops short of a scene that is at once three-dimensional and alive.
A 3D asset holds geometry but not motion; a video holds motion, but no view-consistent geometry one can navigate or simulate.
4D generation is where the two meet: the synthesis of dynamic 3D content whose geometry stays view-consistent as it moves and deforms over time~\cite{singer2023text,zheng2024unified,jiang2023consistent4d,ling2024align,wang2025moving}.
This is the representation an embodied system needs to rehearse how a scene will unfold, and it underwrites photorealistic digital humans~\cite{guo2023vid2avatar,zhao2024sg,zhao2024chase}, interactive AR/VR~\cite{zhao2024hfgs,miao2024pla4d}, and high-fidelity simulation for autonomous driving and embodied AI~\cite{gao2024magicdrive3d,zhao2024efficient}.

Two advances make this tractable.
Temporal 3D representations, deformation-aware fields over NeRFs~\cite{mildenhall2021nerf}, point clouds, and Gaussian splatting~\cite{kerbl20233d,zhao2024efficient}, supply a geometry that can bend and move without losing view consistency, while the extension of generative diffusion into this domain supplies the realism and controllability to drive it.
Together they open dynamic scene synthesis, 4D editing, and interactive agent simulation, at the cost of the difficulty this section returns to throughout: holding long-range temporal coherence while still granting flexible, user-driven control.

Formally, we cast 4D generation as the latent-generator framework of §\ref{subsec:p3_s8_1_1} unrolled along time over a shared 3D representation. Let $\mathcal{R}$ be the 3D-representation space of Eq.~\eqref{eq:p3_s8_2_1_objective} (mesh, NeRF, 3DGS, point cloud, and so on), and let $\kappa:\mathcal{C}\!\to\!\mathcal{C}'$ encode a conditioning signal $\mathbf{c}\in\mathcal{C}$ (text, an image, a video, a static 3D asset, or a combination thereof). 4D generation learns a conditional sampler
\begin{equation}
\label{eq:p3_s8_5_objective}
G_\theta^{\text{4D}}:\mathcal{C}'\!\to\!\mathcal{R}^{[0,T]},\qquad
\hat{\mathbf{r}}^{1:T}\sim p_\theta\!\left(\mathbf{r}^{1:T}\mid\kappa(\mathbf{c})\right),
\end{equation}
that emits a time-indexed sequence of 3D representations $\hat{\mathbf{r}}^{1:T}$ with view-consistent geometry and coherent motion. Two limiting cases anchor the framework: with a single time step ($T=1$) Eq.~\eqref{eq:p3_s8_5_objective} recovers the static 3D generator of Eq.~\eqref{eq:p3_s8_2_1_objective}, and rendering each frame $\mathbf{r}^t$ to a 2D view recovers the video generator of Eq.~\eqref{eq:p3_s8_3_objective}, so 4D generation is precisely the 3D generator of §\ref{sec:p3_s8_2} unrolled in time, or equivalently the video generator of §\ref{sec:p3_s8_3} lifted from per-frame pixels to a shared, renderable geometry.
The subsections are organized by the conditioning input (text, image, video, a static 3D asset, or several at once). Figure~\ref{fig:p3_s8_5_taxonomy} previews this organization together with two methodological strategies that recur across all five, per-scene score distillation and explicit ``generate-then-reconstruct'' supervision; feed-forward generators and rig-based animation of meshes appear alongside them in the subsections where they apply.

\begin{figure}[!ht]
    \centering
    \includegraphics[width=\linewidth]{figures/part3/section8/taxonomy_4d_spatial_generation.pdf}
    \caption{4D spatial generation by conditioning input, (a)--(e), with two representative strategies: (i) score distillation from a frozen diffusion prior (dashed: SDS feedback), and (ii) generating video or multiview targets, then reconstructing a dynamic 3D asset. The strategies are alternatives that may be combined; the chair sequence is illustrative.}
    \label{fig:p3_s8_5_taxonomy}
\end{figure}






\subsection{Text-to-4D Generation}
\label{subsec:p3_s8_5_1}
Text-to-4D generation synthesizes dynamic content directly from natural language, supporting text-driven asset creation and character animation~\cite{HongZPCYL22,chen2024ct4d}. Its progress has been propelled by large language models for semantic understanding and by video foundation models that supply motion priors.


\subsubsection{Task Definition}
\label{subsubsec:p3_s8_5_1_task}
Text-to-4D generation specializes the 4D sampler of Eq.~\eqref{eq:p3_s8_5_objective} to a textual condition $\mathbf{c}=\mathbf{x}$ encoded by the text encoder $\phi$ of Eq.~\eqref{eq:t2i_generator},
\begin{equation}
\label{eq:p3_s8_5_1_objective}
\hat{\mathbf{r}}^{1:T}\sim p_\theta\!\left(\mathbf{r}^{1:T}\mid\phi(\mathbf{x})\right).
\end{equation}
Because language fixes only semantics and leaves geometry and motion grossly under-determined, the sampler must hallucinate both a plausible 3D shape and its temporal evolution from the prompt alone, the source of the semantic-to-motion ambiguity examined in §\ref{subsubsec:p3_s8_5_1_challenges}.



\subsubsection{Methodology}
\label{subsubsec:p3_s8_5_1_method}
Two strategies recur throughout this section, implicit distillation and explicit supervision, implemented across various 4D representations (Figure~\ref{fig:p3_s8_5_taxonomy}).
The more common implicit strategy leverages Score Distillation Sampling (SDS) to indirectly optimize the geometry and motion of a 4D representation, using pre-trained diffusion models as frozen priors~\cite{singer2023text,zheng2024unified}. Formally, it applies a score-distillation update to per-scene parameters $\theta$~\cite{poole2022dreamfusion},
\begin{equation}
\label{eq:p3_s8_5_sds}
\nabla_\theta\mathcal{L}_{\mathrm{SDS}}
=\mathbb{E}_{s,\epsilon,\pi}\!\left[
 w(s)\left(\epsilon_\phi(\mathbf{z}_s;s,\kappa(\mathbf{c}))-\epsilon\right)
 \frac{\partial\mathbf{x}_\theta}{\partial\theta}\right],
\quad \mathbf{z}_s=\alpha_s\mathbf{x}_\theta+\sigma_s\epsilon,
\end{equation}
where $\mathbf{x}_\theta=\mathrm{Render}(\hat{\mathbf{r}}^{1:T}_\theta,\pi)$ is a differentiable rendering at camera $\pi$, $s$ indexes diffusion noise levels, and $\epsilon_\phi$ is a frozen 2D/video diffusion prior. The update differentiates through the renderer while treating the prior noise residual as the supplied gradient signal; this lifts the static score-distillation objective of §\ref{para:p3_s8_2_1_sds} to a temporal sequence.
In contrast, explicitly supervised methods initialize the 4D generation from text-to-3D or text-to-video outputs and use these generated assets as direct supervisory signals~\cite{zhao2024efficient,wu20244d,yang2024deformable,ling2024align,blattmann2023align,xu2024comp4d}. This ``generate-then-reconstruct'' strategy first synthesizes target frames $\hat{\mathbf{y}}$ with a fast generator $\mathcal{G}_{\text{2D}}$ and then fits the 4D asset to them,
\begin{equation}
\label{eq:p3_s8_5_explicit}
\hat{\mathbf{y}}=\mathcal{G}_{\text{2D}}(\mathbf{c}),\qquad
\hat{\mathbf{r}}^{1:T}=\arg\min_{\mathbf{r}^{1:T}}\sum_{\pi,t}\big\|\mathrm{Render}(\mathbf{r}^{t},\pi)-\hat{\mathbf{y}}^{t}_{\pi}\big\|^2,
\end{equation}
replacing score-distillation guidance with explicit pixel targets; reconstruction may itself require per-scene optimization, so a speed advantage depends on the implementation. The two strategies are realized across mesh, NeRF, and 3DGS representations.

\textbf{Mesh-based Methods.} Meshes offer explicit geometric control, which suits precise text-driven shape manipulation and animation.
CT4D~\cite{chen2024ct4d} implements a three-stage ``Generate-Refine-Animate'' pipeline in which a mesh is generated from text, its texture refined, and its vertices animated; TextMesh4D~\cite{dai2025textmesh4d} deforms a generated mesh through a per-face Jacobian deformation field under video-diffusion priors; and AvatarCLIP~\cite{HongZPCYL22} generates a static avatar zero-shot and animates it by registering motion to an SMPL skeleton.

\textbf{NeRF-based Methods.} NeRFs represent continuous geometry and complex temporal variation well.
MAV3D~\cite{singer2023text} introduced multi-stage SDS optimization that refines texture, geometry, and motion with priors from both T2I and T2V diffusion models; 4D-fy~\cite{bahmani20244d} blends 3D-aware, image, and video score signals in a hybrid SDS; and Dream-in-4D~\cite{zheng2024unified} first generates a static asset with a multiview diffusion model and then adds dynamics with a video diffusion model.

\textbf{Gaussian Splatting-based Methods.} 4D Gaussian Splatting (4DGS)~\cite{zhao2024hfgs,wu20244d} renders efficiently and differentiably, which speeds up text-driven dynamic scene synthesis. Early methods followed the two-stage implicit pipeline of Eq.~\eqref{eq:p3_s8_5_sds}: AYG~\cite{ling2024align} generates a static 3D Gaussian model from text via SDS and adds dynamics under a T2V prior, and Comp4D~\cite{xu2024comp4d} uses an LLM to parse a complex prompt into several dynamic objects and their trajectories for a compositional scene.
Later works moved to explicit supervision (Eq.~\eqref{eq:p3_s8_5_explicit}): PLA4D~\cite{miao2024pla4d} aligns the 4D Gaussians pixel by pixel with generated video targets, and 4Real~\cite{yu20244real} uses generated multiview video frames as the reference for photorealistic scene-level generation.
For digital humans, AvatarGO~\cite{cao2024avatargo} generates 4D human--object interactions zero-shot by animating a text-to-3D human with the SMPL-X skinning function.


\subsubsection{Benchmarks}
\label{subsubsec:p3_s8_5_1_benchmarks}
Text-to-4D, like the other subfields of this section, has no standardized benchmark, because large-scale dynamic 3D data with text descriptions is hard to capture and annotate. Training and evaluation instead adapt static 3D repositories (Table~\ref{tab:4d_datasets}): Diffusion4D~\cite{liang2024diffusion4d} filtered Objaverse and Objaverse-XL~\cite{deitke2023objaverse,deitke2023objaversexl} down to 54K high-quality animated assets, and SV4D~\cite{xie2024sv4d} built ObjaverseDy by adaptive sampling and viewpoint adjustment. These are synthetic renderings and lack the multimodal annotations, text in particular, needed to train and evaluate text-conditioned models robustly.

\begin{table}[!ht]
\centering
\fontsize{8}{10}\selectfont
\setlength{\tabcolsep}{1.0mm}
\caption{Datasets used for 4D generation across the subfields of §\ref{sec:p3_s8_5}.}
\label{tab:4d_datasets}
\begin{tabular}{L{0.20\textwidth} L{0.24\textwidth} C{0.20\textwidth} L{0.26\textwidth}}
\toprule
\textbf{Dataset} & \textbf{Content} & \textbf{Scale} & \textbf{Used for} \\
\midrule
Objaverse-XL~\cite{deitke2023objaversexl} & Static and animated 3D objects & 10M+ objects & Source pool for dynamic subsets \\
Diffusion4D~\cite{liang2024diffusion4d} & Animated assets curated from Objaverse / Objaverse-XL & 54K assets & Training image- / video-to-4D \\
ObjaverseDy~\cite{xie2024sv4d} & Animated Objaverse objects, re-rendered & - & Training and evaluating video-to-4D \\
Consistent4D~\cite{jiang2023consistent4d} & Synthetic and real monocular videos & - & Evaluating video-to-4D \\
\bottomrule
\end{tabular}
\end{table}

\paragraph{Evaluation Metrics.} Evaluation across 4D generation combines user studies with metrics borrowed from image and video generation, and the subsections that follow use this same set. Semantic alignment with a prompt is measured by CLIP Score~\cite{radford2021learning} between rendered views and the text; fidelity by PSNR~\cite{hore2010image}, SSIM~\cite{wang2004image}, and LPIPS~\cite{zhang2018unreasonable} where a reference view exists; temporal quality and motion by the Fréchet Video Distance (FVD)~\cite{unterthiner2018fvd}; and overall plausibility and aesthetics by user studies. None of these scores the two properties that make content 4D, cross-view consistency and long-range spatio-temporal coherence.


\subsubsection{Challenges and Trends}
\label{subsubsec:p3_s8_5_1_challenges}
The challenge specific to text-to-4D is semantic-to-motion ambiguity: a prompt such as ``a dancing dog'' admits endlessly many movements, styles, and trajectories, and methods that lean on generic T2V priors struggle to produce diverse, fine-grained motion that matches a nuanced description.
One response uses LLMs for motion planning, as Comp4D~\cite{xu2024comp4d} does when it decomposes a prompt into a structured spatio-temporal plan for several interacting objects; another targets human--object interaction~\cite{cao2024avatargo}, moving toward story-driven dynamic scenes.

% ======================================================================
% ======================================================================


\subsection{Image-to-4D Generation}
\label{subsec:p3_s8_5_2}
Image-to-4D generation synthesizes dynamic content from a single static image, combining the spatial information in the image with temporal priors learned by video models~\cite{yang2024diffusion,sun2024eg4d}. No large image--4D paired dataset exists, which has ruled out direct end-to-end training and pushed the field toward optimization and toward generated supervision.

\subsubsection{Task Definition}
\label{subsubsec:p3_s8_5_2_task}
Image-to-4D generation specializes the 4D sampler of Eq.~\eqref{eq:p3_s8_5_objective} to a single-image condition $\mathbf{c}=I$,
\begin{equation}
\label{eq:p3_s8_5_2_objective}
\hat{\mathbf{r}}^{1:T}\sim p_\theta\!\left(\mathbf{r}^{1:T}\mid\kappa(I)\right),
\end{equation}
where the source view at $t=0$ is fixed by $I$ while every unseen viewpoint and time step must be inferred.
Both the unseen 3D geometry and the entirely absent temporal dynamics have to be inferred from one 2D view, which requires combining the spatial cues of the input with temporal priors from large video models~\cite{ren2023dreamgaussian4d,zhao2023animate124,sun2024eg4d}.



\subsubsection{Methodology}
\label{subsubsec:p3_s8_5_2_method}
Both strategies of §\ref{subsubsec:p3_s8_5_1_method} apply, and most methods use 4D Gaussian Splatting as the representation~\cite{ren2023dreamgaussian4d,yang2024diffusion}.

\textbf{Implicit Optimization with Diffusion Priors.} These methods instantiate Eq.~\eqref{eq:p3_s8_5_sds}, typically combining a multiview diffusion model for geometric consistency with a video diffusion model for temporal coherence. Animate124~\cite{zhao2023animate124} first generates a static NeRF from the image and then refines its dynamics with a video-diffusion SDS loss, and Diffusion\textsuperscript{2}~\cite{yang2024diffusion} reconciles the scores of heterogeneous priors (multiview and video) through Variance-Reducing Sampling, improving both multiview and temporal consistency of the resulting 4D Gaussians.

\textbf{Explicit Supervision with Generated Data.} The slow, iterative nature of SDS motivates the two-stage pipeline of Eq.~\eqref{eq:p3_s8_5_explicit}, in which an image-to-video model first generates an intermediate video that then supervises the 4D reconstruction directly. DreamGaussian4D (DG4D)~\cite{ren2023dreamgaussian4d} generates a static 3D Gaussian model and learns a deformation field to animate it, using the generated video for texture refinement and dynamic supervision, and EG4D~\cite{sun2024eg4d} avoids SDS entirely by chaining image-to-video and image-to-multiview models into coherent image sequences that guide a fast reconstruction.

\textbf{Unified and Tuning-free Generation.} GenXD~\cite{zhao2024genxd} trains one model on both 3D and 4D data, using multiview-temporal modules that separate camera motion from object motion, and Free4D~\cite{liu2025free4d} builds a 4D scene from a single image without fine-tuning any model, distilling pretrained image-to-video and multiview priors into a consistent 4D representation.

\textbf{Specialized Methods for Human Generation.} Dynamic human avatars form a sub-domain in which strong body priors raise fidelity. Human4DiT~\cite{shao2024360} uses a hierarchical 4D transformer that integrates human-specific features, camera parameters, and temporal signals to generate 360-degree human animations from a single image, and Disco4D~\cite{pang2025disco4d} models the body with SMPL-X while representing clothing separately with Gaussian splats, allowing detailed animation of clothed humans.



\subsubsection{Benchmarks}
\label{subsubsec:p3_s8_5_2_benchmarks}
With no image--4D pairs, methods are evaluated on images from Objaverse~\cite{deitke2023objaverse} or on in-the-wild photographs chosen by the authors, which makes comparison across papers unreliable. The metric set is that of §\ref{subsubsec:p3_s8_5_1_benchmarks}, applied to the source view and to pseudo-ground-truth novel views where available.

\subsubsection{Challenges and Trends}
\label{subsubsec:p3_s8_5_2_challenges}
Inferring a complete shape and a plausible motion from one viewpoint is extremely ill-posed, which shows up as geometric inconsistency in unseen views (\eg, the ``Janus problem'') and as unnatural or repetitive motion, since every dynamic detail must be hallucinated.
Domain-specific priors are the main remedy: Human4DiT~\cite{shao2024360} and Disco4D~\cite{pang2025disco4d} regularize geometry and motion with parametric body models and reach much higher fidelity, at the cost of generality.


% ======================================================================
% ======================================================================


\subsection{Video-to-4D Generation}
\label{subsec:p3_s8_5_3}
Video-to-4D generation lifts a single-view video into a dynamic 4D representation that can be viewed from new angles with consistent geometry and motion, turning ordinary video into interactive 4D content. Annotated video-to-4D datasets are hard to curate, so current methods build on pretrained video diffusion priors, either as multi-view supervision or through score distillation.



\subsubsection{Task Definition}
\label{subsubsec:p3_s8_5_3_task}
Video-to-4D generation specializes the 4D sampler of Eq.~\eqref{eq:p3_s8_5_objective} to a monocular video condition $\mathbf{c}=\mathbf{v}^{1:T}$,
\begin{equation}
\label{eq:p3_s8_5_3_objective}
\hat{\mathbf{r}}^{1:T}\sim p_\theta\!\left(\mathbf{r}^{1:T}\mid\kappa(\mathbf{v}^{1:T})\right),
\end{equation}
so that re-rendering $\hat{\mathbf{r}}^{1:T}$ from the input viewpoint reproduces $\mathbf{v}^{1:T}$ while novel views remain geometrically consistent.
Unlike image-to-4D (Eq.~\eqref{eq:p3_s8_5_2_objective}), the temporal dynamics are now \emph{observed} and the residual ambiguity is purely geometric: the unseen off-axis structure.

\subsubsection{Methodology}
\label{subsubsec:p3_s8_5_3_method}
Most methods densify the monocular input into multi-view supervision and reconstruct against it (Eq.~\eqref{eq:p3_s8_5_explicit}), some add score distillation (Eq.~\eqref{eq:p3_s8_5_sds}), and the newest predict the 4D asset feed-forward.

\textbf{Mesh-based Methods.} Meshes are interpretable and efficient to animate. DreamMesh4D~\cite{li2024dreammesh4d} obtains a coarse mesh and drives it with a sparse deformation graph and linear blend skinning (LBS), reaching strong spatio-temporal consistency more cheaply than implicit volumetric methods, and V2M4~\cite{chen2025v2m4} reconstructs an animated, textured mesh asset directly from a monocular video.

\textbf{NeRF-based Methods.} NeRF-based methods model the continuous volumetric content of a video well, and they overcome the single view by generating data. Consistent4D~\cite{jiang2023consistent4d} densifies the monocular input through temporal and multiview interpolation before optimizing a NeRF with SDS, 4Diffusion~\cite{zhang20244diffusion} fine-tunes a video diffusion model with a learnable motion module on synthetic multiview data so that it produces consistent multiview video to guide generation, and SV4D~\cite{xie2024sv4d} completes missing views and frames with a diffusion model and uses the resulting multiview video to supervise the reconstruction.

\textbf{Gaussian Splatting-based Methods.} 4DGen~\cite{yin20234dgen} grounds 4D Gaussian generation in the input video with spatial-temporal pseudo labels; Efficient4D~\cite{pan2024efficient4d} converts the video into a multiview image matrix with a multiview-consistent generator (SyncDreamer) for rapid reconstruction; SC4D~\cite{wu2024sc4d} first fits sparse control Gaussians for coarse shape and motion and then drives dense Gaussians through LBS; and L4GM~\cite{ren2024l4gm} predicts the 4D Gaussians of a whole clip in a single feed-forward pass, trained on rendered animated objects.


\subsubsection{Benchmarks}
\label{subsubsec:p3_s8_5_3_benchmarks}
Video-to-4D has the one dedicated evaluation set in this section, released with Consistent4D~\cite{jiang2023consistent4d}: synthetic and real monocular videos with training and test splits. The metrics are those of §\ref{subsubsec:p3_s8_5_1_benchmarks}, with reconstruction fidelity measured at the input view and novel-view quality measured where synthetic ground-truth multiview video exists.

\subsubsection{Challenges and Trends}
\label{subsubsec:p3_s8_5_3_challenges}
The video constrains motion but not the unseen 3D structure, which leads to geometric distortion, depth ambiguity, and a ``painted-on'' look in novel views.
Explicit geometric structure is the main remedy: DreamMesh4D~\cite{li2024dreammesh4d} and SC4D~\cite{wu2024sc4d} constrain the reconstruction with meshes or sparse control points under LBS, which improves both coherence and efficiency, and combining such structure with the appearance modeling of neural representations is the obvious next step.

% ======================================================================
% ======================================================================

\subsection{3D-to-4D Generation}
\label{subsec:p3_s8_5_4}
3D-to-4D generation animates a static 3D model, so that the geometry is given and only motion and deformation must be synthesized, which makes the large stock of existing 3D assets usable for simulation, virtual reality, and character animation.
Unconditional generators of dynamic neural fields, such as HyperDiffusion~\cite{erkocc2023hyperdiffusion}, which runs diffusion directly on the weights of neural fields, sample 4D content without a given asset; 3D-to-4D instead starts from one, either through explicit mesh deformation and rigging or through neural and physics-based deformation models.

\subsubsection{Task Definition}
\label{subsubsec:p3_s8_5_4_task}
Here the condition is itself a 3D representation $\mathbf{c}=\mathbf{r}_0\in\mathcal{R}$, so only its temporal evolution must be generated. We therefore specialize Eq.~\eqref{eq:p3_s8_5_objective} to a \emph{deformation operator} that takes the static asset and an optional auxiliary driving signal $\mathbf{c}_{\text{aux}}$ (text, audio, or a motion sketch) to a sequence,
\begin{equation}
\label{eq:p3_s8_5_4_objective}
\hat{\mathbf{r}}^{1:T}=\mathcal{D}_\theta\!\left(\mathbf{r}_0,\,\kappa(\mathbf{c}_{\text{aux}})\right),\qquad \hat{\mathbf{r}}^{0}=\mathbf{r}_0,
\end{equation}
with the first frame pinned to the input asset. This places 3D-to-4D in the same family as the editing operators of §\ref{subsec:p3_s8_2_5} and §\ref{subsec:p3_s8_4_4}: it transforms an existing representation rather than synthesizing one from scratch.

\subsubsection{Methodology}
\label{subsubsec:p3_s8_5_4_method}
Methods realize the deformation operator $\mathcal{D}_\theta$ of Eq.~\eqref{eq:p3_s8_5_4_objective} either with explicit mesh structure, increasingly through learned rigging, or with neural and physics-informed deformation fields.

\textbf{Mesh-based Methods.} Meshes give explicit structure and precise control over deformation. Motion3DGAN~\cite{otberdout2022sparse} generates 3D facial animation from a static face mesh by first predicting the displacement of facial landmarks for a target expression and then deforming the full mesh, with a GAN keeping the sequence temporally coherent. A recent line makes arbitrary meshes animatable by predicting the rig itself: MagicArticulate~\cite{song2025magicarticulate} and RigAnything~\cite{liu2025riganything} predict skeletons and skinning weights feed-forward, Puppeteer~\cite{song2025puppeteer} rigs a mesh and then animates it by fitting the rig to videos generated from renderings of the mesh, and AnimateAnyMesh~\cite{wu2025animateanymesh} generates text-driven mesh animation in a single forward pass with a rectified-flow model over a dynamic-mesh latent.

\textbf{NeRF-based Methods.} ElastoGen~\cite{feng2024elastogen} generates physically grounded elastodynamic behavior: it rasterizes a static 3D model and then predicts and refines deformation fields iteratively with recurrent and 3D convolutional modules, so that the motion obeys physical principles.

\subsubsection{Benchmarks}
\label{subsubsec:p3_s8_5_4_benchmarks}
No benchmark with ground-truth dynamic sequences exists; methods are evaluated on a few individual assets, often faces or generic objects, mostly by user study and with the metrics of §\ref{subsubsec:p3_s8_5_1_benchmarks} where a reference animation is available. Physical plausibility, specific to this subfield, can be measured for physics-informed methods such as ElastoGen~\cite{feng2024elastogen} as deviation from a high-fidelity simulator or as violation of conserved quantities such as energy and momentum.

\subsubsection{Challenges and Trends}
\label{subsubsec:p3_s8_5_4_challenges}
Even with the geometry given, the space of plausible animations of a static model is unbounded, so without guidance the task is unconstrained, and current methods specialize to narrow categories such as faces.
Physical priors constrain it (ElastoGen~\cite{feng2024elastogen} grounds motion in elastodynamics), learned rigging makes arbitrary assets animatable, and pairing the asset with text, audio, or motion sketches turns open-ended generation into controllable animation.

% ======================================================================
% ======================================================================


\subsection{Multi-conditional 4D Generation}
\label{subsec:p3_s8_5_5}
Multi-conditional 4D generation steers dynamic scene synthesis with several signals at once: semantic conditions (text), appearance conditions (images, video), structural conditions (skeletal poses, explicit trajectories), and rule-based priors (differentiable physics). The central design problem is to fuse them so that each is respected without one overpowering the others.

\subsubsection{Task Definition}
\label{subsubsec:p3_s8_5_5_task}
We generalize the single-condition cases of Eqs.~\eqref{eq:p3_s8_5_1_objective}--\eqref{eq:p3_s8_5_4_objective} to a set of $K$ conditions, each mapped by its own encoder $\kappa_k$ and fused into the shared embedding:
\begin{equation}
\label{eq:p3_s8_5_5_objective}
\hat{\mathbf{r}}^{1:T}\sim p_\theta\!\left(\mathbf{r}^{1:T}\mid \kappa_1(\mathbf{c}_1),\dots,\kappa_K(\mathbf{c}_K)\right),
\end{equation}
which recovers a single-modality sampler when $K=1$.
Users can then specify not only what is generated (semantics) but how it behaves (motion path, physical behavior, artistic style), which complex applications such as character animation and driving simulation require.

\subsubsection{Methodology}
\label{subsubsec:p3_s8_5_5_method}
Methods fuse the encoders $\kappa_k$ of Eq.~\eqref{eq:p3_s8_5_5_objective}, usually combining score distillation (Eq.~\eqref{eq:p3_s8_5_sds}) with explicit supervision (Eq.~\eqref{eq:p3_s8_5_explicit}) tailored to each condition.

\textbf{Motion and Trajectory Control.} TC4D~\cite{bahmani2024tc4d} lets users define a spline trajectory that guides both the global rigid transformation and the local deformation of an object, with text supplying the semantics. For human animation, skeletons are the natural condition: Beyond Skeletons~\cite{yang2024beyond} integrates text and image conditions with the skeleton to refine motion, and STAR~\cite{chai2024star} uses skeleton-aware SDS and motion retargeting so that avatar movement is both visually coherent and anatomically accurate.

\textbf{Physics-Informed Control.} Sync4D~\cite{fu2024sync4d} and Phy124~\cite{lin2024phy124} integrate differentiable Material Point Method (MPM) simulation, so that the generated content follows realistic dynamics such as elasticity and collisions under video or text guidance.

\textbf{User-Centric and Artistic Control.} Generative Rendering~\cite{cai2024generative} uses a pre-existing dynamic mesh to guide a pretrained text-to-image model, giving precise control over the style and detail of the rendered output, and sketch-based interfaces~\cite{yang2024sketch} let users draw the dynamics they want.

\textbf{Complex, Application-Specific Control.} MagicDrive3D~\cite{gao2024magicdrive3d} generates dynamic driving scenes from 2D bounding boxes, Bird's-Eye-View (BEV) maps, agent trajectories, and text together, producing tightly controlled traffic scenarios.


\subsubsection{Benchmarks}
\label{subsubsec:p3_s8_5_5_benchmarks}
No benchmark exists; building one would require 4D assets annotated with several aligned conditions (\eg, a clip with a description, a skeletal motion sequence, and a trajectory map). Text adherence is scored with the CLIP Score of §\ref{subsubsec:p3_s8_5_1_benchmarks}, while adherence to sketches or trajectories relies on task-specific errors (\eg, distance from the specified trajectory) or on user studies, which remain the only holistic judge of how well a combination of inputs was followed.

\subsubsection{Challenges and Trends}
\label{subsubsec:p3_s8_5_5_challenges}
Semantic guidance from text, geometric constraints from skeletons, and rule-based priors from physics can contradict one another, and keeping each respected without one overpowering the rest while preserving fidelity is the open problem.
Domain-specific stacks such as MagicDrive3D~\cite{gao2024magicdrive3d} sidestep it by fixing the set of inputs for one application, and interactive interfaces such as sketches~\cite{yang2024sketch} point toward iterative creation in which several control channels are adjusted in turn.

% ======================================================================
% ======================================================================


\subsection{Summary, Overarching Challenges and Future Directions}
\label{subsec:p3_s8_5_summary}
Table~\ref{tab:p3_s8_5_subfield_comparison} compares the five subfields as specializations of the 4D sampler of Eq.~\eqref{eq:p3_s8_5_objective}: image- and video-to-4D progressively replace hallucinated structure with observed pixels, 3D-to-4D is given the geometry and generates only motion, and multi-conditional generation fuses several signals at once.

\begin{table}[!ht]
\centering
\fontsize{8}{10}\selectfont
\setlength{\tabcolsep}{1.0mm}
\caption{The five subfields of 4D spatial generation; each sampler or operator specializes Eq.~\eqref{eq:p3_s8_5_objective}.}
\label{tab:p3_s8_5_subfield_comparison}
\begin{tabular}{L{0.13\textwidth} L{0.15\textwidth} L{0.17\textwidth} L{0.25\textwidth} L{0.22\textwidth}}
\toprule
\textbf{Subfield} & \textbf{Conditioning input} & \textbf{Sampler / operator} & \textbf{Dominant method families} & \textbf{Datasets / metrics} \\
\midrule
§\ref{subsec:p3_s8_5_1} Text-to-4D & Text prompt $\mathbf{x}$ & $G_\theta^{\text{4D}}$, $\kappa=\phi$ (Eq.~\eqref{eq:p3_s8_5_1_objective}) & SDS distillation or generate-then-reconstruct; mesh / NeRF / 3DGS; LLM motion planning & Objaverse-derived; CLIP, FVD, PSNR/SSIM/LPIPS, user study \\
§\ref{subsec:p3_s8_5_2} Image-to-4D & Single image $I$ & $G_\theta^{\text{4D}}$, $\kappa(I)$ (Eq.~\eqref{eq:p3_s8_5_2_objective}) & Diffusion-prior SDS; explicit video supervision; unified / tuning-free; human-specific priors & Objaverse / in-the-wild; PSNR/SSIM/LPIPS, FVD, user study \\
§\ref{subsec:p3_s8_5_3} Video-to-4D & Monocular video $\mathbf{v}^{1:T}$ & $G_\theta^{\text{4D}}$, $\kappa(\mathbf{v}^{1:T})$ (Eq.~\eqref{eq:p3_s8_5_3_objective}) & Video-diffusion densification; explicit geometric constraints (mesh / LBS); feed-forward & Consistent4D; PSNR/SSIM/LPIPS, FVD, user study \\
§\ref{subsec:p3_s8_5_4} 3D-to-4D & Static asset $\mathbf{r}_0$ & $\mathcal{D}_\theta$ deformation (Eq.~\eqref{eq:p3_s8_5_4_objective}) & Mesh deformation and learned rigging; physics-informed deformation & Per-asset (faces / objects); FVD, physical plausibility, user study \\
§\ref{subsec:p3_s8_5_5} Multi-conditional & $K$ fused signals & $G_\theta^{\text{4D}}$, $\{\kappa_k\}$ (Eq.~\eqref{eq:p3_s8_5_5_objective}) & Trajectory / skeleton / physics / sketch control; application-specific stacks & No standard suite; CLIP, PSNR/SSIM/LPIPS, FVD, user study \\
\bottomrule
\end{tabular}
\end{table}

\paragraph{Data Scarcity and Generative Priors.}
Large, annotated 4D datasets are the field's main missing ingredient: real dynamic 3D data with multi-view consistency and multimodal annotation (text, motion capture) is expensive and technically hard to collect, which has largely ruled out end-to-end training. The field has therefore converged on distilling pretrained text-to-image, text-to-video, and text-to-3D models, the de facto way to guide 4D synthesis.

\paragraph{From Score Distillation to Explicit Supervision.}
Score distillation~\cite{poole2022dreamfusion} converges slowly, is costly and unstable, and produces multi-faced geometry and flickering textures. The generate-then-reconstruct pipeline of DreamGaussian4D~\cite{ren2023dreamgaussian4d} and EG4D~\cite{sun2024eg4d}, which first generates target videos and then reconstructs against them, is replacing per-prompt SDS, and feed-forward predictors such as L4GM~\cite{ren2024l4gm} and AnimateAnyMesh~\cite{wu2025animateanymesh} take the next step by removing per-scene optimization altogether.

\paragraph{4D Gaussian Splatting as the Default Representation.}
NeRF-based works rendered well but trained and rendered slowly. 3D Gaussian Splatting~\cite{kerbl20233d} and its temporal extension offer real-time, high-fidelity rendering with efficient gradient-based optimization, and their explicit, point-based form is easier to edit and to couple with physical simulation, so 4DGS has become the representation of choice for current methods~\cite{wu20244d, ren2023dreamgaussian4d, ling2024align}.

\paragraph{Consistency, Control, and Physical Plausibility.}
Keeping geometry, appearance, and motion coherent across all viewpoints and over long horizons remains the central difficulty, and control is moving beyond single prompts toward fine-grained manipulation of motion, appearance, and interaction. Physical plausibility is the newest goal: integrating differentiable simulators, as PhysSplat~\cite{zhao2024efficient} and Phy124~\cite{lin2024phy124} do, makes generated dynamics more realistic and predictable. Progress on all three is hard to measure while evaluation relies on small, author-curated test sets and on image and video metrics that ignore cross-view and long-range consistency; 4D-aware benchmarks and metrics are a prerequisite for the rest.

% \newpage


\section{World Modeling}
\label{sec:worldmodeling}
\label{sec:p3_s8_6}

\FramedInitial{D}{espite} its growing popularity across AI research, the notion of a ``world model'' does not yet enjoy a unified or universally agreed-upon definition.
Its conceptual roots trace back to the psychologist Kenneth Craik, who argued in his 1943 monograph \emph{The Nature of Explanation} that the mind ``carries a `small-scale model' of external reality and of its own possible actions within its head''.
In its original form, a world model referred to an individual's internal representation of external reality; contemporary AI has generalized the concept far beyond mental representations to structured models of physical and spatial environments.

Fei-Fei Li has framed the gap between current AI systems and physical understanding in these terms~\cite{Feifei2025spatial}: modern LLMs resemble ``wordsmiths in the dark'', fluent in symbols but divorced from the physical world that gives them meaning, and capable AI must instead understand, act within, and reason about three-dimensional space. On this view world models are new-generation \emph{generative} models that simulate, predict, and imagine the world's dynamics rather than only describe it, which ties world modeling directly to spatial intelligence.



Complementing Li's generative perspective, another influential formulation comes from Yann~LeCun \cite{lecun2022path}. He describes a world model as a system capable of predicting how the world will evolve given its current state and an action.
Formally, it consumes the current state $S(t)$ and an action $A(t)$, and produces a prediction of the next state $S'(t+1)$; in many cases, it also integrates external sensory observations $O(t)$ to update its internal state.
This predictive relation can be written compactly as
\begin{equation}
\label{eq:p3_s8_6_transition}
\hat{S}(t+1) = W_\theta\big(S(t),\, A(t),\, O(t)\big),
\end{equation}
where $W_\theta$ denotes the learned world model parameterized by $\theta$.
Read against the conditional-generation backbone of this chapter (Eq.~\eqref{eq:t2i_generator}), $W_\theta$ occupies the position of the generator, but its decoder is a state-transition predictor rather than a pixel or geometry renderer; this is the sense in which world modeling extends Chapter~\ref{sec:generation} from reconstruction toward prediction.
What exactly constitutes a ``state'' remains open: some work defines it as a compact latent representation in which the world model operates at an abstract semantic level, much as human cognition forecasts events without simulating every physical detail.




Although no single definition of world modeling has been universally accepted, a survey of the literature~\cite{fung2025embodied,li2025comprehensive,ding2025understanding,kong20253d} reveals three recurring paradigms, three realizations of Eq.~\eqref{eq:p3_s8_6_transition} that differ in what they predict.
\textbf{Representation world models} learn abstract latent spaces that capture the state and dynamics of the environment.
\textbf{Generative world models} explicitly simulate or generate future states of the world through data-driven or hybrid mechanisms.
\textbf{Mental world models} aim to infer and represent the mental states of humans or agents, enabling social reasoning and Theory-of-Mind (ToM) capabilities.
Figure~\ref{fig:world-modeling} contrasts them side by side.


 % \cite{}






\begin{figure*}[!t]
    \centering
    \includegraphics[width=\linewidth]{figures/web/world_modeling_tikz.pdf}
    \caption{Three paradigms of world modeling, all instances of the transition $\hat{S}(t{+}1)=W_\theta(S(t),A(t),O(t))$ of Eq.~\eqref{eq:p3_s8_6_transition}, differing in what they predict: (a) a latent state, (b) future observations, and (c) another agent's mental state and next action.}
    \label{fig:world-modeling}
\end{figure*}




\subsection{Representation World Model}
\label{subsec:p3_s8_6_rep}




\subsubsection{Definition}

A representation world model learns an internal latent description of the environment by compressing multimodal observations into structured, predictive embeddings, prioritizing the extraction of physical and semantic regularities over the generation of photorealistic pixels.
Much as LLMs learn abstract linguistic representations through self-supervised pretraining, representation world models are trained on ``world data'' to infer the principles that govern physical and semantic dynamics; the objective is to predict how the world evolves at an abstract level, not to reconstruct pixels.
The latent state is meant to encode objects and their properties (shape, size, appearance, material), their spatial relations (topology, distance, configuration), the dynamics that govern how the scene evolves, and the causal dependence of outcomes on actions. With such a state, an agent can predict over long horizons, reason, and plan without operating in high-dimensional sensory space.





\subsubsection{Multimodal Perception}

A world model draws on several sensory streams, whose encoding is the subject of Chapter~\ref{chapter:3:sec:6}. Vision dominates: images and especially video carry objects, geometry, interaction, and temporal dynamics, and modern systems rely on high-capacity encoders trained with contrastive vision--language objectives for semantic and geometric priors. Audio adds direct evidence about events and the physical state of the surroundings, which language conveys only symbolically. Touch is indispensable where vision is unreliable, occluded, or unavailable (grasping inside a bag, fine manipulation, tool use), yet it remains underexplored because tactile sensors are limited and tactile data are hard to collect at scale.




\subsubsection{World Representation Modeling}

\paragraph{Vision-oriented Modeling.}
A major line of representation world modeling builds latent spaces from visual observations: video frames are encoded into a compact latent space by self-supervised learning, and a transition model predicts future latent states from past latents and actions, which captures the dynamics of the visual world more efficiently than pixel-space prediction.
Concretely, an encoder $E_\phi$ maps observations to latent states and a predictor $P_\theta$ is trained to minimize a prediction loss defined entirely in latent space,
\begin{equation}
\label{eq:p3_s8_6_rep}
\mathcal{L}_{\text{rep}}(\theta,\phi) = \mathbb{E}\!\left[\,\big\lVert P_\theta\big(E_\phi(O_{\le t}),\, A(t)\big) - \operatorname{sg}\!\big(E_{\bar{\phi}}(O_{t+1})\big)\big\rVert^{2}\,\right],
\end{equation}
where $\operatorname{sg}(\cdot)$ denotes the stop-gradient operator and $E_{\bar{\phi}}$ a slowly updated target encoder.
This objective instantiates the transition of Eq.~\eqref{eq:p3_s8_6_transition} with the predicted state $\hat{S}(t+1)$ taken to be an abstract embedding rather than raw pixels (Figure~\ref{fig:world-modeling}(a)), which is precisely what allows representation world models to capture dynamics efficiently while discarding perceptually salient yet predictively irrelevant detail.



The \emph{World Models} architecture~\cite{ha2018world} decomposes the problem into a VAE-based vision model, an RNN-based memory model, and a lightweight controller; the vision and memory components learn latent dynamics from raw pixels without supervision, and the controller then learns its behavior inside the ``dreamed'' environment, which greatly simplifies policy learning.
The Dreamer agents carry this idea into reinforcement learning at scale: Dreamer~\cite{hafner2019dreamer} learns a latent dynamics model from pixels and trains an actor and critic entirely on imagined latent trajectories, and DreamerV3~\cite{hafner2023mastering} applies a single configuration across a wide range of domains. MuZero~\cite{schrittwieser2020mastering} learns a model that predicts only the quantities planning needs (reward, value, and policy) and plans with it by tree search, mastering Atari, Go, chess, and shogi without being given the rules.


Yann~LeCun's autonomous intelligence framework proposes the Joint Embedding Predictive Architecture (JEPA)~\cite{lecun2022path} as the core mechanism for constructing world models.
JEPA learns latent representations and predictions through self-supervised objectives rather than pixel-level reconstruction, and Eq.~\eqref{eq:p3_s8_6_rep} is its generic form.
I-JEPA~\cite{assran2023jepa} first instantiated it on images, predicting the embeddings of masked target regions from a visible context block, and V-JEPA~\cite{bardes2024vjepa} carried the same feature-prediction objective to video, predicting masked spatio-temporal embeddings \emph{in latent space} with no pixel reconstruction, text, or negative examples, and so avoiding the cost and ambiguity of pixel prediction.
V-JEPA~2~\cite{assran2025vjepa} turns this encoder into a world model. Pretrained on more than a million hours of internet video, it is extended with an action-conditioned predictor, V-JEPA~2-AC, post-trained on less than 62 hours of unlabeled robot video; deployed zero-shot on Franka arms, it picks and places objects by planning actions whose predicted embeddings reach a goal image.

The stop-gradient and slowly updated target encoder in Eq.~\eqref{eq:p3_s8_6_rep} exist only to keep the representation from collapsing, and the most recent JEPA work removes them. LeJEPA~\cite{balestriero2025lejepa} shows that an isotropic Gaussian is the optimal embedding distribution for downstream prediction and enforces it with a single sketched regularizer (SIGReg), leaving one trade-off hyperparameter and no teacher--student machinery. LeWorldModel~\cite{maes2026leworldmodel} uses the same regularizer to train an action-conditioned JEPA world model end to end from raw pixels with only two loss terms; at about 15M parameters it trains on a single GPU, plans up to 48 times faster than world models built on foundation-model encoders, and its latents encode physical quantities well enough to flag physically implausible events.
The principle has also reached language: VL-JEPA~\cite{chen2025vl} predicts continuous embeddings of the target text instead of generating tokens, and LLM-JEPA~\cite{huang2026llmjepa} adds an embedding-prediction objective to the next-token training of language models.


DINO-World~\cite{baldassarre2025back} instead freezes DINOv2 as a universal visual encoder, whose features are already semantic, geometric, and object-centric, learns a latent-space predictor from large-scale unlabeled video, and adds action-conditioned prediction to model how interventions change future visual states.



\paragraph{Hybrid World Modeling.}
Another important direction expands representation learning beyond vision by incorporating additional modalities such as LiDAR, camera trajectories, or multi-view geometric signals.
These hybrid models aim to overcome the limitations of purely visual latent spaces by injecting explicit spatial structure, depth cues, and geometric consistency constraints.
For instance, {MUVO}~\cite{bogdoll2025muvo} is a multimodal generative world model that learns jointly from synchronized camera and LiDAR streams.
By integrating LiDAR-derived depth priors into the latent world representation, MUVO achieves more reliable 3D scene prediction, improved object motion modeling, and long-horizon forecasting in autonomous driving environments.
{UniWorld}~\cite{min2023uniworld} introduces a unified pretrained architecture that merges 3D object detection, motion prediction, and scene representation learning.
Through joint training of BEV features, temporal dynamics, and multisensor context, UniWorld produces a structured latent state that simultaneously supports perception and trajectory forecasting.
{Aether}~\cite{zhu2025aether} proposes a geometry-aware unified world model where camera trajectories are treated as explicit action signals.
This allows the model to learn action-conditioned scene transitions and perform predictive visual planning.
By binding egomotion with geometric consistency constraints, Aether generates spatially coherent future predictions and supports counterfactual simulation within a 3D-aware latent space.




\paragraph{Language-oriented World Modeling.}
A final category examines language as a central medium for latent world representation.
This raises an important question: \emph{can LLMs themselves be viewed as world models?}
From a formal perspective, LLMs indeed share structural similarities with classical world models.
They maintain internal states, treat textual context as observations, accept prompts as actions that alter these states, and generate predictions that correspond to inferred future states.
Models such as Dynalang~\cite{lin2023learning} demonstrate that natural language can serve as a universal intermediate representation for encoding environment states, predicting transitions, and expressing causal relations.
LLMs also exhibit emergent knowledge of commonsense physics, intuitive causality, and human behavior derived from massive text corpora.
However, their limitations remain clear: they are trained on symbolic descriptions rather than direct interaction with the physical world; their internal dynamics reflect linguistic statistics more than physical laws; and their world representations tend to be sparse, fragmented, and disconnected from geometric or temporal continuity.
Thus, while LLMs can be interpreted as language-driven world models, their predictive capacity is primarily symbolic and semantic, rather than grounded in spatial or physical understanding.







\subsection{Generative World Model}
\label{subsec:p3_s8_6_gen}

In contrast to representation-based foundation models, generative world models focus on explicitly predicting future states or synthesizing the evolution of the environment over time.
Rather than only learning compact latent encodings, these models aim to simulate how the world might change under different conditions, capturing temporal dynamics and causal structure.
Whereas the representation paradigm of Eq.~\eqref{eq:p3_s8_6_rep} predicts in an abstract latent space, generative world models learn an explicit conditional distribution over future observations and can be unrolled autoregressively over a horizon $H$,
\begin{equation}
\label{eq:p3_s8_6_gen}
p_\theta\big(O_{t+1:t+H} \mid O_{\le t},\, A_{t:t+H-1}\big) = \prod_{k=0}^{H-1} p_\theta\big(O_{t+k+1} \mid O_{\le t+k},\, A(t+k)\big),
\end{equation}
so that sampling from Eq.~\eqref{eq:p3_s8_6_gen} realizes the transition of Eq.~\eqref{eq:p3_s8_6_transition} directly in observation space (Figure~\ref{fig:world-modeling}(b)).
This explicit synthesis is what allows generative world models to be rendered, inspected, and acted within, at the cost of having to model a high-dimensional output distribution.
Broadly, generative world models can be grouped into two main families:
\emph{physics-driven simulation}, which relies on hand-crafted rules and equations,
and \emph{learning-driven world generation}, which learns world dynamics directly from data.
The former can be viewed as the deterministic limit of the latter: it fixes the transition of Eq.~\eqref{eq:p3_s8_6_gen} to a hand-specified rule set instead of a learned distribution, trading sample diversity and coverage for exact reproducibility and verifiability.
A third route, growing quickly since 2024, sits between the two: a learned model writes the rules as code and a simulator executes them (Section~\ref{subsubsec:p3_s8_6_code}).



\subsubsection{Physics-driven Simulation}

Physics-driven simulation builds interactive virtual worlds from human-designed physical, geometric, and logical rules rather than from learned latent spaces. The environment is governed by predefined physical laws, rendering rules, collision logic, and behavioral scripts, specified top-down as explicit mathematical formulations, so the same initial conditions always produce the same outcome; this determinism, together with interpretability and verifiability, is what makes such systems suitable for safety-critical engineering.

Game engines are the canonical example. Physics engines such as Bullet, Havok, PhysX, and Box2D implement Newtonian, rigid-body, soft-body, cloth, particle, and volumetric dynamics, and the engines around them add acceleration structures (\eg, BVH trees), real-time ray tracing, and navigation meshes for path planning, which gives them a natural foothold in spatial-intelligence tasks. Film production relies on the same principle at higher fidelity: pipelines built on Pixar USD, Maya, or Houdini simulate cloth, hair, smoke, fire, and fluids with differential-equation and fluid solvers (\eg, FLIP/MantaFlow) and physically based path tracing. Engineering, architecture, and scientific visualization extend rule-based simulation to structural mechanics, electromagnetics, illumination, and material response through finite element methods, computational fluid dynamics, and photonic simulation.

For embodied AI, physics-driven simulators such as CARLA~\cite{dosovitskiy2017carla}, AirSim, and Isaac Sim~\cite{NVIDIA2021Isaac} remain the standard platforms for testing and training in autonomous driving, navigation, and manipulation, because they model lighting, collisions, dynamics, and sensor noise closely enough to yield data distributions near those of the real world. The same engines underpin digital twins built on platforms such as NVIDIA Omniverse, Unity, and Unreal Engine, which are discussed in Part~IV, Section~\ref{sec:Digital Twins}.





\subsubsection{Learning-driven World Generation}

In contrast to simulation frameworks based on explicit physical rules, learning-driven world generation seeks to acquire world structure, dynamics, and causal regularities directly from large-scale data, in a bottom-up fashion.
These models do not encode Newtonian equations, lighting models, or geometric constraints by hand; instead, world behavior emerges as statistical patterns during training.
For example, a model may never be given the analytic form of the gravity equation, yet by observing countless videos of falling objects, it can still infer the regularity that objects tend to accelerate downward.
The resulting world models are probabilistic and can exhibit complex, emergent behaviors without explicit programming.


Current learning-driven world generation follows two technical directions:
(1) video-centric generative models that learn dynamic world evolution in pixel or latent space,
and (2) 3D/4D asset generation methods that directly model spatially structured worlds in three or four dimensions.




\paragraph{Video-centric Generative World Models.}
Recent advances in video generation have renewed interest in the idea of a ``world simulator''.
Video models are intrinsically suited to encoding spatiotemporal structure, and when they produce sequences that are not only visually coherent but also consistent with basic physical intuition, they begin to resemble implicit world models.
This became especially prominent with the release of OpenAI Sora~\cite{sora2024}, whose generated videos demonstrate strong temporal consistency, long-horizon stability, and apparent physical plausibility, prompting the community to reconsider whether deep video generation might be converging towards a form of physical world modeling.
Subsequently, Google DeepMind's Genie~3~\cite{genie3} extended video generation from passive ``watchable'' content to interactive experiences: the model updates the generated world in response to user actions, enabling agent-in-the-loop, closed-loop world simulation, complementing the action-conditioned scene generators surveyed in §\ref{subsec:p3_s8_2_2}.
World Labs' RTFM~\cite{worldlabs_rtfm} further highlights real-time performance, achieving high frame-rate and low-latency prediction. Systems such as Wan~\cite{wan2025wan}, Emu3.5~\cite{cui2025emu35}, OpenAI Sora~2~\cite{sora2}, and Veo~3~\cite{veo3} continue this trajectory, gradually pushing video generation from mere visual fidelity towards world coherence, physical consistency, and interactivity.
Action-conditioned diffusion has also turned games into testbeds: GameNGen~\cite{valevski2024diffusion} simulates the game DOOM interactively in real time, and DIAMOND~\cite{alonso2024diffusionworldmodelingvisual} trains reinforcement-learning agents entirely inside a diffusion world model of Atari games, while NVIDIA's Cosmos~\cite{nvidia2025cosmos} offers pretrained video world foundation models intended to be post-trained as simulators for robotics and driving.

Video world models are particularly promising in autonomous driving. GAIA-1~\cite{hu2023gaia} proposes a generative world model tailored for driving scenarios, built around a spatiotemporal Transformer that generates future multi-view scenes based on past observations.
By implicitly capturing traffic patterns, pedestrian behavior, scene geometry, and multi-agent motion, GAIA-1 can support closed-loop evaluation, behavior prediction, and synthetic data generation, thereby reducing the cost and risk of traditional real-world testing.



\paragraph{3D/4D Asset Generation as World Modeling.}
A complementary avenue is to build world models that operate directly on structured three- or four-dimensional representations, such as 3D geometry, occupancy fields, or 4D spatiotemporal volumes.
Unlike video-based approaches that primarily model dynamics in pixel space, these methods explicitly parameterize geometry and temporal evolution, making it easier to represent interactive, spatially consistent worlds.


Genesis~\cite{Genesis} exemplifies this direction. It is not only a generative 3D/4D physics simulation platform, but also integrates a vision-language model-driven ``world-building agent'' that uses a tool API to create dynamic 4D environments (geometry plus dynamics).
These generated worlds can be rendered into multiple modalities (RGB, depth, semantics, contact events, etc.), providing rich supervision for training embodied AI systems.
Genesis thus combines a unified physics engine with generative agents for large-scale world construction.


Autonomous driving again serves as a natural application domain for 3D world models.
DriveArena~\cite{yang2025drivearena} builds a closed-loop generative simulation platform, allowing planning modules to iteratively probe and learn within synthesized 3D driving environments.
Its diffusion-model-based multi-agent scene generation can rapidly produce long, interactive driving episodes.
OccWorld~\cite{zheng2024occworld} and OccSora~\cite{wang2024occsora} adopt 4D occupancy grids as the core representation, predicting future 3D occupancy over time from continuous video and sensor streams, thereby providing a unified motion-and-structure world model.
Copilot4D~\cite{zhang2023copilot4d} tokenizes LiDAR scans and forecasts future LiDAR point clouds with discrete diffusion, learning a world model from unlabeled sensor data alone.


Generating \emph{interactive} 3D worlds brings generative world models closest to their conceptual ideal.
Marble~\cite{Feifei2025Marble} from World Labs combines action-conditioned generation, geometric consistency constraints, and multimodal inputs, so that users can act within the generated world and observe the consequences.
Its next-generation model, Atlas~\cite{worldlabs2026atlas}, is an omni world model pretrained from scratch on text, images, video, and 3D: an autoregressive diffusion transformer that conditions every generated image and depth map on an explicit camera pose and grounds them all in a shared 3D context, so that the same model generates camera-controlled video of up to a minute, reconstructs scenes from a handful of images, and exports its output as point clouds or 3D Gaussian splats. World Labs reports that it outperforms specialized open-source reconstruction models, and that it will power future versions of Marble.


\subsubsection{Code-and-Simulator World Models}
\label{subsubsec:p3_s8_6_code}
Between hand-written simulation and learned generation sits a route that has grown quickly since 2024: a language model \emph{writes} the world model as executable code, and a simulator or engine runs it. The learned component is the author of the rules, not the rules themselves,
\begin{equation}
\label{eq:p3_s8_6_code}
\mathcal{P} \sim p_\theta\big(\mathcal{P} \mid c\big), \qquad \hat{S}(t+1) = \mathrm{Sim}\big(S(t),\, A(t);\, \mathcal{P}\big),
\end{equation}
where $c$ is the conditioning context (a description, observations, or a history of interactions), $\mathcal{P}$ the generated program, and $\mathrm{Sim}$ the interpreter or engine that executes it: plain Python, a physics engine such as MuJoCo or PyChrono, or Blender. Compared with Eq.~\eqref{eq:p3_s8_6_transition}, the learned $W_\theta$ is replaced by a program that can be read, tested, and repaired, so the resulting world keeps the determinism and verifiability of physics-driven simulation while gaining the open-endedness of a generative model. The scene programs of Section~\ref{sec:p3_s8_7} rest on the same idea; the difference here is that $\mathcal{P}$ encodes how the world changes under actions, not only how it is laid out.

In one line of work, code is an agent's transition model. Code World Models~\cite{dainese2024codeworld} synthesize Python implementations of reinforcement-learning environments with a generate, improve, and fix search over programs (GIF-MCTS), evaluated on a benchmark of 18 environments, and plan with them faster and more sample-efficiently than with learned dynamics; WorldCoder~\cite{tang2024worldcoder} builds its world model as a program from its own interactions and transfers knowledge to a new environment by editing that program.

In the other, code builds a simulated 3D world. SceneCraft~\cite{hu2024scenecraft} turns a description into Blender Python, planning with a scene graph, converting spatial relations into numerical layout constraints, and checking renders with a vision-language model. GenSim~\cite{wang2023gensim} and RoboGen~\cite{wang2023robogen} have language models write robotic simulation tasks, scenes, and rewards, making the simulator an open-ended source of training data, and the world-building agent of Genesis~\cite{Genesis} follows the same pattern. Recent systems close the loop with the simulator itself. 
Code2Worlds~\cite{zhang2026code2worlds} generates 4D worlds as Blender code refined by a vision-language motion critic; Code2Worlds-2~\cite{zhang2026code2worlds2} compiles such worlds into physically valid MuJoCo scenes and revises them from penetration, contact, and stability signals; and ChronoAgentic~\cite{wang2026chronoagentic} has cooperating agents write and repair PyChrono simulations, meeting the joint semantic and physical correctness criterion on 82.5\% of PhyWorldBench demos against 52.5\% for the best of ten text-to-video models.

The price is coverage. 
A program can express only what its simulator supports, so photorealistic appearance, deformable and fluid materials, and open-ended visual variety still favor learned generators, while code-and-simulator models win where physical correctness has to be checked.


\subsection{Mental World Modeling}
\label{subsec:p3_s8_6_mental}


Physical world models enable agents to understand and predict how the external environment evolves.
However, intelligent behavior in real-world settings also requires an internal model of \emph{other minds}.
Humans naturally construct such models, \ie, often referred to as \emph{mental models}, to reason about others' beliefs, desires, intentions, emotions, and social context.
This ability, commonly associated with \emph{Theory of Mind (ToM)} \cite{astington1995theory}, underlies human social interaction, collaboration, communication, and decision-making.



Craik's ``small-scale models'' of reality are the origin of this notion too~\cite{nersessian1993theoretician,staggers1993mental}, and Johnson-Laird later argued that human reasoning, problem solving, and causal inference rely on abstract internal representations of situations rather than on replaying detailed physical dynamics.
These representations selectively capture the most relevant structure, \ie, objects, relations, causal dependencies, and possible outcomes, while abstracting away low-level perceptual details.
For example, when cycling, a human does not mentally simulate every joint torque or physical trajectory; instead, one imagines ``whether I would fall or not'' when encountering an obstacle.
This high-level, conceptual simulation is a defining property of mental world modeling.



\subsubsection{From Physical to Cognitive Mental World Models}
Human mental simulation operates by combining current observations with stored knowledge to predict plausible futures.
Perception provides the immediate sensory input, which is then transformed into an internal symbolic and probabilistic representation of the world.
Within this representation, the brain runs rapid, approximate simulations that evaluate different hypothetical actions.
These simulations are not physically precise but are cognitively meaningful: they help determine risks, anticipate others' responses, or evaluate alternative courses of action.
In this sense, mental models serve as a compact, abstract, and counterfactual ``world engine'' for decision making.


In artificial systems, the distinction between physical and mental world models mirrors that in human
cognition.
Physical world models predict state transitions in the environment (\eg, dynamics, geometry, object motion), whereas \emph{mental world models} encode the latent mental states of human users or other agents.
These models allow an AI system to reason about what humans know, believe, want, or intend to do, capabilities that embodied agents, collaborative robots, intelligent tutors, and multi-agent planners all need.



Inspired by the Belief-Desire-Intention (BDI)~\cite{rao1995bdi,georgeff1998belief} framework, a mental world model for AI typically includes a small set of structured variables.
\textbf{Beliefs.} These are representations of what the human knows, perceives, or believes about the world.
\textbf{Desires.} These capture user goals, preferences, motivations, and values.
\textbf{Intentions.} These are planned or ongoing actions that the agent expects the human (or another agent) to follow.
Such decomposition enables an AI system to explicitly reason about human rationality, goal pursuit, and action planning.
Formally, a mental world model can be cast as the inference of a latent mental state $M(t)$, comprising the beliefs, desires, and intentions above, from observed behavior, followed by a prediction of the other agent's next action,
\begin{equation}
\label{eq:p3_s8_6_mental}
\hat{M}(t) = \argmax_{M}\, p\big(M \mid O_{\le t}\big), \qquad \hat{A}^{\text{other}}(t+1) = \pi\big(\hat{M}(t)\big).
\end{equation}
This formulation mirrors the transition of Eq.~\eqref{eq:p3_s8_6_transition} but replaces the physical state with the inferred mental state of another agent (Figure~\ref{fig:world-modeling}(c)), so that the predicted quantity is no longer how the environment will evolve but how a human or peer agent is likely to act.

Machine Theory of Mind is now measured directly. ToMi~\cite{le2019tomi} tests false-belief reasoning in short stories with controls against shortcut solutions, BigToM~\cite{gandhi2023bigtom} generates controlled social scenarios with language models to probe whether models infer beliefs from percepts and actions, and MMToM-QA~\cite{jin2024mmtomqa} poses belief and goal questions about a person acting in a household from video and text together. Large language and multimodal models remain unreliable on these tests, whereas MMToM-QA's BIP-ALM, which performs Bayesian inverse planning over a symbolic state with a language model scoring the likelihoods, does better; it is an explicit instance of the inference in Eq.~\eqref{eq:p3_s8_6_mental}.

\paragraph{Dimensions of Mental State Representation.}
Beyond the belief-desire-intention core, a mental world model for embodied agents must span several further layers of internal state and social context. 
It has to read \emph{affect}, recognizing emotional cues and how they shift a person's risk tolerance and willingness to cooperate. 
It has to represent \emph{social dynamics}: the relationships, roles, group structures, and cultural norms that shape behavior. 
And it has to interpret \emph{communication}, both the verbal channel and the non-verbal one carried by tone, gesture, posture, and facial expression. 
These layers are what turn a bare goal estimate into a socially situated account of why a person acts as they do.



Human spatial reasoning is social as well as geometric, which ties mental modeling to spatial intelligence.
To assist or collaborate with humans in homes, factories, classrooms, or vehicles, an agent must infer the intention behind a spatial action (\eg, where a person intends to move next), anticipate trajectories from goals or habits, adapt navigation or manipulation to user preferences, and interpret spatial configurations that reflect what a person believes or misunderstands.


\subsubsection{A Joint Physical-Mental Formulation}
\label{subsubsec:p3_s8_6_mwm}
The benchmarks above still treat the mental state as something inferred next to a world model, not as part of it: Eq.~\eqref{eq:p3_s8_6_mental} estimates another agent's mind, but the environment it acts in is predicted separately. 
Mental World Modeling (MWM)~\cite{fei2026mental} makes the mental state part of the state that the world model of Eq.~\eqref{eq:p3_s8_6_transition} rolls forward. 
We follow its notation, in which subscripts index time, so that $s_t$ and $a_t^{\epsilon}$ play the roles of $S(t)$ and $A(t)$ above.
The case for a joint state is that neither half suffices alone. For a target agent $\epsilon$ whose true action distribution is $p^{\star}$, a physical-only representation fails whenever
\begin{equation}
\label{eq:p3_s8_6_mwm_property}
p^{\star}\big(a_t^{\epsilon} \mid s_t^{\mathrm{phy}}, s_t^{\mathrm{ment}}\big) \neq p^{\star}\big(a_t^{\epsilon} \mid s_t^{\mathrm{phy}}, \tilde{s}_t^{\mathrm{ment}}\big),
\end{equation}
that is, whenever two scenes that are physically identical but differ in what people believe or want ($s_t^{\mathrm{ment}}$ versus $\tilde{s}_t^{\mathrm{ment}}$) lead to different actions; the symmetric condition, with the mental state held fixed and the physical state changed, rules out a mental-only representation.

\paragraph{Joint State.}
The world state is the pair
\begin{equation}
\label{eq:p3_s8_6_mwm_state}
s_t = \big(s_t^{\mathrm{phy}}, s_t^{\mathrm{ment}}\big) \in \mathcal{S}^{\mathrm{phy}} \times \mathcal{S}^{\mathrm{ment}},
\qquad
s_t^{\mathrm{ment}} = \big(\{m_t^{i}\}_{i\in\mathcal{I}},\, \{m_t^{G}\}_{G\in\mathcal{G}},\, R_t^{\mathrm{ment}},\, \alpha_t\big).
\end{equation}
The physical part collects the objects, the embodied characters, their spatial and contact relations, and the environmental conditions of the scene. The mental part holds an individual state $m_t^{i}$ for every agent $i$ in the set $\mathcal{I}$, optional states $m_t^{G}$ for groups $G\in\mathcal{G}$, the social attitudes and role relations $R_t^{\mathrm{ment}}$ between agents, and the atmosphere $\alpha_t$ of the setting. Each individual state extends the belief-desire-intention core introduced above,
\begin{equation}
\label{eq:p3_s8_6_mwm_individual}
m_t^{i} = \big(\mathrm{id}_t^{i},\, b_t^{i},\, q_t^{i},\, g_t^{i},\, \iota_t^{i},\, e_t^{i},\, d_t^{i},\, n_t^{i},\, c_t^{i}\big),
\end{equation}
with identity attributes, beliefs, attention focus, goals, intentions, emotions, dispositions, norms, and behavioral constraints. 
Beliefs, goals, and intentions are the BDI variables; emotions, norms, and the relations $R_t^{\mathrm{ment}}$ turn the affect and social-dynamics layers discussed above into explicit state variables.

\paragraph{Target-Specific Observation.}
The model reasons from the perspective of one target agent $\epsilon\in\mathcal{I}$, which sees only a rendering of the joint state:
\begin{equation}
\label{eq:p3_s8_6_mwm_obs}
o_t^{\epsilon} = \Omega_{\epsilon}(s_t) = \Big(\Omega_{\epsilon}^{\mathrm{phy}}\big(s_t^{\mathrm{phy}};\, \kappa_t^{\epsilon}\big),\; \Omega_{\epsilon}^{\mathrm{ment}}\big(s_t^{\mathrm{phy}}, s_t^{\mathrm{ment}};\, \rho_t^{\epsilon}\big)\Big),
\end{equation}
where $\kappa_t^{\epsilon}$ is the target's perceptual access (position, line of sight, hearing range, available modalities) and $\rho_t^{\epsilon}$ its social and cognitive perspective (role, prior beliefs, relationships to others, awareness of norms, Theory-of-Mind capacity). The mental part of the observation contains the target's view of its own state and its estimates $\widehat{m}^{\,j,(\ell)}_{t,\epsilon}$ of every other agent $j$, nested up to a depth $L$; at the first level they are sampled as
\begin{equation}
\label{eq:p3_s8_6_mwm_tom}
\widehat{m}^{\,j,(1)}_{t,\epsilon} \sim Q^{(1)}_{\epsilon}\big(m_t^{j} \mid o_t^{\epsilon,\mathrm{phy}},\, m_t^{\epsilon},\, R_t^{\mathrm{ment}},\, \alpha_t\big),
\end{equation}
which is the inference of Eq.~\eqref{eq:p3_s8_6_mental}, made stochastic and conditioned on the target's own mind and social context. Spatial access enters the mental model directly through $\kappa_t^{\epsilon}$: what the target can see from where it stands bounds what it can believe, so a poor decision can follow from an object outside its view as much as from a mistaken belief.

\paragraph{Action and Coupled Transition.}
The target proposes $K$ candidate actions $\mathcal{A}_t^{\epsilon} = \Pi_{\epsilon}^{K}(o_t^{\epsilon}) = \{a^{\epsilon}_{t,1},\dots,a^{\epsilon}_{t,K}\}$. Each action has a physical carrier and a mental content, $a_t^{\epsilon} = (a_t^{\epsilon,\mathrm{phy}}, a_t^{\epsilon,\mathrm{ment}})$ (approaching someone, and meaning to comfort them), and the policy factorizes as $\Pi_{\epsilon}(a_t^{\epsilon,\mathrm{phy}}, a_t^{\epsilon,\mathrm{ment}} \mid o_t^{\epsilon}) = \Pi_{\epsilon}^{\mathrm{phy}}(a_t^{\epsilon,\mathrm{phy}} \mid o_t^{\epsilon})\,\Pi_{\epsilon}^{\mathrm{ment}}(a_t^{\epsilon,\mathrm{ment}} \mid a_t^{\epsilon,\mathrm{phy}}, o_t^{\epsilon})$. The world model then predicts the next joint state with a staged transition,
\begin{equation}
\label{eq:p3_s8_6_mwm_transition}
T_{\theta}\big(s_{t+1} \mid s_t, a_t^{\epsilon}\big) = T_{\theta}^{\mathrm{phy}}\big(s_{t+1}^{\mathrm{phy}} \mid s_t^{\mathrm{phy}}, s_t^{\mathrm{ment}}, a_t^{\epsilon,\mathrm{phy}}\big) \cdot T_{\theta}^{\mathrm{ment}}\big(s_{t+1}^{\mathrm{ment}} \mid s_t^{\mathrm{phy}}, s_t^{\mathrm{ment}}, a_t^{\epsilon,\mathrm{phy}}, a_t^{\epsilon,\mathrm{ment}}\big).
\end{equation}
The asymmetry is deliberate. An intention moves nothing until it is carried out physically, so the physical factor ignores $a_t^{\epsilon,\mathrm{ment}}$; the mental factor depends on both components, because the same reassuring words land differently when the speaker is approaching than when walking away. The physical factor still conditions on the mental state, since what objects, space, and visibility allow is joined by what other agents intend to do. Those other agents are not controlled: their actions $a_t^{-\epsilon}$ are marginalized under the simulator's predictive distribution,
\begin{equation}
\label{eq:p3_s8_6_mwm_marginal}
T_{\theta}\big(s_{t+1} \mid s_t, a_t^{\epsilon}\big) = \int T_{\theta}\big(s_{t+1} \mid s_t, a_t^{\epsilon}, a_t^{-\epsilon}\big)\, p\big(a_t^{-\epsilon} \mid s_t\big)\, \mathrm{d}a_t^{-\epsilon},
\end{equation}
with $p(a_t^{-\epsilon} \mid s_t) = \prod_{j\neq\epsilon} \Pi_j\big(a_t^{j} \mid \Omega_j(s_t)\big)$ for simultaneous moves, so that every agent still acts on its own partial observation. The whole pipeline, $s_t \xrightarrow{\Omega_{\epsilon}} o_t^{\epsilon} \xrightarrow{\Pi_{\epsilon}} a_t^{\epsilon} \xrightarrow{T_{\theta}} s_{t+1}$, is Eq.~\eqref{eq:p3_s8_6_transition} with a mental component added to the state and a perspective added to the observation.

\paragraph{Evaluating Candidate Futures.}
Simulating each candidate gives a branch $\hat{s}^{k}_{t+1}$, and the predicted decision is the branch of highest value,
\begin{equation}
\label{eq:p3_s8_6_mwm_select}
\hat{k} = \argmax_{k\in\{1,\dots,K\}}\, r_k .
\end{equation}
In MENTIS, the training-free and fully inspectable baseline of MWM, the branch value $r_k$ is a weighted combination of normalized scores for physical plausibility, mental consistency, and social appropriateness, and a binary safety and legality veto sets it to zero. Each formal component maps to a module: a state parser recovers $\hat{s}_t$, an observation generator implements $\Omega_{\epsilon}$, an action parser structures the candidates, a world-state transitor implements $T_{\theta}$, and an evaluator assigns $r_k$. On situated decision scenarios presented as text, image, and sounding-video stories, experiments across eight LLM-based world models indicate that the explicit mental component is what makes a person's actual decision predictable~\cite{fei2026mental}.


\subsection{Discussions and Challenges}
\label{subsec:p3_s8_6_discussion}

\paragraph{Representation-based World Models as an Alternative Foundation Model Paradigm.}
Representation world models may emerge as a parallel paradigm to current LLMs, functioning as large-scale foundational models trained through massive self-supervised representation learning to capture physical regularities, causal dynamics, and environmental constraints.
Unlike LLMs, which learn the world indirectly through linguistic abstraction, representation world models aim to model the physical world directly from multimodal observations.
Such models could serve video generation, action prediction in VLA systems, 3D and 4D understanding and generation, and decision-making in autonomous driving.



\paragraph{The Necessity of Multi-source Intelligence Integration.}
World models will likely require the integration of heterogeneous informational sources.
On the one hand, LLMs embody language-driven intelligence, while world models emphasize physically grounded intelligence; rather than being mutually exclusive, a truly capable world model must incorporate both forms of reasoning.
On the other hand, accurate world understanding requires far more than visual perception alone: audio, speech, tactile signals, force feedback, and even proprioception provide essential supplementary cues about the dynamics and structure of the physical environment. Yet the current research community focuses overwhelmingly on visual-only modeling.
Building world models that holistically integrate multimodal sensory information remains a major challenge.


\paragraph{Unifying Physical Modeling and Mental Modeling.}
Another promising direction is the integration of physical and mental modeling within a unified world-modeling framework.
Embodied agents operating in real environments must simultaneously reason about the physical evolution of the environment and the mental or behavioral responses of humans (or other agents).
Recent work in multi-agent systems, social navigation, pedagogical reasoning, human-robot collaboration, and interactive LLM agents emphasizes the growing need for AI systems to build internal hypotheses about both external dynamics and human mental states. 
Mental World Modeling~\cite{fei2026mental} is an early formulation of such a joint state (Section~\ref{subsubsec:p3_s8_6_mwm}).
Ultimately, scalable general-purpose world models must simulate not only \emph{what will happen} but also \emph{what others believe, intend, desire, or misunderstand}.


\paragraph{Balance between Generative Fidelity and Physical Correctness.}
A central challenge for learning-driven world models is reconciling generative expressiveness with physical accuracy.
High-capacity video generators often produce visually stunning but physically implausible outcomes, while models optimized for strict physical consistency may underperform in realism or diversity.
Achieving stable long-horizon predictions without temporal drift or compounding errors remains unsolved.
Designing architectures, training objectives, and constraints that jointly enforce physical laws, spatial coherence, and generative flexibility is a key open problem for the next generation of world simulators.


\paragraph{Evaluation, Benchmarks, and Verifiability of World Models.}
Unlike language models, which can be evaluated through perplexity, benchmarks, or task accuracy, world models lack standardized, reliable, and comprehensive evaluation protocols.
Physical correctness, causal consistency, interactive fidelity, multimodal coherence, and long-horizon stability are difficult to quantify in a unified manner.
Furthermore, world models that aim to simulate open-ended environments raise issues of verifiability, reproducibility, and safety: how can we assess whether a world model reliably generalizes to unseen real-world scenarios, or whether its internal dynamics exhibit hidden biases or unsafe behaviors?
Establishing principled evaluation methodologies and shared benchmarks for both representation-based and generative world models is an urgent research direction.

% \newpage
\section{Programmatic and Symbolic Spatial Generation}
\label{sec:p3_s8_7}

\FramedInitial{P}{rogrammatic} and symbolic spatial generation synthesizes 2D and 3D scene structure through explicit rules, executable programs, and symbolic geometric abstractions.
It generates the structural backbone of a scene (spatial layout, hierarchical organization, geometric constraints, executable logic) rather than its raw appearance, which gives it controllability, compositional generalization, global spatial consistency, and interpretability.

Where the largely neural samplers of §\ref{sec:p3_s8_1}–§\ref{sec:p3_s8_6} map a condition directly to pixels, voxels, or a neural latent, programmatic methods emit an intermediate \emph{program} that a deterministic renderer executes into the final scene.
This intermediate layer is what makes the structure auditable, editable, and compositionally reusable, properties that flat neural outputs lack.
Three paradigm families mark the historical and technical progression of the field: \emph{rule- and grammar-based procedural generation}, \emph{symbolic graphics programs and domain-specific languages}, and \emph{neuro-symbolic language-to-program synthesis} driven by large language models.

\subsection{Task Definition}
\label{subsubsec:p3_s8_7_task}

Let $\mathbf{c}\in\mathcal{C}$ denote a conditioning intent, typically a natural-language description, a layout specification, or a partial constraint set, and let $\kappa:\mathcal{C}\!\to\!\mathcal{C}'$ be its encoder.
A programmatic generator $G_\theta$ emits a symbolic program $\rho$ drawn from a graphics language or domain-specific language $\mathcal{L}$, which a deterministic executor $R$ evaluates into a concrete scene $\hat{\mathbf{s}}$ (a raster image, vector graphic, UI specification, or 3D shape):
\begin{equation}
\label{eq:p3_s8_7_objective}
\rho = G_\theta\big(\kappa(\mathbf{c})\big), \qquad \hat{\mathbf{s}} = R(\rho), \qquad \rho \in \mathcal{L}.
\end{equation}
This formulation inherits the conditional-sampler abstraction of Eq.~\eqref{eq:t2i_generator} but replaces its learned decoder $D$ (Eq.~\eqref{eq:t2i_autoencoder}) over a continuous latent space $\mathcal{Z}$ with an interpretable, discrete program $\rho$ and an \emph{exact, non-learned} renderer $R$.
The neural generators of §\ref{sec:generation} are recovered as the limiting case in which $R$ is itself a neural renderer and $\rho$ a continuous latent, whereas the symbolic case constrains $\rho$ to be a human-readable executable artifact.
Two properties follow directly: errors are localizable to specific program statements, and structural constraints can be enforced or verified deterministically on $\rho$ before rendering.



\subsection{Methodology}
\label{subsubsec:p3_s8_7_methodology}

Figure~\ref{fig:p3_s8_7_paradigms} schematizes the three paradigm families that instantiate Eq.~\eqref{eq:p3_s8_7_objective}, contrasting how each obtains the program $\rho$ before the shared deterministic renderer $R$ executes it.

\begin{figure}[t]
    \centering
    \includegraphics[width=\linewidth]{figures/part3/section8/programmatic_spatial_generation.pdf}
    \caption{Three paradigm families for programmatic spatial generation, all producing a program $\rho$ that a deterministic renderer $R$ executes (Eq.~\eqref{eq:p3_s8_7_objective}): (P1) rule- and grammar-based procedural generation, (P2) symbolic graphics programs and domain-specific languages, and (P3) language-to-program synthesis by an LLM planner.}
    \label{fig:p3_s8_7_paradigms}
\end{figure}

\paragraph{Rule- and Grammar-based Procedural Generation.}
The earliest line of work, procedural content generation (PCG)~\cite{parish2001procedural} in computer graphics and game development, synthesizes scenes such as dungeon maps, terrains, architectural plans, or UI layouts from hand-crafted rules or parameterized procedural programs.
Classical formalisms instantiate $G_\theta$ as a fixed production system rather than a learned model.
The Lindenmayer system (L-system)~\cite{lindenmayer1968mathematical}, originally developed to model biological growth, employs parallel rewriting rules to simulate natural growth processes and fractal patterns, making it ideal for self-similar structures such as vegetation and complex terrain.
Formally, the program evolves by repeated parallel application of a production set $\Pi$ to an axiom $\omega$,
\begin{equation}
\label{eq:p3_s8_7_grammar}
\rho_0 = \omega, \qquad \rho_{t+1} = \Pi(\rho_t), \qquad \hat{\mathbf{s}} = R(\rho_T),
\end{equation}
where the generated symbolic string is not the geometry itself but a sequence of Turtle Graphics commands that the interpreter $R$ parses to draw the final shape in 2D or 3D space.
Shape grammars, introduced by \citet{stiny1980introduction}, provide another influential symbolic mechanism by defining geometric primitives and rewrite rules for composing architectural or engineered shapes.
These early paradigms demonstrated a fundamental insight: spatial structures can be synthesized by executing symbolic generative rules, rather than by directly predicting raw geometry.

\paragraph{Symbolic Graphics Programs and Domain-Specific Languages.}
Symbolic graphics programs (SGPs) encode spatial layouts and geometric configurations as formal programs that execute deterministically into images, vector graphics, UI designs, or 3D shapes, bridging abstract scene intent and concrete visual realization. Common languages such as SVG~\cite{bing2025learning}, TikZ/LaTeX~\cite{ellis2018learning}, and HTML/CSS~\cite{chen2025designcoder} already carry rich symbolic layout specifications in real-world use.
In parallel, researchers have developed domain-specific languages (DSLs) tailored for fine-grained spatial and visual synthesis.
Examples include hierarchical GUI-code generators~\cite{zhu2018automatic} that translate visual structures into programmatic UI specifications, and ShapeAssembly~\cite{jones2020shapeassembly}, a 3D assembly language in which objects are represented as executable programs describing part composition, symmetry, and connectivity.
Such DSLs make the advantages of explicit symbolic control concrete: precise geometry, modular design, verifiable constraints, and systematic compositionality.

\paragraph{Neuro-symbolic and Language-to-Program Synthesis.}
A more recent direction couples symbolic programs with learned models, so that $G_\theta$ becomes a trained synthesizer rather than a fixed rule set. The Scene Language~\cite{zhang2025scenelanguage} unifies neural and symbolic components: a program specifying relational and hierarchical layout structure, natural-language labels summarizing entity semantics, and learned embeddings capturing visual identity.
Executable structure, rather than raw geometry, thus becomes a primary representational layer for scenes.
Large language models~\cite{brown2020language,achiam2023gpt} have accelerated the move toward language-to-program spatial generation: instead of predicting images or layouts directly, LLM-based systems act as symbolic planners that synthesize the program by approximately maximizing its conditional likelihood,
\begin{equation}
\label{eq:p3_s8_7_synthesis}
\rho^{\star} = \argmax_{\rho \in \mathcal{L}} \, p_\theta\big(\rho \mid \kappa(\mathbf{c})\big),
\end{equation}
mapping natural-language descriptions to executable layout programs~\cite{aguina2024open, hsu2025programs, yang2024holodeck}.
These programs then guide downstream renderers to produce UI designs, graphic layouts, architectural plans, or 3D scenes with precise structural fidelity, and users can edit or refine the underlying program with fine-grained control.

The three families differ along two cross-cutting axes.
The first is the \emph{representation} of $\rho$: a parallel string-rewriting system (grammar-based generation), an executable program in a general-purpose or domain-specific graphics language (SGPs and DSLs), or a hybrid structure that interleaves code, language, and learned embeddings (neuro-symbolic synthesis).
The second is the \emph{inference procedure} that produces $\rho$: hand-authored production rules, human-written or template-instantiated programs, or learned program synthesis and LLM-based planning. Across all three, the target modality (2D raster, vector graphics, UI, or 3D) is determined solely by the renderer $R$, which keeps the symbolic and rendering concerns cleanly separated.



\subsection{Benchmarks}
\label{subsubsec:p3_s8_7_benchmarks}

Evaluation in this area differs from the distributional metrics of §\ref{sec:generation} because the immediate output is a program, not an image, so quality is assessed along several complementary axes.
\textbf{Program validity} measures whether the emitted $\rho$ parses and executes without error in its target language, a prerequisite that purely neural generators do not face.
\textbf{Rendered-output fidelity} then compares the executed scene $R(\rho)$ against a reference, using pixel- or perceptual-level similarity for raster and vector targets and geometric agreement for 3D shapes.
\textbf{Constraint satisfaction} checks whether explicitly specified spatial relations, counts, alignments, or symmetries hold in the rendered result, a property that symbolic methods are designed to guarantee.
\textbf{Controllability and editability} assess how reliably localized edits to $\rho$ produce the intended localized changes in the scene without unintended global drift. Concretely, these axes are exercised across vector-graphics and SVG synthesis~\cite{bing2025learning}, UI and GUI code generation~\cite{zhu2018automatic, chen2025designcoder}, 3D program induction over part-structured shapes~\cite{jones2020shapeassembly}, evaluated against part-hierarchy benchmarks such as PartNet~\cite{mo2019partnet}, and full scene-program synthesis~\cite{zhang2025scenelanguage, yang2024holodeck}, where the executable program serves simultaneously as the prediction and as the verifiable ground for measurement.

\subsection{Challenges and Trends}
\label{subsubsec:p3_s8_7_challenges}

The field has moved from entirely hand-authored rules toward neuro-symbolic integration with large models, driven by the demand for generators that are more controllable and interpretable than end-to-end neural methods.

\textbf{Generalization.} Automatically synthesized programs still struggle with novel constraints and complex semantic specifications~\cite{mao2025neurosymbolicconcepts}. Program synthesis has to infer the structural logic of a scene rather than memorize surface patterns if the generator is to extrapolate to compositions and constraint combinations absent from training.

\textbf{Efficiency and differentiability.} Neuro-symbolic programs that are both efficient and fully differentiable~\cite{guo2023efficientsymbolicpolicy} would let a rendering-level loss propagate back through the program, so that the whole pipeline, from symbolic planning to neural rendering, could be trained and fine-tuned end to end.

\textbf{Embodied interaction.} Procedurally synthesized interactive environments already supply large-scale embodied worlds~\cite{kolve2017ai2thor, procthor, wang2024pipeline}, and the next step is for agents to generate or validate symbolic scene structure through active exploration rather than from static datasets alone, yielding symbolic representations of physical spaces that can be checked against what the agent observes.



\section{Summary and Outlook}
\label{sec:p3_s8_8}

This chapter has surveyed spatial reconstruction and generation across seven subfields that, taken together, span an axis of increasing output structure and dimensionality: 2D images (§\ref{sec:p3_s8_1}), 3D objects and scenes (§\ref{sec:p3_s8_2}), video (§\ref{sec:p3_s8_3}), embodied motion (§\ref{sec:p3_s8_4}), 4D space-time content (§\ref{sec:p3_s8_5}), world models (§\ref{sec:p3_s8_6}), and interpretable symbolic programs (§\ref{sec:p3_s8_7}). Despite their different output domains, the subfields are unified by a single conditional-generation backbone, established in §\ref{subsec:p3_s8_1_1}: an encoder $\kappa$ maps a condition $\mathbf{c}$ into a latent code, a generator $G_\theta$ samples structure from that code, and a decoder or renderer turns the sampled structure into a concrete artifact (Eq.~\eqref{eq:t2i_autoencoder}, Eq.~\eqref{eq:t2i_generator}). What distinguishes the subfields is not this backbone but the \emph{output space} and the \emph{decoder} that closes it: a pixel decoder for images, a volumetric or splatting renderer for 3D, a spatio-temporal decoder for video, a kinematic sampler for motion, a time-indexed 3D representation for 4D, a state-transition predictor for world models, and an exact, non-learned program executor for symbolic generation. Table~\ref{tab:p3_s8_chapter_comparison} consolidates this view.

\begin{table}[t]
\centering
\caption{The seven spatial-generation subfields of this chapter, all instances of the conditional latent-generator backbone of §\ref{subsec:p3_s8_1_1}.}
\label{tab:p3_s8_chapter_comparison}
\fontsize{8}{10}\selectfont
\setlength{\tabcolsep}{1.0mm}
\begin{tabular}{L{0.17\textwidth} L{0.20\textwidth} L{0.33\textwidth} L{0.22\textwidth}}
\toprule
\textbf{Subfield} & \textbf{Output space} & \textbf{Dominant paradigms} & \textbf{Decoder / renderer} \\
\midrule
§\ref{sec:p3_s8_1} 2D image & Pixel grid (raster) & GAN; autoregressive; latent diffusion; flow matching & Learned pixel decoder \\
§\ref{sec:p3_s8_2} 3D object/scene & Mesh / NeRF / 3DGS / latent & Native 3D diffusion; score distillation; multi-view lift; feed-forward reconstruction & Differentiable / volumetric renderer \\
§\ref{sec:p3_s8_3} Video & Frame sequence (clip latent) & Inflated T2I backbones; spatial / temporal control; editing & Spatio-temporal decoder \\
§\ref{sec:p3_s8_4} Motion & Pose / trajectory sequence & Regression; diffusion samplers; editing operators & Kinematic / skeletal decoder \\
§\ref{sec:p3_s8_5} 4D & Time-indexed 3D representation & Score distillation; generate-then-reconstruct & Deformable / per-frame renderer \\
§\ref{sec:p3_s8_6} World model & Latent state / future observations & Representation; generative; mental world models & Latent transition / observation predictor \\
§\ref{sec:p3_s8_7} Symbolic program & Executable program in a language $\mathcal{L}$ & Procedural grammars; SGP/DSL; LLM program synthesis & Deterministic, non-learned executor \\
\bottomrule
\end{tabular}
\end{table}

\textbf{A shared backbone, specialized by output space.} The progression from §\ref{sec:p3_s8_1} to §\ref{sec:p3_s8_7} is best read not as seven disjoint problems but as one generative abstraction instantiated against progressively richer codomains, and each section casts its own subtasks as conditional specializations of that common sampler. The symbolic case (§\ref{sec:p3_s8_7}) is the limiting instance in which the decoder degenerates to an exact program executor, trading the expressiveness of a learned renderer for interpretability and verifiable control.

\textbf{The 2D prior as a reusable engine.} The image generators of §\ref{sec:p3_s8_1} recur throughout the chapter not as terminal systems but as priors to be distilled or lifted: 3D synthesis distills frozen 2D diffusion through score distillation and multi-view consistency (§\ref{sec:p3_s8_2}), video models inflate text-to-image backbones along the temporal axis (§\ref{sec:p3_s8_3}), and 4D pipelines compose both (§\ref{sec:p3_s8_5}). A single advance in the 2D substrate therefore propagates outward across every higher-dimensional subfield, which concentrates research leverage but also propagates the substrate's failure modes.

\textbf{Physical plausibility and fragmented evaluation.} Across subfields, the recurring failure mode is not low visual fidelity but violated physics and geometry, manifest as multi-view inconsistency in 3D, temporal drift and flicker in video, foot-skating and interpenetration in motion, and dynamic implausibility in 4D. Evaluation remains correspondingly fragmented: each subfield has evolved its own metric tuples, and all are migrating toward learned and vision-language-model judges as scalar similarity scores prove insufficient for spatially structured output. A unified, physically grounded evaluation protocol that spans these subfields remains an open goal.

\textbf{Outlook.} These observations point toward convergence rather than further fragmentation. World modeling (§\ref{sec:p3_s8_6}) already acts as a conceptual capstone that absorbs reconstruction, generation, and dynamics into a single predictive substrate, while programmatic and symbolic generation (§\ref{sec:p3_s8_7}) supplies the controllability and interpretability that purely neural pipelines lack; their integration, neural generators steered and verified by symbolic structure, is a natural trajectory for the field. The capabilities surveyed here also form the generative substrate for the applications of Part~\ref{part:applications}, where embodied agents, digital twins, and interactive environments consume generated spatial content rather than producing it for its own sake.
